% 本文件由 LaTeX 排版 Agent 根据有效 translation.json 自动生成。
% author=Cyprian of Carthage; work_id=0506; title=迦太基的西彼廉书信集

\chapter{迦太基的西彼廉书信集}

% structure: p0001 source=h1 target=omitted reason=page-title-already-rendered-by-parent

\Emph{书信1}\newline{}\Emph{书信2}\newline{}\Emph{书信3}\newline{}\Emph{书信4}\newline{}\Emph{书信5}\newline{}\Emph{书信6}\newline{}\Emph{书信7}\newline{}\Emph{书信8}\newline{}\Emph{书信9}\newline{}\Emph{书信10}\newline{}\Emph{书信11}\newline{}\Emph{书信12}\newline{}\Emph{书信13}\newline{}\Emph{书信14}\newline{}\Emph{书信15}\newline{}\Emph{书信16}\newline{}\Emph{书信17}\newline{}\Emph{书信18}\newline{}\Emph{书信19}\newline{}\Emph{书信20}\newline{}\Emph{书信21}\newline{}\Emph{书信22}\newline{}\Emph{书信23}\newline{}\Emph{书信24}\newline{}\Emph{书信25}\newline{}\Emph{书信26}\newline{}\Emph{书信27}\newline{}\Emph{书信28}\newline{}\Emph{书信29}\newline{}\Emph{书信30}\newline{}\Emph{书信31}\newline{}\Emph{书信32}\newline{}\Emph{书信33}\newline{}\Emph{书信34}\newline{}\Emph{书信35}\newline{}\Emph{书信36}\newline{}\Emph{书信37}\newline{}\Emph{书信38}\newline{}\Emph{书信39}\newline{}\Emph{书信40}\newline{}\Emph{书信41}\newline{}\Emph{书信42}\newline{}\Emph{书信43}\newline{}\Emph{书信44}\newline{}\Emph{书信45}\newline{}\Emph{书信46}\newline{}\Emph{书信47}\newline{}\Emph{书信48}\newline{}\Emph{书信49}\newline{}\Emph{书信50}\newline{}\Emph{书信51}\newline{}\Emph{书信52}\newline{}\Emph{书信53}\newline{}\Emph{书信54}\newline{}\Emph{书信55}\newline{}\Emph{书信56}\newline{}\Emph{书信57}\newline{}\Emph{书信58}\newline{}\Emph{书信59}\newline{}\Emph{书信60}\newline{}\Emph{书信61}\newline{}\Emph{书信62}\newline{}\Emph{书信63}\newline{}\Emph{书信64}\newline{}\Emph{书信65}\newline{}\Emph{书信66}\newline{}\Emph{书信67}\newline{}\Emph{书信68}\newline{}\Emph{书信69}\newline{}\Emph{书信70}\newline{}\Emph{书信71}\newline{}\Emph{书信72}\newline{}\Emph{书信73}\newline{}\Emph{书信74}\newline{}\Emph{书信75}\newline{}\Emph{书信76}\newline{}\Emph{书信77}\newline{}\Emph{书信78}\newline{}\Emph{书信79}\newline{}\Emph{书信80}\newline{}\Emph{书信81}\newline{}\Emph{书信82}\par

\SourceRecord{Translated by Robert Ernest Wallis.；Ante-Nicene Fathers；Vol. 5.；Edited by Alexander Roberts, James Donaldson, and A. Cleveland Coxe.；Buffalo, NY: Christian Literature Publishing Co.,；1886.；<http://www.newadvent.org/fathers/0506.htm>.}

% structure: p0001 source=h1 target=section reason=page-title

\section{书信一}

致\Person{多纳徒}\par

\WorkTitle{论证}——\Person{西彼廉}曾应许\Person{多纳徒}，将与他谈论神圣之事，如今因提醒而履行诺言。他首先称赞在圣洗中所赐予的天主恩宠，并宣告自己如何因此改变；最后指出世界的错误，劝勉人轻视世界，并致力于读经与祈祷。\par

1. \Person{凯基利乌·西彼廉}致\Person{多纳徒}，祝安。你最亲爱的\Person{多纳徒}，你正确地提醒了我：我不但记得我的承诺，而且承认现在是履行它的适当时候，因为丰收节庆邀请心灵在安息中放松，并享受岁末每年指定的休息。此外，这个地方与季节相宜，花园的宜人景色与温和秋日的柔风相协调，使感官得到抚慰和愉悦。在这样的地方，整日谈论是令人愉快的，并且通过（神圣的）比喻来训练内心的良知以领受神圣的诫命。为使任何不敬的闯入者不打断我们的交谈，也不让吵闹家庭的无节制喧哗干扰我们，让我们寻找这个凉亭。邻近的树丛确保我们独处，而藤蔓枝条在支撑它们的芦苇间蜿蜒下垂的漫游路径，为我们造了一个藤蔓的门廊和绿叶遮蔽所。我们愉快地在这里用言语装扮我们的思想；当我们用眼睛欣赏树木和藤蔓的悦目景色时，心灵同时因所听到的而受教导，因所看到的而得滋养，尽管此时你唯一的乐趣和兴趣在于我们的谈论。轻视视觉的快乐，你的眼现在注视着我。你以心灵和耳朵全然是一位聆听者；而且是一位热切与爱成正比的聆听者。\par

2. 然而，我的心灵所能传达给你的任何东西，是怎样的或多少呢？我浅薄悟性的贫乏平庸只产生非常有限的收获，并不以丰饶的沉积使土壤富饶。尽管如此，我将凭我所有的能力着手此事；因为我即将谈论的题目本身将帮助我。在法庭上，在公共集会中，在政治辩论中，滔滔不绝的口才可能是善辩野心的荣耀；但在谈论主天主时，一种纯洁简朴的表达力求以实质，而非以口才的力量，促成信德的坚信。因此，从我这里接受并非机巧而是沉重的话语，并非以修辞装饰以取悦群众听众的言词，而是简单且因其无粉饰的真诚而适于宣告天主慈悲的言词。接受那些在说出之前已被感受到的东西，那些并非经过多年缓慢劳苦积累而成，而是一口气吸入成熟恩宠的东西。\par

3. 当我仍躺在黑暗和阴沉的夜里，飘摇不定，在这自夸时代的泡沫上被抛来抛去，对我游荡的脚步毫无把握，不知我真正的生命为何，远离真理与光明时，我曾认为这是一件难事，尤其是相对于我那时的品性而言，即一个人能够重生——这是天主的慈悲为我的救恩所宣告的真理——并且在拯救的水泉中被赋予新生命的人能够脱下他先前所是；并且，虽然保留了他全部的身体结构，却在内心和灵魂上被改变。我说：\Quote{这样的转变如何可能，即所有那些要么与生俱来、在我们物质本性的腐败中硬化，要么后天获得、因长期习惯而根深蒂固的东西，能被突然而迅速地脱去？这些东西已深深且彻底地植根于我们。习惯于盛宴和豪宴的人，何时学会节俭？那闪耀于金银和紫袍中、以昂贵服饰闻名的人，何时将自己降至普通和简单的衣着？曾感受过权杖（\OriginalTerm{法西斯}{fasces}）和公民荣誉的魅力的人，厌恶成为单纯的平民和不光彩的公民。那被成群的食客簇拥、因众多好管闲事的随从而显赫的人，当独处时视之为惩罚。酒之爱引诱，骄傲膨胀，愤怒燃烧，贪婪扰心，残忍刺激，野心喜悦，欲望加速毁灭，这些诱惑从不放松束缚，正如过去必然如此。}\par

4. 这些是我常有的思绪。因为我自己也曾被先前生活中无数的错误束缚，以为绝无可能从中解脱，于是甘心沉溺于纠缠的恶习；由于对更好的事绝望，我便纵容自己的罪，仿佛它们是我的一部分，与我同根。但此后，藉新生之水的帮助，往日污秽被洗去，从上而来的静谧纯净之光注入我重修和好的心中——此后，藉从天上吹来的圣神，第二次出生使我恢复为新人——于是，奇妙地，疑惑之事立刻开始向我证实，隐藏之事被揭示，黑暗之事被光照，先前看似困难的事开始显明可行之道，曾被以为不可能的事竟能成就；我得以承认，先前生于肉身、活在罪恶中的，是属土的、属地的，但现在已开始属于天主，并由圣神所活化。你自己必定知道，也与我一样记得，那罪恶之死与德行之生从我们身上夺走了什么，又给了我们什么。你无需我告知便知此事。任何自夸的言谈都是可憎的，虽然我们\Emph{实际上}不能自夸，只能感恩，因我们凡不归于人德而宣认为天主恩赐的一切，都当感恩；以至于如今我们不犯罪，是信德之工的开始，而我们从前犯罪，则是人类谬误的结果。我们一切能力来自天主；我说，来自天主。我们从他获得生命，从他获得力量，藉从他汲取并领受的能力，我们虽仍在世上，却能预知未来之事的征兆。唯愿恐惧成为无罪的守护者，使那因仁慈而藉天上恩宠的临在流入我们心中的主，因正义的顺从得以存留于感恩之心的客舍，免得我们所得的确信滋生疏忽，致使旧敌再次潜近我们。\par

5. 但若你持守无罪之路、正义之路，若你以坚定稳健的步伐行走，若你全心全力依赖天主，只\Emph{是}你开始成为的样貌，那么，随着你超性恩宠的增加，行事的自由与能力必赐予你。因为属天恩赐的分施，不像属世好处那样有度量或限制。那自由流溢的圣神不受任何界限约束，不被某些有限空间内的封闭栅栏阻碍；它永流不息，丰盈洋溢。唯愿我们的心干渴，并预备领受：我们以何等宽大的信德接近它，便从它汲取何等丰沛的恩宠。由此赐予能力，以贞洁的节制、健全的心智、纯朴的话语、无瑕的德行，能消灭毒液以医治病人，以恢复的健康清除愚昧灵魂的污秽，向敌对者宣告和平，向暴烈者赐予安宁，向悖逆者赐予温良——以惊人的恐吓迫使那些潜入人体意图毁灭人的不洁游魂承认自己，以重击驱赶它们出来，使它们挣扎、号叫、呻吟，因不断加增的痛苦而伸展，以鞭笞击打它们，以烈火炙烤它们：这事在暗中进行，却不可见；所受的击打是隐藏的，但惩罚是明显的。因此，就我们已开始成为的而言，我们所领受的圣神拥有其自由行动之权；而就我们尚未改变身体与肢体而言，属肉体的目光仍被此世的云雾所遮蔽。这心灵的统辖何其伟大，其能力何其强大，不仅它能脱离世界的邪恶结盟，如同被洁净之人不受敌对入侵的玷污，而且它的力量变得更大更强，以致能以它的统辖制服那攻击者的全部傲慢大军！\par

6. 但为使天主的事因真理的开展而更显光辉，我将赐你光明以领悟它，罪的昏暗被擦去。我将从此隐藏世界的黑暗上揭去幔帐。请你短暂地设想自己被提到某座不可攀登的高山之巅，从那里俯瞰下方事物的景象，向四方转目，观看这波涛世界的漩涡，而你自己远离尘世——你必立刻开始同情这世界，并因自省与对天主的日益感恩，以更大的喜乐欢欣，因你已逃脱了它。请看道路被强盗阻断，海洋被海盗围困，战争与营中血腥的恐怖散布全地。整个世界浸透了彼此的鲜血；而杀人，若是个别行为，被视为罪行，但若大规模施行，则被称为美德。对邪恶行为要求免罚，并非因其无罪，而是因残暴大规模地施行。\par

7. 如今，你若将目光和注意转向这些城邑本身，你将看到一群比任何孤寂更令人悲伤的聚集。角斗表演正在准备，好让鲜血取悦残酷眼目的贪欲。身体以更强烈的食物滋养，四肢的有力体魄以筋肉和肌肉充实，好使这为刑罚而养肥的可怜虫死得更艰难。人被屠杀，好让人得到满足；最能杀人的技能成为一种练习和艺术。罪行不仅被犯下，而且被教导。还有什么比这更不人道——还有什么比这更令人厌恶？训练是为了获得杀人的能力，而杀人的成就就是它的荣耀。我请问，那是怎样的一种状况，又可能像什么——在其中，无人定罪的人们却将自己献给野兽——他们年纪成熟，容貌足够美丽，穿着昂贵的衣服？活着的人，为自愿的死亡而装饰；可怜的人，夸耀自己的苦难。他们与野兽搏斗，不是为了他们的罪行，而是为了他们的疯狂。父亲看着自己的儿子；兄弟在竞技场上，姐妹就在近旁；尽管更华丽的排场增加了表演的代价，但可耻啊！连母亲也会支付那增加的费用，好让她在场观看自己的苦难。目睹如此可怕、如此不敬、如此致命的情景，他们似乎没有意识到自己是用眼睛弑亲的人。\par

8. 由此，转目看向另一种表演的可憎之事，同样不可不悲叹。在剧场中，你也会看到足以令你悲伤和羞愧之事。那是悲剧厚底靴，以诗句讲述古代日子的罪行。古老的弑亲和乱伦的恐怖，以旨在表达真理形象的动作展现出来，以便随着岁月流逝，任何从前犯下的罪行都不会被遗忘。每一代人都由所听到的事提醒：凡曾做过的事，都可能再做。罪行从不因岁月流逝而灭绝；邪恶从不因时间过程而消除；不敬从未被遗忘埋葬。那些已不再是实际恶行之事，成为榜样。此外，在滑稽戏中，通过丑行的教导，观众被吸引要么重新思考他可能在暗中所行之事，要么听到他可能行的事。奸淫在被观看时就被学习；当那有公共权威的恶行纵容罪恶时，那主妇，也许原本去观看表演时还是一位端庄的妇人，回来时却变得放荡。更进一步，在违背自己出生的盟约和律法的情况下，被戏子的姿态玷污，细致观看乱伦的可憎之事，那是何等道德败坏，何等可憎之事的刺激，何种罪恶的食粮！男人被阉割，他们性别的所有骄傲和活力在其衰弱身体的耻辱中变得柔弱；那里最令人愉悦的人，就是最彻底把男人打破成女人的人。他因自己的罪行而成长于赞美中；他越是被贬低，就越被认为是技艺高超。这样的人被观看——可耻啊！而且被愉快地观看。这样的生物还能暗示什么？他煽动感官，他讨好情欲，他驱走善良胸怀那更有力的良心；而且，这诱人的可憎之事不乏权威，好让恶行以更不易察觉的方式潜入人民之中。他们描绘不洁的维纳斯、奸淫的马尔斯；而他们的那位朱庇特，不仅在统治上至高无上，在罪恶上也是如此，在雷霆之中为尘世之爱所燃，时而变白为天鹅的羽毛，时而如金雨倾泻，时而借助飞鸟来侵犯男孩的纯洁。现在请提问：观看这些事物的人，能有健康的心智或端庄吗？人们仿效他们崇拜的神明，对这些可怜人而言，他们的罪行就成了他们的宗教。\par

9. 唉！若能置身于那高耸的瞭望塔上，以目光窥探隐秘之所——若能打开卧室紧闭的门，将其中幽暗的角落召回知觉的视线——你将会看见不洁之人所行之事，没有贞洁的眼睛能够观看；你将看见那些连观看本身都是罪行的事；你将看见那些被罪恶的疯狂所兽化的人否认他们做过、却又急忙去做的事——男人以疯狂的情欲扑向男人，行那些连行此事的人本身也不得满足的事。若有人犯下这等事，却仍指控他人，那我便是被欺骗了。堕落者诋毁堕落者，认为自己虽知罪，却已逃脱，仿佛良心不足以为定罪。那些公开作控告者的人，私下却是罪犯，在定罪罪犯的同时也定自己的罪；他们在外谴责自己在家中所行之事，甘愿行他们所指控之事——这种大胆确与罪恶相配，这种厚颜与无耻之人完全相称。我恳请你不要奇怪这类人所说的话：他们口舌的冒犯，在言语上还是他们最轻的罪过。\par

10. 但在考虑了布满陷阱的公共道路、遍布全世界的各种战斗、血腥或可耻的表演、以及情欲的憎恶——无论是在妓院中公开出售还是隐藏在家庭墙壁之内——这些憎恶，其肆无忌惮程度与犯罪的隐秘性成正比——之后，也许你会认为至少广场上不会有这些事情，它既不会受到令人愤怒的冤屈，也不会因罪犯的关联而被玷污。那么，请将你的目光转向那里：在那里你会发现比以往更可憎的事情，以至于你会更渴望移开你的眼睛，尽管法律被刻在十二铜表上，法令被公开规定在铜板上。然而，就在法律之中，冤屈仍在发生；就在法令面前，邪恶仍在进行；即使在正义受到捍卫的地方，无辜也得不到保护。争辩者的怨恨交替爆发；当和平在托加袍中被打破时，广场上回荡着纷争的疯狂。近在咫尺的是长矛和刀剑，还有刽子手；还有撕裂的利爪、拉伸的刑架、燃烧的火焰——一个可怜的人体所遭受的折磨比它的肢体还要多。在这种情况下，有谁能帮忙呢？他的保护人？他假装帮忙，却欺骗了他。法官？但他出卖了他的判决。坐在那里审判罪行的人自己就在犯罪，法官变成了罪犯，以便被告无辜地死去。罪行到处都是；在罪的多重形式中，有害的毒药通过堕落的心灵起作用。一个人伪造遗嘱，另一个人通过致命的欺诈作出虚假证词；一方面，孩子们被欺骗失去了继承权，另一方面，陌生人被赋予了他们的财产。对手提出指控，诬告者攻击，证人诽谤，四面都是雇佣来的声音的谄媚无耻，开始伪造指控，与此同时，有罪者甚至不和无辜者一起灭亡。没有人惧怕法律；没有人在乎调查官或法官；当判决可以用金钱买通时，人们就不在乎它了。如今，在有罪者中，无辜反而成了罪行；凡是不效法恶人的人，就对他们构成了冒犯。法律与罪行妥协了，凡是公开的事情都开始被允许。那里还能有什么谦逊，什么正直呢？既然没有一个人谴责恶人，而你所遇到的都是那些自己本应被谴责的人。\par

11. 但为了不让我们显得是在挑选极端的例子，并且为了贬低而试图吸引你的注意力到那些更良心的目光可能被悲伤和反感所冒犯的事物上，我现在将引导你去看世界上那些无知者认为是美好的事物。在这些事物中，你也会看到令你震惊的事情。至于你所认为的荣誉、你所认为的权杖、你所认为的财富富足、你所认为的军营中的权力、长官职务中的紫色荣耀、指挥官职位中的放纵权力——其中隐藏着诱惑性祸害的毒液，以及一种微笑的邪恶之貌，确实是欢快的，但却是隐藏灾难的奸诈欺骗。正如某些毒药，其味道被甜味所药化，在致命的毒汁中狡猾地混合，饮用时看起来像普通的饮料，但当它被喝下去时，你所吞下的毁灭就会袭击你。你看，确实，那个穿着华丽衣服的人，在他自己看来，在他紫色的衣服中闪闪发光。然而，他付出了多么卑劣的代价才获得了这闪耀！他首先必须忍受骄傲者的多少蔑视！他作为一个早期的朝臣，围攻了多少傲慢的门槛！他必须在傲慢的大人物们多少轻蔑的脚步之前，拥挤在门客的人群中，以便不久之后，类似的行列可以带着问候跟随他并走在他前面——一个队伍等候的不是他本人，而是他的权力！因为他不应该因其品格而受尊重，而应因其权杖而受尊重。最后，你可以看到这些的堕落结局，当趋炎附势的谄媚者离开时，食客们抛弃了他，玷污了退隐到私人状态的人暴露的一侧。正是那时，对挥霍家庭财产所造成的不法行径打击着良心，那时，耗尽财富的损失才被知晓——这些开销是用来购买民众的青睐，并以善变的虚空恳求来求得人民的赞许。诚然，这是一种虚荣和愚蠢的夸耀，想在令人失望的奇观中展示那些人民不会接受、并且会毁掉官员的东西。\par

12. 此外，你认为富足的那些人，他们又将森林加到森林上，并将穷人从邻近之处排除出去，把他们的田地无边无际地远远延伸开去，据有巨堆的金银和大量的钱财，要么堆积成堆，要么埋藏起来——即使在他们财富之中，这些人也被模糊思绪的焦虑撕碎，唯恐强盗劫掠，唯恐凶手攻击，唯恐某个更富的邻人的嫉妒变为敌意，并以恶意的诉讼来烦扰他们。这样的人无论在进食还是在睡眠中都得不到安全。在宴席中他叹息，尽管他用宝石镶嵌的杯饮酒；当奢华的床榻用其柔软的怀抱包裹他那因宴饮而倦怠的身体时，他躺在羽绒中却无法安眠；这个可怜的人并未察觉，这些不过是镀金的折磨，他被自己的黄金束缚，他是其奢华和财富的奴隶而非主人。啊，可憎的感知盲目，和愚蠢贪婪的深沉黑暗！尽管他可以卸下重担并摆脱负担，他却继续盘踞在他那令人烦恼的财富之上——他顽固地紧握着他那折磨人的积存。他对依附者没有慷慨，对穷人没有施舍。然而这样的人却称那些钱财为自己的，他们忌妒地劳苦守护在家中，仿佛那是别人的钱财，他们从中既不为朋友、也不为儿女、甚至也不为自己得到任何益处。他们的占有仅仅在于：他们能阻止别人占有它；啊，何等惊人的名称颠倒！他们称那些东西为\Emph{货物}，而他们却绝对只用于\Emph{坏}的用途。\par

13. 或者你以为，那些在荣誉花冠和巨大财富中，在皇家宫殿的辉煌中，被警觉的武装卫士环绕的人——他们也安全吗？——他们至少有一些稳定的持久吗？他们比其他人有更大的恐惧。一个人必然既害怕别人，也被人害怕。显贵同样向更有权势者索取代价，尽管他可能被成群随从包围，并用众多随从的围护和保护来保卫自己。正如他不允许下属感到安全，他自己也必然缺乏安全感。那些使别人恐惧之人的权力，首先对自己是可怕的。它微笑以发怒，它谄媚以欺骗，它引诱以杀害，它高举以摔落。带着某种祸害的利息，所达到的尊严和荣誉越高，所需的惩罚利息就越大。\par

14. 因此，唯一宁静而可信的平安，唯一坚固、稳定且恒常的安全，就是一个人从这纷扰世界的漩涡中退出，锚定在救恩港口的基石上，将目光从地上举向天上；并且，既已被接纳于天主的恩赐，心灵已非常接近他的天主，他可以夸耀，人间事务中别人视为崇高伟大的，全然在他的意识之下。那真正比世界更大的人，不能从世界渴求任何东西，也不能从世界欲求任何东西。那保障何等稳固、何等不受所有冲击；那保护在天赐福佑中何等属天——即从这纠缠世界的网罗中解脱，从世俗渣滓中被净化，并为永生不朽之光而被装备！他将看见先前攻击我们的仇敌那狡猾的祸害如何在我们身上进展。我们被迫更爱我们将成为的，因为我们得以认识并谴责我们曾经是的。为此目的，无需花费贿赂或劳力的代价，以使人的升迁、尊严或权力通过精心努力在他内产生；但这是天主白白赐予的恩赐，且人人可得。如同太阳自发照耀，如同白天赐予光明，如同泉水涌流，如同雨水带来湿润，同样，天上的圣神将自己注入我们内。当灵魂举目望天而认出其造主时，它便升得比太阳更高，远远超越这地上一切权势，并开始成为它相信自己是的。\par

15. 然而，你——天国的战事已征召你进入属灵营垒——只须遵守在宗教德行中纯洁且受训练的纪律。在祈祷和阅读中都要恒心；时而与天主交谈，时而让天主与你交谈，让祂以祂的诫命教导你，让祂引导你。祂所使富足的人，无人能使之贫穷；因为，事实上，凡胸怀一次被天上食粮所供应的人，不可能有贫穷。用黄金装饰的天花板，以及用昂贵大理石马赛克装饰的房屋，在你看来将是鄙陋的，既然你知道你自己更应该被成全，你自己更应该被装饰，并且那作为天主居住其内的殿宇、圣神已开始在其中安家的居所，比所有其他居所更重要。让我们用纯洁的颜色美化这房屋，用正义之光光照它：这房屋永不会因年久朽坏而衰败，也不会被墙壁或黄金颜色的暗淡所玷污。凡人为美化的东西都会朽坏；这样的东西不包含真实拥有的实质，不能给其拥有者提供持久的保证。但这房屋保持在永恒生动的美丽中，在完全的尊荣中，在持久的光辉中。它既不能衰败，也不能被毁灭；当身体回归它时，它只能被塑造成更大的完美。\par

16. 最亲爱的\Person{多纳图}，这些事暂且简略提及。因为你所受益听闻的，既悦纳你的忍耐（以你的良善宽容）、你平衡的心智、你坚定的信德——没有什么比天主内令你喜悦的事更悦耳了——然而，我们既是邻居，时常交谈，谈话也当有节制；既然这是一段休憩假期，闲暇时光，趁夕阳西斜，让我们以喜乐度过余下的白昼，连宴饮时刻也不失天上的恩宠。愿节制的宴席回响着圣咏；你记忆牢固、嗓音优美，请照常担任此职。你将为你最亲爱的朋友提供更佳的款待，若我们在聆听属灵之事时，宗教音乐的甜美令我们耳悦。\par

\SourceRecord{Translated by Robert Ernest Wallis.；Ante-Nicene Fathers；Vol. 5.；Edited by Alexander Roberts, James Donaldson, and A. Cleveland Coxe.；Buffalo, NY: Christian Literature Publishing Co.,；1886.；<http://www.newadvent.org/fathers/050601.htm>.}

% structure: p0001 source=h1 target=section reason=page-title

\section{第二封书信}

\Emph{从\Place{罗马}神职到\Place{迦太基}神职，关于真福\Person{西彼廉}的隐退。}\par

\Emph{论点。——\Place{罗马}神职从副执事\Person{克雷门提乌斯}得知，在迫害时期\Person{西彼廉}已隐退自己。因此，以他们惯常对信仰的热忱，他们提醒\Place{迦太基}神职他们的职责，并指示他们在主教缺席的间隔中，在背教者的情况下该做什么。}\par

1. 我们已由你们的副执事\Person{克雷门提乌斯}得知，真福神父\Person{西彼廉}已因某一理由隐退；\Quote{他这样做完全正确，因为他是一位卓越的人，且因为一场冲突即将来临}，这是天主为与祂的仆人们一同对抗仇敌而在世上允许的，并且祂还愿意这场冲突向天使和人类显示，胜利者将获得冠冕，而失败者将亲自领受那已向我们显明的惩罚。此外，既然我们这些似乎居于高位、在牧人位置上的人有责任看守羊群；若我们被发现疏忽，将如同我们的前任也被如此说，他们同样因疏忽而被置于职责中，说：\BibleQuote{我们没有寻求那迷失的，没有纠正那游荡的，没有包扎那折断的，却吃了它们的奶，穿了它们的羊毛}；\BibleRef{则 34:3}然后主自己，成全律法和先知所写的，教导说：\BibleQuote{我是善牧，我为羊舍命。但雇工，羊不是他自己的，看见狼来，就撇下羊逃跑，狼就分散了它们}。对\Person{西满}，祂也如此说：\BibleQuote{你爱我吗？他回答：我爱祢。祂对他说：你喂养我的羊}。\BibleRef{若 21:17}我们知道这句话正是出于他隐退的情况，其余的门徒也如此做了。\par

2. 因此，亲爱的弟兄们，我们不愿你们被发现是雇工，却愿意你们是善牧，因为你们知道，若你们不劝诫我们的弟兄在信德上站稳，就必有极大的危险威胁你们，以免弟兄团体被彻底根除，如同那些冲入偶像崇拜的人一般。我们不仅仅在言语上劝勉你们这样做；你们将从许多来到你们这里的我们的人那里确知，天主降福我们，我们自己也怀着所有焦虑和世俗的风险，做了并仍在做这一切事，我们眼中更看重对天主的敬畏和永恒的苦楚，而非对人的恐惧和短暂的难受，没有离弃弟兄们，而是劝勉他们坚守信德，并准备好与主同行。我们甚至召回了那些正上去做他们所被迫做的事的人。教会屹立在信德中，尽管有些人因极大的恐惧而跌倒，无论是那些卓越的人，或是那些被捕后因怕人而害怕的人：然而，我们没有遗弃他们，虽然他们与我们分离了，但我们曾劝勉他们，并且仍劝勉他们悔改，或许他们能从那位能赐予宽恕者那里获得赦免；免得若他们被我们遗弃，就变得更糟。\par

3. 因此，弟兄们，你们看，你们也应该这样做，好使那些跌倒的人能因你们的劝诫而改正心意；如果他们一旦再次被捕，可以宣认，并如此弥补他们先前的罪过。还有其它你们应尽的本分，我们在这里也附加了，即若有人陷入这诱惑而开始生病，并为自己所做的事悔改，且渴望共融，应无论如何赐予他们。或者，若你们有寡妇或卧床不起的人，无法维持生计，或那些在监狱中或被驱逐出自己住所的人，这些人都应当有人照管他们。此外，望教者患病时，不应被欺骗，而应给予他们帮助。并且，最重要的是，若殉道者和其他人的尸体未被埋葬，那些有责任做此事的人将面临相当大的危险。因此，无论你们中谁，也不论在什么场合履行了这项职责，我们确信他被视为好仆人——如同那在小事上忠信的人，并被指派管理十座城市。然而，愿那将一切赐予仰望祂者的天主，赐予我们，使我们能在这些善工中被发现。被囚禁的弟兄们问候你们，长老们和整个教会也问候你们，教会自己也以最深切的焦虑看守所有呼求主名的人。我们也同样请求你们记念我们。此外，要知道\Person{巴西雅努斯}已来到我们这里；我们请求你们这些对天主有热忱的人，将这封信的副本寄给你们所能寄到的任何人，看时机便利，或自己创造机会，或传递消息，好使他们能坚定不移地站稳信德。祝你们，亲爱的弟兄们，永远安康。\par

\SourceRecord{Translated by Robert Ernest Wallis.；Ante-Nicene Fathers；Vol. 5.；Edited by Alexander Roberts, James Donaldson, and A. Cleveland Coxe.；Buffalo, NY: Christian Literature Publishing Co.,；1886.；<http://www.newadvent.org/fathers/050602.htm>.}

% structure: p0001 source=h1 target=section reason=page-title

\section{第三封书信}

\Emph{致住在罗马的司铎与执事。公元250年。}\par

\Emph{说明——这是一封亲切友好的书信；无需正式说明，尤其因为从标题本身即可充分领会。罗马神职人员的回信已遗失。}\par

1. \Person{西彼廉}致留在罗马的司铎与执事，弟兄们，安好。我亲爱的弟兄们，当我仍在犹疑不定时，关于我那杰出的同僚离世的传闻仍不确定，我便收到了你们由副执事\Person{克雷孟提乌斯}带给我的信函，其中极为详尽地告知了我他光荣的结局；我甚为欣喜，因为与他治理的无瑕相称，他也获得了光荣的结局。此外，我尤为祝贺你们，因你们以如此公开而辉煌的见证纪念他，以致经由你们，我不仅得知了你们纪念你们的主教的光荣，也得知了应成为我信德与德行之榜样的事。因为正如一位主教的跌倒倾向于导致其追随者毁灭性的跌倒，反之亦然，当一位主教以其信德的坚定将自己作为效法的榜样展示给弟兄们时，这乃是一件有益而有助的事。\par

2. 此外，我还阅读了另一封信函，其中既未明确写明写信人，亦未写明收信人；由于在同一封信中，笔迹、内容乃至纸张本身都让我觉得其中有些内容被删改或与原文不同，我便将原样收到的信函寄还给你们，以便你们核查这是否就是你们交给副执事\Person{克雷孟提乌斯}的那封。因为若神职人员的信函的真实性被任何虚假或欺骗所败坏，那是一件极为严重的事。因此，为使我们知晓此事，请确认笔迹和签名是否出自你们，并就事实真相再回信给我。至亲爱的弟兄们，我衷心祝愿你们平安。\par

\SourceRecord{Translated by Robert Ernest Wallis.；Ante-Nicene Fathers；Vol. 5.；Edited by Alexander Roberts, James Donaldson, and A. Cleveland Coxe.；Buffalo, NY: Christian Literature Publishing Co.,；1886.；<http://www.newadvent.org/fathers/050603.htm>.}

% structure: p0001 source=h1 target=section reason=page-title

\section{第四封书信}

\Emph{致司铎与执事。}\par

\Emph{论述。——\Person{西彼廉}从隐退之处劝勉他的神职人员，在他的缺席期间应团结一致；不可让囚犯和其他穷苦人有所缺乏；还应保持人民安静，以免若他们成群涌向监狱探望殉道者，最终导致这一特权被禁止。公元250年。}\par

1. \Person{西彼廉}致以他亲爱的弟兄司铎与执事，问候。承蒙天主的恩宠，我最亲爱的弟兄们，我平安无恙，向你们请安。由于此地的情况不容我现在与你们同在，我恳请你们，凭着你们的信德和虔诚，在那里履行你们和我的职分，使纪律和勤勉都一无所缺。此外，关于筹款以应付开支，无论是为了那些以光荣的声音宣认了他们的主而被投入监狱的人，还是为了那些在贫困和匮乏中辛劳、却仍然坚立在主内的人，我恳求不要有任何缺乏，因为在当地收集的全部小额款项已分发给神职人员用于此类情况，以便许多人能有手段来帮助个人的需要和负担。\par

2. 我也恳请你们，不要缺乏智慧和谨慎以维护和平。因为虽然弟兄们出于爱情渴望接近和探望那些良好的宣信者——神圣的眷顾已藉他们光荣的开端将光明照耀在他们身上——然而，我认为这种渴望必须谨慎地满足，不要成群结队——不要一下子聚集众多人数，免得因此激起恶意，进入的途径被禁止，这样，当我们贪得无厌地想要得到一切时，反而失去一切。因此，请商议并设法以节制更安全地处理此事，以便那些在那里与宣信者一同献祭的司铎，也能一次一位地与执事轮流替换；因为这样更换人员、使聚集的人群有所不同，猜疑就会减少。因为我们温柔谦卑地对待一切，正如天主仆人所当行的，我们应当适应时世，提供安宁，并关怀人民。弟兄们，我亲爱的、渴望已久的，祝你们永远诚心如意；并请纪念我。向全体弟兄问安。执事\Person{维克多}以及与我同在的人向你们致安。祝好！\par

\SourceRecord{Translated by Robert Ernest Wallis.；Ante-Nicene Fathers；Vol. 5.；Edited by Alexander Roberts, James Donaldson, and A. Cleveland Coxe.；Buffalo, NY: Christian Literature Publishing Co.,；1886.；<http://www.newadvent.org/fathers/050604.htm>.}

% structure: p0001 source=h1 target=section reason=page-title

\section{第五封书信}

\Emph{致长老与执事。}\par

\Emph{内容提要。——本书信的内容与前一封大致相同，不同之处在于作者指示宣信者也应由圣职人员告诫其本分，注意谦卑，并服从长老与执事。他自己的退隐恰为提供了这个机会。}\par

1. 西彼廉致各位长老与执事，他的弟兄们，问安。亲爱的弟兄们，我原本希望借这封信向我的全体圣职人员问安，祝他们健康平安。但由于这风暴时期已在很大程度上压垮了我的百姓，更令我悲伤的是，它还波及了一部分圣职人员，我祈求主，藉着天主的仁慈，我今后至少能见到你们平安，因我得知你们在信德和德行上都坚定不移。虽然有些理由似乎促使我赶去与你们相会，首先，例如，出于我对你们的渴望和爱慕，这是我祈祷中的首要考虑；其次，我们能够共同商议那些为了教会的治理所需要的普遍利益之事，经过充分的商议仔细审查后，智慧地安排它们——但我仍认为，暂时保持我的退隐和宁静更为妥当，以利于我们所有人的平安和安全所需的其它益处。这些益处将由我们亲爱的弟兄特尔土禄向你们报告，他除了热诚献身于神圣劳苦之外，还是这个建议的发起者：要我谨慎、节制，不可鲁莽地出现在公众面前；尤其要提防那个我多次被查问和搜索的地方。\par

2. 因此，我信赖你们的爱德和虔诚，对此我已深知，在这封信中我既劝勉又命令你们，你们中那些出现最少嫌疑、最少危险的人，应代我履行我的职责，处理那些宗教管理所需要的事。与此同时，要尽可能多地照顾好穷人；尤其是那些以坚定不移的信德站立、没有离弃基督羊群的人，藉着你们的勤勉，供应他们所需，使他们能承受贫困，这样，艰难时期未能动摇他们信德的事，不至于因匮乏而击垮他们的痛苦。此外，要更加热切地照顾那些光荣的宣信者。虽然我知道他们中已有许多人靠弟兄们的誓愿和爱心得到维持，但若有人缺乏衣服或生计，应供应他们一切必需品，正如我先前在你们仍被囚禁在狱中时写信给你们的——只要让他们从你们那里知道并接受教导，学习圣经的权威所要求的教会纪律对他们有何要求：他们应当谦卑、温和、和平，持守他们名号的荣誉，以便那些因见证而获得荣耀的人，也藉着他们的品行获得荣耀，并在一切事上寻求主的悦纳，在完成他们的赞美时，证明自己配得天上的冠冕。因为还有比已经看到的更多的事要完成，因为经上记载说：\Quote{人未死以前，不要赞美他；} 又说：\BibleQuote{你应当至死忠心，我就赐给你那生命的冠冕。} \BibleRef{默 2:10} 主也说：\BibleQuote{惟有忍耐到底的，必然得救。} \BibleRef{玛 10:22} 让他们效法主，祂在受难之时不是更骄傲，而是更谦卑。因为那时祂洗门徒的脚，说：\Quote{我是你们的主和老师，尚且洗你们的脚，你们也当彼此洗脚。我给你们作了榜样，叫你们照着我向你们所做的去做。} 也让他们效法宗徒保禄的榜样，他在多次被囚、受鞭打、被野兽追赶之后，在一切事上仍然温柔谦卑；甚至在升到第三层天和乐园之后，他也没有骄傲地自夸，而是说：\BibleQuote{我们也未尝白吃人的饭，倒是辛苦劳碌，昼夜做工，免得叫你们一人受累。} \BibleRef{得后 3:8}\par

3. 我请求你们，将这些事提醒我们的弟兄们。正如经上说：\BibleQuote{凡自高的，必降为卑。} \BibleRef{路 14:11} 现在正是他们更当惧怕那设陷阱的仇敌的时候，他更急切地攻击最强壮的人，并且正是因为被击败而变得更恶毒，竭力战胜他的征服者。愿主使我尽快再次见到他们，并通过有益的劝勉坚固他们的心，以保持他们的荣耀。因为我听到他们中有些人邪恶而骄傲地到处乱跑，沉溺于愚昧或纷争，感到悲伤；基督的肢体，甚至是已经宣认基督的肢体，竟被非法的姘居玷污，不能被执事或长老管理，反而因少数几个人的邪恶品格，使许多善良的宣信者的光荣蒙受玷污；这些人应当害怕，免得被他们的见证和审判定罪，被排除在他们的团契之外。最后，那才是卓越而真实的宣信者，教会后来不因他羞愧，反而夸耀。\par

4. 关于我们的同司铎\Person{多纳图}、\Person{福图纳图斯}、\Person{诺瓦图斯}和\Person{戈尔狄乌斯}写信给我之事，我无法独自答复，因为自从我开始担任主教职以来，我就下定决心，没有你们的意见和民众的同意，绝不凭自己的私人见解行事。但是，一旦藉着天主的恩宠，我能与你们相聚，那时我们将按照各自应有的尊严，共同讨论那些已经完成或将要进行的事。亲爱的、深切渴望的弟兄们，我衷心祝愿你们永远平安，并请记念我。请替我热忱地问候你们那里的弟兄团体，并告诉他们记念我。祝平安。\par

\SourceRecord{Translated by Robert Ernest Wallis.；Ante-Nicene Fathers；Vol. 5.；Edited by Alexander Roberts, James Donaldson, and A. Cleveland Coxe.；Buffalo, NY: Christian Literature Publishing Co.,；1886.；<http://www.newadvent.org/fathers/050605.htm>.}

% structure: p0001 source=h1 target=section reason=page-title

\section{第六封书信}

\Emph{给司铎\Person{罗加提亚努斯}及其他宣信者。公元二五〇年。}\par

\Emph{论证。——他劝勉\Person{罗加提亚努斯}及其他宣信者持守纪律，免得那些口里宣认基督的人，行为上却似乎否认了他；并顺便责备他们中的一些人，他们因信德被流放，却不怕未经允许就返回本国。}\par

\Person{西彼廉}致司铎\Person{罗加提亚努斯}及其他宣信者，他的弟兄们，问候。我之前已经给你们写了一封信，以洋溢的言辞祝贺你们的信德和勇毅；现在，我的声音别无他事，首先就是以喜乐之心，一再地宣布你们之名的荣耀。因为在我的祈祷中，我还能祈求什么比看到基督的羊群因你们宣信的尊荣而受到光照更大更好呢？因为虽然所有弟兄都应当为此喜乐，但在共同的喜乐中，主教的分内是最大的。因为教会的荣耀就是主教的荣耀。正如我们为那些被敌对迫害击倒的人悲伤，同样我们也为你们欢喜，因为魔鬼未能胜过你们。\par

然而，我藉着我们共同的信德，并藉着我对你们衷心真实而纯朴的爱，劝勉你们，既然在这初次交锋中胜过了仇敌，就要以勇敢而持久的勇毅持守你们的光荣。我们仍在世界上；我们仍处身战场；我们天天为生命而战。必须留心，使你们在如此开端之后，也能有所长进，并使你们以如此蒙福的开端所开始的事，能在你们内得以完成。能够获得什么，这是小事；能够保住所得，这才是大事；正如信德本身和救恩的诞生，不是由于接受，而是由于保存才使人活。实际上，使人保持在天主面前的，不是获得，而是成全。主在他的教导中教训了这事，他说：\Quote{看，你已痊愈了；不要再犯罪，免得你遭遇更糟的事。}你们要想象他也这样对宣信者说：\Quote{看，你已成为宣信者；不要再犯罪，免得你遭遇更糟的事。}同样，\Person{撒罗满}和撒乌耳，以及许多其他的人，当他们行走在主的道路上时，能够保住赐给他们的恩宠。当他们离弃了主的纪律时，恩宠也离弃了他们。\par

我们必须持守在赞美和荣耀的直而窄的路上；既然平安、谦逊和良好生活的宁静适合于所有基督徒，按照主的话，他只垂顾这样的人：\BibleQuote{贫苦和心灵痛悔的人，以及因我的话战兢的人}\BibleRef{依 66:2}，那么，你们这些已成为其余弟兄榜样的宣信者，就更应当遵守并实行此事，因为你们的品德应当激发所有人效法你们的生活和行为。就如犹太人被天主疏远了，他们使\BibleQuote{天主的名在外邦人中受亵渎}\BibleRef{罗 2:24}，同样另一面，那些因遵守纪律而得以用赞美为天主的名作证的人，则为天主所喜爱，正如经上所载，主自己预先警告说：\Quote{你们的光当这样照在人前，叫他们看见你们的善行，光荣你们在天上的父。}宗徒\Person{保禄}说：\BibleQuote{你们在世界中应该发光}\BibleRef{斐 2:15}。同样，\Person{伯多禄}也劝告说：\BibleQuote{亲爱的，你们是外方人和旅客，要戒绝与灵魂作战的肉欲；在外邦人中要常持良好的品行，好使那些毁谤你们为作恶的人，因见到你们的善行，而光荣主。}\BibleRef{伯前 2:11}事实上，你们中的大部分，我高兴地说，都留意此事，并且因宣信本身的荣誉而改善，以宁静和德行高尚的生活守卫并保存其光荣。\par

但我听说有些人玷污了你们的团体，并以腐败的言谈毁灭了卓越之名的赞美；你们自己，作为本身赞美的爱护者和守护者，应当斥责、制止和纠正他们。因为当有人终日酗酒荒淫，另有人返回他被放逐的祖国，一旦被捕就灭亡，如今不是作为一个基督徒，而是作为一个罪犯时，你们的名声遭受了何等耻辱！我听说有些人自高自大，虽然经上写着：\Quote{不要心高气傲，要害怕；因为天主既然没有怜悯自然的枝条，你们也要小心，免得他连你们也不怜悯。}我们的主\BibleQuote{被牵去如同待宰的羊，又像羔羊在剪毛者面前无声，他也没有开口。}\BibleRef{依 53:7}\Quote{我没有叛逆，也没有反驳。我将我的背转给击打者，我的面颊转给掌掴者；我没有掩面躲避唾污。}如今，凡生活在此一位内并因他而活的人，竟敢自高自大，忘记他所行的事，也忘记他亲自或藉他的宗徒留给我们的命令吗？但如果\Quote{仆人不能大于他的主人}，那么，那些谦卑、平安、沉默地跟随主的人，就当踏着主的足迹而行，因为人越卑微，就越被高举；正如主说：\BibleQuote{你们中间最小的，将成为大的。}\BibleRef{路 9:48}\par

5. 那么，我所怀着极度的痛苦和悲伤得知的那件事——它在你眼中应显得多么可憎——就是：不乏有人以可耻而臭名昭著的姘居玷污天主的殿宇，以及在告解后成圣并得荣耀的肢体，将他们的床榻与女人混杂！在这些事中，即使良心没有受到污染，单凭他们的过错为他人提供毁灭的榜样，这本身已是极大的罪债。你们中间也不应有争论和嫉妒，因为主将祂的平安留给了我们，并且经上记载：\BibleQuote{你应爱人如己。} \BibleRef{肋 19:18} \BibleQuote{但你们若彼此咬伤、指责，要当心不要被彼此消灭。} \BibleRef{玛 22:39} 我也恳求你们远离辱骂和诽谤，因为\Quote{辱骂者不能承受天主的国；} 那宣认过基督的舌头，应当保持健全、纯洁，配得上它的尊荣。因为那依照基督的诫命、谈论和平、善良和公义的人，每日都在宣认基督。我们受洗时已经弃绝了世界；但现在，当我们受天主考验和认可、舍弃一切所有、跟随了主、并站立和活在祂的信德与敬畏中时，我们确实弃绝了世界。\par

6. 让我们以互相劝勉来坚固彼此，并在主内不断前进；这样，当祂按祂的仁慈赐予祂所应许的平安时，我们就能焕然一新、几乎成为新人回到教会，并被我们的弟兄或外邦人接纳，在一切事上得到改善和更新；那些先前在我们勇敢中钦佩我们光荣的人，如今将钦佩我们生活里的纪律。亲爱的弟兄们，祝你们常怀热忱，多多保重，并请记念我。\par

\SourceRecord{Translated by Robert Ernest Wallis.；Ante-Nicene Fathers；Vol. 5.；Edited by Alexander Roberts, James Donaldson, and A. Cleveland Coxe.；Buffalo, NY: Christian Literature Publishing Co.,；1886.；<http://www.newadvent.org/fathers/050606.htm>.}

% structure: p0001 source=h1 target=section reason=page-title

\section{\WorkTitle{第七封书信}}

\Emph{致圣职人员，论向天主的祈祷。}\par

\Emph{论据——本封书信的论点与前两封大致相同，只是在此劝勉要勤于祈祷。}\par

1. 西彼廉致诸位司铎与执事，弟兄们安。虽然我知道，亲爱的弟兄们，因我们对天主所怀的敬畏，你们也在不断的祈求与热切的祈祷中常向祂恳求，但我仍提醒你们虔敬的焦虑：为平息并恳求主，我们必须不仅以言语哀恸，还要以禁食、眼泪及各种迫切的方式。因为我们必须认识到并承认，那因这场灾祸而极其混乱的摧毁——它已大大蹂躏，甚至仍在蹂躏我们的羊群——是依照我们的罪降临在我们身上，因为我们没有遵守主的道，也没有遵守为我们的救恩所赐予的天上诫命。我们的主遵行了祂父的旨意，而我们却不遵行我们主的旨意；我们急于追求产业与利益，寻求满足我们的骄傲，完全投身于嫉妒与纷争，不关心纯朴与信德，只在言语上、而非行动上弃绝世俗，我们每个人都只取悦自己，而令所有人不悦——因此我们受到应得的打击，因为经上记载：\Quote{那知道主人旨意，却不遵行的仆人，必受许多鞭打。} \BibleRef{路 12:47} 然而，我们岂不更该受何等的鞭打与击打，因为连那些本该是他人德行生活榜样的宣信者，也不遵守纪律？因此，当一些人对自己的宣信过分自夸、不谦卑时，酷刑便临到他们身上，且是无间断的折磨、无终结的定罪、无死亡的安慰——这些酷刑并不轻易让他们抵达荣冠，反在刑架上摧残他们，直至使他们背弃信仰，除非有人蒙天主怜悯，在酷刑中离世，不是因酷刑终止而得荣耀，而是因死亡迅速而获荣光。\par

2. 这些事是因我们自己的过错与应得而承受的，正如天主的审判曾预先警告我们，说：\Quote{如果他们离弃我的法律，不遵行我的典章，亵渎我的律例，不守我的诫命，我就要用杖责罚他们的过犯，用鞭责罚他们的罪孽。} 正因如此，我们感到杖与鞭的责罚，因为我们既不以善行取悦天主，也不为我们的罪作补赎。让我们从内心深处、以整个心神恳求天主的仁慈，因为祂自己也接着说道：\Quote{然而，我的慈爱必不远离他们。} 让我们祈求，就必领受；若我们的领受有所延迟与缓慢，因为我们严重冒犯了祂，就让我们敲门，因为\Quote{凡敲门的，就给他开门} \BibleRef{路 11:10}，只要我们的祈祷、我们的叹息与眼泪，扣响那门；我们当以这些迫切并坚持，即使祈祷是同心合意地献上。\par

3. 因为——这更促使并催迫我写这封信给你们——你们应当知道（主曾屈尊显示并启示这事），在异象中曾说道：\Quote{祈求，就必得着。} 随后，在场的群众被吩咐为某些指定的人祈祷，但在他们的祈求中却有分歧的声音与不一致的意愿，这大大触怒了那曾说\Quote{祈求，就必得着} 的一位，因为百姓的不和谐、未能有一致而简单的弟兄同意与合一的和谐；正如经上记载：\Quote{天主使同心合意的人住在一家中；} 我们在宗徒大事录中读到：\Quote{众信友都同心合意，没有一个人说他所拥有的任何东西是自己的。} \BibleRef{宗 4:32} 主曾亲自吩咐我们，说：\Quote{这是我的命令：你们要彼此相爱。} \BibleRef{若 15:12} 又说：\Quote{我告诉你们：若你们中有两人在地上同心合意地祈求任何事，我在天之父必为他们成就。} \BibleRef{玛 18:19} 如果两人同心能成就如此大事，那么若全体都和睦，会成就何事？但如果我们领受主所赐的平安，众弟兄都同心，我们早已从天主的仁慈中求得我们所寻求的；也不致在此危及我们救恩与信德的时刻长久摇摆不定。是的，确实，这些祸患本不会临到弟兄们身上，如果弟兄们曾以同一心神同心合意。\par

4. 因为异象中还显示，有一位家主坐着，一个年轻人坐在他右边，他面带焦虑与些许不满，带着一种愤慨，右手托腮，神情忧伤地占据他的位置。但另一个人站在左边，拿着一面网，威胁要撒网，以捕捉站在周围的群众。看见的人惊奇这预示什么，有人告诉他：那坐在右边的青年之所以忧伤痛苦，是因为他的诫命没有被遵守；而左边的那个则在得意，因为他得到了机会，从家主那里获得毁灭的权力。这一切在这场灾害的风暴来临之前很久就已显明。我们已看见那曾显示的事成就了：当我们轻视主的诫命、不遵守祂所赐的救恩律法时，那仇敌便获得了行恶的能力，并用他的网网住了那些装备不全、过于粗心而无力抵抗的人。\par

5. 让我们热切祈祷，不断呻吟恳求。因为要知道，亲爱的弟兄们，不久前我在异象中也因此受到责备，说我们在祈祷中昏睡，没有警醒祈祷；而天主，祂责备祂所爱的人\BibleRef{希 12:6}，当祂责备时，责备是为了改正，改正是为了保存。因此，让我们击碎并摆脱睡眠的束缚，以紧迫和警醒祈祷，正如宗徒保禄命令我们所说：\Quote{不断祈祷，并在此事上警醒。}\BibleRef{哥 4:2}因为宗徒们也不分昼夜地祈祷；主自己，我们纪律的老师，我们榜样的道路，也常常警醒祈祷，正如我们在福音中所读：\Quote{祂出去，上山祈祷，整夜向天主祈祷。}\BibleRef{路 6:12}确实，祂所祈求的，是为我们而求，因为祂不是罪人，而是担当了别人的罪。但祂为我们祈祷，以至于在另一处我们读到：\Quote{主对伯多禄说：看，撒殚想要筛你们像筛麦子一样；但我已为你祈祷，使你的信德不致丧失。}\BibleRef{路 22:31}如果祂为我们和我们的罪劳苦、警醒并祈祷，我们更应该恒心祈祷；首先，向主自己祈祷和恳求，然后透过祂，向天主圣父赔补！我们有一位为我们罪过的辩护者和中保——耶稣基督，主和我们的天主，只要我们悔改以往的罪过，承认并认明我们的罪，即我们如今冒犯主的罪，并承诺今后遵行祂的道路，敬畏祂的诫命。圣父纠正并保护我们，只要我们仍在患难和困苦中坚守信德，也就是说，紧紧依附于祂的基督；正如经上所载：\Quote{谁能使我们与基督的爱隔绝呢？是患难吗？是困苦吗？是迫害吗？是饥饿吗？是赤身露体吗？是危险吗？是刀剑吗？}\BibleRef{罗 8:35}这些事都不能分离信徒，没有什么能撕裂那些紧依附于祂身体和血的人。这种迫害是对心灵的考验和察验。天主愿意我们被筛净和考验，正如祂一向考验祂的子民；然而在祂的考验中，帮助从未在任何时候缺少给信徒。\par

6. 最后，祂恩赐给我们这些最卑微的仆人，尽管我们置身于许多罪过中，不配祂的俯就，但祂出于对我们的仁慈，命令说：\Quote{告诉他，}祂说，\Quote{要平安，因为和平即将来临；但在此期间，有一点延迟，以使一些仍存留的人得以考验。}我们被这些神圣的俯就所提醒，要节制饮食，适量饮酒；免得世俗的诱惑削弱那已因天上活力而振奋的心胸，或免得因过于丰盛的宴饮而沉重的心智，在祈祷和恳求上不够警醒。\par

7. 我有责任不隐瞒这些特别的事，也不独自藏于自己的意识中——这些事可以使我们每个人受教和受引导。你们也不要将此信藏在自己中间，而要让弟兄们阅读。因为想要自己的弟兄不受警告和教导的人，就会拦截主用以屈尊警诫和教导我们的这些话。让他们知道我们被主考验，并愿他们永不失去那使我们一度相信祂的信德，即在此当前患难的冲突中。让每个人承认自己的罪，现在就脱去旧人的行为。\Quote{因为手扶着犁向后看的人，不配进天主的国。}\BibleRef{路 9:62}最后，罗特的妻子，当被救出时，违背命令向后看，便失去了逃脱的益处。\BibleRef{创 19:26}让我们不看那后面的事，魔鬼在那里叫我们回转，而要看前面的事，基督在那里召唤我们。让我们举目望天，免得大地及其欢愉和诱惑欺骗我们。让我们每人不仅为自己，也为所有弟兄向天主祈祷，正如主教导我们祈祷时，祂命令每人不要私下祈祷，而是吩咐他们祈祷时，为所有人以共同的祈祷和和谐的恳求祈祷。如果主看见我们谦卑和平安；如果祂看见我们彼此联合；如果祂看见我们畏惧祂的愤怒；如果我们因当前的磨难而受纠正和改正，祂将保护我们免受敌人骚动的侵害。惩戒先行，宽恕也必随后。\par

8. 让我们只有不停地祈求，并满怀信心必将领受，以纯朴和同心恳求主，不仅以呻吟，也以眼泪恳求，正如那些处于哀哭者之废墟与畏惧者之残存之间的人所应当恳求的；处于屈服者的多重杀戮与仍然站立者的微小坚定之间。让我们祈求和平早日恢复；使我们能在隐藏和危险中迅速得到帮助；使主屈尊向祂仆人显示的那些事得以实现——教会的复兴，我们救恩的保障；雨后的晴朗，黑暗后的光明，风暴与旋风后的平静安宁；父爱慈爱的援助，天主尊威惯常的伟大，以使迫害者的亵渎得以抑制，失足者的悔改得以更新，坚守者的坚定信德得以夸耀。亲爱的弟兄们，我衷心问候你们，永远平安，并记念我。以我的名义问候弟兄们；提醒他们记念我。再见。\par

\SourceRecord{Translated by Robert Ernest Wallis.；Ante-Nicene Fathers；Vol. 5.；Edited by Alexander Roberts, James Donaldson, and A. Cleveland Coxe.；Buffalo, NY: Christian Literature Publishing Co.,；1886.；<http://www.newadvent.org/fathers/050607.htm>.}

% structure: p0001 source=h1 target=section reason=page-title

\section{\WorkTitle{第八封书信}}

\Emph{致殉道者与宣信者}\par

\Emph{论点。—— \Person{西彼廉}，极其惊叹非洲殉道者的坚毅，以他们的同侪\Person{玛帕利库斯}为例，敦促他们坚持到底。}\par

\Person{西彼廉}致殉道者以及在基督我们的主和天主父内的精修者，永远的救恩。我最勇敢而蒙福的弟兄们，听到你们的信德和英勇，我们的慈母教会以此为荣，我由衷欢喜并感恩。最近，她的确曾荣耀，因为因着持久的宣认，那将基督的精修者驱逐流放的刑罚被承受了；然而，当前的宣认在受苦中更为英勇，因此更加显赫和尊荣。战斗加剧了，战斗者的荣耀也增加了。你们并未因酷刑的恐惧而退缩，反而被酷刑本身更加激励去战斗；你们勇敢而坚定地以准备好的热忱返回，要在极端的争战中战斗。我得知你们中有些已戴上了冠冕，有些则已接近胜利的冠冕；但所有被危险关在光荣团体中的人，都以同样共通的德性热情被激励去持续战斗，正如基督的士兵在神圣军营中所应做的：好使任何诱惑都不能欺骗你们信德的不朽坚定，任何威胁都不能恐吓你们，任何苦难或酷刑都不能战胜你们，因为\Quote{那在我们内的，比那在世界上的更大}；地上的刑罚不能胜过那使人降下的，正如天主的保护胜过那使人升起的。弟兄们光荣的战斗证明了这一真理，他们成为其余人克服酷刑的引领者，提供了德性和信德的榜样，在斗争中争战，直到斗争屈服，被克服。最勇敢的弟兄们，我能用什么赞美来称许你们？用什么言语来颂赞你们心中的力量和信德的坚毅？你们忍受了最严酷的酷刑拷问，直到光荣的终点，没有向苦难屈服，反而苦难向你们屈服了。酷刑的终结，即酷刑本身未能给予的，冠冕给予了。更加严厉的酷刑拷问持续了长时间，得到这样的结果：不是要推翻坚定的信德，而是要更快地将天主的人送到主那里。在场的人群惊叹地观看这场天上的比赛——天主的比赛、属灵的比赛、基督的战争——看到祂的仆人们以自由的声音、不屈的意志、神圣的德性站立——固然赤手空拳，没有世俗的武器，却满怀信心，以信德为武装。受酷刑者比施刑者更为勇敢；被击打撕裂的肢体战胜了弯曲撕裂的铁钩。带着全部怒火的多次鞭打不能征服无敌的信德，即使包裹内脏的薄膜破裂，受酷刑的不再是天主仆人的肢体，而是他们的伤口。鲜血流淌，足以熄灭迫害的烈焰，以其光荣的血腥制服地狱的火焰。哦，主面前是何等的景象——在祂士兵的忠诚和奉献中，是何等崇高、伟大、悦乐天主的眼睛！正如圣咏所载，圣神同时向我们说话和警告：\Quote{上主看圣者的死，是珍贵的。}珍贵的死是以其鲜血买得了不朽，从其德行的圆满获得了冠冕。基督在其中何等喜乐！祂何等甘愿在如此忠信的仆人内战斗并得胜，作为他们信德的保护者，且给予信徒如同那接受者所相信他领受的一样！祂亲临自己的比赛；祂扶持、坚固、激励祂圣名的斗士和捍卫者。那一次为我们战胜死亡的祂，总是在我们内战胜死亡。祂说：\BibleQuote{当他们}，祂说，\BibleQuote{把你们交出时，不要思虑你们要说什么：因为在那时刻，自会给你们要说什么。因为说话的不是你们，而是你们父的圣神在你们内说话。}\BibleRef{玛 10:19}当前的斗争证明了这话。当最蒙福的\Person{玛帕里库斯}在酷刑中对总督说\Quote{你们明天将见一场比赛}时，充满圣神的声音从殉道者口中迸发。那他以德性和信德作证所说的话，主使之成就了。一场天国的比赛展示了，天主的仆人应在应许的战斗比赛中戴上冠冕。这就是先知\Person{依撒意亚}昔日所预言的比赛，说：\Quote{你们与人相争，并非轻易的事，因为天主安排了这场争斗。}为了表明这争斗将是什么，他接着说了这话：\BibleQuote{看，一位贞女将怀孕生子，你要称祂的名字为厄玛奴耳。}\BibleRef{依 7:14}这就是我们信德的战斗，在其中我们参战，在其中我们得胜，在其中我们受冠冕。这就是蒙福的宗徒\Person{保禄}向我们指明的战斗，我们应当奔跑并获得光荣的冠冕。他说：\BibleQuote{你们岂不知道吗？}他说，\BibleQuote{在赛场上赛跑的，大家都跑，但只有一个得奖？你们也该这样跑，好能得到奖赏。}\BibleQuote{他们这样做，为得到可朽坏的冠冕，而我们要得到不朽的。}\BibleRef{格前 9:24}此外，他陈述自己的战斗，并宣告自己即将为主成为祭品，说：\BibleQuote{我已被奠祭，我离世的时期到了。这场好仗，我已打完；这场赛跑，我已跑到终点；这信德，我已保持了。从今以后，正义的冠冕已为我预备下了，就是主，正义的审判者，到那一日必要赏给我的；不但赏给我，也赏给一切爱慕祂显现的人。}\BibleRef{弟后 4:6}因此，这场战斗，昔日由先知预言，由主发起，由宗徒进行，\Person{玛帕里库斯}以其本人及其同伴的名义再次向总督应许了。那忠信的声音在它的应许上没有欺骗；他展示了他所承诺的战斗，并获得了应得的奖赏。我不仅恳求，而且劝勉其余的人，你们都当跟随那位如今已极为蒙福的殉道者，以及那场战斗的其他伙伴——士兵和同道，信德坚定，受苦忍耐，在酷刑中得胜——使那些因宣认的纽带和牢狱的款待而联合的人，也得以在德行的圆满和天上的冠冕中联合；使你们以喜乐擦干我们慈母教会的眼泪，她为许多人的沉船和死亡而哀恸；并且使你们以榜样激励其他站稳之人的坚定。如果战斗呼唤你们，如果你们比赛的日子来到，要勇敢参战，坚定战斗，因为知道你们是在一位临在的主眼前战斗，你们是以宣认祂的圣名达到祂自己的光荣；祂不仅这样观看祂的仆人，而且祂自己也在我们内战斗，祂自己参战——祂自己在我们的战斗比赛中不仅加冕，而且被加冕。但如果在你们比赛的日子之前，因天主的仁慈，和平降临，仍要让你们保有健全的意志和光荣的良知。你们中没有人应因仿佛不如那些在你们之前忍受了酷刑、征服并践踏了世界、从而以光荣的道路到达主那里的人而感到悲伤。因为主是\BibleQuote{洞察肺腑和心肠的。}\BibleRef{默 2:23}祂看透隐秘之事，并注视那隐藏的。为从祂那里获得冠冕，祂自己的见证就足够了，祂将审判我们。因此，亲爱的弟兄们，两种情况同样崇高而显赫——前者更为稳妥，即以我们胜利的圆满迅速到达主那里；后者更为喜乐，即在获得荣耀之后，得到休假，在教会的赞美中繁荣。哦，我们有福的教会，被天主屈尊的恩宠光照，在我们时代被殉道者的光荣鲜血所彰显！她先前因弟兄们的行为而洁白，如今因殉道者的血而变成紫红。在她的花朵中，不缺玫瑰，也不缺百合。现在，愿每个人竭力争取两种荣誉中最大的尊严。愿他们接受冠冕，或是白色的，如劳苦的冠冕，或是紫红的，如受苦的冠冕。在天上的军营中，和平与斗争各有其花朵，基督的士兵可为此受冠冕得光荣。我向你们，最勇敢且亲爱的弟兄们，在主内衷心告别；并请记念我。再会。\par

\SourceRecord{Translated by Robert Ernest Wallis.；Ante-Nicene Fathers；Vol. 5.；Edited by Alexander Roberts, James Donaldson, and A. Cleveland Coxe.；Buffalo, NY: Christian Literature Publishing Co.,；1886.；<http://www.newadvent.org/fathers/050608.htm>.}

% structure: p0001 source=h1 target=section reason=page-title

\section{第九封书信}

\Emph{致圣职人员，论及某些司铎在迫害平息之前且未经主教知情，贸然将平安赐予背道者之事}\par

\Emph{谕旨。——本书信之谕旨见于第十四封书信的以下文字：——他说：\Quote{致司铎与执事：司铎职之刚毅未曾缺失，致使某些人，既对纪律不甚留心，又因鲁莽急躁而仓促行事，已经开始与背道者共融者，受到了制止。}}\par

西彼廉致诸位司铎与执事，他的弟兄们，安好。亲爱的弟兄们，我长久忍耐，希望我含默的容忍有助于平静。但因某些人无理而鲁莽的僭越，以其胆大妄为企图扰乱殉道者的尊荣、宣信者的谦逊以及全体子民的安宁，我不可再保持沉默，唯恐过度的缄默会给子民和我们自己带来危险。因为当某些司铎，既不记得福音也不记得自己的身份，而且既不思虑主将来的审判，也不思虑此时被置于他们之上的主教，竟将全部权柄归为己有——这在我们的前任治下是绝对不曾有过的事——并以此羞辱和藐视主教时，我们岂不应惧怕主的不悦所带来的一切危险？\par

我固然希望，若能不牺牲我们弟兄的安全，他们就能得偿所愿，获得一切；我本可装作不知，忍受我主教权位所受的羞辱，如同我素常所忍受的那样。但现在不是容我们姑息的时候，因为我们的弟兄团体正被你们当中的一些人欺骗；这些人虽没有能力恢复救恩，却想要讨人喜欢，反而成了背道者更大的绊脚石。因为那件迫害所迫使他们犯下的滔天大罪，犯下者自己也知道其严重性。因为我们的主和审判者曾说：\BibleQuote{凡在人前承认我的，我在天上的父面前也必承认他；但凡在人前否认我的，我也必否认他。} \BibleRef{玛 10:32} 祂又说：\BibleQuote{人子的一切罪和亵渎，都可得赦免；但那亵渎圣神的，永不得赦免，而应受永罚。} \BibleRef{谷 3:28} 而且蒙福的宗徒也曾说：\BibleQuote{你们不能喝主的杯，又喝魔鬼的杯；不能共享主的筵席，又共享魔鬼的筵席。} \BibleRef{格前 10:21} 凡将这番话隐瞒不向我们弟兄说的人，就是在欺骗他们，使这些可怜人本可真心悔改，以祈祷和行为满足天主——祂既为父又仁慈——却被更深地诱入丧亡；而那些本可振作起来的人，反而跌得更深。因为在较小的罪中，罪人通常须按规定时间做补赎，按纪律的规则前来公开告解，并经主教及圣职人员覆手而获得共融权；但现在，在规定的时期尚未届满之时，迫害仍在肆虐，教会本身的平安尚未恢复，他们就被准予共融，他们的名字被提呈；而补赎尚未完成，告解尚未举行，主教及圣职人员的覆手尚未施行，圣体却被赐予他们；尽管经上记载：\BibleQuote{凡不配地吃主的饼、喝主的杯的，就是干犯主的身体和主的血。} \BibleRef{格前 11:27}\par

但现在，那些很少遵守圣经律法的人并不有罪；有罪的是那些身居职位却未将此事告知弟兄的人，好使弟兄们在上级的教导之下，以敬畏天主之心，并依照祂所赐予和规定的诫命行事。再者，他们还使蒙福的殉道者受人非议，并让天主光荣的仆人与天主的司铎一同卷入；因此，尽管殉道者念及我的位分，已致信于我，要求查验他们的意愿，并赐予平安——但应先等我们的慈母教会本身，因主的仁慈，以及天主的保护，使我返回祂的教会，而获得平安——然而这些人，无视蒙福的殉道者与宣信者为我所保持的尊荣，藐视主的律法和那同一殉道者与宣信者所要求遵守的诫命，在迫害的恐惧尚未平息之前，在我回去之前，几乎甚至在殉道者离世之前，就与背道者共融，并为他们奉献和赐予圣体：即便殉道者在其荣耀的热忱中，可能对圣经不够审慎，并有所额外要求，也应由司铎和执事加以规劝，正如过去一向所行的那样。\par

为此，天主的责罚昼夜不止地惩治我们。因为除了夜间的异象，白昼间，在我们中间，那纯真年龄的孩童充满圣神，在神魂超拔中以眼看见，以耳听见并说出那些主屈尊来警告和教导我们的事。当那命我退避的主再使我回到你们中间时，你们必会听见一切。与此同时，你们当中那些鲁莽、不慎、自夸、不尊重人的人，至少该敬畏天主，知道若他们仍坚持同一行径，我将行使主命我运用的劝诫权柄；好使他们暂时不得奉献，并须在我、宣信者本人以及全体子民面前辩护自己的理由，直到蒙天主允许，我们重新聚于慈母教会的怀抱之时。关于此事，我已致信殉道者、宣信者及子民；两封信我都已命人读给你们听。亲爱的弟兄们，我衷心祝愿你们在主内永远安好，并请记得我。安好。\par

\SourceRecord{Translated by Robert Ernest Wallis.；Ante-Nicene Fathers；Vol. 5.；Edited by Alexander Roberts, James Donaldson, and A. Cleveland Coxe.；Buffalo, NY: Christian Literature Publishing Co.,；1886.；<http://www.newadvent.org/fathers/050609.htm>.}

% structure: p0001 source=h1 target=section reason=page-title

\section{\WorkTitle{信札第十通}}

\Emph{致寻求为背教者赐予平安的殉道者与宣信者。}\par

\Emph{论据。——此信的缘由如下文第十四信所述：当我发现那些以亵渎的接触玷污了手与口，或同样以邪恶的证书玷污了良心的人，等等。}\par

1. 西彼廉致他所爱的弟兄，殉道者与宣信者，问安。我处境中的焦虑和对主的畏惧，我勇敢而亲爱的弟兄们，促使我以书信劝告你们：那些如此虔诚而勇敢地持守主之信德的人，也当持守主的法律与纪律。因为所有基督的士兵都当遵守他们统帅的命令，而你们更当特别遵行祂的命令，因为你们已成了他人勇气与敬畏天主的榜样。我原以为，与你们同在的司铎和执事们，会按照往昔在我前任时的惯例，更充分地以福音的法律劝告和教导你们，使出入监狱的执事们以他们的忠告和圣经的诫命来节制殉道者的愿望。但现在，我怀着极大的忧伤得知，他们不仅没有向你们提示神圣的诫命，反而加以限制，以致某些司铎既不畏惧天主，也不尊重主教，竟松弛了你们自己出于谨慎而遵行的事，即在天主面前和尊重天主的司铎的事。虽然你们曾致信于我，请求审查你们的愿望，并在迫害结束后，当我们与神职人员重新会聚时，为某些背教者赐予平安；但这些司铎违背福音的法律，也违背你们恭敬的请求，在补赎尚未完成，甚至最严重、最可耻的罪过尚未告明之前，在主教和司铎尚未为悔罪者覆手之前，竟胆敢以他们的名义奉献并给予他们圣体圣事，即亵渎主的圣体，尽管经上记载：\Quote{无论何人，若不配地吃主的饼，喝主的杯，就是干犯主的身、主的血。}\BibleRef{格前 11:27}\par

2. 对于背教者，此事或可予以宽恕。因为哪个死人不想赶快复活？谁不渴望获得自己的救恩？但管理他们的人有责任遵守规定，教导那些急躁或无知的人，以免本应作群羊牧人的人反倒成了屠夫。因为让步于那些导致毁灭的事，乃是欺骗。背教者并不能因此被提升，反而因得罪天主，更被推向灭亡。因此，让他们甚至从你们那里学习他们本应教导的事；让他们将你们的请求和愿望留给主教，并等候成熟和平的时机，按你们的请求赐予平安。首先，母亲教会应从主那里先获得平安，然后，按照你们的愿望，再考虑她子女的平安。\par

3. 我最勇敢而亲爱的弟兄们，既然我听说你们受到一些人的无耻逼迫，你们的谦逊受到强暴；我以我所能的恳求，求你们铭记福音，并思量以往你们的先辈殉道者曾如何让步，他们在各方面何等谨慎，你们也当焦虑而审慎地衡量那些向你们请愿者的愿望，因为你们作为主的朋友，将来要与祂一同审判，必须检查每个人的行为、事迹和功德。你们也要考虑他们罪过的种类和性质，以免任何事被你们仓促而不配地允诺，或被我施行，以致我们的教会甚至在外邦人面前蒙羞。因为我们常受眷顾和管教，且被提醒，要无玷污、无违犯地遵守主的诫命；我发现这在你们那里并未停止，以至神圣的审判也教导你们中许多人在教会的纪律上受训。这一切若能成就，就在于你们以审慎的宗教意识调节那些向你们所求的事，洞察并抑制那些因徇情面而施恩惠，或企图以不法交易牟利的人。\par

4. 关于此事，我已致信给神职人员和教友，并吩咐将这两封信读给你们听。但你们也当以你们的勤勉加以纠正和改善，按名指定你们愿意赐予平安的人。因为我听说，证书是这样给予某些人的，以致说：\Quote{愿某某人与他的朋友们一起被接纳进入共融。}殉道者从未做过这样的事，以使一个模糊而盲目的请求后来归咎于我们。因为说\Quote{某某人与他的朋友们}打开了宽阔的门路；于是可能有二十、三十或更多的人被呈报给我们，他们被说成是接受证书者的邻居、亲戚、自由人和仆人。因此我请求你们，在证书上按名指定那些你们亲眼所见、你们所认识、其补赎已近乎完全满足的人，并按着信德与纪律将信件寄给我们。我祝愿你们，最勇敢而亲爱的弟兄们，在主内永远安康；并请记念我。祝好。\par

\SourceRecord{Translated by Robert Ernest Wallis.；Ante-Nicene Fathers；Vol. 5.；Edited by Alexander Roberts, James Donaldson, and A. Cleveland Coxe.；Buffalo, NY: Christian Literature Publishing Co.,；1886.；<http://www.newadvent.org/fathers/050610.htm>.}

% structure: p0001 source=h1 target=section reason=page-title

\section{书信十一}

\Emph{致人民}\par

\Emph{论点——这封信的要旨亦见于\WorkTitle{第十四封信}，\Quote{在人民中间，}他说，\Quote{我已尽力平息他们的心，并教导他们要持守在教会的纪律中。}}\par

西彼廉向坚守的人民中的弟兄们问安。亲爱的弟兄们，我从自身知道你们为我们的弟兄的跌倒而哀恸悲伤，我也与你们一同为每个人哀恸悲伤，并感受那蒙福的宗徒所说的话：\BibleQuote{有谁软弱，我不软弱呢？有谁跌倒，我不焦急呢？}\BibleRef{格后 11:29}他又在他的书信中规定说：\BibleQuote{若一个肢体受苦，所有的肢体都一同受苦；若一个肢体得荣耀，所有的肢体都一同欢乐。}\BibleRef{格前 12:26}因此，我同情你们的痛苦和悲伤，为我们的弟兄们，他们因迫害的严厉而失足倒下，用他们的创伤给我们带来了同样的痛苦，因为他们带走了我们的一部分心肠。天主的怜悯能够医治他们。然而，我认为不可急躁，也不可轻率地做任何事，免得在寻求和平时更严重地招致天主的义怒。蒙福的殉道者曾就某些人写信给我，要求审查他们的愿望。当主赐给我们大家和平，我们开始回到教会时，那时我们将在你们面前，凭你们的判断来审查每个人的愿望。\par

2. 然而我听说某些司铎，既不记念福音，也不考虑殉道者写给我的信，也不为主教保留其司祭职和尊严的荣誉，已经开始与失足者共融，并为他们献祭，给他们圣体，而他们本应按时得到这些。因为，在那些不是得罪天主的较小罪过中，补赎可以在一段时期内完成，并通过审查履行补赎者的生活来告解，且除非主教和圣职人员的双手先按在他身上，无人能领受共融；那么在这样重大和极端的罪过中，按照主的纪律，更应当谨慎和适度地遵守这一切！确实，我们的司铎和执事本应警告你们，使他们可以照顾托付给他们的羊群，并藉着天主的权威教导他们通过祈祷获得救恩的方式。我深知我们人民的和平与敬畏，他们本会在补赎和祈求天主息怒时保持警醒，除非某些司铎为了讨好他们而欺骗了他们。\par

3. 因此，你们自己也要引导他们每个人，并按照天主的诫命，以你们的劝告和你们的节制来引导失足者的心。不要让任何人过早地摘取未熟的果实。不要让任何人将他破碎的船匆忙重新抛入深海，除非他仔细修好了它。不要让任何人急于接受和穿上撕裂的长袍，除非他看到它被巧匠补好，并经过漂洗者整理。请他们耐心听从我的劝告。请他们等待我的回来，以便当天主仁慈地让我来到你们那里时，我能够与许多共同主教一起，按照主的纪律召集起来，并在宣信者和你们的意见面前，审查蒙福的殉道者的信件和愿望。关于此事，我已经写信给圣职人员和殉道者及宣信者，这两封信我都已指示向你们宣读。亲爱的、最渴望的弟兄们，我衷心祝你们在主内永远平安，并请记念我。再见。\par

\SourceRecord{Translated by Robert Ernest Wallis.；Ante-Nicene Fathers；Vol. 5.；Edited by Alexander Roberts, James Donaldson, and A. Cleveland Coxe.；Buffalo, NY: Christian Literature Publishing Co.,；1886.；<http://www.newadvent.org/fathers/050611.htm>.}

% structure: p0001 source=h1 target=section reason=page-title

\section{第十二封书信}

\Emph{致圣职人员，论及背道者与慕道者：他们不应被置于无人监管之下。}\par

\Emph{论纲——本书信以及其后书信的要点，见于下述第十四封书信。\Quote{但后来，}他说，\Quote{有些背道者，无论出于自愿或受他人怂恿，竟大胆要求，仿佛试图以暴力强夺殉道者与宣信者所应许予他们的平安，}等等。}\par

1. 西彼廉问候他的弟兄——各位长老与执事。我亲爱的弟兄们，我多次写信给你们，却不见你们回信，令我诧异；虽然我们兄弟团体的利益与需要，若能从你们那里获得消息，我便能准确查考并商议事务的管理，那自是最好的。但是，既然我看出目前尚无法前往你们那里，而夏季已经开始——这个季节常有持续且严重的疾病扰乱——我认为必须这样处理我们的弟兄：那些从殉道者获得证书、并能借其在天主前的特权得助的人，若遭遇任何不幸或疾病危险，不必等待我的到来，可在任何在场的长老面前，若找不到长老而死亡开始临近，甚至可在执事面前，宣认他们的罪过，以便藉为他们忏悔而覆手后，他们能带着殉道者通过致我们的书信希望赐予他们的平安，回归主前。\par

2. 也以你们在场关怀其余背道的子民，用你们的安慰鼓励他们，使他们不致丧失信德和天主的仁慈。因为那些温和、谦卑、以真诚忏悔持守善工的人，绝不会被主的助佑和扶持所遗弃；神圣的救治也将赐予他们。至于慕道者，若有人遭遇危险、濒临死亡，你们的警惕不可缺席；不可拒绝向哀求天主恩宠的人施予主的仁慈。我祝愿你们，亲爱的弟兄们，始终在主内衷心安康，并记念我。请以我的名义问候全体兄弟团体，提醒他们并请他们记念我。再见。\par

\SourceRecord{Translated by Robert Ernest Wallis.；Ante-Nicene Fathers；Vol. 5.；Edited by Alexander Roberts, James Donaldson, and A. Cleveland Coxe.；Buffalo, NY: Christian Literature Publishing Co.,；1886.；<http://www.newadvent.org/fathers/050612.htm>.}

% structure: p0001 source=h1 target=section reason=page-title

\section{\WorkTitle{书信第十三}}

\Emph{致圣职人员，论及那些急于获得平安的人。公元二五〇年。}\par

\Emph{论据——平安必须通过补赎获得，而补赎藉遵守诫命得以实现。那些患病的人，若因殉道者的代祷而得缓解，可被接纳获享平安；但其他人须被阻止，直至教会的平安得到保障。}\par

1. 西彼廉致诸位长老与执事，他的弟兄们，安好。我已阅读你们的信函，亲爱的弟兄们，信中提到你们并未缺乏对弟兄们的有益劝告，即要摒弃一切鲁莽急躁，向天主显示敬虔的忍耐，好使我们在祂的怜悯下聚集时，能按照教会纪律讨论各类事宜，尤其因为经上记载：\BibleQuote{要回想你是从哪里跌落的，并要悔改。}\BibleRef{默 2:5} 那悔改的人，是记念神圣诫命，以温柔和忍耐，并服从天主的司铎，藉其服从和正义的行为，在主面前蒙恩的人。\par

2. 然而，由于你们提到有些人无礼，急切地要求被接纳进入共融，并希望在此事上由我给你们一些规则，我认为在上一封给你们的信中已充分写明了此事：那些从殉道者获得证书，并能藉其帮助在主面前解决其罪债的人，若开始遭受任何疾病或危险，当他们做了告解，并由你们在他们身上实行覆手以承认其补赎后，应被交托给主，并享有殉道者所应许的平安。至于其他人，他们没有从殉道者获得任何证书，心中嫉妒（因为这不单是少数人、一个教会、一个行省的事，而是全世界的事），必须仰赖主的保护，等候教会本身的公共平安。因为这对我们所有人的谦逊、纪律乃至生命都是适宜的，即当首领们与圣职人员会面，在那些站立的教友面前（他们本身也因其信德和敬畏而应受尊敬），我们能够以共同商议的虔诚来安排一切。但是，对那些急切者而言，这是何等不虔敬和有害的事啊——当那些流亡者、被逐出家乡、被剥夺一切财产的人尚未回到教会时，一些背教者却急于抢先于宣信者本人，在他们之前进入教会！如果他们如此过度焦虑，他们已在自己的权力中拥有他们所要求的，因为时机本身比他们所求的更为慷慨地提供给他们。战斗仍在进行，争斗每日都在举行。如果他们真正并恒心地对自己所做的事忏悔，且其信德的热情得胜，那么那不能被延迟的人将可获赐冠冕。我祝你们，亲爱的弟兄们，永远在主内安康；并请记得我。请以我的名义问候全体弟兄，告诉他们纪念我。再会。\par

\SourceRecord{Translated by Robert Ernest Wallis.；Ante-Nicene Fathers；Vol. 5.；Edited by Alexander Roberts, James Donaldson, and A. Cleveland Coxe.；Buffalo, NY: Christian Literature Publishing Co.,；1886.；<http://www.newadvent.org/fathers/050613.htm>.}

% structure: p0001 source=h1 target=section reason=page-title

\section{第十四封书信}

\Emph{致聚集在罗马的长老与执事。}\par

\Emph{主旨——他陈述了自己退隐的原因及在此期间所行之事，并将他写给其子民的书信副本寄往罗马以作辩护；不仅如此，他还引用了那些书信中的原话。}\par

\Person{西彼廉}致他聚集在罗马的弟兄们，即长老与执事，问候。得知亲爱的弟兄们，我所做和正在做的事以一种有些歪曲不实的方式被传告给你们，我认为有必要写这封信给你们，以便向你们陈述我的行为、我的纪律和我的勤勉；因为，正如主的诫命所教导的，骚乱刚一爆发，百姓便以激烈的呼喊再三要求我出面，我考虑到与其说是自身的安危，毋宁说是弟兄们的公共安宁，便暂时退隐，以免因我过于鲁莽的出现，使已开始的骚动进一步激化。然而，虽身体缺席，我却在精神、行动和劝告上毫无欠缺，以致未能尽我微薄的能力，按照主的规诫，在一切事上以我的忠告为弟兄们提供任何我能给予的益处。\par

我所做的事，这十三封在不同时间发出的书信已向你们宣告，我已将它们转寄给你们；在这些信中，按照信德之律和敬畏天主，靠着主的助佑，我微薄的能力所能尽力而为的，对圣职界的劝告、对宣信者的鼓励、对流放者必要时给予的斥责、以及对整个弟兄团体恳求天主仁慈的呼吁与劝勉，都一无所缺。后来，当酷刑来临时，我的话传到了那些受刑的弟兄以及尚在监禁中将要受刑的弟兄那里，以坚强和安慰他们。此外，当我发现那些以亵渎性的接触玷污了手脚和口唇，或是以邪恶的证书同样感染了良心的人，四处恳求殉道者，并以纠缠不休和过分的请求腐化宣信者，以致不分青红皂白地、不对个别人进行审查，每天发出成千上万违反福音律法的证书时，我便写信，尽可能以我的劝告将殉道者和宣信者召回到主的诫命上。对长老和执事，司祭职的活力也未欠缺；以致一些不太留意纪律、仓促鲁莽、已开始与背教者共融的人，被我的干预所制止。此外，在百姓中，我也尽力安抚他们的心，并教导他们维护教会纪律。\par

后来，当一些背教者，或出于自愿，或受他人煽动，以大胆的要求冲出，似乎要竭力强求殉道者和宣信者所许诺给他们的平安时，关于此事我也两次写信给圣职界，并命令读给他们听；即为了暂时以任何方式缓和他们的暴力，但凡那些从殉道者处领受证书的人，若在离世前做了告解并接受了覆手以行补赎，他们应带着殉道者所许诺的平安，被交托于主。我并非在此给他们制定法律，或鲁莽地自命为指令的制定者；而是因为既应尊敬殉道者，又应遏制那些急于扰乱一切之人的激烈行为；此外，当我读到你们最近通过副执事\Person{克莱孟提乌斯}寄给我圣职界的信函，大意是应给予那些背教后患病并忏悔地渴望共融者援助时，我认为最好持守你们的判断，以免我们本应在一切事上合一并一致的程序，在任何方面有所分歧。其余人的案件，即使他们从殉道者处领受了证书，我也下令完全推迟，保留至我亲临之时；这样，当主赐予我们和平，数位主教开始聚集一处时，我们便能安排并改革一切，也得到你们忠告的助益。亲爱的弟兄们，我衷心祝愿你们永远安好。\par

\SourceRecord{Translated by Robert Ernest Wallis.；Ante-Nicene Fathers；Vol. 5.；Edited by Alexander Roberts, James Donaldson, and A. Cleveland Coxe.；Buffalo, NY: Christian Literature Publishing Co.,；1886.；<http://www.newadvent.org/fathers/050614.htm>.}

% structure: p0001 source=h1 target=section reason=page-title

\section{第十五封书信}

\Emph{致莫依斯与玛克西穆，以及其余各位宣信者。}\par

\Emph{本信主旨见于下文第三十一封书信，其中罗马圣职团这样说：\Quote{为此我们深感并致以丰厚的谢意，因你们以书信照亮了他们监牢的幽暗。}}\par

1. \Person{西彼廉}致\Person{莫依斯}与\Person{玛克西穆}司铎及其他宣信者，诸位弟兄，安好。\Person{塞莱里诺}，你们信德与勇德的同伴，在光荣战事中天主的士兵，来到我这里，亲爱的弟兄们，他强有力地使我心中浮现你们全体及每一位。我见他在我面前，便看见了你们全体；当他亲切而屡次谈及你们对我的爱时，我在他的言语中听到了你们。当有像他这样的人从你们那里给我带来这样的消息时，我极为喜乐。在某种意义上，我也与你们同在牢中。我想，我既与你们的心如此相连，便同享天主嘉许的喜乐。你们每个人的爱使我有份于你们的荣光；圣神不允许我们的爱分离。宣认将你们关在监牢，而情爱将我也关在那里。我确实日夜记念你们，无论在多人同献的祭献中祈祷，还是独自以私下的恳求祈祷时，我都向主祈求对你们冠冕与赞美的圆满承认。然而我微薄的能力不足以回报你们；你们在祈祷中记念我时，给的更多，因为你们已只呼吸天上的事，只默想神圣的事，甚至因受苦的延迟而升到更高之处；并且因时间的漫长，你们的荣耀非但不减，反而增长。第一次且唯一的宣认便使人有福；你们每当被要求离开牢狱时，却以信德和勇德选择留在牢中，你们便是在宣认；你们的赞美如同日子一样多；月份流转，你们的功勋常增。那一时受苦而得胜者只胜一次；但不断与刑罚搏斗、不被苦难击败者，则日日加冕。\par

2. 如今，让长官与执政官或总督去吧！让他们以一年一度的尊荣标志和十二束权杖夸耀吧！看，你们身上的天上尊荣已被一年荣耀的光辉所印证，并且已在得胜荣耀的延续中，越过了回归之年的循环。旭日与残月照亮世界；但在你们那里，那造日月的主在你们的囚室里是更大的光，基督的光辉在你们心中、在你们思想中闪耀，以那永恒明亮的光照亮了牢狱的幽暗——这对别人是那样可怖与致命。冬季经过月份的交替而逝去；但你们被关在牢中，所经历的并非冬日的严寒，而是逼迫的寒冬。冬天之后是温和的春天，以玫瑰为乐，以花朵为冠；但在你们面前，有乐园欢愉中的玫瑰与花朵，天上的花环编在你们的额上。看，夏天丰收累累，打谷场上堆满谷粒；但你们既播种了荣耀，便收获荣耀的果实，且被放在主的打谷场上，看那糠秕被不灭的火烧尽；你们自己像麦粒，经过簸扬的珍贵谷粒，如今洁净并收在仓里，将牢狱的居所视为谷仓。秋天也不缺少履行季节职务的神恩。葡萄在外被压榨，那将来流入杯中的葡萄在压榨机中被踩踏。你们，主葡萄园中丰硕的串串，果已成熟的枝条，被世俗压力的磨难受挤压，在折磨人的牢狱中充满你们的酒榨，流出你们的血代替酒；勇敢忍受苦难，你们甘愿饮下殉道的杯。如此，岁月在主仆身上流转——如此，季节的交替以神功与天上赏报而庆祝。\par

3. 极有福的是你们中那些踏着这些荣耀足迹、已离开世界的人；他们走完了德性与信德的旅程，获得了主的拥抱与亲吻，获得了主本身的喜乐。然而你们的荣耀并不少，你们仍在战斗，将要跟随同伴们的荣耀，长久作战，以不动摇、不摇撼的信德站稳，日日以你们的德行在天主眼前展示一幕景象。你们的斗争愈长，冠冕便愈高。这斗争是一个，却充满多种多样的战斗；你们以勇气的力量征服饥饿，轻视干渴，践踏牢狱的污秽与刑罚居所的恐怖。刑罚在那里被制伏；酷刑被耗尽；死亡不是被惧怕，而是被渴望，被永生的赏报所克服，以致得胜者被永恒的生命加冕。你们的心志如今该是怎样，何等高举，何等宽广，既然如此重大之事已被决定，除了天主的诫命与基督的赏报，别无他念！那时，意志唯独是天主的旨意；虽然你们仍在肉身中，现在所活的不是今世的生命，而是来世的生命。\par

现在，亲爱的弟兄们，你们还应当记念我；在你们伟大而神圣的思虑中，也要在心意和心神中想起我；当那声音——因告解的净化而辉煌，因其尊荣的持续旋律而可赞——上达天主的耳中时，我应常在你们的祈祷和恳求中；天门为之敞开，它从此被征服的世界区域上升到上面的领域，并从主的良善中获得它所求的一切。因为你们从主的仁慈中求什么，而你们不配获得呢？——你们如此遵守了主的诫命，以纯朴信德的活力保持了福音的纪律，以未受玷污的德行之荣耀，勇敢地站于主的诫命和祂的宗徒旁边，并以你们殉道的真理坚定了许多人的动摇信德。诚然，你们是福音的见证人，确实是基督的殉道者，安息于祂的根上，建立在坚固的磐石基础上，你们将纪律与德行结合，引导他人敬畏天主，并将你们的殉道作为榜样。我愿你们，极其勇敢而亲爱的弟兄们，永远衷心告别，并记念我。\par

\SourceRecord{Translated by Robert Ernest Wallis.；Ante-Nicene Fathers；Vol. 5.；Edited by Alexander Roberts, James Donaldson, and A. Cleveland Coxe.；Buffalo, NY: Christian Literature Publishing Co.,；1886.；<http://www.newadvent.org/fathers/050615.htm>.}

% structure: p0001 source=h1 target=section reason=page-title

\section{第十六封书信}

\Emph{宣信者致\Person{西彼廉}。}\par

\Emph{内容提要——\Person{路西安}以殉道者之名所写的证书。}\par

所有宣信者致\Person{西彼廉}神父，安好。凡在您看来其犯罪后所作交代令人满意者，我们均已给予平安；我们亦希望您将这封回函告知其他主教。我们吩咐您与圣殉道者共享平安。\Person{路西安}写此信，在场的有圣职人员，一位驱魔者和一位读经师。\par

\SourceRecord{Translated by Robert Ernest Wallis.；Ante-Nicene Fathers；Vol. 5.；Edited by Alexander Roberts, James Donaldson, and A. Cleveland Coxe.；Buffalo, NY: Christian Literature Publishing Co.,；1886.；<http://www.newadvent.org/fathers/050616.htm>.}

% structure: p0001 source=h1 target=section reason=page-title

\section{第十七封书信}

\Emph{关于前述及后续书信——致司铎与执事。}\par

\Emph{论旨——在教会恢复平安之前，不应重视来自殉道者的证书。}\par

\Person{西彼廉}问候司铎与执事诸位弟兄。主说：\BibleQuote{我垂顾谁呢？岂非那谦卑安静、敬畏我言语的人吗？}\BibleRef{依 66:2} 虽然我们都当如此，但那些在严重失足后，必须藉着真诚的补赎和完全的谦卑来赢得主恩宠的人，更当特别如此。如今我已读过全体宣信者的一封信，他们希望我将其告知所有同仁；在信中，他们请求确保那些在我们看来，对其犯罪后的行为已有满意交代之人，享有他们所给予的平安。此事有待我们众人商议和判断，我不敢妄下断言，将自己的决定强加于共同事务。因此，我们暂时应遵守我近日致你们的信件，我已将副本送交许多同仁，他们回信表示赞同我的决定，并说在主赐我们平安、得以聚集一处、逐一审查个案之前，不可偏离此决定。但为使你们知道我的同仁\Person{加尔多尼乌斯}致信与我以及我的回复，我已将两封信的副本附在本函中。恳请你们向我们的弟兄们宣读全部内容，使他们越来越安于忍耐，不要再因不愿听从我或福音，也不允许按其所有宣信者的信件审查各自案件，而将另一过失加于先前的过失。亲爱的弟兄，我衷心祝你们永远安好，并请记得我。问候全体弟兄。再会！\par

\SourceRecord{Translated by Robert Ernest Wallis.；Ante-Nicene Fathers；Vol. 5.；Edited by Alexander Roberts, James Donaldson, and A. Cleveland Coxe.；Buffalo, NY: Christian Literature Publishing Co.,；1886.；<http://www.newadvent.org/fathers/050617.htm>.}

% structure: p0001 source=h1 target=section reason=page-title

\section{第十八封书信}

\Emph{卡尔多纽斯致西彼廉。}\par

\Emph{辩论——当新迫害紧迫之际，某些背教者已宣认基督，故在流放之前请求平安，卡尔多纽斯就是否应赐予平安一事咨询西彼廉。}\par

卡尔多纽斯致西彼廉及同留居迦太基的同僚司铎们，安好。时势所迫，我们不可贸然赐予平安。但有必要写信给你们，因为那些曾献祭后又受试炼的人，现已成了流放者。因此在我看来，他们已弥补了昔日的罪过，因他们如今舍弃了财产和家园，悔改并追随基督。故此，费利克斯——他曾在德基慕斯手下协助执行司铎职务，与我一同被囚时十分亲近（我深知这位费利克斯的为人）——他的妻子维多利亚，以及信友路济乌斯，均遭放逐，并留下了他们的财产，现由国库保管。此外，有一位名叫波娜的妇人，被其夫拖去献祭，虽良心无犯罪之咎，但因执其手者行了祭献，她便向他们喊说：\Quote{我没有做，是你们做的！}——她也遭放逐。因此，既然这众人都请求平安，说：\Quote{我们恢复了曾失去的信德，我们悔改了，并公开宣认了基督}——尽管在我看来，他们理应获得平安——但我仍将他们的事提交你们的裁断，以免我显得轻率妄为。所以，若你们愿意我凭共同决议行事，请写信告知。请代我们问候诸位弟兄；诸位弟兄也问候你们。亲爱的弟兄们，我衷心祝愿你们永远安康。\par

\SourceRecord{Translated by Robert Ernest Wallis.；Ante-Nicene Fathers；Vol. 5.；Edited by Alexander Roberts, James Donaldson, and A. Cleveland Coxe.；Buffalo, NY: Christian Literature Publishing Co.,；1886.；<http://www.newadvent.org/fathers/050618.htm>.}

% structure: p0001 source=h1 target=section reason=page-title

\section{第十九封书信}

\Emph{\Person{西彼廉}答复\Person{卡尔多纽}。}\par

\Emph{论点——\Person{西彼廉}在此信中没有处理任何特殊问题，只是赞同\Person{卡尔多纽}的意见，即：对于那些因真正的悔改和宣认基督之名而配享平安，并因此回归基督的背教者，不应拒绝他们的平安。}\par

\Person{西彼廉}致其弟兄\Person{卡尔多纽}，安。我们收到了你的信，亲爱的弟兄，这封信极其明智，满盈诚实与信德。我们并不惊异，因为你在主圣经中娴熟且多受锻炼，你凡事都谨慎而聪明。你关于将平安赐予我们弟兄们的判断十分正确，他们藉着真正的悔改和宣认主的荣耀，已为自己恢复了那平安，因他们那之前曾定自己罪的言语而成义。既然他们已洗去了自己所有的罪，他们先前的污点，靠着主的助佑，已被一种更强的德能除去，他们不应再如同一度匍匐在地的，躺卧在魔鬼的权下；因为他们虽被放逐，被剥夺了一切财产，却已站立起来，开始与基督一同站立。而且我也希望其余的人在他们跌倒之后能悔改，并恢复到从前的情形；而且为使你知道我们是如何处理这些人的——他们带着迫切而热切的鲁莽和强求来索取平安——我已寄给你一卷书，连同五封书信，那是我写给圣职人员和教友，并也写给殉道者和宣信者的。这些信已经寄给我们许多同僚，并使他们满意；他们回信说，他们也依据公教信仰，同意我同样的意见。但愿你也将这同一件事传达给你尽可能多的同僚，好使众人之间，按照主的诫命，能保持同一行动方式和同一共识。愿你，亲爱的弟兄，常时衷心安康。\par

\SourceRecord{Translated by Robert Ernest Wallis.；Ante-Nicene Fathers；Vol. 5.；Edited by Alexander Roberts, James Donaldson, and A. Cleveland Coxe.；Buffalo, NY: Christian Literature Publishing Co.,；1886.；<http://www.newadvent.org/fathers/050619.htm>.}

% structure: p0001 source=h1 target=section reason=page-title

\section{Epistle 20}

\Emph{\Person{Celerinus} 致 \Person{Lucian}}。\par

\Emph{论点：—— \Person{Celerinus} 代其在罗马失足的姐妹们，向迦太基的宣信者们恳求和平。}\par

1. \Person{Celerinus} 向 \Person{Lucian} 致候。写这封信给你，我的主和兄弟，我既喜且忧——喜，是因我听说你为我们的主耶稣基督救主之名受试炼，并在世上的官长面前宣认了祂的名；忧，是因自从与你相处以来，我从未能收到你的信。而近来又有一重忧愁加于我：虽然你知道我们共同的兄弟 \Person{Montanus} 从监牢里由你处来我这里，你却未向我提及你的安康，也未提及任何与你相关的事。然而，这不断地发生在天主的仆人身上，尤其是那些被拣选为基督作证的人。因为我知道，如今谁都不再挂虑世上的事，而是盼望天上的冠冕。此外，我曾想，也许你忘了写信给我。因为即使我从最卑微的地位被你称为\Quote{你的}或\Quote{兄弟}，若我配得被称为 \Person{Celerinus}，那么，当我也处于如此血色的宣认中时，我仍记得年长的弟兄们，并在信中问候他们，因为旧日的情谊仍环绕我和他们。然而，我恳求主所爱的人，首先，若你已在那圣血中洗净，并为我们的主耶稣基督之名受苦，在我的信于今世找到你之前，或它们如今到达你处，愿你回信给我。愿那因你所宣认之名而加冕祂的人如此待你。因为我相信，虽在今世我们彼此不见，但在将来，我们将在基督面前拥抱。请祈求，使我也配与你们一同被加冕。\par

2. 然而，你要知道，我处在极大的患难中；仿佛你与我同在，我日夜记念你从前的爱，只有天主知道。因此我请求你成全我的愿望，并为我的姐妹（属灵）的死亡与我一同悲伤，她在这毁灭之时从基督面前跌倒了；因为她献了祭，惹怒了我们的主，这在我们是显而易见的。因她的行为，在这逾越节的喜乐之日，我日夜哭泣，在泪、麻衣和灰烬中度日，并且至今仍在度过，直到我们的主耶稣基督的帮助，以及借着你们，或那些已得了冠冕的我的主们（你将向他们请求）所显出的慈爱，来救助这可怕的船难。因为我记得你从前的爱，你会与我们其余的人一同为我们的姐妹们悲伤，你也熟知她们——即 \Person{Numeria} 和 \Person{Candida}。因她们的罪，既然视我们为兄弟，我们应当警醒。因为我相信，基督会按她们的忏悔和她们对从你们那里来的被放逐的同工们所行的善工（你们将亲自从他们那里听到她们的善工）——我相信，基督必会怜悯她们，当你们这些祂的殉道者向祂恳求之时。\par

3. 因为我听说，你已领受了穿紫袍者的职务。哦，你有福了，即使睡在地上，也实现了你一向渴望的愿望！你曾渴望为祂的名被投入狱中，如今已成事实；正如经上所载：\BibleQuote{主必照你心之所愿赐给你。} 如今你被立为天主的司鐸治理他们，而同一人，他们的服务者，也承认了这事。因此，我请求，我的主，并因我们的主耶稣基督恳求，请你将此事转告其余的同工们，你的弟兄们，我的主们，并请求他们：你们中谁先被加冕，就当赦免我们的姐妹们 \Person{Numeria} 和 \Person{Candida} 如此重大之罪。这后者，我常称她为 \Person{Etecusa}——天主是我的见证——因为她为自己献了礼物，以免献祭；但她似乎只是上到 \Place{Tria Fata}，再从那里下来。因此，我知道她没有献祭。她们的案子最近被听审后，首领们命令她们暂且保持现状，直到一位主教被指派。但，尽可能借着你们神圣的祈祷和恳求，我们信赖这些，因为你们是基督的朋友兼证人，（我们祈求）你们在这些事上都肯宽仁。\par

4. 因此，我恳求，亲爱的 \Person{Lucian} 主，请记念我，应允我的请求；愿基督赐你那神圣的冠冕，不但借着宣认，也借着圣洁，你常在其中行走，并常作圣徒们的榜样和证人，愿你向我的所有主们，你的弟兄宣信者们，述说这一切事，使他们从你得到帮助。因为，我的主和兄弟，你应当知道，为她们请求此事的不只我一人，还有 \Person{Statius} 和 \Person{Severianus}，以及所有从你们那里来到此地的宣信者们；正是这些姐妹们下到港口，接他们进城，并照料了六十五人，直至今日仍在一切事上服事他们。因为众人都与她们同在。但我不应再加重你神圣的心，因我知道你必欣然劳苦。\Person{Macharius} 与他的姐妹们 \Person{Cornelia} 和 \Person{Emerita} 问候你，因你血腥的宣认以及所有弟兄们的宣认而喜乐；还有 \Person{Saturninus}，他本人也与魔鬼角力，英勇地宣认了基督的名，并且在铁爪的酷刑下英勇宣认，也极力恳求此事。你的弟兄们 \Person{Calphurnius} 和 \Person{Maria} 以及所有圣洁的弟兄们问候你。因为你也应当知道，我已写信给我的主们，你的弟兄们，并请你屈尊为他们诵读。\par

\SourceRecord{Translated by Robert Ernest Wallis.；Ante-Nicene Fathers；Vol. 5.；Edited by Alexander Roberts, James Donaldson, and A. Cleveland Coxe.；Buffalo, NY: Christian Literature Publishing Co.,；1886.；<http://www.newadvent.org/fathers/050620.htm>.}

% structure: p0001 source=h1 target=section reason=page-title

\section{书信二十一}

\Emph{路齐安答复色勒林。}\par

\Emph{提要——路齐安同意色勒林的请求。}\par

1. 路齐安致他的主人、若堪称为在基督内的同工色勒林，问候。最亲爱的弟兄与主人，我已收到你的信，你在信中满溢善意之言，使我因这重负几乎被如此大的喜乐压倒；我欣喜地阅读，因你如此伟大的谦卑，使我得以读到那封我也渴望已久能读到的信，信中你竟愿纪念我，在信中你向我说：\Quote{若我堪当你称为弟兄，}——像我这样曾在下级官员面前战战兢兢宣认天主之名的人。因为你，因天主的旨意，当你宣认时，不仅吓退了那大蛇本身、反基督的先锋，而且以那声音和那些神圣的言辞征服了他。由此我知道你多么热爱信德，对基督的纪律多么热忱，我深知且欣喜你正积极投身其中。如今，已在殉道者中受尊敬的亲爱的弟兄，你愿以你的信使我负担过重，你在信中告诉了我们关于我们姊妹们的事。为了她们，我愿我们有可能提及她们而不至于也记起如此重大的罪行；的确，我们那时就不会像现在这样流着许多眼泪想念她们了。\par

2. 你应当知道关于我们所发生的事。当有福的殉道者保禄尚在肉身时，他召叫我并对我说：\Quote{路齐安，我在基督面前对你说：若有人在我被召去后，向你求平安，你要因我的名赐予他。}此外，我们所有主屈尊在这般磨难中召去的人，藉我们的书信，经共同同意，将平安赐予了所有人。所以，弟兄，你看出我如何部分地执行了保禄所吩咐的，也正如我们在这磨难中所一致决定的。在这磨难中，我们按皇帝的谕令被判处饥渴而死，并被关在两间牢房里，为要藉饥渴削弱我们。此外，因酷刑所致的火焰如此难忍，无人能承受。但现在我们已达到了光明本身。因此，亲爱的弟兄，请向努梅里亚和坎迪达致意，她们将按照保禄的训令得享平安；还有以下列出的其他殉道者：在伪证者地牢中的巴苏斯，在酷刑中的马帕利库斯，在监狱中的福尔图尼奥，受刑后的保禄，福尔图纳塔，维克托里努斯，维克托，赫伦尼乌斯，尤利娅，马尔提亚尔，以及阿里斯托，他们藉天主的旨意，在狱中死于饥饿。几天后你将听说我也成为他们的同伴。因为现在从我被再次关押之日到给你写信之日，已有八天。在这八天之前，连续五天我按定量只得到一小块面包和水。因此，弟兄，正如在这里，既然主已开始将平安赐予教会本身，按照保禄的训令和我们的论述，在主教前陈述案情并作出宣认之后，我请求不仅这些人可得平安，而且所有你知道与我们心意相近的人也都可以。\par

3. 我所有的同工问候你。请代我问候在你那里的主的宣信者们，你已告知了他们的名字，其中有撒图尔尼努斯和他的同伴——他也是我的同工——以及玛里斯、科莱克塔、埃梅里塔、卡尔普尔尼乌斯和玛利亚、萨比娜、斯佩西娜，还有众姊妹雅努阿里亚、达提瓦、多纳塔。我们问候撒图鲁斯及其家人、巴西安努斯和全体圣职人员、乌拉尼乌斯、阿莱克修斯、昆泰努斯、科洛尼卡，以及所有我没有写出名字的人，因为我已经疲倦了。所以，请他们宽恕我。我衷心向你和阿莱克修斯、格图利库斯、兑换银钱者及众姊妹告别。我的姊妹雅努阿里亚和索菲亚——我把她们托付给你——向你问候。\par

\SourceRecord{Translated by Robert Ernest Wallis.；Ante-Nicene Fathers；Vol. 5.；Edited by Alexander Roberts, James Donaldson, and A. Cleveland Coxe.；Buffalo, NY: Christian Literature Publishing Co.,；1886.；<http://www.newadvent.org/fathers/050621.htm>.}

% structure: p0001 source=h1 target=section reason=page-title

\section{\WorkTitle{第二十二封书信}}

\Emph{致住在罗马的圣职人员，论及许多宣信者，以及宣信者\Person{路西安}的冒进和宣信者\Person{塞莱里努斯}的谦逊。}\par

\Emph{摘要。——在这封信中，\Person{西彼廉}告知罗马圣职人员关于失足者要求恢复和平的煽动性请求，以及\Person{路西安}的冒进。为使罗马圣职人员更了解这些事，\Person{西彼廉}确保不仅他自己的信件，而且\Person{塞莱里努斯}和\Person{路西安}的信件，也送达给他们。}\par

1. \Person{西彼廉}致住在罗马的司铎和执事，他的兄弟们，安好。我先前写给你们的信函中，已说明我所行之事，并略述了我的纪律和勤勉。但随后又发生了一件事，也不应隐瞒你们。因为我们的兄弟\Person{路西安}，他自己也是宣信者之一，确实在信德上热诚，在美德上坚强，但在阅读主的话上却不甚稳固。他做了一些事，一时自充为愚人的权威，以致以\Person{保禄}之名由他亲手写成的证明书不加区别地分发给了许多人；而殉道者\Person{玛帕利库斯}则谨慎谦逊，铭记法律和纪律，没有写过任何违背福音的信，只是出于对跌倒的母亲的家庭情感，才命令给予她和平。\Person{撒图尼努斯}在受酷刑后，仍留在狱中，也没有发出这类信函。但\Person{路西安}不仅在\Person{保禄}仍在狱中时以他的名义到处分发由他自己亲手写成的证明书，甚至在\Person{保禄}去世后仍继续以他的名义做同样的事，说这是\Person{保禄}命令他的，却不知道他必须服从主而非同仆。同样，以受酷刑的青年\Person{奥勒留}的名义，也由同一\Person{路西安}亲手写了多份证明书，因为\Person{奥勒留}自己不会写字。\par

2. 为在一定程度上制止这种做法，我写信给他们，并附在前一封信中寄给你们。在信中，我恳切要求并劝勉他们考虑主和福音的法律。但在我把信寄给他们，以期事情得以更温和适度地处理之后，同一\Person{路西安}以全体宣信者的名义写了一封信，信中几乎每一信德的纽带、对天主的敬畏、主的命令以及福音的神圣和真诚都被解散了。因为他以全体名义写道，他们已给予所有人和平，并希望这道谕令通过我传达给其他主教。我将那封信的副本寄给了你们。其中确实加了一句：\Quote{那些在犯罪后已做了满意补赎的人}——这件事更激起对我的憎恨，因为当我开始听取每个人的案情并予以查究时，我似乎拒绝了许多人，而他们现在都夸耀自己从殉道者和宣信者那里得到了（和平）。\par

3. 最后，这种煽动性的做法已开始显现；因为在我们行省的一些城市里，群众已攻击他们的管理者，并迫使他们立即给予和平，而群众都叫喊说他们早已从殉道者和宣信者那里得到了和平。这些管理者因恐惧而屈服，几乎没有力量抵抗他们，无论是凭心灵的刚毅还是信德的力量。此外，在我们这里，一些不安分的人，过去我很难管束他们，并拖延到我回来，他们被这封信煽动得像火把一样，开始更加激烈，并强索已经给予他们的和平。我把我关于这些事写给圣职人员的信的副本寄给你们，此外，还有我的同僚\Person{卡尔多尼乌斯}凭其正直和忠诚所写的信，以及我的回信。我都寄给你们阅读。另外，我也寄给你们善良而勇敢的宣信者\Person{塞莱里努斯}写给同为宣信者的\Person{路西安}的信的副本，以及\Person{路西安}的回信；好让你们知道我所有方面的工作和勤勉，并了解真相本身：宣信者\Person{塞莱里努斯}是多么温和谨慎，在谦卑和对我们信德的敬畏中是多么恭敬；而\Person{路西安}，如我所说，在对主的话的理解上不太熟练，并且由于他的轻率，因他对我的敬虔处理所引起的不悦，而造成了损害。因为主曾说过：万民要在圣父、圣子、圣神之名内受洗，他们过去的罪要在圣洗中除去；但此人，对训诫和法律一无所知，竟命令以\Person{保禄}之名给予和平并除去罪过；他还说这是\Person{保禄}命令他的，正如你们将在同一\Person{路西安}写给\Person{塞莱里努斯}的信中所见，他很少考虑到：不是殉道者创造了福音，而是殉道者被福音所造就；因为\Person{保禄}宗徒，主所召选的器皿，也在他的书信中写道：\BibleQuote{我惊奇你们竟这么快就离开了那藉基督的恩宠召叫你们的那一位，而归向另一个福音；这福音并不是另一个，只是有些人扰乱你们，并企图改变基督的福音。但即使是我们，或是一位从天上来的天使，若向你们宣讲我们所宣讲的福音以外的任何福音，当受诅咒。正如我们先前说过的，现在我再重复：若有人向你们宣讲你们所领受的福音以外的任何福音，当受诅咒。}\par

4. 然而，我收到你写给我神职人员的信，来得正是时候；同样，那些蒙福的宣信者，\Person{莫伊塞斯}、\Person{马克西穆斯}、\Person{尼各斯特拉图斯}以及其他，写给\Person{撒图尔尼努斯}、\Person{奥勒里乌斯}和其他人的信，也来得正好，这些信中包含着福音的圆满力量与主法律的坚定纪律。当我在此地辛劳并竭尽信德之力对抗恶意的攻击时，你的话语极大地帮助了我，以致我的工作从上头被缩短，以至于在我最后寄给你的信件到达你之前，你就向我声明，按照福音的法律，你的判断也强烈而一致地与我相合。亲爱的、深切的弟兄们，我愿你们永远衷心珍重。\par

\SourceRecord{Translated by Robert Ernest Wallis.；Ante-Nicene Fathers；Vol. 5.；Edited by Alexander Roberts, James Donaldson, and A. Cleveland Coxe.；Buffalo, NY: Christian Literature Publishing Co.,；1886.；<http://www.newadvent.org/fathers/050622.htm>.}

% structure: p0001 source=h1 target=section reason=page-title

\section{\WorkTitle{第二十三封书信}}

\Emph{致圣职界，论寄往罗马的信函，以及关于任命\Person{撒图鲁斯}为读经员、\Person{奥普塔图斯}为副执事之事。公元二五〇年。}\par

\Emph{论旨：圣职界藉此信得知\Person{撒图鲁斯}与\Person{奥普塔图斯}的圣职任命，以及\Person{西彼廉}寄往罗马的信函。}\par

\Person{西彼廉}致他的弟兄司铎与执事，问候。亲爱的弟兄，为使你们知晓那写给某人和某人所答复的，我已将每封信的副本寄给你们，我相信我的回复不会使你们不悦。但我应当在此信中将此事告知你们：因一极为紧急的理由，我已寄信给住在城里的圣职人员。既然我应当通过圣职人员写信，而我知道我们中的许多人都不在，留在那里的少数人几乎不足以承担日常职责的服事，就有必要任命一些新人，好能派遣出去。那么，你们要知道，我已任命\Person{撒图鲁斯}为读经员，并任命宣信者\Person{奥普塔图斯}为副执事；我们已依共同商议，将二人列为圣职之类，因我们曾在逾越节再三委托\Person{撒图鲁斯}诵读，并当我们与教师司铎一起仔细考验读经员时，从读经员中任命\Person{奥普塔图斯}为听道者的教师——首先察验他们身上是否一切合宜，即那些在准备担任圣职的人身上应具备的一切。因此，在你们缺席时，我并未做什么新事；而是那依我们众人的共同商议已经开始的事，基于紧急需要而完成了。亲爱的弟兄，我祝你们永远由衷安康，并记念我。再会。\par

\SourceRecord{Translated by Robert Ernest Wallis.；Ante-Nicene Fathers；Vol. 5.；Edited by Alexander Roberts, James Donaldson, and A. Cleveland Coxe.；Buffalo, NY: Christian Literature Publishing Co.,；1886.；<http://www.newadvent.org/fathers/050623.htm>.}

% structure: p0001 source=h1 target=section reason=page-title

\section{第二十四封书信}

\Emph{致莫依塞、玛克西穆及其余宣信者。}\par

\Emph{论点：此信为祝贺罗马宣信者而作。}\par

1. 西彼廉向莫依塞、玛克西穆二位司铎，并致其余最亲爱的弟兄宣信者，安。英勇有福的弟兄们，我早已从传闻中得知你们信德与勇毅的荣光，心中大喜，并深切祝贺你们，因我主耶稣基督至高的屈尊俯就为你们预备了因宣认祂的圣名而得的花冠。因为你们已在我们这时代的战斗中成为首领和先锋，竖起了天战之旗；你们开始了属灵的竞赛，这是天主决意藉你们的勇毅如今要进行的；你们以不可动摇的力量和不屈的坚定，击破了初起战争的第一波冲击。由此产生了战斗的喜乐开端；由此开始了胜利的吉祥预兆。在这里，殉道者们藉着酷刑得以完成。但那在斗争中先行、成为弟兄们勇毅榜样的人，在光荣上与殉道者同等。因此你们将亲手编织的花冠交付给我们，并以救恩之杯为兄弟作保。\par

2. 在这荣耀的宣认开端和得胜战争的预兆之上，又加上了纪律的持守，这是我从你们最近写给与你们一同在宣认中归于主的同僚们的书信的严厉告诫中所观察到的，你们殷切劝诫，要坚定而严格地遵守那一次交付给我们的福音神圣诫命和生活的律令。看哪，你们荣耀的另一崇高境界；看哪，与宣认并行的，是天主喜悦的双重功勋——以坚定的脚步站立，并在这斗争中藉你们信德的力量驱逐那些试图破坏福音、将不虔之手伸向颠覆主之诫命的人：——先前已显示出勇气的迹象，如今又提供了生活的教训。主在复活后派遣宗徒时，吩咐他们说：\Quote{天上地下的一切权柄都赐给了我。所以你们要去，使万民成为门徒，因父、及子、及圣神之名给他们授洗，教训他们遵守我所吩咐你们的一切。}\BibleRef{玛 28:18} 宗徒若望记念这嘱咐，随后在他的书信中规定：\Quote{我们若遵守祂的诫命，由此就知道我们认识祂。那说\InnerQuote{我认识祂}却不遵守祂诫命的，是撒谎者，真理不在他内。}\BibleRef{若一 2:3} 你们推动这些诫命的持守；你们遵守神圣属天的命令。这就是做主的宣信者；这就是做基督的殉道者——在所有邪恶中保持自己宣认的坚定并稳固无虞。因为要为主成为殉道者，却试图推翻主的诫命；要用祂赐予你的屈尊俯就反过来攻击祂——仿佛用从祂领受的武器成为叛逆——这就是想宣认基督，却又否认基督的福音。因此，最英勇最忠信的弟兄们，我为你们欢喜；如同我祝贺那里因勇力光荣而受尊崇的殉道者一样，我也同样祝贺你们获得主之纪律的冠冕。主以多种多样的慷慨施予祂的屈尊俯就。祂以丰盛的种类分派了善战士兵的赞美和他们属灵的光荣。我们也共享你们的尊荣；我们算你们的荣耀为我们的荣耀，因为我们的时代被如此幸福所照亮，使我们这个日子得以看见经考验的天主仆人和基督的士兵戴上冠冕。最英勇有福的弟兄们，我衷心祝你们永远安好，并请记念我。\par

\SourceRecord{Translated by Robert Ernest Wallis.；Ante-Nicene Fathers；Vol. 5.；Edited by Alexander Roberts, James Donaldson, and A. Cleveland Coxe.；Buffalo, NY: Christian Literature Publishing Co.,；1886.；<http://www.newadvent.org/fathers/050624.htm>.}

% structure: p0001 source=h1 target=section reason=page-title

\section{第二十五封书信}

\Emph{摩西、马克西穆斯、尼科斯特拉图斯及其他宣信者答复上述书信。公元250年。}\par

\Emph{论旨。—— 他们感激地承认罗马宣信者从西彼廉的书信中所获得的安慰。殉道不是惩罚，而是幸福。福音之言是燃起信德的火种。关于背教者，他们同意西彼廉的判断。}\par

致迦太基教会主教凯基利乌斯·西彼廉：摩西、马克西穆斯，司铎；尼科斯特拉图斯、鲁菲努斯，执事；以及其他在天主圣父、祂的圣子耶稣基督我们的主、并圣神内，持守信德真理的宣信者，致候。兄弟，我们身处各种多重忧伤之中，由于目前几乎全世界许多兄弟的荒凉，这主要的安慰已达于我们：我们因收到你的书信而振作，并从我们忧伤心灵之悲痛中获得了些许缓解。由此我们已能看出，神圣眷顾的恩宠愿我们如此长久地被囚禁在牢狱的锁链中，或许别无他因，只为使我们受你的书信教导并更热切地激励后，能以更诚恳的意志达到预定的冠冕。因为你的书信照耀我们，如同暴风雨中的平静，如同汹涌大海中期盼的安宁，如同劳作中的安息，如同危险与痛苦中的健康，如同最浓密黑暗中的明亮炽光。因此我们以渴慕的心灵畅饮它，以饥饿的渴望领受它；我们欢喜地发现，由此我们得到充分的饱足和坚强，足以迎战仇敌。主必会酬报你的那份爱，并会归还你这如此善行应得的果实；因为劝勉者并不比受难者更少配得冠冕的赏报；教导者并不比实行者更少配得赞美；警告者并不比战斗者更少受尊敬；只不过有时功荣的份量更多地归于训练者，而非那显为可教之学徒的人；因为后者若没有前者的教导，也许就不会拥有他所实行的。\par

2. 因此，我们再要说，兄弟西彼廉，我们获得了极大的喜乐、极大的安慰、极大的振作，尤其因你以荣耀而应得的赞美描述了殉道者的荣耀——我不说是死亡，而是不朽。因为这样的离去应当用如此的话语来宣报，以便所叙述的事情能按它们发生的方式被述说。如此，从你的书信，我们看见了殉道者那些荣耀的凯旋；并仿佛以我们的眼睛追随他们升天，凝视他们坐在天使、天上的掌权者和宰制者中间。此外，我们仿佛以耳朵感知到主在圣父面前赐予他们那应许的证言。正是这，每日提升我们的精神，并激励我们追随如此尊严的踪迹。\par

3. 因为从天的屈尊，有什么更荣耀、更幸福的事能临到任何人，比在死亡本身中，在行刑者面前宣认主天主？比在世界权力的狂怒、多样、精妙的酷刑中，即使身体被拉扯、撕裂、切割成碎片，仍以自由的精神（虽将离去）宣认基督是天主子？比撇下世界，升到天上？比离弃世人，站在天使中间？比冲破一切世俗障碍，已在天主眼前自由站立？比毫无迟延地享受天国？比因基督的名成为基督苦难的同伴？比因天主的屈尊成为自己审判官的审判官？比因宣认祂的名而保有纯洁的良心？比拒绝服从反对信德的人类亵圣法律？比以公开的证言为真理作证？比借着死亡，征服那为众人所畏惧的死亡本身？比借着死亡本身，获得不朽？比当被撕成碎片，被一切残忍刑具折磨时，以折磨本身战胜折磨？比以心灵的力量与残破身体的一切剧痛搏斗？比不因自己流血而颤抖？比在拥有信德忍受惩罚后，开始爱自己的惩罚？比认为不离开生命是对生命的伤害？\par

4. 因为我们的主，如同以祂福音的号角，激励我们参加这场战斗，祂说：\BibleQuote{爱父亲或母亲超过我的，不配属于我；爱自己灵魂超过我的，不配属于我。不背起自己的十字架跟随我的，也不配属于我。}\BibleRef{玛 10:37} 又说：\BibleQuote{为义而受迫害的人是有福的，因为天国是他们的。人若迫害你们、仇恨你们，你们是有福的。你们欢喜踊跃罢！因为你们的祖先也这样迫害过你们以前的先知。}\BibleRef{玛 5:10} 又说：因为你们要站在君王和权贵面前，兄弟要把兄弟置于死地，父亲要把儿子，那坚持到底的，必要得救；并且\BibleQuote{胜利的，我要赐他坐在我宝座上，就如我得了胜，在我父的宝座上坐下一样。} 此外，宗徒说：\BibleQuote{谁能使我们与基督的爱隔绝？是困苦吗？是窘迫吗？是迫害吗？是饥饿吗？是赤身露体吗？是危险吗？是刀剑吗？（正如经上记载：为了你，我们终日被置于死地，被看作待宰的羊。）然而，靠着那爱我们的，我们在这一切事上大获全胜。}\BibleRef{罗 8:35}\par

5. 当我们读到福音中所收集的这些以及类似的事情，并感受到，就好像有火炬放在我们下面，藉着主的话点燃我们的信德，我们不仅不惧怕，反而甚至激怒真理的敌人；我们藉着不向它们屈服这一事实，已经战胜了天主的反对者，并征服了它们反对真理的邪恶律法。虽然我们尚未倾流我们的血，但我们已准备好倾流。谁也不可认为我们延迟离世是出于任何慈悲；因为这种延迟阻碍我们，妨碍我们的光荣，推迟天堂，扣留对天主的荣耀瞻仰。因为在这种竞赛中，并且在信德在交锋中奋力争战的那类竞赛中，藉着延迟而推迟殉道者，并非真正的慈悲。因此，亲爱的\Person{西彼廉}，请恳求主，愿祂以其仁慈，每日更加丰富地、更加准备充足地用更大的丰盈和准备武装并装饰我们每一位，并以祂大能的力量坚定并坚强我们；并且，作为一位好统帅，终将带领祂的士兵们——他们迄今已在我们的监狱营中被训练和考验——出来，进入为他们摆在前面的战场。愿祂向我们展示神圣的武器，那些不败的兵器——正义的胸甲，它从未惯于被破碎——信德的盾牌，它不能被刺穿——救恩的头盔，它不能被打碎——圣神的剑，它从未惯于受损伤。因为我们应将这些事情交给谁，请他为我们祈求，胜过交给我们如此可敬的主教呢？如同预定的祭品向司铎请求援助一般。\par

6. 请看我们的另一喜乐：您虽因时代的状况暂时与弟兄们分离，却常在主教职务的职责中以书信多次坚固了宣信者；您曾从自己的正义所得提供必要的供应；您总在一切事上以某种方式显示自己亲临；在您的职责中，您从未像逃兵一样退缩。但我们不能缄默那更强烈地激励我们获得更大喜乐的事，而必须用我们全部的见证声音来描述。因为我们注意到，您以适当的责斥恰当地责备了那些人：他们不顾自己的罪，在您不在时以仓促和急切的心愿从司铎那里强索了和平；以及那些不尊重福音、竟以亵渎的轻率将主的\Emph{圣物}赐给狗，将珍珠给猪的人；虽然一个大罪——它以难以置信的破坏性几乎蔓延到整个大地——只应，如您自己所写，谨慎而有节制地处理，并经过所有主教、司铎、执事、宣信者、甚至那些坚定站立的平信徒的咨议，正如您在您的书信中亲自证明的一样；这样，当我们不合时宜地想要修复废墟时，我们就不至于看起来造成了另一更大破坏，因为如果如此轻易地给予罪人赦免，神圣之言还剩下什么呢？当然，他们的心灵要得到鼓舞和滋养，直到他们成熟的时候，并且要藉着圣经教导他们犯下了何等重大和超越的罪。也不要因为他们是多数而受鼓舞，而更要因为他们是不少而受约束。不知羞耻的数目从未惯于在减轻罪行方面有分量；然而，羞愧、谦逊、忍耐、纪律、谦卑与服从——等待他人对自己的判断，并承受他人对自己判断的判决——这正是证明补赎的；这正是使深伤口愈合的；这正是使跌倒精神的废墟升起、并修复它们，平息并抑制他们狂怒罪恶的炽热蒸汽的。因为医生不会将健康身体的食物给予病人，唯恐不合时宜的营养不仅不能抑制，反而刺激狂怒疾病的力量——也就是说，唯恐那原可通过禁食更快消减的，因着不耐烦，反而因消化不良而延长了。\par

因此，被不虔敬祭献所玷污的手，必须以善工来洁净；被诅咒食物所玷污的可怜嘴唇，必须以真诚忏悔的言语来清洁；灵魂必须在信者之心的隐秘处被更新和祝圣。愿常听到忏悔者不断的呻吟；愿从眼中流下信实的泪水，不只一次，而是再三再四，好使那些曾邪恶注视偶像的眼睛，以悦乐天主的泪水，洗去它们所做的不法之事。对于疾病，除了忍耐，别无所需：疲惫虚弱的人与痛苦角力；他们若因容忍而能克服痛苦，最终才能盼望康复；因为医生过于仓促造成的伤疤是不可靠的；若不用药从极慢处忠实使用，任何小小的意外都会毁掉治疗。火焰很快又会燃起大火，除非整堆火的材料连最细微的火星都被熄灭；好使这类人公正地知道，甚至连他们自己从这拖延中反而更受益，而且通过必要的延缓应用了更可靠的疗法。此外，哪里能说那些承认基督的人被囚禁在肮脏的监狱里，而否认祂的人却毫无信德之险？哪里能说他们为天主的缘故被枷锁束缚，而那些没有保持对天主的承认者却未被剥夺共融？哪里能说被囚的殉道者献出了光荣的生命，而那些背弃了信德的人却没有感觉到危险和罪过的严重性？但如果他们显得过于急躁，以无法容忍的急切要求共融，他们便是以鲁莽无羁的舌头徒然说出那些怨怼和敌意的指责，这些指责在真理面前毫无效力，因为他们本可以凭自己的权利保留那如今因一种他们自愿寻求的必要性而被迫乞求的东西。因为那曾能承认基督的信德，也本可以被基督保存在共融中。我们视您，蒙福且最荣耀的父，在主内永远衷心告别；并请记念我们。\par

\SourceRecord{Translated by Robert Ernest Wallis.；Ante-Nicene Fathers；Vol. 5.；Edited by Alexander Roberts, James Donaldson, and A. Cleveland Coxe.；Buffalo, NY: Christian Literature Publishing Co.,；1886.；<http://www.newadvent.org/fathers/050625.htm>.}

% structure: p0001 source=h1 target=section reason=page-title

\section{Epistle 26}

\Person{Cyprian}致失足者。\par

\Emph{本书信的论旨见下第二十七封信。他说：\Quote{他们写信给我，不是请求赐予平安，而是声称平安已归他们所有，因为他们说\Person{保禄}已赐予众人平安；正如你将在他们信中读到的那样，我已将那信的副本连同我对他们的简短答复寄给你。}但失足者所写的信（他答复的那封）现已佚失。}\par

1. 我们的主，我们应当遵守其诫命与训诲，在福音中论及主教的尊荣与祂教会的秩序时，对\Person{伯多禄}说：\BibleQuote{我告诉你，你是\Person{伯多禄}，在这磐石上，我要建立我的教会，地狱的门不能胜过她。我要把天国的钥匙交给你；凡你在地上捆绑的，在天上也要被捆绑；凡你在地上释放的，在天上也要被释放。}从此，经由时代变迁与继承，主教的序列与教会的计划得以延续；因此教会建立在主教之上，教会的一切行动都由这些统治者管理。既然这事基于神圣的法律，我感到惊奇的是，有些人竟胆大妄为地选择以教会的名义写信给我，仿佛他们是以教会的名义写的；而教会是建立在主教与圣职人员和所有\Emph{持守信德}的人身上。因为天主的仁慈和祂不受约束的大能，决不容许失足者的群体被称为教会；因为经上记载：\BibleQuote{天主不是死人的天主，而是活人的天主。}\BibleRef{玛 22:32} 我们确实希望所有人都能获得生命；我们祈祷，藉着我们的恳求与哀叹，他们能恢复原状。但如果某些失足者声称自己是教会，如果教会是在他们中间、在他们内，那么除了请求这些人恩准我们进入教会之外，还有什么剩下呢？因此他们应当谦卑、安静、温良，如同那些因记念自己的罪过而应平息天主的人，而不应以教会的名义写信，因为他们更应意识到自己是在写信给教会。\par

2. 但有些失足者最近写信给我，他们谦卑、温良、战战兢兢、敬畏天主，并且一向在教会中荣耀而慷慨地工作，从未向主夸耀自己的劳动，因为知道祂说过：\BibleQuote{你们做完了一切所吩咐的，只当说：我们是无用的仆人，所做的本是我们应当做的。}\BibleRef{路 17:10} 他们想起这些事，虽然从殉道者那里得到了证明书，但仍为使他们的补赎能被主接纳，这些人在恳求中写信给我，承认自己的罪过，真心悔改，并不轻率或强求地急于得到平安；而是等待我的到来，说甚至平安本身，若在我亲临的时候领受，对他们也会更加甘甜。主是我的见证，祂曾屈尊指示这样的仆人和这类仆人应得祂怎样的慈恩，我多么为他们庆幸啊！这些信我最近收到，而现在读到你们写的信大不相同，我请求你们将自己的意愿分辨清楚；无论你们是谁寄了这封信，请将你们的名字添加到证明书上，并将附有各人名字的证明书传达给我。因为我必须先知道我要答复谁；然后我将按照我的地位与行为的平庸，对你们所写的每件事作出回应。我祝愿亲爱的弟兄们，永远在主内衷心平安，并按照主的纪律安静而安宁地生活。祝好。\par

\SourceRecord{Translated by Robert Ernest Wallis.；Ante-Nicene Fathers；Vol. 5.；Edited by Alexander Roberts, James Donaldson, and A. Cleveland Coxe.；Buffalo, NY: Christian Literature Publishing Co.,；1886.；<http://www.newadvent.org/fathers/050626.htm>.}

% structure: p0001 source=h1 target=section reason=page-title

\section{书信第二十七}

\Emph{致司铎与执事。}\par

\Emph{主旨——本书信的主旨与前一封颇相符，看来正是他在下列信中提及的那一封；因为他赞扬其圣职人员将\Place{狄达}的\Person{加约}及其执事（二人因轻率地与背教者共融）逐出共融，并劝勉他们对其他人也如此行。}\par

1. 西彼廉致司铎与执事，他的弟兄，祝安。你们行事正直且合于纪律，亲爱的弟兄们，因为又我在场同僚的建议，你们决定不与\Place{狄达}的司铎\Person{加约}及其执事共融；他们因与背教者共融并献上祭品，常被逮住他们的邪恶错误中；且他们一再（如你们所写信告知我的），当被我的同僚警告不要如此行时，仍固执于他们的狂妄和鲁莽，也欺骗了我们百姓中的某些弟兄，对他们的益处我们愿以全谦卑来照料，对他们的得救我们关心，不是出于虚伪的奉承，而是出于真诚的信德，使他们能以真实的补赎、呻吟和悲伤恳求主，因为经上记载：\BibleQuote{你要回想你是从哪里坠落的，并要悔改。}\BibleRef{默 2:5}再者，神圣的圣经说：\BibleQuote{主这样说：当你们归依并哀悼时，你们将得救，并知道你们曾在哪里。}\BibleRef{依 30:15}\par

2. 然而那些如何能哀悼和悔改呢？那些司铎轻率地认为可与他们共融，从而阻挠了他们的呻吟和眼泪，却不知道经上记载：\BibleQuote{他们称你们有福的，使你们迷失，毁灭你们脚下的路。}\BibleRef{依 3:12}自然，我们健全而真实的忠告无法成功，因为有害的谄媚和奉承阻碍了救人的真理，而背教者受伤不健康的灵魂遭受肉体患病者常遭受的：即当他们拒绝健康的食物和有益的饮品，视之为苦和厌恶，而渴望那些当下似乎令他们愉悦和甘甜之物时，他们因自己的鲁莽和不节制招致自身的灾祸和死亡。有技巧的医生的真实疗法也无法有益于他们的安全，因为甜蜜的诱惑以其魅力欺骗着他们。\par

3. 所以，你们要按我的书信，忠信而健全地商议此事，不要背离更好的忠告；并留心将这些书信也读给我的同僚们，若有任何在场，或若有任何会到你们那里；使我们在同心合意中，维持一个健康的计划，以安抚和医治背教者的创伤，并打算在主仁慈恩准下，当我们开始聚集时，对所有人做非常充分的处理。与此同时，若有任何不受约束和急躁的人，无论是我方的司铎或执事或外人，敢在我们的法令之前与背教者共融，让他被逐出我们的共融，并在主许可下，当我们再次聚集时，在我们所有人面前为自己的鲁莽辩护。此外，你们希望我回复我对副执事\Person{斐禄默诺}和\Person{福尔图纳图斯}，以及辅祭\Person{法沃里努斯}的看法，他们在考验时期中途退去，现已返回。关于此事，我不能自任唯一的审判者，因为许多圣职人员仍缺席，甚至至今仍未考虑应返回其岗位；而且每个人的案件必须分别考虑和充分调查，不仅与我的同僚，也与全体百姓自己。因为一件事将来可能作为教会服务人员的榜样，必须以审慎的考量来权衡和判决。与此同时，让他们只暂停每月的分配，不是要显得被剥夺了教会的职务，而是等到所有事情处于健全状态，保留至我到来。亲爱的弟兄们，我衷心祝你们常安。问候所有弟兄，再会。\par

\SourceRecord{Translated by Robert Ernest Wallis.；Ante-Nicene Fathers；Vol. 5.；Edited by Alexander Roberts, James Donaldson, and A. Cleveland Coxe.；Buffalo, NY: Christian Literature Publishing Co.,；1886.；<http://www.newadvent.org/fathers/050627.htm>.}

% structure: p0001 source=h1 target=section reason=page-title

\section{书信第二十八封}

\Emph{致驻罗马的众司铎与执事。}\par

\Emph{论点——罗马圣职人员获悉背教者要求平安之鲁莽。}\par

\Person{西彼廉}致驻罗马的诸司铎与执事，他的弟兄们，安。我们共同的爱与事理要求，亲爱的弟兄们，我不应将我们中间所发生的事有任何隐瞒，不让你们知道，好使我们能为教会的管理益处而有一致的计划。因为在我写了那封由读经员\Person{撒图鲁}与副执事\Person{奥普塔特}送交给你们的信之后，某些背教者的共同鲁莽（他们拒绝悔改并向天主做补赎）给我写信，不是请求将平安赐予他们，而是声称平安已赐予他们；因为他们说\Person{保禄}已将平安给了所有人，正如你们将在他们的信中读到的那样，我给你们寄去一份副本，同时也附上我暂时对他们所作的简短回复。但为使你们也知道我后来写给圣职界的是怎样一封信，我还寄给你们一份副本。但若他们的鲁莽终究不能被我的或你们的信函所抑制，也不听从有益劝告，我将按照主依其福音所命令的那样采取行动。我祝你们，亲爱的弟兄们，永远在衷心问候中。告别。\par

\SourceRecord{Translated by Robert Ernest Wallis.；Ante-Nicene Fathers；Vol. 5.；Edited by Alexander Roberts, James Donaldson, and A. Cleveland Coxe.；Buffalo, NY: Christian Literature Publishing Co.,；1886.；<http://www.newadvent.org/fathers/050628.htm>.}

% structure: p0001 source=h1 target=section reason=page-title

\section{\WorkTitle{第二十九封书信}}

\Emph{\Place{罗马}的诸位长老与执事，致\Person{西彼廉}。}\par

\Emph{论据——\Place{罗马}教会声明其对背教者的判决与迦太基法令一致。任何对背教者施予的宽恕，必须符合福音的法律。由宣信者赐予的平安，仅基于恩宠与善意，这从背教者被提交给主教这一事实可见一斑。\Person{费利奇西穆斯}对和平的煽动性要求，应归因于派系之争。}\par

\Place{罗马}的诸位长老与执事，致父亲\Person{西彼廉}安。亲爱的弟兄，当我们仔细阅读你托副助祭\Person{福图纳图斯}送来的信函时，我们陷入双重悲伤，并为双重忧虑所困扰：一则在如此迫害的困境中，你得不到任何安息；二则背教弟兄们不理智的鲁莽行事，竟至发出危险的狂妄言论。虽然这些事令我们及我们的精神深感痛苦，但你依据合宜的纪律所展现的严厉与严格，缓和了我们沉重的忧伤；因为你正确地制止了一些人的恶意，并通过劝勉悔改，指明了合法的救恩之道。他们竟会如此急切地走到这一步，我们确实相当惊讶：在如此巨大的过度罪恶中，他们竟以如此迫切、不合时宜且痛苦的方式，不是请求，而是要求平安，甚至声称他们在天上已经拥有平安。若他们已拥有平安，为何还要请求他们所拥有的？但若正因为他们请求，证明他们并未拥有，为何他们不接受那些他们以为应当向其请求平安之人的审判？若他们以为从别处获得了共融的特权，就让他们试着将其与福音相比，这样，如果它与福音律法不相违背，它最终才能充分有益于他们。但若它似乎建立在与福音真理相悖的基础上，它又如何能给予福音共融呢？因为既然每一项特权都考虑到团体的特权，且其前提是与它所寻求联合的那一位的旨意不相冲突；那么，由于它与其所寻求联合者的旨意相疏离，它必然失去团体的宽恕与特权。\par

2. 那么，让他们看看他们在这件事上试图做什么。因为若他们说福音确立了一项法令，而殉道者确立了另一项，那么，他们使殉道者与福音对立，就会处于双重危险中。因为一方面，福音的尊威若能被另一项新法令所征服，便似乎已破碎倾倒；另一方面，若发现殉道者并非通过遵守那使他们成为殉道者的福音而获得荣耀的宣信冠冕，那么这冠冕便从他们头上被夺去。因此，没有人比那努力从福音中获得殉道者名号的人，更有理由谨慎地不判定任何与福音相反的事。此外，我们还想知道这一点：如果殉道者成为殉道者，无非是因为他们不献祭，从而保持教会的平安，甚至流自己的血，以免在酷刑的折磨下被击败，失去平安而丧失救恩；那么，他们认为那若献祭便以为无法得到的救恩，应赐给那些据说已献祭的人，其理由何在？尽管他们应当维护他人遵守那似乎他们自己也摆在眼前的律法。在此事中，我们观察到，他们恰恰提出了那他们认为有利于他们却反而反对他们自身理由的东西。因为若殉道者认为应赐予他们平安，为何他们自己不去赐予？为何他们认为，如他们自己所说，应将此事提交给主教？因为那命令某事去做的人，必定能做他所命令的事。但据我们理解，确切地说，正如事情本身所言明并宣告的，最神圣的殉道者认为必须在两方面保持适度的谦逊与真理。因为当许多人催促他们时，他们将这些人交给主教，认为这样既能顾及自己的谦逊，不再被打扰，又能通过不与他们共融，判断福音律法的纯洁应保持不受损害。\par

3. 但凭借你的爱德，兄弟，永不停止安抚失足者的心灵，向迷途者提供真理的良药，虽然病人的性情常常拒绝治疗者的友善帮助。这失足者的伤口尚新，疮伤仍在肿起；因此我们确信，当经过更长时间，他们的急切消退时，他们就会喜爱那将他们指向可靠良药的迟延本身；只要没有人为了自己的危险而武装他们，悖逆地教导他们，替他们要求过早共融的致命毒药，而非迟延的有益疗法。因为我们不相信，没有某些人的煽动，他们会全都这样鲁莽地为自己要求平安。我们知道迦太基教会的信德，我们知道她的培育，我们知道她的谦卑；因此我们也感到惊奇，我们竟会在信中看到某些针对你有些粗鲁的暗示，尽管我们常常在许多相互之爱的例证中了解到你们彼此的友爱和爱德。因此，现在他们应当悔改自己的过犯，他们应证明自己对失足的忧伤，他们应显示谦逊，表现谦卑，展示一些羞愧，通过他们的顺服，为自己恳求天主的仁慈，并通过给天主的司铎应有的尊荣，为自己引来天主的怜悯。如果那些站稳者的祈祷能得到他们自己谦卑的帮助，这些人自己的信函将会好多少倍！因为所求的更容易获得，当所求的对象配得所应获得的之时。\par

4. 然而，对于\Place{兰贝萨}的\Person{普里瓦图斯}，你一如既往地行事，希望通知我们此事，因为它令人忧虑；因为我们都应当为整个普世教会的身体警醒，其成员分散在各个行省。但即使在没有收到你的信之前，那个狡猾之人的欺诈也未能逃过我们的眼睛；因为先前，当某一位\Person{富图鲁斯}——\Person{普里瓦图斯}的旗手——从那个非常邪恶的团体中前来，希望以欺诈手段从我们这里获得信件时，我们既非不认识他是谁，他也没有得到他想要的信件。我们在主内衷心向你告别。\par

\SourceRecord{Translated by Robert Ernest Wallis.；Ante-Nicene Fathers；Vol. 5.；Edited by Alexander Roberts, James Donaldson, and A. Cleveland Coxe.；Buffalo, NY: Christian Literature Publishing Co.,；1886.；<http://www.newadvent.org/fathers/050629.htm>.}

% structure: p0001 source=h1 target=section reason=page-title

\section{第三十封书信}

\Emph{罗马神职致西彼廉}\par

\Emph{论述——罗马神职在前一封信中已谈及的事，在本信中更全面、更实质地加以深入；此外，他们还答复了西彼廉的另一封信（该信被认为已佚失，他们从中引用了几句话）。他们感谢西彼廉寄给罗马宣信者和殉道者的信件。}\par

1. 致父亲西彼廉，常驻罗马的司铎与执事们，愿你们平安。虽然一个意识到自身正直、依靠福音纪律的力量并在天上谕令中作自己真实见证的心灵，习惯于以天主为唯一的审判者而满足，既不寻求他人的称赞，也不畏惧他人的指责，然而那些人更值得双重的赞美，他们知道自己只对天主作为审判者负责，却仍希望自己的行为也得到弟兄们的认可。西彼廉兄弟，你这样做并不奇怪，因为你以惯有的谦逊和天生的勤勉，希望我们发现我们是你计划的分享者而非评判者，这样我们因赞同你的行为而与你同得称赞，并且因我们完全同意你而能与你同为你良善计谋的共同继承人。同样，我们都认为自己在其中劳苦，因为在其中我们都被视为在同样的谴责与纪律的合一中联合。\par

2. 因为在和平中有什么比保持神性严厉的适当严格更为适宜？在逼迫的战争中又有什么比这更为必要？抗拒这严厉的人，必然会飘摇在事务不稳定的航程中，被各种不确定的风暴抛来抛去；而仿佛从手中夺走指南的舵，他将教会救恩的船驶向礁石；以至于教会的救恩似乎只能这样得到保障：即像驱散逆浪一样击退任何反对者，并像在风暴中紧握救恩的舵一样维持纪律本身一向守护的规则。这计谋并非今日才被我们考虑，这些对付恶人的突然措施也并非最近才想到；在我们中间，这被认为是古老的严厉、古老的信仰、古老的纪律，因为宗徒不会如此赞扬我们，当他说\Quote{你们的信德传遍了全世界} \BibleRef{罗 1:8}，除非从那时起，那力量就已经从那些时代汲取了信德的根基；从这赞扬和荣耀中堕落是极大的罪。因为从未达到赞美的宣告是较小的耻辱，而从赞美的高度跌落是更大的；没有得到美好见证的尊荣是较小的罪过，而失去美好见证的尊荣是更大的；未曾因德行而闻名、默默无闻而无可赞美是较小的不光彩，而被剥夺了信德、失去了应有的赞美则是更大的。因为那些为某人荣耀而宣告的事，除非以忧虑和谨慎的辛劳加以持守，否则会膨胀为极大罪过的反感。\par

3. 我们这样说并非不诚实，我们以前的信函已经证明，在其中我们以非常清楚的陈述向你们宣告了我们的意见：既针对那些因非法提交邪恶的证书而自证为不忠信的人，仿佛他们以为能逃脱魔鬼的网罗；其实，他们与接近邪恶祭台的人一样，因向魔鬼作证而被牢牢抓住；也针对那些在证书制作后使用它们的人，尽管他们不在场，却因命令如此书写而确实宣称了他们的在场。因为吩咐作恶的人并非无罪；同意公开谈论罪行的人，即使自己没有犯下，也与罪行有关。既然信德的全部奥秘被认为包含在宣认基督圣名之中，那么寻求狡猾托词为自己开脱的人，就是否认了基督；而想要显得已经满足对抗福音而颁布的谕令或法律的人，就是因他想要显得服从的行为而服从了那些谕令。此外，我们也宣告了我们的信德和同意，对抗那些用非法祭献玷污了自己手和口的人，他们的心早已玷污；因此他们的手和口也玷污了。罗马教会绝不可用如此亵渎的轻率来松懈她的力量，通过推翻信德的威严而松驰她严格的神经；以致当你们堕落弟兄的残骸不仅仍然躺着，而且还在周围倒下时，应该提供过于仓促的补救措施，而肯定无济于事，并借虚假的仁慈，在他们旧伤痕上又加上新伤痕；以至于连忏悔也被从这些可怜的人身上夺去，使他们更加败亡。因为如果连医生本人也因截断忏悔而为新危险开辟道路，只隐藏伤口，不让时间的必要治疗来愈合疤痕，那么宽恕的医药在哪里有益呢？这不是医治，而是——若我们说实话——杀害。\par

4. 然而，你们有从那些精修圣人那里来的信件，与我们的信件一致。他们因宣认的身份而仍被囚禁在此监狱中，并因福音的竞赛，他们的信德已在光荣的宣认中为他们加冕。在那些信件中，他们坚持了福音纪律的严厉，并撤销了不合法的请求，以免成为教会的耻辱。除非他们这样做了，否则福音纪律的废墟就不容易修复，尤其是因为保持福音活力的完整和尊严，最适合那些为了福音而将自己交付给狂怒者折磨和切割的人，这样他们就不会在殉道的时候，若想成为福音的背叛者而理应失去殉道的荣誉。因为一个人若不守护他所拥有的，在拥有它的条件下，通过违反他拥有它的条件，便会失去他所拥有的。\par

5. 在这件事上，我们也应当，并且确实给予你们丰厚的感谢，因为你们用信件照亮了他们监狱的黑暗；你们以一切可能的方式进入监狱去探望他们；你们通过讲话和信件，更新了他们那因自身信德和宣认而坚强的心灵；你们以配得上的赞美追随他们的幸福，激励他们更加热切地渴求天上的荣耀；你们敦促他们前进；你们用话语的力量鼓舞了那些我们相信并希望即将成为胜利者的人；这样，尽管一切似乎都来自那些精修圣人的信德和天主的仁慈，但他们在殉道中似乎在某种程度上成了你们的债务人。但再次回到我们似乎已经离题的话语出发点，你们将会看到附在后面的我们也寄往\Place{西西里}的那种信件；虽然我们有更大的必要推迟此事；自从最可纪念的\Person{法比安}离开后，由于事务和时代的困难，我们至今没有任命主教来妥善安排这类事务，并以权威和智慧处理那些跌倒者的事。然而，你们自己在如此重要的事上已经宣告的，也令我们满意，即必须首先维护教会的和平；然后，召集主教、司铎、执事、精修圣人以及坚定的平信徒举行会议，共同处理跌倒者的案件。因为由许多人商议后，去审查似乎由许多人犯下的罪行，似乎极为不公和沉重；或者说，当如此大的罪行已知已扩散到许多人中时，一个人就应作出判决；因为不是由许多人同意的法令不会是稳固的。请看几乎整个世界都被蹂躏，观察到跌倒者的残余和废墟到处散落，并思考因此需要商议的范围，与罪行似乎广泛传播的程度成正比。让药物不小于伤口，让补救措施不少于死亡，这样，正如那些跌倒者因盲目鲁莽而过于不慎而跌倒，那些努力纠正这一祸害的人也应当在商议中运用一切节制，以免任何做得不当的事被所有人视为无效。\par

6. 因此，让我们——那些至今似乎已逃脱这个时代灾难的人，以及那些似乎已陷入这些时代灾难的人——以同样的商议，同样的祈祷和泪水，恳求天主的尊威，为教会的名字祈求和平。让我们用彼此的祈祷，轮流珍爱、守护、武装彼此；让我们为跌倒者祈祷，使他们被扶起；为站立者祈祷，使他们不致被诱惑到毁灭的地步；为那些据说已跌倒的人祈祷，使他们认识到自己罪过的严重性，并明白这需要不是短暂或仓促的医治；让我们祈祷，悔改也伴随着跌倒者宽恕的果实；如此，当他们理解了自己的罪行时，他们愿意暂时对我们有耐心，不再搅扰教会的动荡状态，以免他们自己似乎为我们点燃了内部的迫害，并且他们的不安被加增到他们罪过的堆中。因为谦逊非常适合那些其罪过中受谴责的是无耻之心的人。让他们确实敲门，但肯定不要破门而入；让他们出现在教堂的门槛前，但肯定不要越过它；让他们在天上营地的门口守望，但让他们以谦逊武装自己，通过谦逊他们认识到自己是逃兵；让他们重新吹起祈祷的号角，但不要以此吹响战争的点；让他们确实以谦逊的武器武装自己，并重新拿起信德的盾牌，他们因对死亡的恐惧而通过否认将其放下；但那些已经武装起来的人要相信，他们是对抗他们的仇敌魔鬼，而不是对抗为他们的跌倒而哀伤的教会。一个谦逊的恳求将大大有益于他们；一个羞怯的祈求，必要的谦卑，一种不轻率的忍耐。让他们以眼泪作为他们苦难的使者；让他们从内心深处发出的叹息履行辩护者的职责，并证明他们对所犯的罪过的悲痛和羞愧。\par

7. 不仅如此，若他们因所犯之罪的重大而战栗；若他们以真正医治之手处理心中和良心上那致命的创伤以及那隐秘而微妙的恶毒的深层处所，让他们甚至羞愧于祈求；除非，再次强调，若不恳求和平的帮助，那将是更大的危险和耻辱。但让这一切都在圣事中进行；在他们恳求的法律中，应顾及时机；祈求应当谦卑，恳求应当克制，因为被恳求者应当被软化，而非激怒；正如应仰望天主的仁慈，同样也应仰望天主的责罚；正如经上所载：\BibleQuote{我豁免了你全部的债，因为你恳求了我。} \BibleRef{玛 18:32} 同样也记载着：\Quote{凡在人前否认我的，我也将在我父及祂的天使前否认他。} 因为天主既是仁慈的，祂也要求遵守祂的诫命，并且确实严格要求；正如祂邀请人赴宴，那没有穿婚宴礼服的人，祂捆起他的手和脚，把他丢到圣人集会之外。祂预备了天堂，但祂也预备了地狱。祂预备了安息之所，但祂也预备了永罚。祂预备了人不可接近的光明，但祂也预备了永夜那广阔而永恒的黑暗。\par

8. 在这些事上，我们渴望持守这条中庸之道的节制，我们长期，实际上许多人，而且，加上我们附近一些可及的主教，以及那些被逼迫的烈火从其他省份驱赶出来的远处的主教，共同认为在任命一位主教之前不应做任何新事；但我们相信，对于失足者的照顾必须适度处理，以便在此期间，当主教的赋予被天主对我们保留时，那些能够忍受延迟之拖延者的案件应被悬置；但对于那些迫近的死亡不容忍延迟的人，如果他们已悔改并多次宣认憎恨自己的行为；如果他们以眼泪、以叹息、以哭泣显露出悲伤和真正痛悔心灵的标志，当就人力所能判断而言，不再有生存的希望时；最后，应给予他们这种谨慎和小心帮助，天主自己知道祂将如何对待他们，以及祂将如何衡量祂审判的天平；然而我们却焦虑地注意，既不让不虔诚者赞扬我们的圆滑便利，也不让真正痛悔者指责我们的严苛为残忍。至福和光荣的父，我们祝愿您在主内永远衷心安康；并请记念我们。\par

\SourceRecord{Translated by Robert Ernest Wallis.；Ante-Nicene Fathers；Vol. 5.；Edited by Alexander Roberts, James Donaldson, and A. Cleveland Coxe.；Buffalo, NY: Christian Literature Publishing Co.,；1886.；<http://www.newadvent.org/fathers/050630.htm>.}

% structure: p0001 source=h1 target=section reason=page-title

\section{第三十一封书信}

\Emph{致\Place{迦太基}圣职界，论寄往\Place{罗马}及从该处收到的信件。}\par

\Emph{提要——\Place{迦太基}圣职界受请留意，务使\Place{罗马}圣职界的信件及\Person{西彼廉}的答复得以传达。}\par

\Person{西彼廉}致司铎及执事，他的弟兄们，安好。为使你们，我亲爱的弟兄，知道我曾向在\Place{罗马}任职的圣职界发出了什么信件，以及他们如何答复了我，此外，\Person{莫伊塞斯}和\Person{马克西穆}（司铎）、\Person{鲁菲努斯}和\Person{尼科斯特拉图斯}（执事）以及与他们一同被囚禁在狱中的其余宣信者，也如何答复了我的信件，我已将副本寄给你们阅读。你们要尽最大努力，务使我所写的和他们的答复，让我们的弟兄们知晓。此外，若有任何来自外地的主教，我的同僚，或司铎、执事，在你们当中或到你们那里，让他们从你们那里听到这一切；若他们希望抄写信件副本并带给他们的教友，就给他们抄写的机会。尽管我也已吩咐读经员、我们的弟兄\Person{撒图鲁斯}，准许任何愿意抄写的人自由抄写。这样，在暂时以任何方式整顿教会状况时，众人能持守一个且忠信的共识。至于其他需处理的事，正如我也写信给几位同僚，待我们蒙主准许，开始聚集在一处时，将在公议会中更充分地商议。祝你们，弟兄们，心爱的和渴慕的，永远衷心地平安。请代问候弟兄团体。祝安。\par

\SourceRecord{Translated by Robert Ernest Wallis.；Ante-Nicene Fathers；Vol. 5.；Edited by Alexander Roberts, James Donaldson, and A. Cleveland Coxe.；Buffalo, NY: Christian Literature Publishing Co.,；1886.；<http://www.newadvent.org/fathers/050631.htm>.}

% structure: p0001 source=h1 target=section reason=page-title

\section{第三十二封书信}

\Emph{致圣职人员和民众：关于\Person{奥勒留}被任命为读经员}\par

\Emph{论题：\Person{西彼廉}告知圣职人员和民众，他已任命宣信者\Person{奥勒留}为读经员，并顺便赞扬他品德与心志的恒毅，这使他甚至配得教会中更高的品级。}\par

1. \Person{西彼廉}致诸位长老与执事，并全体民众，安好。亲爱的弟兄们，在圣职人员的任命中，我们通常事先与你们商议，并以共同的意见衡量个别人的品格与功绩。然而，当主的认可先于人的见证时，就不必等待人的见证。我们的兄弟\Person{奥勒留}，一位杰出的青年，已获主认可并为天主所爱，年纪虽轻，但在德行与信德上却已成熟；他在本性的能力上逊于年龄，但在\Emph{他所应得的}尊荣上卓越——他在这里进行过双重的战斗，两次宣认，两次在宣认的胜利中光荣，既在赛程中征服并被流放，又最终在更激烈的战斗中奋战，在受苦的争战中得胜凯旋。每当仇敌想召出天主的仆役时，这位敏捷勇敢的兵士便一面战斗，一面征服。此前，当他被流放时，在少数人面前作战尚且是小事一桩；他更配得在法庭上以更显赫的德行作战，以至在征服长官之后，也战胜了总督，并在流放之后，接着克服了酷刑。我无法发现应更多称道他身上的什么——他伤口的荣耀，还是他品格的谦逊；他是因德行的尊荣而杰出，还是因谦逊的可敬而值得赞扬。他既在尊严上如此卓越，又在谦卑中如此卑微，以至于他似乎是被天主特意保留，作为其余人学习教会纪律的榜样，即天主的仆人们应如何在宣认中以勇气获胜，并在宣认之后以其品格引人注目。\par

2. 这样一个人，不应按年龄而应按功绩来判断，他配得圣职任命中更高的品级与更大的提升。但在此期间，我认为他应从读经的职务开始；因为对于一个以光荣的言语宣认过主的人来说，最适合莫过于以庄严重复神圣读经来宣扬他；莫过于在宣告基督见证的高超话语之后，诵读那产生殉道者的基督福音；莫过于在离开刑台后到读经台前；莫过于在那里曾为外邦人群所瞩目，在这里则为弟兄们所看见；莫过于在那里曾使周围民众惊叹地聆听，在这里则使兄弟团体喜悦地听闻。那么，最亲爱的弟兄们，要知道此人已由我和当时在场的同事们一起任命。我知道你们会欣然接受这个消息，并渴望我们教会中能有尽可能多这样的人被任命。而且，因为喜乐总是急切的，欢愉不能忍受迟延，他目前就在主日为我诵读；也就是说，他以庄严开始他的读经职务，开启了和平的序幕。请你们恒常地恳切祈祷，并用你们的祈祷协助我的祈祷，使主的仁慈眷顾我们，早日将司铎平安地还给民众，并将殉道者作为读经员与司铎一同赐还。亲爱的在圣父天主及耶稣基督内的弟兄们，我衷心祝你们常常安好。\par

\SourceRecord{Translated by Robert Ernest Wallis.；Ante-Nicene Fathers；Vol. 5.；Edited by Alexander Roberts, James Donaldson, and A. Cleveland Coxe.；Buffalo, NY: Christian Literature Publishing Co.,；1886.；<http://www.newadvent.org/fathers/050632.htm>.}

% structure: p0001 source=h1 target=section reason=page-title

\section{\WorkTitle{Epistle 33}}

\Emph{致圣职人员与全体信众，关于擢升\Person{塞莱里努斯}为宣读员}\par

\Emph{说明——这封信与前信意旨大致相同，只是他极力称赞\Person{塞莱里努斯}在宣认信仰时的坚毅。此外，这两封信均写于他退隐期间，这一点从上下文的情况足以表明。}\par

1. \Person{西彼廉}致诸位司铎、执事及全体子民，他的在主内的弟兄们，问安。亲爱的弟兄们，我们应当承认并拥抱天主的恩惠，主屈尊在当今时代以祂的眷顾赐予祂的善作证者和光荣的殉道者一段喘息之机，以装饰和彰显祂的教会，好让那些曾英勇宣认基督的人，后来在教会的职务中装饰基督的圣职。因此，当你们收到我的信时，请与我一同欢欣喜乐；在这封信中，我和当时在场的同僚们向你们提及我们的弟兄\Person{塞莱里努斯}，他以勇气和品德同样光荣，被加入我们的圣职行列，并非出于人的推荐，而是出于\OriginalTerm{天主的恩宠}{divine condescension}；他原对顺服教会有所迟疑，却因教会自身的劝诫和鼓励，在夜间的异象中被催促，不得拒绝我们的劝说；而教会更有力量约束了他，因为他不应没有教会的尊荣，主既以天上荣耀的尊位尊崇了他。\par

2. 此人在我们这世代斗争中首当其冲；他是基督士兵中的先锋；他在迫害初燃之际，便与骚乱的元凶和创始者交锋，以不可战胜的坚韧征服了他的对手。他为他人开辟了得胜之路；虽受重创，却因漫长争战的持久痛苦而奇迹般得胜。他被囚禁在狭窄的牢房十九天，受刑且戴着镣铐；但身体虽被锁链束缚，灵魂却依然自由无拘。他的肉体因长期忍饥挨渴而消瘦；但天主以属灵的滋养喂养他那活在信德与德行中的灵魂。他卧于刑罚中，刑罚反使他更为坚强；被囚禁，却比囚禁他的人更伟大；俯伏在地，却比站立的人更高昂；被捆绑，却比捆绑他的锁链更牢固；受审判，却比审判他的人更崇高；尽管双脚在刑架上被缚，但那蛇却被践踏、被碾碎、被击败。他光荣的身体上闪耀着伤口的明亮证据；那些明显的痕迹显明出来，出现在这人因长期消耗而疲惫的筋脉与四肢上。伟大的事——奇妙的事——是弟兄们所能听到的，关于他的德行和赞美。若有人像\Person{多默}那样，对所闻之事信心不足，也不缺少眼见之信，好叫人既听到也能看到。在天主的仆人身上，伤口的荣耀成就了胜利；伤痕的记忆保存了那荣耀。\par

3. 在我们亲爱的\Person{塞莱里努斯}身上，这种荣耀头衔并非陌生新事。他正追随着亲族的足迹；他以同等的天主的恩宠的尊荣与他的父母和亲属竞争。他的祖母\Person{塞莱里娜}不久之前已戴上殉道的冠冕。此外，他的叔父和舅父\Person{劳伦蒂乌斯}与\Person{埃格纳提乌斯}，他们自己也曾一度在世俗的阵营中作战，但他们是天主真实的、属灵的士兵，藉着宣认基督击倒了魔鬼，从而藉他们卓越的苦难从主那里获得了棕榈枝与冠冕。我们常为他们奉献祭献，正如你们所记得的，每当我们以年度纪念庆祝殉道者的苦难与日子时。因此，他不能是何等卑下低劣之辈，因为这一家族的尊荣和高贵的血统，由家中德性与信德的榜样所激发。倘若在世俗家庭中，成为贵族是纹章与赞美之事，那么在天上纹章中成为高贵等级，其赞美和尊荣是何等之大！我不知该称谁为更蒙福——是那些先辈，因如此显赫的后裔，还是他，因如此光荣的起源。天主的恩宠就这样平均地流淌于二者之间，并来回传递，正如后代的尊荣灿烂了他们的冠冕，同样，祖先的崇高照亮了他的荣耀。\par

4. 当此人在如此的天主的恩宠下，因迫害者本人对他的见证与惊叹而闻名，来到我们这里时，亲爱的弟兄们，除了将他安置在讲坛上，即教会的审判座上，使他以更高地位的崇高为依托，并因尊荣的光辉而为全体百姓所瞩目，好让他宣读他所如此勇敢而忠信地追随的主的诫命与福音，还能做什么呢？让那每日宣认主的嘴唇，在主的言语中被听见。让我们看看他在教会中是否能再被提升到更高的等级。作证者所能为弟兄们带来的益处，莫过于当福音的诵读从他口中被听到时，每个听见的人都效法诵读者的信德。他本应在诵读上同\Person{奥勒留}联合；在此事上，他也在天主尊荣的联盟中与他联合；在一切德行与赞美的标记上，他已与他结合。两人平等，彼此相似，他们崇高于荣耀的程度，恰与谦卑于谦逊的程度相当。他们因天主的恩宠被高举，也同样因自身的平安与宁静而自卑，且同样为每一个人提供了德性与品行的榜样，既适合战斗也适合和平；前者因力量而受赞美，后者因谦逊而受赞美。\par

5. 在这样的仆人身上，主欢悦；在这样的明认者身上，祂得以光荣——他们的生活方式与品行如此有利于彰显他们的光荣，以致为他人提供了纪律的教导。为此，基督愿意他们长久留在教会内；为此，祂保全了他们，将他们从死亡中夺回——可以说是为他们行了某种复活；好使弟兄们看到，在光荣上没有什么高于他们，在谦卑上没有什么低于他们，而弟兄团体的生活之道应与这些人相伴。因此，你们要知道，这些人目前被任命为读经员，因为灯应当放在灯台上，好照亮众人，并且他们光荣的容貌应当安置在更高的地方，在那里，被周围所有的弟兄团体看见，就能给观看者带来光荣的激励。但是你们要知道，我已经打算赐予他们司铎的尊荣，好使他们与司铎受到同样的荣耀，获得均等的月度份额，日后在他们年长力强时与我们同坐；虽然就年岁特质而言，没有人可以被视为低下，因为他已经以自己光荣的尊贵完成了他的年岁。我祝你们，弟兄们，可爱且深切渴望的，永远衷心平安。\par

\SourceRecord{Translated by Robert Ernest Wallis.；Ante-Nicene Fathers；Vol. 5.；Edited by Alexander Roberts, James Donaldson, and A. Cleveland Coxe.；Buffalo, NY: Christian Literature Publishing Co.,；1886.；<http://www.newadvent.org/fathers/050633.htm>.}

% structure: p0001 source=h1 target=section reason=page-title

\section{第三十四封书信}

\Emph{致同一人，关于立\Person{努米迪库斯}为司铎。}\par

\Emph{绪论——\Person{西彼廉}告知圣职团与信众，他已立\Person{努米迪库斯}为司铎，并简要赞扬其功绩。}\par

\Person{西彼廉}致司铎们、执事们及全体信众——我极亲爱且长久思念的弟兄们，问安。最亲爱的弟兄们，那属于我们教会的共同喜乐与最大光荣的事，理当告知你们；因为你们必须知道，我受天主的俯就所劝告和指教，\Person{努米迪库斯}司铎应被列入\Place{迦太基}司铎之列，并与我们同坐于圣职人员之中——他是一位因宣认之光而极其显赫、因德行和信德双重荣耀而得高举的人；他藉着自己的劝勉，先于自己派遣了大量的殉道者，他们被石头和火焰所杀；他又喜悦地看见他的妻子与他同在，与其余人一同被焚烧（我应说，被保全）。他本人半被焚毁，被石头掩埋，被当作死人留下；后来他的女儿，出于忧心的慈爱，寻找她父亲的尸体——他被发现半死，被拖出并复苏，但他不情愿地留在了他自己先前派遣的同伴之外。然而他留下的原因，如我们所见，是这样：为使主能将他添入我们的圣职团，并用光荣的司铎来妆饰我们因某些人的跌倒而荒凉的司铎数目。当天主允许时，他将在其地区被提升至更大的职务，那时，藉主的保护，我们将再次来到你们面前。与此同时，愿所启示的事得以成就，使我们以感恩之心领受天主的这份恩赐，并期望从主的仁慈获得更多同类光荣，好使祂教会的力量得以更新，祂能造就如此温良谦逊的人，在我们集会的尊荣中繁盛。我祝你们，极亲爱且长久思念的弟兄们，永远衷心平安。\par

\SourceRecord{Translated by Robert Ernest Wallis.；Ante-Nicene Fathers；Vol. 5.；Edited by Alexander Roberts, James Donaldson, and A. Cleveland Coxe.；Buffalo, NY: Christian Literature Publishing Co.,；1886.；<http://www.newadvent.org/fathers/050634.htm>.}

% structure: p0001 source=h1 target=section reason=page-title

\section{第三十五封书信}

\Emph{致圣职人员，论照顾穷人与旅客。}\par

\Emph{论点——他警告他们不可忽视寡妇、病患、穷人或旅客。}\par

\Person{西彼廉}致他亲爱的弟兄——司铎与执事，问候。因天主的恩宠，我平安地致候你们，亲爱的弟兄们，渴望很快来到你们中间，满足我自己、你们以及所有弟兄的心愿。然而，我也必须顾全共同的和平，同时，尽管心中疲惫，却要远离你们，以免我的出现激起外邦人的嫉妒和暴行，使我成为破坏和平的原因，而我本应顾全众人的安宁。因此，当你们写信告知诸事已安排妥当，我应当前往时，或主若俯允事先告知我时，我就会到你们那里去。因为哪里能比主愿意我信从并成长的地方更美好、更喜悦呢？我请求你们殷勤照顾寡妇、病患和所有穷人。此外，你们可以从我自己的份额中为旅客提供费用，若有穷困者；我已将这笔份额留给了我们的同僚司铎\Person{罗加提亚努斯}。为避免此份额被全部挪用，我已通过辅祭\Person{纳里库斯}另外寄送一份给他，以便穷人能得到更充足和及时的照顾。亲爱的弟兄们，我衷心祝愿你们永远平安；请纪念我。代我向你们的弟兄们致意，并告诉他们要纪念我。\par

\SourceRecord{Translated by Robert Ernest Wallis.；Ante-Nicene Fathers；Vol. 5.；Edited by Alexander Roberts, James Donaldson, and A. Cleveland Coxe.；Buffalo, NY: Christian Literature Publishing Co.,；1886.；<http://www.newadvent.org/fathers/050635.htm>.}

% structure: p0001 source=h1 target=section reason=page-title

\section{\WorkTitle{书信第三十六}}

\Emph{致司铎，命他们向狱中的宣信者显示各种慈爱。}\par

\Emph{论点。——他劝勉其司铎，应向狱中的宣信者，无论生者或死者，行使一切慈爱与关怀；仔细记下他们的忌日，以便每年纪念；最后，也不应忘记穷人。}\par

1. \Person{西彼廉}致司铎及执事，他的弟兄们，问候。虽然我知道，最亲爱的弟兄，你们在我多次的书信中已受劝勉，要向那些以光荣之声宣认了主、被囚在狱中的人显示一切关怀；然而我一再劝勉你们，切莫缺少对他们的关怀，他们的光荣毫无缺失。我多么希望这地方和我职位的状况允许我此刻与他们同在；我会立即并乐意在我被指定的职务中，向我们最勇敢的弟兄们履行一切爱德的责任。但我恳求你们，让你们的勤勉代表我的职责，并做一切应当为那些因天主的眷顾而在信德与美德上如此卓越的人所做的事。对那些虽然未在狱中受酷刑，却以光荣的死亡离开那里的人，也要施加更热心的看顾和关怀。因为他们的美德和光荣并不少，也应被列于真福殉道者之中。就他们所能的，他们承受了他们准备并装备好要承受的一切。在天主眼中将自己献给酷刑和死亡的人，已经承受了他愿意承受的一切；因为不是他缺少酷刑，而是酷刑缺少他。\BibleQuote{凡在人前承认我的，我也必在我天上的父前承认他} \BibleRef{玛 10:32}，主说。他们已承认了他。\BibleQuote{那坚持到底的，必得救} \BibleRef{玛 10:22}，主说。他们坚持了，并将他们美德的纯洁无瑕的功绩保持到底。再者，经上记载：\BibleQuote{你当至死忠心，我就赐给你那生命的冠冕} \BibleRef{默 2:10}。他们忠信、坚定、不可征服地坚持到底，直至死亡。当在狱中和锁链中的意愿和宣认之名加上死亡的结局时，殉道者的光荣就圆满了。\par

2. 最后，也要记下他们离去的日子，使我们能在殉道者的纪念中庆祝他们的纪念，尽管我们最忠信而热心的弟兄\Person{特土禄}，他除了在其他事务上以各种劳作服事弟兄们之外，在照顾他们遗体的任何事上也毫不疏忽，已经写信并将写信向我告知那些日子，我们的真福弟兄们在狱中通过光荣死亡之门进入永生；而且我们在这里为他们举行奉献和祭献以示纪念，这些事，在天主的护佑下，我们不久将与你们一同举行。你们的关怀（如我常所写的）和勤勉也不应忘记穷人——我是指那些在信仰中坚定不移、与我们勇敢战斗、没有离开基督军营的人；事实上，我们现在应向他们显示更大的爱和关怀，因为他们既不被贫困所迫，也不被迫害的风暴击倒，而是忠诚信实地事奉主，并为其他穷人树立了信德的榜样。我祝你们，亲爱的弟兄，极切思念的，永远衷心告别；并请记念我。请以我的名义问候弟兄们。再见。\par

\SourceRecord{Translated by Robert Ernest Wallis.；Ante-Nicene Fathers；Vol. 5.；Edited by Alexander Roberts, James Donaldson, and A. Cleveland Coxe.；Buffalo, NY: Christian Literature Publishing Co.,；1886.；<http://www.newadvent.org/fathers/050636.htm>.}

% structure: p0001 source=h1 target=section reason=page-title

\section{\WorkTitle{第三十七封书函}}

致\Person{卡尔多尼乌斯}、\Person{赫尔库拉努斯}及其他人士：关于绝罚\Person{费利奇西穆斯}\par

论点——\Person{费利奇西穆斯}及其叛乱同伙应被禁止参与全体共融。\par

1. \Person{西彼廉}致\Person{卡尔多尼乌斯}与\Person{赫尔库拉努斯}，他的同僚，也致\Person{罗加提亚努斯}与\Person{努米迪库斯}，他的司铎同伴，问安。最亲爱的弟兄们，收到你们的来信，我深感痛心，因为我一直对自己期望并希望保全所有弟兄的安全，并按照爱德的要求保持羊群不受伤害，而你们现在告诉我，\Person{费利奇西穆斯}以邪恶和诡计图谋了许多事；以致除了他旧日的欺诈和掠夺（我先前已颇为了解）之外，他现在甚至还试图分裂主教与一部分人群；也就是将羊群与牧人分离，将儿子与父母分离，并分散基督的肢体。尽管我派遣你们作为我的代表，用资金去处理我们弟兄们的急需，并且若有任何人还希望操练其手艺，也要用足够的增额协助他们的愿望，同时还要记录他们的年龄、状况和功过——好让我，这个负有完全认识他们所有人职责的人，可以将那些配得、谦卑和温良的人擢升到教会管理的职务中——他却横加干涉，命令不得周济任何人，并且我原本希望经过仔细查考才能确定的事，也不得查明；他还以邪恶的权势和暴力的恐吓威胁我们那些首先前来接受救济的弟兄，说那些愿意听从我的人，在死亡中不得与他共融。\par

2. 既然在这一切之后，他既没有被我的职位之尊所感动，也没有被你们的权威和在场而动摇，反而出于自己的冲动，扰乱弟兄们的平安，带着更多人冲了出来，以急躁的疯狂自命为党派的领袖和叛乱的魁首——在这方面，我确实为几位弟兄感到庆幸，他们已远离这种狂妄，宁愿与你们一致，以便能留在他们的母亲教会内，并从施予的主教那里领取他们的俸禄；我确切知道，其他人也将和平地这样做，并会很快从他们鲁莽的错误中退出——与此同时，既然\Person{费利奇西穆斯}威胁说，那些听从我们的人，即那些与我们共融的人，在死亡中不得与他共融，就让他领受他首先宣布的判决吧，好让他知道自己已被我们绝罚；因为他除了以最清楚的事实所知道的欺诈和掠夺之外，还加上了奸淫的罪行，这是我们那些庄重的弟兄们所宣称他们已经发现并断言将要证明的；这一切我们将在以后，在主的许可下，与许多同僚聚集一处时，进行司法审理。但是\Person{奥根都斯}，他既不顾他的主教也不顾他的教会，同样在这阴谋和叛乱中与他联合，如果他进一步与他坚持，就让他承担那好斗和暴躁的人所惹来的判决。此外，凡与他的阴谋和叛乱联合的人，要知道他不得在教会中与我们共融，因为他已自愿选择与教会分离。请将我这封信读给我们的弟兄们听，并转交至迦太基的圣职人员，附上那些与\Person{费利奇西穆斯}联合者的名字。亲爱的弟兄们，我祝愿你们永远在主内衷心安康，并请记念我。别了。\par

\SourceRecord{Translated by Robert Ernest Wallis.；Ante-Nicene Fathers；Vol. 5.；Edited by Alexander Roberts, James Donaldson, and A. Cleveland Coxe.；Buffalo, NY: Christian Literature Publishing Co.,；1886.；<http://www.newadvent.org/fathers/050637.htm>.}

% structure: p0001 source=h1 target=section reason=page-title

\section{书信三十八}

\Emph{\Person{Caldonius}、\Person{Herculanus}及其他人论将\Person{Felicissimus}及其同伙开除出教的书信。}\par

\Emph{摘要。——\Person{Caldonius}、\Person{Herculanus}及其他人执行前封书信所命之事。}\par

\Person{Caldonius}同其同僚\Person{Herculanus}和\Person{Victor}，以及司铎\Person{Rogatianus}和\Person{Numidicus}。我们将\Person{Felicissimus}和\Person{Augendus}逐出共融；也将流放者中的\Person{Repostus}，血染者中的\Person{Irene}，以及女裁缝\Person{Paula}逐出；你们应从我的签署得知此事；我们还拒绝了\Person{Sophronius}和\Person{Soliassus}\OriginalTerm{卖香肠的}{budinarius}——他自己也是流放者之一。\par

\SourceRecord{Translated by Robert Ernest Wallis.；Ante-Nicene Fathers；Vol. 5.；Edited by Alexander Roberts, James Donaldson, and A. Cleveland Coxe.；Buffalo, NY: Christian Literature Publishing Co.,；1886.；<http://www.newadvent.org/fathers/050638.htm>.}

% structure: p0001 source=h1 target=section reason=page-title

\section{书信三十九}

\Emph{致民众：关于\Person{非利奇西穆斯}派系的五位分裂司铎}\par

\Emph{要旨：正如在此前一封书信中，\Person{西彼廉}告诉圣职人员，现在他告诉民众：\Person{非利奇西穆斯}连同其派系的五位司铎应被远离，他们不仅毫无辨别地给予失足者平安，还煽动造反和分裂反对他自己。}\par

1. \Person{西彼廉}致全体民众，问安。最亲爱的弟兄们，虽然最忠实和正直的司铎\Person{维尔提乌}，以及司铎、宣信者、因天主降尊纡贵的荣耀而闻名的\Person{罗加提亚努斯}和\Person{努米迪库斯}，还有执事们，都是好人，专注于教会行政的一切职责，连同其他服务人员，以他们的临在给予你们充分的关注，并且不停止以勤勉的劝诫坚固个人，而且以他们健康的劝告治理和改造失足者的思想，然而，我尽我所能地劝诫，并尽我所能地以我的书信探望你们。我说，最亲爱的弟兄们，我以我的书信前来；因为某些司铎的恶意和背信弃义造成了这一点，即不允许我在复活节前到你们那里去；因为他们念念不忘他们的阴谋，并保留着对我主教职位的古老毒液，即对你们的投票和天主的审判，他们重新开始对我的旧攻击，并以他们惯常的诡计再次开始他们的亵渎阴谋。而且，的确，出于天主的眷顾，既非出于我们的愿望或希望，不，尽管我们宽容沉默，他们已遭受了他们应得的惩罚；以至于，不是被我们驱逐，他们自动地驱逐了自己。他们自己，在各自的良心面前，已根据你们的投票和天主的判决对自己宣判了刑罚。这些阴谋家和恶人自动地将自己逐出了教会。\par

2. 现在已显出\Person{非利奇西穆斯}的派系从何而来；它立于何根，凭何力。这些人从前给某些宣信者提供鼓励和劝诫，不要同意他们的主教，不要按照主的诫命以信德和安静持守教会的纪律，不要以无可指责和无玷污的品行保持他们宣信的荣耀。并且，恐怕腐蚀某些宣信者的思想，以及希望武装我们破碎的兄弟会的一部分反对天主的司祭职还嫌不够，他们现在以他们恶毒的诡计将注意力转向失足者的毁灭，使患病和受伤者远离他们伤口的医治，以及那些因跌倒的不幸而较不适宜、较不坚强以采纳更坚强劝告的人；并以一种虚假和平的谎言邀请他们走向致命的鲁莽，放弃祈祷和恳求，而通过这些，以长久连续的补赎，主才得以平息。\par

3. 但我祈求你们，弟兄们，要谨防魔鬼的陷阱，并关心你们自己的救恩，勤勉地警惕这导致死亡的谬误。这是另一种迫害和另一种试探。那五位司铎正是最近与官吏联合颁布法令的五位首领，他们为要推翻我们的信德，以真理的悖谬将弟兄们软弱的心转向他们致命的网罗。现在，同样的阴谋，同样的颠覆，又由五位司铎与\Person{非利奇西穆斯}联结而行，为毁灭救恩，使天主不被祈求，使那否认基督的人不得向同一基督求怜悯，就是他曾经否认的；使罪过之后，悔改也被夺走；使主不得通过主教和司铎被平息，而是主的司铎被遗弃，一种亵渎任命的传统应产生，违反福音的纪律。并且，虽然这曾经由我们以及宣信者、城市圣职人员、而且由在我们省或海外任命的所有主教共同安排，即有关失足者的案件不应有任何新奇之处，除非我们全部聚集在一处，比较我们的意见，决定一个温和的判决，既合纪律也合慈悲——他们反叛了我们这个意见，而所有司铎的权威和能力被派系的阴谋摧毁。\par

4. 我现在承受着怎样的痛苦，最亲爱的弟兄们，就是我自己不能在此时到你们那里去，我自己不能接近你们每一位，我自己不能按照主及其福音的教导劝诫你们！至今两年的流放还不够，与你们、与你们的容貌和视线悲伤的分离——持续的忧伤和哀悼，在我没有你们的孤独中，以我不断的哀悼将我撕碎，而且我的眼泪日夜流淌，以至于连司铎——就是你们以如此爱和热忱所造立的那一位——都没有机会向你们问安，或被拥入你们的怀抱。这更大的忧伤加在我疲惫的精神上，即在如此多的忧虑和急需中，我自己不能赶到你们那里去，因为由于背信之人的威胁和陷阱，我们担心我们到来时那里可能不会产生更大的骚动；因此，虽然主教应当在一切事上注意和平与安宁，他自己却似乎为叛乱提供了材料，并重新激起了迫害。然而，因此，亲爱的弟兄们，我不但劝诫而且忠告你们，不要轻率地信任有害的话语，不要轻易同意欺骗性的言论，不要以黑暗为光明，以夜为昼，以饥饿为食物，以口渴为饮料，以毒药为药物，以死亡为安全。不要让那些人的年龄或权威欺骗你们，他们对应于两位长老的古代邪恶；正如他们试图败坏和侵犯贞洁的\Person{苏撒纳}，同样他们现在也以他们淫乱的教义试图败坏教会的贞洁和侵犯福音的真理。\par

5. 主大声呼喊说：\Quote{不要听从假先知的话，因为他们心中的异象欺骗了他们。他们说话，但不是出于主的口。他们向藐视主话的人说：你们必得平安。}\BibleRef{耶 23:16} 他们现在在赐平安，而他们自己却没有平安。他们许诺将失足者带回并召回教会，而他们自己却已离开了教会。天主是一位，基督是一位，教会是一个，有一个由主的话建立在磐石上的座位。不能另设祭台，也不能另立新司祭职，除非同一祭台和同一司祭职。凡在别处聚集的，就是分散。凡由人的疯狂所制定，以致违反天主安排的事，都是不贞的、不敬的、亵圣的。要远远离开这类人的传染，逃避他们的话语，视之如毒疮和瘟疫，正如主警告你们说：\Quote{他们是盲人领盲人。但如果盲人领盲人，两人都要掉进坑里。}\BibleRef{玛 15:14} 他们拦截你们的祈祷，就是你们日夜同我们一起向天主倾倒的祈祷，为以正义的赔补平息祂；他们拦截你们的眼泪，就是你们用来洗去所犯之罪孽的眼泪；他们拦截你们真正而忠信地从主的仁慈中求得的平安；他们不知道经上记载：\Quote{那说预言的先知或作梦的，若使你们离开主你们的天主，必被处死。}\BibleRef{申 13:5} 可爱的弟兄们，不要让人使你们偏离主的道路；不要让任何人从基督的福音中夺走你们基督徒；不要让任何人从教会中夺走教会的子女；让那些愿意灭亡的人独自灭亡吧；让那些离开教会的人独自留在教会之外吧；让那些反叛主教的成为没有主教的人吧；让那些从前依照你们的投票、如今依照天主的审判，应当承受自己阴谋和恶毒的判决的人，独自承受他们结党的惩罚吧。\par

6. 主在他的福音中警告我们说：\Quote{你们废弃天主的诫命，为要遵守你们的传统。}\BibleRef{谷 7:9} 愿那些废弃天主的诫命并致力于遵守自己传统的人，被你们勇敢而坚定地拒绝；让失足者的一次跌倒就足够了；不要让任何人用诡计把那些想要起来的人推倒；不要让任何人把那些跌倒的人更深地推下和压制，我们正是为这些人祈祷，愿他们被天主的手和臂扶起；不要让任何人使那些半死不活、恳求恢复昔日健康的人完全失去获救的希望；不要让任何人熄灭了那些在失足的黑暗中摇摆不定的人得救道路上的一切光明。宗徒教训我们说：\Quote{若有人讲异端，不赞同我们主耶稣基督的健全言语及其教义，他因愚昧而自高自大；这样的人你要躲开。}\BibleRef{弟前 6:3} 他又说：\Quote{不要让人以空言欺骗你们，因为为了这些事，天主的义怒降在背命之子身上。所以你们不要作他们的同伙。} 你们没有理由被空言欺骗，并开始作他们堕落的同伙。我恳求你们离开这样的人，听从我们的劝告，我们每天为你们向主献上不断的祈祷，愿你们因主的仁慈被召回教会，我们为你们祈求天主赐予最圆满的平安，首先是母亲，然后是她的子女。也要将你们的祈求祈祷与我们的祈祷祈求联合；将你们的眼泪与我们的哀哭混合。要躲避那些把羊从牧人那里分离的狼；躲避那从世界之初就一直诡诈说谎的魔鬼的毒舌，他撒谎为要欺骗，谄媚为要伤害，许诺好为要给予恶，许诺生命为要置之于死。现在他的话语也显而易见，他的毒害也清清楚楚。他许诺平安，为使平安不可能得到；他许诺救恩，为使犯罪之人不能得到救恩；他许诺教会，同时却如此谋划，使相信他的人完全丧亡于教会之外。\par

7. 如今正是时机，亲爱的弟兄们，你们这些站立稳固的人要勇敢地坚持，并保持你们在迫害中以持续的坚定所保持的光荣稳定；如果你们中有人因仇敌的欺骗而跌倒，那么在这次第二次诱惑中，你们要为你们的希望和你们的平安忠信地寻求劝告；并为了使主宽恕你们，你们不应离开主的司铎，因为经上记载：\Quote{那行事傲慢，不听从当时司祭或审判官的人，那人必被处死。}\BibleRef{申 17:12} 对于这次迫害，这是最后最末的诱惑，它本身也因主的保护将要很快过去；使我在复活节后再次与我的同僚一同出现在你们面前，他们到场时，我们将能够按照你们的判断和我们大家事先已决定的共同意见来安排并完成需要做的事。但如果有人拒绝悔改和向天主作赔补，而屈服于费利西西莫的党派及其随从，并加入异端派别，当知道此人以后不能再返回教会与基督的主教和子民共融。我最亲爱的弟兄们，我祝你们永远衷心安康，并请你们与我一同不断祈祷，使天主的仁慈得以被祈求。\par

\SourceRecord{Translated by Robert Ernest Wallis.；Ante-Nicene Fathers；Vol. 5.；Edited by Alexander Roberts, James Donaldson, and A. Cleveland Coxe.；Buffalo, NY: Christian Literature Publishing Co.,；1886.；<http://www.newadvent.org/fathers/050639.htm>.}

% structure: p0001 source=h1 target=section reason=page-title

\section{第四十封书信}

\Emph{致\Person{科尔乃略}，论其拒绝接受\Person{诺瓦提安}的祝圣。}\par

\Emph{\Person{诺瓦提安}派往\Place{迦太基}教会通报其祝圣的使者被\Person{西彼廉}拒绝。}\par

\Person{西彼廉}致其弟兄\Person{科尔乃略}，请安。亲爱的弟兄，\Person{诺瓦提安}差遣了长老\Person{玛克西穆}、执事\Person{奥森都}，以及一位叫\Person{玛开乌}的和\Person{隆吉诺}，来到我们这里。但我们从他们带来的信件以及他们的言谈和声明中发现，\Person{诺瓦提安}已被祝圣为主教；我们因这违背天主教会的不合法祝圣之邪恶而深感不安，遂立即决意必须阻止他们与我们共融；同时，在驳斥并回绝了他们顽固坚持的论点之后，我和几位聚到我这里的同僚，正在等待我们的同僚\Person{卡尔多纽}和\Person{福尔图纳}的到来——他们最近被我们派往你处，以及出席你祝圣典礼的诸位同僚主教那里——以便当他们来到并报告事情的真相时，藉着他们更大的权威和明确的证据，反对派的邪恶能被平息。此外，我们的同僚\Person{庞培}和\Person{斯德望}也来了，他们为教导我们，提出明确的证据和见证，与他们的庄重和忠信相符，以致于根本无须再听那些\Person{诺瓦提安}所差来之人的话。当他们在我们庄严的集会中闯入，带着恶意的辱骂和喧嚣的吵闹，要求我们和民众公开审查他们所提出并声称将要证明的指控时，我们说：我们的庄重不容许我们容忍同僚的荣誉——他已经被许多人可称赞的判断所拣选、祝圣和认可——再受到竞争者辱骂之声的质疑。由于要将他们被驳斥和压制、以及被显明因非法企图而引发异端的种种事宜汇集成信，将是一件冗长的工作，你将从我们的同僚长老\Person{普里米提乌}那里听到一切详情，当他到你那里时。\par

而且，唯恐他们的狂暴胆量永不停止，他们在此地也竭力将基督的肢体分散到分裂的党派中，切割并撕裂天主教会唯一的身体，以致于他们挨家挨户、穿街走巷，或从一城到另一城、经过某些地区，为他们的顽固和错误寻找同伴，加入他们的分裂。我们已对他们作出这样的答复，并且不会停止命令他们抛弃有害的纷争和争论，并认识到离弃他们的母亲是不虔敬的；要承认并明白，一旦一位主教被其同僚和民众的见证与判断所任命并认可，就绝不可能再另立一位。因此，如果他们平平安安、忠信地顾及自己的利益，如果他们承认自己是基督福音的维护者，就必须回到教会。最亲爱的弟兄，我衷心祝你永远安康。\par

\SourceRecord{Translated by Robert Ernest Wallis.；Ante-Nicene Fathers；Vol. 5.；Edited by Alexander Roberts, James Donaldson, and A. Cleveland Coxe.；Buffalo, NY: Christian Literature Publishing Co.,；1886.；<http://www.newadvent.org/fathers/050640.htm>.}

% structure: p0001 source=h1 target=section reason=page-title

\section{第四十一封书信}

\Emph{致\Person{高内略}：论\Person{西彼廉}赞同其祝圣，以及关于\Person{斐理基息茂}}\par

\Emph{论据——\Person{西彼廉}为自己最初没有毫不犹豫地相信\Person{高内略}的祝圣而辩解，直到收到同僚\Person{加尔多尼}和\Person{福都纳}的信件，这些信件充分证明其合法性；并且顺便重复，关于\Person{诺瓦西安}派的对立派别，他并非在第一次就支持该派，而是支持\Person{高内略}，甚至达到拒绝接受针对他的指控的程度。}\par

\Person{西彼廉}致其兄弟\Person{高内略}，问安。最亲爱的兄弟，正如天主仆人所应做的，尤其是正直和平的司铎所应做的，我们最近派遣了我们的同僚\Person{加尔多尼}和\Person{福都纳}，使他们不仅通过我们信件的劝说，而且通过他们本人的在场和你们所有人的建议，竭尽全力努力劳动，将分裂团体的成员带回到天主教会合一之中，并将他们联系到基督徒爱德的纽带里。但是，由于对立派别的顽固和固执的倔强不仅拒绝了其根源和母亲的怀抱和拥抱，而且甚至随着一种更糟更糟的纷争扩散和复活，为自己任命了一位主教，并且违背一次交付的、关乎上主定立和天主教合一的圣事，在教会之外制造了一个奸淫且敌对的首领；在收到你的信以及我们同僚的信之后，也在我们的同僚、优秀且为我们所爱的\Person{庞培}和\Person{斯德望}到来之时，通过他们，所有这些事情毫无疑问地被提出并带着普遍的喜乐向我们证明，符合神圣传统和教会体制的圣德与真理的要求，我们已经将我们的信件寄给了你。此外，将这些同样的事情提请全省各地同僚的注意，我们也命令我们的弟兄们，带着他们的信件，寄给你。\par

\Emph{此事已完成}，尽管我们的心思和意图已经清楚地向这里的弟兄们和全体民众宣布，当时我们最近收到双方的信件，我们读了你的信，并当众宣布了你被祝圣为主教。此外，铭记共同荣誉，并尊重司铎的庄重和圣洁，我们拒绝了对方团体寄给我们的文件中用刻毒语言堆砌的那些内容；既考虑又衡量，在如此庞大和如此虔诚的弟兄们的集会中，天主的司铎们同坐一堂，祂的祭台设立其间，这些东西既不应该被宣读也不应该被听见。因为这些内容不应轻易提出，也不应粗心大意地粗鲁公开，它们可能通过好争吵的笔触在听众心中引起丑闻，并用不确定的意见扰乱那些远隔重洋、分隔居住的弟兄们。让那些人小心，他们顺从自己的愤怒或私欲，忘记天主的法律和圣洁，乐于暂时传播他们无法证明的事情；尽管他们可能无法成功地摧毁和毁灭无辜，却满足于用谎言和虚假的谣言玷污无辜。确实，我们应该努力，如同主教和司铎所应做的，对于任何人所写的此类事情，我们应拒绝它们。因为否则，我们所学习并宣告的圣经教导会成为什么？\BibleQuote{要谨守你的舌头，不出恶言，你的嘴唇，不说诡诈的话？}\BibleRef{咏 34:13}。以及别处：\BibleQuote{你的口充满邪恶，你的舌头编造诡诈。你坐着毁谤你的兄弟，诬蔑你母亲的儿子}\BibleRef{咏 50:19-20}。还有宗徒所说：\BibleQuote{一切坏话都不可从你们的口中出来，但要说合于建树信德的话，使听者获得恩宠}\BibleRef{厄 4:29}。此外，我们表明了正确的行为方式：如果这些内容是由某些人的诽谤性鲁莽写成的，我们不允许在我们中间宣读它们。因此，最亲爱的兄弟，当这样的信件送到我这里来反对你时，尽管它们是你同坐的你的同司铎的信件，但既然它们散发着宗教纯朴的语调，没有回响任何咒骂和辱骂的吠声，我就命令在圣职人员和民众面前宣读它们。\par

但在渴望得到我们出席你祝圣仪式的同僚的信件时，我们并未忘记古老的习惯，也未寻求任何新奇。因为你通过信件宣布自己被立为主教已经足够，除非在另一方有异议派别，他们通过诽谤和造谣的捏造扰乱了同僚以及许多弟兄们的心思和困惑。为了平息此事，我们认为有必要从我们那些写信给我们的同僚那里获得强大而断然的权威；他们宣布他们的信件证据完全符合你的品格、生活和教导，从而剥夺了你的对手以及那些喜爱新奇或邪恶之人的任何疑虑或分歧的顾虑；根据我们以健全理性权衡的建议，在这个海洋中漂泊的弟兄们的心灵真诚而坚定地认可了你的司祭职。因为，我的兄弟，我们特别为之奋斗并应当为之奋斗的是，要谨慎地尽可能维护主所交付并通过祂的宗徒、即我们这些继承者所交付的合一，并且尽我们所能，将那些被某些人的任性派别和异端诱惑从母亲教会中分开的分散和迷失的羊群聚集到教会中；只有那些因顽固和疯狂而坚持、不愿回到我们这里的人被留在外面；他们自己将不得不向主交代他们所造成的分裂和分离，以及他们所离弃的教会。\par

4. \Person{比尔}，至于某些司铎和\Person{斐里奇西默}在此地的事端，为使你知晓此次所作所为，我们的同仁已寄给你亲笔签署的信函，好让你在听取双方陈述后，从他们的信件中得知他们的意见和宣判。然而，兄弟，你若也吩咐将我最近通过同仁\Person{卡尔多尼}和\Person{福尔图纳图斯}寄给你的信件的副本宣读出来，以共慰藉，那便更好；这些信件是我就同一\Person{斐里奇西默}及其司铎团的事写给那边圣职人员的，也写给那边的教友，要在弟兄们面前宣读；表明你的祝圣以及整个事件的经过，好使那边的弟兄团体和此地的弟兄团体都能通过我们知晓一切。此外，我亦通过我所派遣的副执事\Person{梅蒂乌斯}和辅祭\Person{尼塞福鲁斯}，随信附上这些信件的副本。至爱的兄弟，我愿你常怀喜乐安康。\par

\SourceRecord{Translated by Robert Ernest Wallis.；Ante-Nicene Fathers；Vol. 5.；Edited by Alexander Roberts, James Donaldson, and A. Cleveland Coxe.；Buffalo, NY: Christian Literature Publishing Co.,；1886.；<http://www.newadvent.org/fathers/050641.htm>.}

% structure: p0001 source=h1 target=section reason=page-title

\section{第四十二封书信}

\Emph{致同一人，关于他已发送信件给被\Person{诺瓦提安}诱惑的宣信者一事。}\par

\Emph{论题——本信之论题已充分见于标题。显然，本信与次信由同一位信使送出。}\par

\Person{西彼廉}致其兄弟\Person{高尔乃略}安好。至爱的弟兄，我以为我有义务且你必须收到一封短信，致那些与你同在、被\Person{诺瓦提安}与\Person{诺瓦图斯}的固执和邪恶所诱惑、已离开教会的宣信者；在此信中，我可劝勉他们，为着我们彼此的爱，回归他们的母亲，即天主教会。我首先将此信托付给副执事\Person{梅提乌斯}送交你审阅，以免有人声称我写的与信中内容不符。此外，我已吩咐由我派往你处的同一\Person{梅提乌斯}，他应受你决断的指引；你若认为此信应交给宣信者，则他应交付。至爱的弟兄，祝你始终由衷平安。\par

\SourceRecord{Translated by Robert Ernest Wallis.；Ante-Nicene Fathers；Vol. 5.；Edited by Alexander Roberts, James Donaldson, and A. Cleveland Coxe.；Buffalo, NY: Christian Literature Publishing Co.,；1886.；<http://www.newadvent.org/fathers/050642.htm>.}

% structure: p0001 source=h1 target=section reason=page-title

\section{第四十三封书信}

\Emph{致罗马的宣信者：他们应回归合一。}\par

\Emph{要旨。——他劝勉那些被\Person{诺瓦提安}与\Person{诺瓦图斯}的派别诱惑的罗马宣信者们回归合一。}\par

\Person{西彼廉}致\Person{玛克西穆}、\Person{尼科斯特拉图}及其他宣信者，问安。正如你们经常从我的书信中所领会的那样，亲爱的诸位，我一向在言语中表达对你们\OriginalTerm{宣认}{confession}的何等敬意，以及对你们所属的兄弟团体何等的爱；我恳求你们相信并顺从这封信，我写此信是本着纯朴与忠诚为你们、为你们的行为、为你们的赞誉而考虑。因为有一件事使我心力交瘁且忧伤难忍，一种被打击的、几乎跌倒的无法忍受的悲伤攫取了我，就是我发现你们在那里，违背教会的秩序，违背福音的法律，违背\OriginalTerm{公教会体制}{Catholic institution}的合一，竟然同意另立一位主教。那是既不义也不允许的事；另立一个教会；撕裂基督的肢体；以分裂的嫉妒撕裂主羊群的一心一体。我恳求至少在你们当中，那不合法的弟兄分裂不要再继续；但要记住你们的宣认和神圣的传统，回到你们所离开的母亲那里；你们从她那里出来，在同一个母亲的欣喜中达到了宣认的光荣。不要以为你们这样就是持守基督的福音，当你们离开基督的羊群、离开祂的平安与和好时；因为更合适的是，光荣而优秀的士兵们坐在自己的营中，这样在营中共同处理并照管那些需要共同处置的事情。因为我们的同心同德决不可分裂，且因为我们不能离弃教会而到外面去找你们，我们以尽可能的劝勉恳求你们，宁愿回到你们的母亲教会和我们兄弟团体。我祝愿你们，最亲爱的弟兄们，永远热忱地安好。\par

\SourceRecord{Translated by Robert Ernest Wallis.；Ante-Nicene Fathers；Vol. 5.；Edited by Alexander Roberts, James Donaldson, and A. Cleveland Coxe.；Buffalo, NY: Christian Literature Publishing Co.,；1886.；<http://www.newadvent.org/fathers/050643.htm>.}

% structure: p0001 source=h1 target=section reason=page-title

\section{\WorkTitle{书函四十四}}

\Emph{致科尔乃略，论阿德鲁梅提内的颇利卡普。}\par

\Emph{论点：——他在此书中为自己所发生之事辩护，即当他在阿德鲁梅提乌姆期间，颇利卡普的圣职人员曾从该地致信，并非写给科尔乃略，而是写给罗马的圣职人员；尽管此前颇利卡普本人更常写信给科尔乃略。从上下文本身已足够清楚地看出，此书是在前述诸书之后写成的。}\par

1. 西彼廉致其兄弟科尔乃略安好。我已阅读了你亲爱的兄弟，由我们的同司铎普利米提乌斯带来的书信，在其中我察觉到你因如下事而感到不快：尽管来自阿德鲁梅提内殖民地的信件是以颇利卡普的名义寄给你的，但在利贝拉利斯与我来到那地方之后，信件便开始从那里寄给司铎和执事们。\par

2. 关于此事，我希望你知道，并且确实相信，这并非出于轻率或藐视。然而，当我们的几位同僚聚集在一处，决定在派遣我们的同主教卡尔多尼乌斯和福图纳图斯作为使节到你那里去的同时，所有事情都应暂时维持现状，直到那些同僚们使那里的事情归于和平，或查明其真相后回到我们这里；留在阿德鲁梅提内殖民地的司铎和执事们，在我们的同主教颇利卡普缺席的情况下，并不知道我们所共同决定的事。但当我们来到他们面前，我们的意图被了解后，他们自己也开始遵行其他人所做的事，以至于在那里存留的教会之间的合一毫无损伤。\par

3. 然而，有些人有时会以他们的言语扰乱人们的心思和精神，因为他们所叙述的事情与事实不符。因为我们，凡从此地航行者，无不提供计划，使他们能不冒任何风险地航行，我们知道我们曾劝勉他们承认并持守天主教会的根源和基质。但由于我们的行省幅员辽阔，且努米底亚和毛里塔尼亚与之相连；唯恐在城市中发生的分裂会以不确定的意见扰乱远方之人的心思，我们决定——通过主教们获得事情的真相，并为你的祝圣得到更大的权威证明，并最终为消除每个人心中的任何疑虑——由所有在该行省任何地方的人寄信给你，正如事实上所做的，使我们所有的同僚都能明确地认可并维护你和你的共融，亦即同时维护天主教会的合一及其爱德。这一切已借着天主的引导成就，我们的计划在天主眷顾下得以推进，我们为此欢欣。\par

4. 因为如此，你主教职的真理和尊贵在最公开的光明中，并以最显明和实质的认可得以确立；以致从我们同僚的回函中（他们曾从那里写信给我们），以及从我们的同主教庞培、斯特凡努斯、卡尔多尼乌斯和福图纳图斯的叙述和见证中，你祝圣的必要原因、正确秩序，以及你那令人荣耀的无辜，都得以被众人知晓。为使我们可以与其余同僚坚定而稳固地履行此职，并在天主教会和谐一致中保持合一，天主的俯就将成就此事；好让那屈尊在祂的教会中拣选并任命司铎的主，也以祂的善意和助佑保护那些已被拣选和任命的人，激励他们治理，并供给他们力量以约束恶人的悖逆，提供温柔以眷顾失足者的悔改。亲爱的兄弟，愿你衷心平安。\par

\SourceRecord{Translated by Robert Ernest Wallis.；Ante-Nicene Fathers；Vol. 5.；Edited by Alexander Roberts, James Donaldson, and A. Cleveland Coxe.；Buffalo, NY: Christian Literature Publishing Co.,；1886.；<http://www.newadvent.org/fathers/050644.htm>.}

% structure: p0001 source=h1 target=section reason=page-title

\section{第四十五封书信}

\Emph{科尔乃略致西彼廉：论宣信者回归合一。}\par

\Emph{主旨。——科尔乃略告知西彼廉关于宣信者庄严回归教会，并加以描述。}\par

1. \Person{科尔乃略}致其兄弟\Person{西彼廉}，问安。鉴于我们因那些宣信者所承受的忧虑与不安——他们曾被那狡诈诡谲之人的诡计与恶意所蒙蔽，几乎被欺骗、被离间，远离了教会——当我们看到他们承认错误，认清那恶人如同毒蛇般的毒计，并一心一意地以纯朴的意志（正如他们所一致承认的）回到他们所离开的教会时，我们心中充满了喜乐，并向全能的天主和我们的主基督献上感谢。首先，我们那些信德得到认可、热爱和平、渴望合一的弟兄们宣告，这些人的傲慢已平息；然而，没有足够的把握让我们轻易相信他们已彻底改变。但随后，宣信者\Person{乌尔巴诺}和\Person{西多尼乌}来到我们的司铎那里，坚称宣信者兼司铎\Person{马克西穆}同样渴望回到教会；但由于此前发生过许多他们策划的事，我们的同牧主教们和我的信件也已告知你，因此我们不能轻易信任他们；我们决定亲耳从他们口中听到他们通过信使所传告的那些事。当他们来到，司铎们要求他们说明所做之事，并指控他们不久前以他们的名义通过所有教会多次发送充满诽谤和责难的信件，几乎扰乱了所有教会；他们声称自己被欺骗，不知道那些信件的内容；只是由于被误导，他们也犯了分裂的行为，成了异端的始作俑者，以至于他们让人按手在他身上如同按手在主教身上一样。当这些和其他事项被指控后，他们恳求这些事情能被取消并完全从记忆中抹去。\par

2. 因此，我将整件事提交到我面前，决定召集司铎会议（当时有五位主教在场，今日也同样在场），以便通过深思熟虑的商议，取得一致同意，决定如何处理他们个人之事。为使你知晓所有人的意见和每个人的建议，我也决定将我们各种意见告知你，你将在下文读到。这些事情完成后，\Person{马克西穆}、\Person{乌尔巴诺}、\Person{西多尼乌}以及几位与他们联合的弟兄来到司铎会议，恳切祈求之前的事能被遗忘，不再提及；并\Emph{承诺}今后，如同什么事也未曾发生或说过，双方互相原谅，他们将向天主呈献洁净纯全的心，遵循福音的话说：\BibleQuote{心里洁净的人是有福的，因为他们要看见天主。}\BibleRef{玛 5:8}剩下的事是通知全体教友这一整件事，使他们看到这些人——他们曾长久哀叹为迷途离散的人——如今在教会中站稳了。他们的意愿为人所知后，弟兄们的大集会聚集起来。众人同声感谢天主；大家以眼泪表达心中的喜乐，拥抱他们如同他们今日才从牢狱的刑罚中被释放。引用他们自己的话说：\Quote{我们，他们说，知道\Person{科尔乃略}是由全能天主和我们的主基督所选立的至圣天主教会的主教。我们承认我们的错误；我们被欺骗；我们被诡诈的背信和饶舌所蒙骗。因为虽然我们似乎曾与一个分裂者和异端者有过某种共融，但我们的心在教会中始终是纯正的。因为我们并非不知道只有一个天主；只有一位我们曾承认的主基督；只有一位圣神；而且在天主教会中应该只有一位主教。}我们难道不正是因他们的这一宣认而合理地允许，他们在世俗权力面前所宣认的，如今在教会中得以肯定吗？因此，我们命令司铎\Person{马克西穆}回到他的位置；其余的人我们以极大的民众赞许接纳。但我们将一切交托给全能的天主，一切都在祂的权能之下。\par

3. 因此，兄弟，我们于同一时辰、同一时刻将这一切写成信寄给你；我已立即派遣辅祭\Person{尼塞福禄}，迅速赶去登船，以便毫不耽搁，你犹如亲临该司铎会议和民众集会，向全能天主和我们的主基督谢恩。但我们相信——不，我们确信——其他同样陷入此错误的人不久也会回到教会，当他们看见他们的首领与我们一同行动。兄弟，我认为你也应当将这些信函寄给其他教会，使众人皆知这位分裂者和异端者的狡诈与欺骗日渐归于无有。亲爱的兄弟，祝安好。\par

\SourceRecord{Translated by Robert Ernest Wallis.；Ante-Nicene Fathers；Vol. 5.；Edited by Alexander Roberts, James Donaldson, and A. Cleveland Coxe.；Buffalo, NY: Christian Literature Publishing Co.,；1886.；<http://www.newadvent.org/fathers/050645.htm>.}

% structure: p0001 source=h1 target=section reason=page-title

\section{第四十六封书信}

\Emph{西彼廉答复高内略，祝贺宣信者从裂教中回归。}\par

\Emph{论据：——他祝贺宣信者回归教会，并提醒他们这一回归对天主教会何等有益。}\par

1. 西彼廉致弟兄高内略安。至亲爱的弟兄，我承认我不停地向全能的天主圣父、并向祂的基督、我们的主、天主与救主，献上最大的感谢，因为教会如此受到天主的保护，其合一与圣洁并未因背信的顽固和异端的邪恶而不断或完全败坏。我们读了你的书信，因着我们共同愿望的\Emph{实现}而满心喜乐地领受了最大的欢欣：即司铎马克西穆、宣信者乌尔巴诺，与息多尼奥及玛加略，已重新进入天主教会；也就是说，他们抛弃了错误，放弃了裂教的——不，是异端的——疯狂，以健全的信德重新寻求合一与真理的家园；他们从那里出发走向光荣，如今光荣地返回；那些曾宣认基督的人，此后不应再抛弃基督的营帐，那些在力量与勇气上未曾被征服的人，不应再试探其爱德与合一的信德。看，他们赞颂的无瑕纯洁；看，这些宣信者未受玷污的稳固尊严：他们离开了背弃者与逃亡者，他们离开了信仰的背叛者和天主教会的攻击者。有理由，如你所写，人民与弟兄团体以最大的喜乐迎接他们回归；因为在那些保持了光荣并回归合一的宣信者的荣耀中，没有人不把自己视为同伙与分享者。\par

2. 我们可以从自身的情感估量那日的喜乐。因为如果在此地，全体弟兄因你寄来的关于他们宣信的书信而喜乐，并以最大的热忱接受了这共同喜乐的消息，那么当事情本身与普遍的喜悦在众人眼前展开时，那里的喜乐该是何等大呢！因为主在福音中说：\BibleQuote{对于一个悔改的罪人，在天上也有极大的喜乐。}那么，对于那些带着光荣与赞美回归天主教会、并以自己信德与榜样的认可为他人开辟回归之路的宣信者，在地上，不仅在天上，该是何等更大的喜乐呢！因为这错误曾误导我们某些弟兄，使他们以为自己在追随宣信者的共融。当这错误被除去后，光明注入了众人的心中，天主教会被显明为唯一的，既不能被切割也不能被分裂。如今，再没有人容易被一个狂怒的裂教者的空谈所欺骗，因为已经证明，基督良善而光荣的兵士，不能长久地被他人的诡诈与背信拘留在教会之外。至亲爱的弟兄，祝你永远安康。\par

\SourceRecord{Translated by Robert Ernest Wallis.；Ante-Nicene Fathers；Vol. 5.；Edited by Alexander Roberts, James Donaldson, and A. Cleveland Coxe.；Buffalo, NY: Christian Literature Publishing Co.,；1886.；<http://www.newadvent.org/fathers/050646.htm>.}

% structure: p0001 source=h1 target=section reason=page-title

\section{\WorkTitle{第四十七书}}

\Emph{\Person{科尔乃略}致\Person{西彼廉}，论\Person{诺瓦提安}及其党羽的分裂。}\par

\Emph{内容提要：——\Person{科尔乃略}向\Person{西彼廉}报告\Person{诺瓦提安}分裂之事。}\par

\Person{科尔乃略}致其兄弟\Person{西彼廉}，问安。为使这可怜之人受天主权势推倒后（在你们驱逐\Person{玛克西木}、\Person{隆基努}和\Person{玛凯乌}之后），他再度兴起；正如我在前信（由宣信者\Person{奥根都}带给您）中所暗示，我想\Person{尼科斯特拉图}、\Person{诺瓦都斯}、\Person{艾瓦里斯特}、\Person{普里穆斯}和\Person{狄约尼削}已经到达那里。因此应当注意，要将此事告知我们所有的同僚主教和弟兄们：\Person{尼科斯特拉图}被指控诸多罪行，不仅在他管理其世俗女恩主的事务时犯下欺诈和掠夺，而且还（这为他保留为永罚）盗走了教会不少存款；\Person{艾瓦里斯特}是分裂的始作俑者；\Person{泽图斯}已被任命为主教接替他，成为他先前所牧之百姓的继任者。但那个分裂者和异端者，凭其恶毒和贪得无厌的邪恶，策划了比他在自己人中一直实行的事更大更坏的事，好叫你们知道，那分裂者和异端者始终有怎样的领袖和保护者附在其侧。亲爱的兄弟，我愿您永远衷心安康。\par

\SourceRecord{Translated by Robert Ernest Wallis.；Ante-Nicene Fathers；Vol. 5.；Edited by Alexander Roberts, James Donaldson, and A. Cleveland Coxe.；Buffalo, NY: Christian Literature Publishing Co.,；1886.；<http://www.newadvent.org/fathers/050647.htm>.}

% structure: p0001 source=h1 target=section reason=page-title

\section{第四十八封书信}

\Emph{\Person{西彼廉}答\Person{科尔乃略}，论\Person{诺瓦图斯}之罪行。}\par

\Emph{论旨。—— \Person{西彼廉}称赞\Person{科尔乃略}，因\Person{科尔乃略}及时警告了他：就在有罪派系前来见他的次日，他就收到了\Person{科尔乃略}的信。然后他详细描述了\Person{诺瓦图斯}的罪行，以及此前\Person{诺瓦图斯}在\Place{阿非利加}挑起的裂教。}\par

1. \Person{西彼廉}致他的弟兄\Person{科尔乃略}，安好。亲爱的弟兄，你既勤勉又友爱，速派辅祭\Person{尼塞福鲁}来我们这里，他既告诉了我们关于宣信者返回的光荣喜讯，又最充分地指导我们抵御\Person{诺瓦西安}和\Person{诺瓦图斯}为攻击基督的教会而设的新而恶毒的诡计。因为就在前一天，那异端邪恶的恶毒派系已抵达这里，它本身已经丧亡，并准备毁灭那些加入它的人；而在第二天，\Person{尼塞福鲁}带着你的信到达了。从这信我们既自己得知，也开始教导和训导别人：\Person{厄瓦里斯托}从主教的身份，现已连平信徒也不如；反而被逐出教座和民众，自基督的教会流放，在远处各省份漂泊；他自己既已使真理和信仰沉船，就为那些像他的人预备可怕的沉船。此外，\Person{尼科斯特拉托}失掉了神圣职务的执事职，因他以亵圣的欺诈窃取了教会的钱财，并否认了寡妇孤儿的存款；他并不那么想来到\Place{阿非利加}，而是想从城中逃往那里，出于对自己劫掠和可怕罪行的意识。如今，他作为教会的逃兵和流亡者，仿佛改变气候就能改变人似的，继续夸耀并自称为宣信者，虽然\Emph{他}再也不能，也不可被称为基督的宣信者，因他否认了基督的教会。因为当宗徒\Person{保禄}说：\BibleQuote{为此，人要离开父亲和母亲，依附自己的妻子，二人成为一体。这奥秘真是伟大，但我是指基督和教会说的。}\BibleRef{弗 5:31}——我说，当有福的宗徒说这话，并以他的圣声为基督与教会的合一作证，彼此以不可分割的纽带相依附时，那不与基督的净配及祂的教会同在的人，怎能与基督同在呢？那抢夺并亏负了基督教会的人，又怎能自任治理教会的职责呢？\par

2. 至于\Person{诺瓦图斯}，你们无需向我们提到什么，因为\Person{诺瓦图斯}更应由我们向你们指出：他总贪图新奇，以无尽贪欲的掠夺而暴怒，因膨胀骄傲的傲慢和愚蠢而自大；在那里的主教中名声狼藉；总被所有司铎齐声谴责为异端者和背信之人；总好打听，以便出卖；他阿谀以欺骗，从不忠信以便爱；是点燃叛乱火焰的火炬和火；是使信仰沉船的旋风与风暴；是安静的敌人、宁静的对手、和平的仇敌。最后，当\Person{诺瓦图斯}从你们那里离开，也就是说，当那风暴和旋风离去，那里部分地恢复了平静；因他煽动而离开教会的荣耀而良善的宣信者，在他离开城市后，回到了教会。正是这个\Person{诺瓦图斯}，首先在我们中间播下了纷争和裂教的火焰；他使这里的弟兄们与主教分离；在迫害本身中，他对于我们的人民来说，仿佛是另一场迫害，为了颠覆弟兄们的心智。正是他，未经我的许可和知情，出于他的派系性和野心，任命他的随从\Person{费利西西莫}为执事，并以他自己的风暴航行到\Place{罗马}去颠覆教会，在那里试图做类似和同样的事，强行将一部分民众与神职人员分开，并分裂那紧密结合、彼此相爱的兄弟情谊的和谐。既然\Place{罗马}因其伟大而显然应优先于\Place{迦太基}，他在那里犯下了更重大、更严重的罪行。他在一处违反教会任命了执事，在另一处却任命了主教。任何人不应对这类人如此感到惊讶。恶人总是被他们自己狂暴的情欲疯狂地冲昏头脑；犯下罪行之后，他们被败坏良心的意识所激动。那些无论是在生活的言谈还是品格的平安中都没有持守教会神圣纪律的人，也不能留在天主的教会中。被他掠取的孤儿、被他欺骗的寡妇，以及被扣留的教会钱财，正从他身上追讨我们所看见在他疯狂中降临的惩罚。他的父亲也饿死在街头，死后甚至也没有被他埋葬。他妻子的子宫被他的脚跟击中；在随之而来的流产中，胎儿产出，成了父亲谋杀的果实。而今，他竟敢谴责那些献祭之人的手，他自己却在脚上更加有罪，因为即将出生的儿子正是被那脚所杀？\par

3. 他早已惧怕这罪行的意识。因此他确信，他不仅会被逐出司铎团，还会被禁止共融；由于弟兄们的紧迫要求，调查的日期即将到来，届时他的案件将要在我们面前处理，若不是迫害阻止的话。他欢迎这一点，带着一种逃避和规避定罪的欲望，犯下了所有这些罪行，并引起了所有这些骚动；以至于那个将被逐出教会、被排除教会的人，以自愿离开来抢先司铎的审判，仿佛抢先判决就能免于刑罚似的。\par

4. 至于其他弟兄，我们为他们被他欺骗而忧伤，我们努力使他们能避开那狡猾骗子的有害邻舍，逃脱他诱惑的致命网罗，再次寻求那他被天主的权威逐出的教会。确实，赖主的助佑，我们相信他们能因祂的慈悲而归回，因为一个人除非显然必须丧亡，否则不会丧亡，因为主在祂的福音中说：\Quote{凡不是我天主圣父所栽种的，必要被拔除。}\BibleRef{玛 15:13} 唯独那没有被天主圣父的诫命和警告所栽种的人，才能离开教会；唯独他能离弃主教们，与分裂分子和异端者一同固守于疯狂之中。但天主圣父的慈悲、我主基督的宽仁、以及我们自己的忍耐，必将其他人与我们联合。亲爱的弟兄，祝你永远安康。\par

\SourceRecord{Translated by Robert Ernest Wallis.；Ante-Nicene Fathers；Vol. 5.；Edited by Alexander Roberts, James Donaldson, and A. Cleveland Coxe.；Buffalo, NY: Christian Literature Publishing Co.,；1886.；<http://www.newadvent.org/fathers/050648.htm>.}

% structure: p0001 source=h1 target=section reason=page-title

\section{Epistle 49}

\Emph{马克西穆斯与其他宣信者致西彼廉，论他们脱离分裂归来。}\par

\Emph{论述。——他们告知西彼廉，他们已返回教会。}\par

\Person{马克西穆斯}、\Person{乌尔巴诺斯}、\Person{西多尼奥斯}、\Person{马卡里乌斯}，致他们的弟兄\Person{西彼廉}，安。\newline{}我们确信，最亲爱的弟兄，你也同样真诚地与我们一同喜乐，因为我们经过商议，尤其顾全教会的利益与平安，搁置一切其他事项，将它们留待天主的审判，与我们的主教\Person{科尔内略}以及全体圣职人员重归于好。你务须从我们这封信中得知，此事是在整个教会的喜乐中，甚至以弟兄们踊跃的爱心成就的。我们祈祷，最亲爱的弟兄，愿你长久安康。\par

\SourceRecord{Translated by Robert Ernest Wallis.；Ante-Nicene Fathers；Vol. 5.；Edited by Alexander Roberts, James Donaldson, and A. Cleveland Coxe.；Buffalo, NY: Christian Literature Publishing Co.,；1886.；<http://www.newadvent.org/fathers/050649.htm>.}

% structure: p0001 source=h1 target=section reason=page-title

\section{\WorkTitle{第五十封书信}}

\Emph{西彼廉致宣信者，祝贺他们从分裂中归来。}\par

\Emph{论点——西彼廉祝贺罗马的宣信者回归教会，并回复他们的信件。}\par

\Person{西彼廉}致司铎\Person{马克西莫}，以及他的弟兄\Person{乌尔巴诺}、\Person{西多尼乌}和\Person{马图雷斯特}，问候。最亲爱的弟兄们，当我读到你们写给我的关于你们回归、教会和平及弟兄和解的信件时，我承认，我内心的喜悦与我之前得知你们宣信的荣耀、感恩地听到你们属天和属灵争战声望时的喜悦一样巨大。因为这也是你们信德和赞美的另一项宣信：承认教会是唯一的，不参与他人的错误，更确切地说是邪恶；从同一营地再次出发，你们曾以最强劲的力量从那里跃出以进行战斗并征服对手。因为战场的战利品应当带回那里，那里曾接收了战斗的武器，免得基督的教会不保留这些为荣耀而装备的光荣战士。然而现在，你们在主的平安中保持了信德的适当准则和不可分割的爱德与和谐的法则，并以你们的行径为他人树立了爱德与和平的榜样；如此，教会的真理和我们所持守的福音奥迹的合一，也因你们的同意和纽带而联系在一起；基督的宣信者不会成为错误的领袖，因为他们曾作为美德与荣誉的可称赞的创始者挺身而出。\par

让别人去考虑他们能多么祝贺你们，或每个人多么为自己夸耀吧：我承认，对于你们这次和平的回归和爱德，我更祝贺你们，也更向他人夸耀你们。因为你们应当单纯地听听我内心的想法。我曾十分悲伤，极其痛苦，因为我不能与那些我曾开始爱过的人共融。当分裂和异端的错误在你们出狱后抓住你们时，你们的荣耀似乎留在了地牢里。因为当基督的士兵没有从监狱回到教会，尽管他们曾带着教会的赞美和祝贺进入监狱时，你们名字的尊严似乎留在了那里。\par

因为虽然在教会中似乎有莠子，但我们的信德和爱德不应受阻，以至于我们因看到教会中有莠子就自己离开教会：我们只应努力成为麦子，当麦子开始被收进主的仓库时，我们就能为我们的劳苦和工作收获果实。宗徒在他书信中说：\BibleQuote{在大户人家，不但有金器银器，也有木器瓦器；有作为贵重的，有作为卑贱的。}\BibleRef{弟后 2:20} 让我们努力，最亲爱的弟兄们，尽我们所能地劳苦，使我们成为金器或银器。但只有主有权打破瓦器，因为祂也被赐予了铁杖。仆人不能大过主人，也没有人能为自己索取父唯独赐给子的东西，以至于认为他可以拿起簸箕来扬净场，或凭人的判断把所有的莠子从麦子中分离出来。那是骄傲的顽固和亵渎的僭越，是堕落疯狂自取的。而有些人总是为自己索取比温和正义的要求更多的权柄，他们从教会中灭亡；当他们傲慢地自高自大，被自己的膨胀蒙蔽时，他们失去了真理的光。为此，我们保持节制，考虑主的天平，思念天主父的慈爱和怜悯，长期仔细地在心里掂量，并权衡以适度的节制该做什么。\par

所有这些事你们可以深入察看，如果你们愿意阅读我最近在这里宣读的论著，我也为了我们彼此的爱德，将它们转交给你们阅读；其中对于失足者，既不缺少责备的惩戒，也不缺少医治的药方。此外，我微薄的能力也尽其所及表达了天主教会的合一。我现在更加相信这篇论著会使你们喜悦，因为你们现在以如此方式阅读它，既赞同又喜爱；因为我们在文字中所写的，你们在行为上成就了，当你们在爱德与和平的合一中回到教会时。最亲爱的弟兄们，我衷心诚挚地祝愿你们健康，深切思念你们。\par

\SourceRecord{Translated by Robert Ernest Wallis.；Ante-Nicene Fathers；Vol. 5.；Edited by Alexander Roberts, James Donaldson, and A. Cleveland Coxe.；Buffalo, NY: Christian Literature Publishing Co.,；1886.；<http://www.newadvent.org/fathers/050650.htm>.}

% structure: p0001 source=h1 target=section reason=page-title

\section{书信五十一}

\Emph{致\Person{安托尼亚努斯}：论\Person{科尔乃略}与\Person{诺瓦替安}}\par

\Emph{论述。——当\Person{安托尼亚努斯}收到\Person{诺瓦替安}的信，开始倾向他的一党时，\Person{西彼廉}坚定他原有的意见，即继续与他的主教并因此与天主教会共融。他为自身在关于失足者一事上意见的改变辩护，最后他阐明\Person{诺瓦替安}异端的本质所在。}\par

\Person{西彼廉}致其弟兄\Person{安托尼亚努斯}，问安。至爱的弟兄，你坚定地维护司铎团的和谐，并依附天主教会，你告知我你并未与\Person{诺瓦替安}共融，而是遵从我的劝告，与我们的同主教\Person{科尔乃略}保持同一共识。此外，你要求我将同一封信的副本转交给我们的同僚\Person{科尔乃略}，以便他能放下一切忧虑，并立即知道你是与他——即与天主教会——共融的。\par

2. 但随后我们同司铎\Person{昆图斯}带来了你的另一封信，我在其中观察到你的心意受到\Person{诺瓦替安}的信影响，开始动摇。因为虽然你先前已坚定地确定了你的意见和同意，你在这封信中却要求我再写信告诉你，\Person{诺瓦替安}引入了什么异端，或者\Person{科尔乃略}是凭什么理由与\Person{特洛斐摩}及献祭者共融。在这些事上，如果你确实出于信仰的热忱而焦虑谨慎，并勤勉探求一件可疑之事的真相，那么一个在敬畏天主中尚未决定的犹豫不安的心意，是不应受谴责的。\par

3. 然而，我看到你在信中的第一个意见之后，后来受到\Person{诺瓦替安}的信的搅扰，我首先声明这一点，至爱的弟兄：庄重的人，以及一旦以坚固的坚毅建立在稳固磐石上的人，不会，我说，被微风摇动，甚至被狂风暴雨所动，免得他们善变而不定的心意被各种意见——就像一阵阵狂风冲击他们那样——反复动摇，因而带着些许轻率的责备改变主意。为了使\Person{诺瓦替安}的信不会对你或任何人造成这种情况，我将如你所愿，我的弟兄，以寥寥数语向你陈述此事。而且首先，既然你似乎也对我所做的事感到困扰，我必须在你眼前澄清我本人和我自己的原因，免得有人以为我轻率地改变了我的初衷，并且在起初的开始我保持了福音的勇毅，但后来似乎使我的心意转变，脱离了纪律和从前判断的严厉，以致认为那些用证书玷污了自己良心，或献了可憎祭的人，应该轻易地得到平安。这两件事我都做了，并非没有经过长期权衡和深思熟虑的理由。\par

4. 因为当战斗仍在进行，荣美争战的冲突在迫害中炽烈时，必须以各种劝勉和充分的紧迫激励士兵的勇气，尤其是要以我声音的号角唤醒失足者的心意，使他们不仅以祈祷和哀恸走上悔改的道路；而且，由于有机会重复争战并重获救恩，使他们受到我声音的斥责，并被激发而趋向于更热烈的宣认和殉道的光荣。最后，当司铎和执事关于某些人写信给我，说他们毫无节制、急切地要求领受共融时，我在我现存的那封信中答复他们，当时还加上了这句话：\Quote{如果这些人如此过分急切，他们自己有权得到他们所要求的，因为时机本身为他们提供了比他们要求的更多的东西：战斗仍在进行，争战天天举行；如果他们真正而实在地为所作的事悔改，并且他们信德的热忱占上风，那不能被耽误的人可以得到荣冠。}但我推迟决定如何处理失足者的事，为的是当安宁和宁静被赐予，当天主的上智允许主教们聚在一处时，然后经过传阅和权衡从所有意见中收集的建议，我们确定必须要做什么。但如果任何人在我们的会议之前，并且在全体建议所决定的意见之前，鲁莽地愿意与失足者共融，他自己应被阻止领受共融。\par

5. 我也将这十分详尽地写信到罗马，给那时仍在没有主教的情况下行事的圣职人员，以及给宣信者、当时仍被关在监狱但现在已在教会内、与\Person{科尔乃略}联合的司铎\Person{玛克西穆}和其他人。你可以从他们的回信知道我写了这些，因为他们信上这样写：\Quote{然而，你所亲自就如此重要之事宣布的，令我们满意，即必须先维护教会的和平；然后，聚集会议，包括主教、司铎、执事、宣信者以及站立的平信徒，我们应处理失足者的事。}此外还加上——当时\Person{诺瓦替安}撰写并亲口宣读他所写的，而司铎\Person{莫伊塞斯}（当时仍是宣信者，但现在是殉道者）签字——应当将平安赐给那些患病并离世在即的失足者。这封信被传送到整个世界，并被所有教会和所有弟兄所知悉。\par

6. 然而，依照早先的决定，当迫害平息、聚会的机会得以恢复时，我们聚集了许多主教——他们的信德与天主的保护使他们保持健全与安全。双方陈明了圣经，我们以健全的节制平衡了决议：既不将共融与和平的希望完全拒绝给跌倒者，以免他们因绝望而更加堕落，并因教会对他们关闭而如外邦人一样生活在世俗中；也不放松福音的责戒，以致他们轻率地冲向共融，而是应延长补赎期，以哀伤的心情恳求慈父的宽恕，并根据一本小册子（我相信它已送达你们）中所载的内容，逐一审查每个人的情况、愿望和需要，其中汇集了我们决议的若干要点。此外，恐怕非洲主教的人数似乎不足，我们也就此写信给我们在罗马的同僚主教\Person{科尔乃略}，他本人也与众多主教举行了会议，以同样的庄重和健全的节制，认同了我们所持的相同意见。\par

7. 关于此事，如今有必要写信给你们，使你们知道我没有轻率行事，而是按照我之前在信中所包含的，将一切推迟到我们会议的共同决定，并且确实尚未与任何跌倒者共融，只要仍有使他们不仅获得宽恕、而且获得冠冕的途径。但后来，由于我们同僚的一致意见以及聚集弟兄并医治他们创伤的利益需要，我顺应了时势的必要，认为必须为许多人的得救着想；我现在并不退离我们会议中已共同决定的事，尽管有许多声音散布各种言论，从魔鬼的口中说出反对天主司铎的谎言，并且到处流传，以破坏大公合一的和谐。但身为好兄弟和同心合意的司铎，你们不宜轻易听信恶人和背教者可能说的话，而应通过考察我们的生活与教导，仔细衡量我们那些谦逊庄重的同僚所做之事。\par

8. 现在，最亲爱的弟兄，我要谈到我们同僚\Person{科尔乃略}的品格，以便你们能更公正地认识他——不是根据恶人和毁谤者的谎言，而是根据上主天主的判断（祂立他为主教），以及他全体同僚主教的见证——全世界的主教以绝对一致的态度赞同了这一点。因为——这件事以可赞美的宣告将我们最亲爱的\Person{科尔乃略}推荐给天主、基督、祂的教会以及全体同僚司铎——他并非突然获得主教职位；而是经历了教会所有职务的晋升，在神圣的管理中多次为主立功，藉着宗教服务的所有阶梯登上了司祭职的高峰。而且，他既没有要求主教职位，也不渴望它；也非像其他人在骄傲与自大的膨胀中所做的那样，强行夺取；相反，他素来安静、温和，正如那些被天主选召担任此职的人所惯有的，他顾念自己童贞节制的谦逊以及天生且持守的恭敬谦卑，并没有像某些人那样强求被立为主教，而是他自己忍受了强迫，被迫接受了主教职务。他是被当时在罗马城的许多同僚主教立为主教的，他们写信给我们论及他的祝圣，信件充满荣誉与赞美，并以突出的见证来宣告他。此外，\Person{科尔乃略}是由天主及其基督的判断、几乎全体圣职人员的见证、当时在场的民众的投票以及古代司铎与善人的集会而立为主教的——在他之前无人如此被立，当\Person{法比盎}的位置，即\Person{伯多禄}的位置与司祭宝座的等级空缺之时；这一位置既被天主的旨意占据，且由我们全体的一致同意所确立，如今任何想要成为主教的人，必须从外部产生；凡不持守教会合一的人，不能拥有教会的祝圣。无论他是谁，即使大大自夸并为自己要求甚多，他是亵渎的，是外方人，是在外的。正如在第一位之后不能有第二位，任何在应当唯一的人之后被立的人，不是他的第二，而实际上什么都不是。\par

9. 后来，当他接受了主教之职——不是通过恳求或强求，而是通过那设立司铎的天主的旨意——在他履行主教职务时，有着何等的德行，何等的刚毅，何等坚定的信德——这是我们应当以纯朴之心仔细审视并加以赞美的事——他无畏地坐在\Place{罗马}的司铎椅上，当时，一位令天主的司铎憎恶的暴君，正威胁着那些可说的和不可说的事，因为他会更耐心、更容忍地听到一位敌对君王被兴起反对他，而不是听到一位天主的司铎在\Place{罗马}被设立。最亲爱的弟兄，难道不应该以最高的德行和信德的见证来称赞这个人吗？难道他不该被列于光荣的宣信者和殉道者之中吗？他长久地坐着，等待那些要撕裂他身体的人，以及那残暴暴君的复仇者；当\Person{科尔乃略}抵抗他们的致命谕令，并以他信德的力量践踏他们的威胁、痛苦和酷刑时，他们要么用剑攻击他，要么将他钉在十字架上，要么用火焚烧他，要么以某种闻所未闻的刑罚撕裂他的内脏和肢体。虽然保护的主的尊严和良善保护了那成为司铎的人——祂愿意他成为司铎；然而，就\Person{科尔乃略}的虔诚和敬畏而言，他承受了所能承受的一切，并首先以他的司铎职务战胜了那暴君——后来这暴君在武装和战争中被击败。\par

10. 至于关于他的一些不名誉和恶意的传言，当你知道这总是魔鬼的工作时，你就不会感到奇怪：魔鬼用谎言伤害天主的仆人，用虚假的意见玷污光荣的名字，使那些在自己的良心之光中明亮的人被别人的传言所玷污。此外，你应当知道，我们的同僚已经调查过，并确实发现他并没有像某些人所暗示的那样，被证书的污点所玷污；他也没有与那些献过祭的主教进行亵渎的共融，而只是将那些案件已被审理、无辜性已被认可的人与我们一起接纳。\par

11. 至于\Person{特洛斐摩}，你也希望得到关于他的消息，情况并不像恶意之人的报告和谎言所传达给你那样。因为，正如我们的前任经常做的，我们最亲爱的弟兄在召集弟兄时，出于必要而让步；既然一大部分群众已随\Person{特洛斐摩}离开，现在\Person{特洛斐摩}回到教会，并以祈祷的忏悔弥补并承认他先前的错误，并以完全的谦卑和补赎召回了他前不久带走的弟兄们，他的祈祷得到了垂听；不仅\Person{特洛斐摩}，而且与他一起的极多弟兄，都被接纳进入主的教会——这些人若没有随\Person{特洛斐摩}一起来，就不会全部回到教会。因此，与几位同僚在那里商议后，\Person{特洛斐摩}被接受了，因为弟兄们的回归和许多人的救恩的恢复替他弥补了。然而，\Person{特洛斐摩}被接纳的方式只是作为平信徒共融，而不是像恶意之人的信函所告知你的那样，让他占据司铎的位置。\par

12. 而且，关于你所被告知的，\Person{科尔乃略}到处与那些献过祭的人共融，这个情报也来自背信者的虚假报告。因为那些离开我们的人不能赞美我们，我们也不应期待取悦那些背叛我们、反抗教会、并强行坚持从主教会中拉走弟兄的人。因此，最亲爱的弟兄，不要轻易听信或相信任何反对\Person{科尔乃略}和反对我的流言蜚语。\par

13. 因为如果有人患病，就在危险中给予帮助，这是已决定的。然而，在他们被帮助后，并在危险中获得了平安，我们不能窒息他们，或毁灭他们，或通过我们的武力或双手催促他们走向死亡的结局；仿佛因为平安赐给了临终者，那些接受了平安的人就必须死；尽管这里更显示出神圣的爱和父般的仁慈：那些在给予他们的平安中接受了生命保证的人，更因所接受的平安而被束缚于生命。因此，如果接受了平安之后，天主给了他们缓期执行，没有人应当为此抱怨司铎，因为已经决定要在危险中帮助弟兄。最亲爱的弟兄，你也不可像某些人那样认为，那些接受证书的人应与那些献过祭的人同等看待；因为即使在献过祭的人中，情况和案情也常常不同。我们不可将一个人与另一个人等同看待：一个人立刻欣然投入那可憎的祭献，而另一个人在长期挣扎和抵抗后才在强迫下达到那致命的结局；一个人出卖了自己和所有亲属，而另一个人自己为众人奔赴试炼，保护了他的妻子、儿女和整个家庭，亲自承担危险；一个人强迫他的家人或朋友犯罪，而另一个人宽恕了家人和仆人，甚至接待了许多弟兄——他们因放逐和逃亡而离开——到他家中并予以款待；向主显明并献上许多活着且安全的灵魂，以祈求一个受伤的灵魂。\par

14. 由于献祭的人之间存在很大区别，把那些领了证书的人与那些献了祭的人等同起来，是多么缺乏怜悯，何等痛苦的困难啊！因为那领了证书的人可以说：\Quote{我以前读过，也因主教的训诲而知道，我们不可向偶像献祭，天主的仆人不该崇拜图像；因此，为了不做出这不合法的事，当获得证书的机会被提供时（我自己本也不该接受，除非这机会摆在我面前），我就去，或者委托别人去见长官，说我是基督徒，我不允许献祭，我不能到魔鬼的祭坛前，并且为此我付出代价，以便我不去做那对我不合法的事。}可是现在，连那因领了证书而被玷污的人，——在从我们的劝诫中得知他甚至连这事也不该做，并且虽然他的手是纯洁的，没有致命的食物接触玷污他的嘴唇，但他的良心仍然被玷污之后，他听到我们时哭泣，哀叹，并且现在被提醒他犯罪的事，与其说是出于罪恶，不如说是出于错误而被欺骗，他作证说他已为将来受训并准备好了。\par

15. 如果我们拒绝那些在一种可容忍的良心上有些自信之人的忏悔，那么立刻，与他们的妻子、他们所保全的儿女一起，他们被魔鬼的邀请匆匆带入异端或分裂；在审判之日，将会归咎于我们，说我们没有关心受伤的羊，并且因为一只受伤的羊而失去了许多健康的羊。然而主留下那九十九只整全的，去寻找那一只迷失和疲乏的，并在找到后亲自把它扛在肩上；我们不仅不去寻找失足者，反而当他们来到我们这里时驱赶他们；而假先知们不停地在糟蹋和撕裂基督的羊群，我们却给狗和狼机会，以致那些可憎的迫害没有毁灭的人，我们因自己的刚硬和不人道而毁灭了他们。而且，最亲爱的弟兄，宗徒所说的话将会怎样呢？\BibleQuote{我在一切事上使众人喜悦，不求自己的利益，只求许多人的利益，为叫他们得救。你们该效法我，如同我效法基督一样。}又说：\BibleQuote{对软弱的人，我就成为软弱的，为赢得那软弱的人。}\BibleRef{格前 9:22}又说：\BibleQuote{若一个肢体受苦，所有的肢体都一同受苦；若一个肢体受尊荣，所有的肢体都一同欢乐。}\par

16. 最亲爱的弟兄，哲学家和斯多葛派的原则是不同的，他们宣称所有罪过都是相等的，一个庄重的人不该轻易被触动。但是基督徒与哲学家之间有很大差别。当宗徒说：\BibleQuote{你们要谨慎，免得有人用哲学和虚妄的诡诈把你们掳去，}\BibleRef{哥 2:8}我们应当避免那些并非来自天主的仁慈，而是出自一种过于严苛的哲学的傲慢所生的事物。此外，关于梅瑟，我们在圣经中读到：\BibleQuote{梅瑟为人十分谦和；}\BibleRef{户 12:3}主在福音中说：\BibleQuote{你们应当慈悲，就像你们的父也慈悲了你们一样；}又说：\BibleQuote{健康的人不需要医生，有病的人才需要。}\BibleRef{玛 9:12}那说\Quote{我只医治健康的人，他们不需要医生}的人，能施展什么医术呢？我们应当向受伤的人提供我们的援助、我们的医治；不要看他们如同死人，反而要看他们如同半死不活地躺着，我们看见他们在那次致命的迫害中受了伤，如果他们完全死了，就绝不会后来由同一些人成为宣认者和殉道者。\par

17. 但是，既然在他们身上有一种东西，能藉着后来的忏悔而坚强起来成为信德；并且藉着忏悔，力量被武装成德行，这力量本来是不能被武装的，如果一个人因绝望而跌倒；如果，被严厉而残酷地与教会分离，他转向外邦人的方式和世俗的工作，或者，如果被教会拒绝，他转到异端者和分裂者那里去；在那里，虽然他后来因那名而被处死，但是，由于被置于教会之外，并与合一及爱德分离，他在死亡中不能获得冠冕。因此，最亲爱的弟兄，经对每个人的情况加以审查后，决定暂时接纳那些领了证书的人，而那些献了祭的人在死亡时应得到帮助，因为在逝者的地方没有告解，也没有人能因我们而被强迫去忏悔，如果忏悔的果实被夺走了。如果战斗先来临，被我们强化后，他将被发现为战斗准备好了武装；但如果在战斗之前疾病紧逼他，他就带着平安和共融的安慰离去。\par

18. 此外，我们并不预判主将是审判者；除非祂发现罪人的忏悔是圆满和真诚的，那时祂才会批准我们在此所决定的事。然而，若有人用虚伪的忏悔欺骗我们，那不被欺骗、洞察人心的天主，将审判我们不完全看清的事，主会修正祂仆人的判决。然而，最亲爱的弟兄，我们应当记住经上记载：\Quote{扶助兄弟的兄弟，必被高举。} 宗徒也曾说：\Quote{你们各人应当各自顾念自己，免得你们也被诱惑。你们要彼此分担重担，这样就满了基督的法律。} \BibleRef{迦 6:1} 此外，他在书信中斥责骄傲，挫败他们的狂妄说：\Quote{那自以为站得稳的，要小心，免得跌倒。} 在另一处又说：\Quote{你是谁，竟论断别人的仆人？他或站住或跌倒，自有他的主人；而且他必站住，因为主能使他站住。} 若望也证明主耶稣基督是我们在罪中的\OriginalTerm{辩护者}{Advocate}和\OriginalTerm{中保}{Intercessor}，说：\Quote{我的孩子们，我写这些话给你们，是要你们不犯罪。若有人犯罪，我们在父那里有一位\OriginalTerm{辩护者}{Advocate}，就是耶稣基督，扶持者；祂是为我们的罪作\OriginalTerm{赎罪祭}{propitiation}。} 保禄宗徒也在书信中写道：\Quote{当我们还是罪人的时候，基督为我们死了；何况现在，因祂的血而成义，我们更要藉着祂免去天主的义怒。} \BibleRef{罗 5:8}\par

19. 考虑到祂的爱和慈悲，我们不该在照顾弟兄时如此刻薄、残忍、不人道，而要与哀哭的人同哀哭，与哭泣的人同哭泣，并尽我们所能藉着爱的帮助和安慰扶持他们；既不要过于苛刻固执地拒绝他们的忏悔；也不要过于松散轻率地贸然给予共融。看！一位受伤的弟兄在战场上被敌人击倒。魔鬼正试图杀死他所受伤的人；而基督在敦促，祂所救赎的人不可完全灭亡。我们协助哪一方？我们站在哪边？我们是站在魔鬼一边，任由他毁灭，并像福音中的司祭和肋未人一样经过我们倒卧无生命的弟兄身边？还是，作为天主和基督的司祭，我们效法基督所教导和所做的，从敌人嘴里夺回受伤的人，好让他痊愈并保存给审判的天主？\par

20. 而且，最亲爱的弟兄，不要认为因为放宽了对失足者的忏悔，并向忏悔者提供了和平的希望，弟兄们的勇气就会减弱，殉道就会减少。真正信德坚定的人不动摇；在那些全心敬畏和爱慕天主的人中，他们的正直持续坚定。因为我们对奸淫者也给予忏悔的时间，并赐予和平。然而，教会中并没有因此缺乏童贞，贞洁的崇高志向也不会因他人的罪恶而萎靡。教会因如此多的童贞女而繁荣昌盛；贞洁和端庄保持着它们荣耀的基调。也不会因为奸淫者的忏悔和宽恕变得容易而打破节制的活力。站立以求得宽恕是一回事，达到荣耀是另一回事；被投入监狱，直到还清最后一分钱才能出来是一回事，立即获得信德和勇气的赏报是另一回事。因罪恶而长期受折磨，被火炼净并长时间净化是一回事，藉着苦难洗净一切罪恶是另一回事。最后，在审判之日等候天主的判决是一回事，立即被主加冕是另一回事。\par

21. 事实上，在我们的前任中，我们省的一些主教认为不应将和平赐予奸淫者，并完全关闭了奸淫的忏悔之门。然而，他们并没有退出同僚主教的集会，也没有因他们坚持严厉或谴责而破坏天主教会合一的纽带；因此，因为有些人将和平赐予了奸淫者，那未赐予的人也不应与教会分离。在和谐的纽带保持，天主教会不可分割的圣事持续的同时，每位主教安排和指导自己的行为，并且要向主汇报自己的意图。\par

22. 但我感到惊讶，有些人如此顽固，认为悔改不可赐予失足者，或以为宽恕应拒绝于悔罪者，尽管经上记载：\BibleQuote{你该回想你是从哪里跌落的，你该悔改，行你最初所行的事}\BibleRef{默 2:5}，这话当然是向那显然已跌倒的人说的，主劝勉他藉自己的行为重新站起来，因为经上记载：\BibleQuote{施舍能救人脱离死亡}\BibleRef{多 4:10}，当然不是指那曾被基督的血所消灭、且圣洗和我们的救赎者的救恩恩宠已救我们脱离的死亡，而是指后来因罪过潜入的死亡。此外，在另一处经文也给予悔改的时间；主威胁那不悔改的人说：\BibleQuote{我有许多事反对你，}祂说，\BibleQuote{因为你容忍你的妻子依则贝尔，她自称女先知，教导并引诱我的仆人行淫乱，吃祭偶像之物；我曾给她悔改的时间，但她不肯悔改她的淫乱。看哪，我要把她抛在床上，与她行淫的人也要遭受大苦难，除非他们悔改自己的行为；}\BibleRef{默 2:20} 主当然不会劝勉他们悔改，倘若祂不向悔改者应许慈悲。在福音中祂说：\BibleQuote{我告诉你们，对于一个罪人悔改，在天上也要这样欢乐，比对于九十九个无须悔改的义人更为欢乐。}\BibleRef{路 15:7} 因为经上记载：\BibleQuote{天主并没有制造死亡，也不喜欢生灵丧亡}\BibleRef{智 1:13}，那不愿任何人丧亡的天主，当然愿意罪人悔改，藉悔改重新回归生命。同样，祂藉\Person{岳厄尔}先知呼喊说：\BibleQuote{现在，上主你们的天主这样说：你们应全心归向我，禁食、哭泣、哀悼；要撕裂你们的心，不要撕裂你们的衣服，归向上主你们的天主，因为祂宽仁慈悲，缓于发怒，富于慈爱，对所定的灾祸回心转意。}\BibleRef{岳 2:12} 在\WorkTitle{圣咏集}中，我们也读到天主的斥责与仁慈，祂在宽容的同时也恐吓，祂惩罚好使人改正，改正后则保全。\BibleQuote{我要用杖惩罚他们的过犯，}祂说，\BibleQuote{用鞭笞责罚他们的罪孽。但我的慈爱仍不彻底收回。}\par

23. 主在福音中，阐述天主圣父的爱，说：\BibleQuote{你们中间有哪个人，儿子向他求饼，反而给他石头？或者求鱼，反而给他蛇？你们虽然邪恶，尚且知道把好东西给你们的儿女，何况你们在天之父，岂不更要把好东西赐给求祂的人？}\BibleRef{玛 7:9} 主在这里比较血肉之父和天主圣父永恒而慷慨的爱。如果那地上的恶父，因一个罪恶败坏的儿子深感冒犯，但若后来看到同一个儿子改过自新，弃绝从前生活的罪过，因悔改的忧伤而恢复清醒、道德和纯洁的纪律，他岂不既欢喜又感恩，以父亲般的喜悦急切拥抱那个曾被他赶出的回头浪子？何况那唯一真父，良善、慈悲、仁爱——乃至祂本身就是良善、慈悲与爱——岂不更欢喜祂自己儿子的悔改！祂不威胁那些现在悔改、哀伤痛哭的人，反而应许宽恕和仁慈。因此主在福音中称哀恸的人是有福的，因为哀恸的人必蒙怜悯。那顽固骄傲的人是为自己积蓄愤怒和将来审判的惩罚。因此，最亲爱的弟兄，我们决定：那些不悔改，也不以全心和明确的哀恸表白来显示为罪忧伤的人，若在疾病和危险中开始祈求共融与平安，必须完全剥夺他们的希望；因为驱使他们祈求的不是对罪的悔改，而是迫近死亡的警告；那不曾深思自己将要灭亡的人，不配在死亡中得到安慰。\par

24. 然而，关于诺瓦提安的性格，最亲爱的弟兄，你曾要求将消息写成关于他引入了何种异端；要知道，首先，我们甚至不应该探究他所教导的，只要他在合一之外教导。无论他是谁，无论他是什么，凡不在基督的教会内的，就不是基督徒。虽然他可能自夸，并以高傲的言辞宣扬他的哲学或口才，然而，凡没有保持弟兄之爱或教会合一的，连他先前所是的也失去了。除非在你看来，他是一个主教——当一位主教已在教会中由十六位共主教选出时——却试图通过贿赂，由叛离者之手而成为一个奸淫的、外来的主教；并且，虽然只有一个教会，由基督分散到整个世界成为许多肢体，并且也只有一个主教职，通过许多主教的和谐群体而扩散；尽管有天主的传统，尽管有天主教会的联合且处处凝聚的合一，他仍试图建立一个属人的教会，并派遣他的新宗徒到许多城邑，以便建立他自己任命的新基础。并且，虽然各城和各省已有年迈、信德纯全、在考验中被证明、在迫害中被放逐的主教被祝圣，此人竟敢在这些之上另立虚假的主教：仿佛他既能以他新努力中的坚持走遍整个世界，又能通过传播他自己的不睦而拆散教会身体的结构，不知道分裂分子起初总是狂热，但他们不能增加或增添非法开始的事物，而是立即与他们的恶竞争一同衰落。但是，他不能持有主教职，即使他先前已被立为主教，因为他已自绝于同主教们的身体和教会的合一；因为宗徒告诫我们要互相扶持，不要背离天主所指定的合一，并说：\BibleQuote{彼此在爱中担待，尽力保持圣神的合一在和平的联系中。}\BibleRef{弗 4:2} 那么，凡不保持圣神的合一和和平的联系，并自绝于教会的队伍和司铎的集会的，不能拥有主教的权力或尊荣，因为他拒绝保持主教职的合一或和平。\par

25. 再者，那是何等傲慢的肿胀，何等谦卑与温良的遗忘，何等自傲的夸耀，竟有人敢或以为自己能够做连主都没有赐予宗徒们的事；他以为自己能辨别莠子与麦子，或者，仿佛他已获准拿着簸箕并清理禾场，就试图将糠秕与麦子分开；既然宗徒说：\BibleQuote{但在大屋里，不但有金器银器，也有木器和瓦器，}\BibleRef{弟后 2:20} 却以为可选择金器银器，鄙视、丢弃并谴责木器和瓦器；而木器只在主的日子由神圣火焰的烈焰才被烧毁，瓦器则只有那位获得铁杖的才可打破。\par

26. 或者，如果他自命为心与肾的探究者和审判者，就让他对一切案件同样判断。既然他知道经上记载：\BibleQuote{看，你已痊愈了；不要再犯罪，免得你遭遇更坏的事，}\BibleRef{若 5:14} 就让他从自己身边和同伴中分离欺诈者和奸淫者，因为奸淫者的案件远比取得证书者更严重、更糟糕；因为后者因需要而犯罪，前者因自由意志；后者认为不祭拜已足够，而受到错误的欺骗；前者，侵犯他人的婚姻纽带，或进入妓院，进入平民的污秽和肮脏的深坑，以可憎的不洁玷污了圣化的身体和天主的殿，正如宗徒所说：\BibleQuote{凡人所犯的罪都在身体以外，但那犯奸淫的，却是冒犯自己的身体。}\BibleRef{格前 6:18} 然而，即使对这些人也赐予了悔改，并留下了哀恸和赎罪的希望，按同一位宗徒的话：\BibleQuote{我害怕，当我来到你们那里时，我要为许多从前犯了罪而没有悔改的人哀恸，因为他们行过不洁、奸淫和放荡。}\BibleRef{格后 12:21}\par

27. 不可让新异端者因声称不与崇拜偶像者共融而自满；尽管在他们中间既有奸淫者，也有欺诈者，这些人被视为犯崇拜偶像之罪，正如宗徒所说：\BibleQuote{你们要明白：凡是淫乱的、不洁的、贪财的——贪财者即是崇拜偶像者——在基督和天主的国内，都得不到产业。}\BibleRef{弗 5:5} 又说：\BibleQuote{所以，你们要致死属于地上的肢体，戒绝淫乱、不洁、邪情、恶欲和贪财——贪财即是崇拜偶像——因这些事，天主的愤怒才降来。}\BibleRef{哥 3:5} 因为我们的身体是基督的肢体，我们每个人都是天主的宫殿；谁若以奸淫玷污天主的宫殿，便是玷污天主；谁若犯罪，遵行魔鬼的旨意，就是服事魔鬼和偶像。因为恶行并非来自圣神，而是来自仇敌的煽动，由不洁之神所生的私欲迫使人们反抗天主、顺从魔鬼。因此，若他们说一人因他人的罪而被玷污，并坚持声称犯罪者的崇拜偶像之罪转移到他们自己所说无罪的人身上，他们便不能免除崇拜偶像之罪，因为宗徒的证明清楚指出，与他们共融的奸淫者和欺诈者都是崇拜偶像者。然而，按照我们的信德和所赐予的神圣宣讲准则，真理的原则与我们一致：每人自己承担自己的罪；一人不能为另一人成为有罪，因为主预先警告我们说：\BibleQuote{义人的义必归于义人，恶人的恶必归于恶人。}\BibleRef{则 18:20} 又说：\BibleQuote{父亲不可因儿女而死，儿女不可因父亲而死；每人要因自己的罪而死。}\BibleRef{申 24:26} 阅读并遵守这些，我们确实认为不应阻止任何人获得补赎的果实和平安的希望，因为我们知道，按照神圣圣经的信德——天主自己是圣经的作者，并在其中劝勉——罪人被召回悔改，且悔改者不会得不到宽恕和慈悲。\par

28. 啊，被欺骗的弟兄们的嘲笑！啊，可怜而无知的哀悼者虚妄的欺骗！啊，异端制度的无效而无益的传统！劝勉人作赔补的悔改，却又从赔补中除去医治；对我们的弟兄说：\Quote{你们要哀悼、流泪，日夜叹息，为洗除和洁净你们的罪而大力并频繁地劳作；但做完这一切之后，你们要在教会之外死去。凡为平安所必需的事，你们都要做，但你们所求的平安，一点也得不到！} 谁不会立刻丧亡？谁不会因彻底绝望而跌倒？谁不会转而放弃一切哀悼的意图？你以为农夫会劳作，如果你对他说：\Quote{你要用一切农艺技能耕种田地，勤勉坚持耕作；但你将收不到庄稼，压不出葡萄汁，得不到橄榄园的果实，也摘不到树上的苹果；} 或者，如果劝勉某人拥有并使用船只，你对他说：\Quote{弟兄，你购买上等木材的原料；用最坚固而精选的橡木编织你的龙骨；致力制作船舵、绳索、风帆，使船建造并装配好；但做完这些之后，你永远看不到它航行和航行的结果吗？}\par

29. 这乃是关闭和截断忧伤与悔改的道路；以致虽然在整个圣经中，主上主抚慰那些归向祂并悔改的人，但悔改本身却被我们的硬心和残忍夺去，这残忍阻断了悔改的果实。然而，若我们发现不应阻止任何人悔改，并且主的司铎们可以将平安赐给那些恳求并哀求主怜悯的人——因为祂是仁慈和慈爱的——那么，哀悼者的呻吟就应当被接纳，悲伤者的悔改果实也不应被拒绝。因为在死者的地方没有告解，也不能在那里做告解，所以那些全心悔改并祈求了平安的人，应当被接纳入教会，并在教会内为主保存；主必定会在祂来到教会时，审判那些在教会内的人。但是背教者和叛逃者，或敌对者和仇敌，以及那些破坏基督教会的人，即使他们在教会外为祂的名而被杀，也不能被接纳进入教会的平安，按照宗徒的话，因为他们既没有保持精神的合一，也没有保持教会的合一。\par

30. 最亲爱的弟兄，以上我从许多事中简略地略述了这几件，为要借此满足你的愿望，并将你更紧密地联系到我们团体和身体的共融中。但若你得到机会和力量前来我们这里，我们将能更充分地彼此商议，更有效地更广泛地考虑那些有益于拯救性合议的事宜。最亲爱的弟兄，祝你永远衷心平安。\par

\SourceRecord{Translated by Robert Ernest Wallis.；Ante-Nicene Fathers；Vol. 5.；Edited by Alexander Roberts, James Donaldson, and A. Cleveland Coxe.；Buffalo, NY: Christian Literature Publishing Co.,；1886.；<http://www.newadvent.org/fathers/050651.htm>.}

% structure: p0001 source=h1 target=section reason=page-title

\section{第52封书信}

\Emph{致\Person{福图纳图斯}及其他同僚：关于那些因受酷刑而屈服的人}\par

\Emph{论题——\Person{西彼廉}受同僚咨询：某些在酷刑中被压倒的失足者是否应准予共融？他答复道，鉴于他们已忏悔三年之久，他认为应予接纳；但在逾越节后将有主教会议与他共议，届时将再行商议。}\par

1. \Person{西彼廉}致其兄弟\Person{福图纳图斯}、\Person{阿希姆努斯}、\Person{奥普塔图斯}、\Person{普里瓦提亚努斯}、\Person{多纳图卢斯}及\Person{斐理克斯}，安好。最亲爱的兄弟们，你们来信告知，当你们在\Place{卡普萨城}为祝圣一位主教时，我们的兄弟及同僚\Person{苏佩里乌斯}向你们提出：我们的兄弟\Person{尼努斯}、\Person{克莱门提亚努斯}和\Person{弗洛鲁斯}，先前在迫害中被捕，曾承认主的名，并战胜了官长的暴力和狂乱群众的攻击；其后在总督面前受重刑折磨，被酷烈的痛苦所征服，因长时间的折磨而从他们以全备信德之力所趋向的荣耀等级上坠落；在这一次并非出于自愿而是出于必要的严重失足之后，他们连续三年未曾中止忏悔。你们以为应当咨询：现在是否适宜接纳他们进入共融。\par

2. 实则，依我之见，我认为主的仁慈不会吝于那些已知曾站在战斗行列中、曾承认主名、以坚定不移的信德战胜官长暴力和狂乱群众的冲击、曾遭受监禁、曾在总督的恐吓和周围民众的喧嚣中长期抵抗、并以持久的反复折磨而撕裂他们的酷刑；以致最后时刻因肉身的软弱而被征服，可由先前功绩的辩护而减轻。对于这样的人，失去荣耀或许已足够，但我们不应再向他们关闭宽恕之门，剥夺他们父的爱和我们共融；我们认为，如你们所写，他们以过度的忏悔哀恸连续三年忧伤哭泣，足以恳求主的仁慈。我确实不认为，对那些在作战的勇敢中并未在战斗前退缩，且若新的斗争来临可能重获荣耀的人，轻率而过急地赐予和平是不谨慎的。因为当会议决定，处于疾病危险中的忏悔者应得到帮助并赐予和平，那么那些我们看见并非因心灵软弱而跌倒，而是因肉身的软弱在战斗中受伤后无法承受他们宣信的荣冠的人，岂不更应优先获得和平？尤其因为，他们虽渴望死亡，却不被允许被杀害，而是酷刑足够长久地折磨他们疲惫的躯体，并非为了征服那不可征服的信德，而是为了耗尽软弱的肉身。\par

3. 然而，既然你们来信要我与我许多同僚充分商议此事，而如此重大的主题需要更广泛、更谨慎的会议讨论；且如今几乎全体同僚在首次庆祝逾越节期间，都与各自的兄弟住在家里；待他们完成在各自人群中所举行的庆典，并开始到我这里来时，我将与每一位更广泛地考虑此事，以便在众多司铎的会议中权衡出一个确定的意见，关于你们所咨询的主题，在我们中间确立，并写信给你们。最亲爱的兄弟们，我衷心祝你们永远安康。\par

\SourceRecord{Translated by Robert Ernest Wallis.；Ante-Nicene Fathers；Vol. 5.；Edited by Alexander Roberts, James Donaldson, and A. Cleveland Coxe.；Buffalo, NY: Christian Literature Publishing Co.,；1886.；<http://www.newadvent.org/fathers/050652.htm>.}

% structure: p0001 source=h1 target=section reason=page-title

\section{第五十三封书信}

\Emph{致\Person{科尔内略}：论给予背教者平安}\par

\Emph{摘要：\Person{西彼廉}以整个主教会议的名义向\Person{科尔内略}神父宣布这项法令；因此这封信与其说是\Person{西彼廉}本人的信，不如说是整个非洲主教会议的信。}\par

\Signature{\Person{西彼廉}、\Person{利贝拉里斯}、\Person{卡尔多纽}、\Person{尼各美德}、\Person{凯奇利乌}、\Person{尤尼乌}、\Person{马鲁提乌}、\Person{费利克斯}、\Person{苏克塞苏斯}、\Person{福斯提诺}、\Person{福尔图纳图}、\Person{维克多}、\Person{撒图尼努}、另一位\Person{撒图尼努}、\Person{罗加提安}、\Person{特图卢}、\Person{路西亚诺}、\Person{欧提克斯}、\Person{安普卢}、\Person{萨提乌}、\Person{塞孔迪诺}、另一位\Person{撒图尼努}、\Person{奥雷略}、\Person{普里斯库}、\Person{赫尔库拉诺}、\Person{维克托里库}、\Person{昆图}、\Person{奥诺拉图}、\Person{蒙塔诺}、\Person{霍尔滕西亚诺}、\Person{维里亚诺}、\Person{扬布斯}、\Person{多纳图}、\Person{庞培}、\Person{波利卡普}、\Person{德默特里}、另一位\Person{多纳图}、\Person{普里瓦提安}、另一位\Person{福尔图纳图}、\Person{罗加图}和\Person{莫努卢}向他们的弟兄\Person{科尔内略}致敬。}\par

1. 至亲的弟兄，我们确实已在早先共同商议后决定：那些在迫害的猛烈中被仇敌击倒、背教，并以非法祭献玷污自己的人，应进行长期而完全的补赎；若疾病危险紧迫，应在死亡关头领受平安。因为教会不应对敲门者关闭，也不应将救恩希望的帮助拒绝给那些哀哭恳求的人，这不合理，圣父的爱和天主的仁慈也不允许；这样，当他们离开这个世界时，不应在没有共融与平安的情况下被遣送到他们的主那里；因为那颁布法律的那位亲自规定，凡在地上捆绑的，在天上也要被捆绑；并且他还允许，那些在地上首先在教会里被释放的，在那里也可以被释放。但现在，当我们看到另一场患难的日子又开始临近，并被频繁重复的暗示提醒：我们应为主所宣告的争战做好准备并武装起来，也应以我们的劝诫准备那些由天主眷顾托付给我们的子民，并从各处聚集所有渴望武器、急切在主的营中战斗的基督士兵；因这必要性的迫使，\Emph{我们决定}：要给那些没有退出主的教会，但从背教第一天起就没有停止悔改、哀哭并恳求主的人，给予平安；并决定他们应当为迫近的争战被武装和装备。\par

2. 因为我们必须遵从合适的暗示和劝诫，使羊群不至于在危险中被牧人遗弃，而全体羊群能聚集一处，主的军队为天上战争的竞赛而集合。因为当和平与平安盛行时，哀悼者的悔改被合理地延长了更久的时间，帮助只给予病危者离世之时，这允许长时期推迟哀悼者的眼泪和在病危时给予临终者迟来的帮助。但现在，平安确实必要，不是为了病人，而是为了强壮者；我们也不应给予垂死者共融，而应给予活着的人，免得我们使那些我们激励和劝勉去战斗的人赤手空拳、毫无装备，而应以基督圣体圣血的保护来坚固他们。并且，既然圣体圣事被设立正是为了成为领受者的保障，\Emph{就有必要}以主的丰盛恩宠的保护来武装那些我们愿意使之安全抵抗仇敌的人。因为我们若拒绝给那些将要投入战争的人基督的血，又怎能教导或激励他们为承认祂的名而流血呢？我们若不先允许他们在教会中因共融之权饮主的杯，又怎能使他们适于殉道之杯呢？\par

3. 至亲的弟兄，我们应当分辨：有些人已经背教，回到他们所弃绝的世界，过着外邦人的生活；或者已沦为叛徒到异端者那里，每日拿起弑亲的武器攻击教会；而有些人并不离开教会的门槛，不断悲伤地恳求天主和父的安慰，宣示他们现在已为争战准备就绪，并预备好为他们的主的名和自身的救恩勇敢站立争战。在这时期，我们赐予平安，不是给睡觉的人，而是给警醒的人。我们赐予平安，不是在放纵安闲中，而是在武器之中。我们赐予平安，不是为了休息，而是为了战场。如果我们按照所听闻、所期望、所相信他们的那样，他们将勇敢站立，并在交锋中与我们一同击倒仇敌，我们将不会后悔赐予如此勇敢的人平安。是的，赐予殉道者平安，是我们主教职的伟大光荣和荣耀，使我们作为司铎，每日举行天主的祭献，能为天主预备祭品和牺牲。但是，如果——愿主从我们的弟兄们中除去这事——任何一个背教者欺骗，以诡计寻求平安，在迫近争战之时领受平安而无战斗之志，他就背叛并欺骗了自己，心中藏着一件事，口中说出另一件事。我们，就我们所能看见和判断的，只看每个人的外貌；我们不能窥察人心，洞察思想。关于这些，那辨别和搜查隐藏事物的那一位会审判，祂将迅速来临，审判心中的秘密和隐藏之事。但恶人不应当阻挡善人，反而恶人应当被善人帮助。因此，平安也不应被拒绝给那些将要忍受殉道的人，因为有些人会拒绝它，因为为了这个目的，平安应当赐给所有将要投入战争的人，免得因我们的无知，那在争战中应被加冕的人反而被疏忽。\par

4. 谁也不要说：\Quote{接受殉道者是在自己的血中受洗，并且他不需要从主教那里得到平安，因为他即将拥有自己荣耀的平安，并因主的屈尊俯就而领受更大的赏报。}首先，若未受教会装备以应战，他就不能适于殉道；若未领受圣体以振奋激励，他的精神就是软弱的。因为主在福音中说：\BibleQuote{但当人把你们交出时，不要思虑你们该说什么；因为在那个时刻，必会赐给你们该说什么。因为不是你们说话，而是你们父的圣神在你们内说话。}\BibleRef{玛 10:19}如今，既然祂说父的圣神在那些被交出并置于宣认祂圣名的人内说话，那么，一个人若未先在领受平安时领受了父的圣神——父的圣神加给祂仆人们力量，亲自在我们内说话和宣认——他又怎能被视为预备好或适于那宣认呢？此外，若有人舍弃一切所有而逃，住在隐密处和独处中，偶然落入盗匪之手，或死于热症和衰弱，难道我们不会被指控：这样一位好兵士，舍弃一切所有，轻视自己的房屋、父母和子女，宁愿追随他的主，却死于无平安、无共融吗？在审判之日，难道不会将懈怠疏忽或残酷无情归咎于我们，就是我们这些牧人，既不愿在平安中照顾那托付给我们、交托给我们的羊群，也不愿在战斗中武装他们吗？难道主不会藉祂的先知向我们提出控诉说：\BibleQuote{看，你们喝牛奶，穿羊毛，宰杀肥壮的，却不牧放我的羊群。你们没有扶助软弱的，没有医治有病的，没有包扎受伤的，没有领回迷失的，没有寻找失落的，却用强力劳役那强壮的。我的羊因没有牧人而四散，成了田野一切野兽的食物；没有人去寻找，也没有人领回。因此，主这样说：看，我必与牧人为敌，我必向他们追讨我的羊，使他们不再牧放我的羊；他们也不再牧放，我要从他们口中救出我的羊，我要以公义牧养他们。}\par

5. 因此，为了不使主托付给我们的羊群从我们口中被追讨——我们竟藉此拒绝平安，竟对人施行残忍的严厉而非天父之爱的\Emph{仁善}——我们已决定，藉圣神的提示和主的劝诫——这些劝诫通过许多明显的神视传达给我们——因为敌人已被预告并显示已迫近，要将基督的兵士聚集到营中，审查每个人的情况，并赐平安给失足者，确切地说，要给那些即将作战的人装备武器。我们相信，在默想天父的慈悲时，这事必蒙你们喜悦。但若我们同人中有谁在这争战紧迫之时，认为不应将平安赐予我们的弟兄姐妹，他将在审判之日向主交账，或为他严厉的苛刻，或为他非人性的硬心。我们，按合於我们的信德、爱德和挂虑，已将我们自己心中所想陈明於你们，就是：争战之日已近，凶暴的敌人将很快兴起攻击我们，一场斗争即将来临，不像从前那样，而是更严重、更剧烈。这事常常从上面显示给我们；关于这事，我们常受主眷顾和慈悲的劝诫，我们信赖祂的人，可确知祂的帮助和慈爱，因为祂在平安中预告祂的兵士战斗将要来临，也必在争战中赐给他们胜利。我们祝你，最亲爱的弟兄，永远衷心平安。\par

\SourceRecord{Translated by Robert Ernest Wallis.；Ante-Nicene Fathers；Vol. 5.；Edited by Alexander Roberts, James Donaldson, and A. Cleveland Coxe.；Buffalo, NY: Christian Literature Publishing Co.,；1886.；<http://www.newadvent.org/fathers/050653.htm>.}

% structure: p0001 source=h1 target=section reason=page-title

\section{书信五十四}

\Emph{致科尔乃略，论福尔图纳图斯和费利奇西穆斯，或驳斥异端者。}\par

\Emph{纲目。——西彼廉在此信中主要警告科尔乃略，不要听信费利奇西穆斯和福尔图纳图斯对他所作的诽谤，也不要被他们的威胁吓倒，而要勇敢无畏，正如天主的司铎在对抗异端者时所当行的；即是那些追随异端者中通行的习惯，以蔑视教会中的一位主教开始他们的异端和分裂的人。}\par

1. 最亲爱的兄弟，我已读了你的来信，是由我们的兄弟辅祭撒杜鲁斯带来的；信中充满了兄弟之爱与教会纪律及司铎的责备；你表示，费利奇西穆斯——他并非基督的新仇敌，而是早已因诸多严重的罪行被绝罚，且不仅由我的审判，也由众多同为主教的审判所定罪——已被你在那里拒绝；当他带领一群亡命之徒和派系前来时，你以主教应尽之全力将他逐出教会。这教会早已被天主之威严及我主基督审判者之严厉，与像他一样的人一同驱逐；这分裂与不和的肇始者、骗取托付钱财者、侵犯贞女者、破坏与败坏许多婚姻者，不应再因他无耻及乱伦的接触而玷污基督迄今无玷、圣洁、端庄的新娘。\par

2. 但是，兄弟，当我读了你的另一封信（你附在第一封信之后）时，我颇为惊讶地发现，你竟在某种程度上被那些前来之人的威胁和恐吓所困扰；据你所写，他们以极端绝望之态攻击并威胁你，说若你不接受他们带来的信件，他们将公开宣读，并说出许多卑鄙无耻、有辱其口的话。但若事情如此，最亲爱的兄弟，即最邪恶之人的胆大妄为竟值得畏惧，且恶人不能以正当公平的方式成就之事，却可凭鲁莽与绝望达成，那么主教之权威及治理教会之崇高神圣权力便告终结；我们也不能继续存在，或实际上不再是基督徒，若此事已发展到我们竟要害怕被逐之徒的威胁或诡计的地步。因为外邦人和犹太人都威胁，异端者及所有心中被魔鬼占据的人，每日都以狂怒之声证明其毒恶疯狂。因此，我们不应因其威胁而屈服；仇敌与敌人也并不因此大于基督，因其在世界中自许与承担如此之多。最亲爱的兄弟，我们应持守不可动摇的信德力量；面对一切向我们咆哮之浪潮的冲击与侵袭，应如磐石般坚韧不拔，以坚定不屈的勇气屹立。对于一位主教而言，他生活在恐惧与危险之中，却正因这些恐惧与危险而光荣，恐惧或危险来自何处并不重要。因为我们不应只考虑和关注外邦人或犹太人的威胁，当我们看到主自己也被祂的兄弟抛弃，被祂亲自从宗徒中所拣选的人出卖；在世界的开始，也是兄弟杀了义人\Person{亚伯尔}，愤怒的兄弟追赶逃走的\Person{雅各伯}，年幼的\Person{若瑟}因兄弟的行为而被卖。在福音中我们也读到预言，我们的仇敌将是我们自己家里的人，首先在合一的圣事中结合的人，将彼此出卖。是谁出卖或是谁施暴并不重要，因为天主容许那些祂指定要加冕的人被出卖。因我们遭受基督所受之苦并非耻辱，他们做出犹大所行之事也非荣耀。但是这些威胁者何等傲慢，何等膨胀、浮夸、虚妄的夸耀，竟然在我缺席时\Emph{那里}威胁我，而在这里他们却在我手中！我不怕他们的辱骂，他们每日以此伤害自己和自己的生命；我不因他们的棍棒、石头和刀剑而战栗，他们以杀亲之词挥舞这些；尽其所能在天主面前他们是杀人者。然而除非主允许他们杀人，他们并不能杀人；尽管我必死一次，但他们每日以仇恨、言语和恶行杀我。\par

3. 然而，最亲爱的弟兄，不可因此放弃教会的纪律，也不可放松司铎的谴责，因为我们被辱骂所困扰，或被恐惧所动摇。因为圣经迎面对我们提出警告说：\BibleQuote{但那傲慢自大、夸口自诩、使灵魂如地狱一样膨胀的人，必一事无成。}\BibleRef{哈 2:5}又说：\BibleQuote{不要畏惧罪人的言语，因为他的荣耀必成为粪土和蛆虫。今日他高升，明日就寻不见，因为他已化为尘土，他的思念也必消灭。}\BibleRef{加上 2:62}又说：\BibleQuote{我曾见恶人趾高气扬，高耸如\Place{黎巴嫩}的香柏；我经过，看哪，他已不在；我寻找他，却寻不着他的踪迹。}这种高抬自己、自大、傲慢夸口的言语，不是来自教导谦卑的基督，而是来自反基督的精神。主曾借先知斥责他说：\BibleQuote{因为你心里说：我要升到天上，我要将我的宝座高举在天主的众星之上；我要坐在北方的极山上，我要升到云中，我要与至高者相似。}\BibleRef{依 14:13}他又接着说：\BibleQuote{然而你必下到地狱，到地极之处；凡看见你的，都必惊奇。}\BibleRef{依 14:15}因此，圣经在另一处也以类似的惩罚警告这样的人说：\BibleQuote{因为万军的上主的日子，要临到一切骄横傲慢的人，临到一切自高自大的人。}\BibleRef{依 2:12}所以，人因自己的口舌和言语，立时就被暴露；他心中怀有基督还是反基督，可由他的言语辨别，正如主在福音中所说：\BibleQuote{毒蛇的种类！你们既是恶的，怎能说出善来？因为心里充满什么，口里就说什么。善人从善库中发出善来，恶人从恶库中发出恶来。}\BibleRef{玛 12:34}因此，那因大罪而受苦的富人，曾向当时躺在\Person{亚巴郎}怀中、安居于安慰之地的\Person{拉匝禄}求援；而他自己在煎熬的痛苦中，被燃烧的烈焰所吞噬，他身体中最受折磨的，正是他的口舌，因为无疑他在口舌上犯罪最多。\par

4. 因为经上记载：\BibleQuote{辱骂人的不能承受天主的国}\BibleRef{格前 6:10}，主又在福音中说：\BibleQuote{凡向弟兄说\InnerQuote{你是愚人}的，凡说\InnerQuote{拉卡}的，应受地狱的火。}\BibleRef{玛 5:22}那么，那些不仅对弟兄，而且对司铎堆积这类言语的人，如何能逃脱主报复的谴责呢？司铎被赋予了天主如此屈尊的尊荣，以至于凡不听从他的司铎、以及当时审判他的人，立时被处死？在\BibleBook{}{申命纪}中，上主天主说：\BibleQuote{若有人擅自行事，不听从司铎或审判官——无论当时是谁——那人必被处死；众百姓听见，就必害怕，不敢再行恶。}\BibleRef{申 17:12}再者，当\Person{撒慕尔}被犹太人轻视时，天主说：\BibleQuote{他们不是轻视你，而是轻视了我。}\BibleRef{撒上 8:7}主也在福音中说：\BibleQuote{听从你们的，就是听从我；拒绝你们的，就是拒绝我；拒绝我的，就是拒绝那派遣我的。}\BibleRef{路 10:16}当祂洁净了癞病人后，说：\BibleQuote{去吧，将身体给司铎察看。}\BibleRef{玛 8:4}后来在祂受难时，祂曾受一位司铎仆人的掌掴，那仆人对祂说：\BibleQuote{你这样回答大司铎吗？}\BibleRef{若 18:22}主没有说出任何侮辱大司铎的话，也没有减损司铎的尊荣。反而维护自己的无辜，表明说：\BibleQuote{我若说得不对，你指证不对之处；若说得对，你为什么打我？}\BibleRef{若 18:23}随后在\BibleBook{宗徒}{大事录}中，当有人对蒙福的宗徒\Person{保禄}说：\BibleQuote{你竟辱骂天主的大司铎吗？}\BibleRef{宗 23:4}——尽管那些人已经开始亵圣、不虔敬、血腥，主已经被钉十字架，他们不再保留任何司铎的尊荣和权柄——然而\Person{保禄}顾念这头衔本身，即使是空虚的、司铎的影儿，说：\BibleQuote{弟兄们，我不知道他是大司铎；因为经上记载：\InnerQuote{不可毁谤你百姓的首领。}}\BibleRef{宗 23:5}\par

5. 既然如此重大、如此众多的榜样，以及其他许多先例，藉此司祭的权威与能力因天主的屈尊就卑而得以确立，你们认为那些敌视司祭、反抗天主教会，既不被预告之主的警告所吓倒，也不为将来审判的报复所畏惧的人，究竟是些什么样的人呢？因为异端的产生和分裂的起源，无不出于这一原因：不服从天主的司祭；也不考虑在教会中有一人暂时担任司祭，暂时代替基督作审判者；如果整个弟兄团体按照神圣的教导服从此人，就没有人会挑起任何反对司祭团的事端；没有人会在神圣的审判、人民的投票、同僚主教的同意之后，使自己成为审判者——不是审判主教，而是审判天主。没有人会因分裂基督的合一而撕裂教会。没有人会自我陶醉、傲慢膨胀，创立一个分离在外的新异端，除非有人怀着如此亵渎的胆量、如此堕落的头脑，以致认为司祭可以在没有天主审判的情况下被任命，而主在福音中说：\Quote{两只麻雀不是卖一个铜钱吗？但若没有你们父的旨意，它们连一只也不会掉在地上。}\BibleRef{玛 10:29} 既然祂说连最小的事没有天主的旨意也不会发生，难道有人会认为在祂的教会中，那些最高、最大的事是在天主不知情或不允许的情况下发生的吗？而且司祭——即祂的管家——不是由祂的命令任命的吗？这便不是那使我们生活的信德；这不是给予天主尊荣，我们知道并相信万事都由祂的指引和决定所统驭和管理。毫无疑问，确实有主教并非因天主的旨意而被设立，而是那些在教会之外被设立的——那些违反福音的诫命和传统而设立的；正如主自己在十二先知中所宣称的：\Quote{他们为自己立了君王，却不是出于我。}\BibleRef{欧 8:4} 又说：\Quote{他们的祭品如同丧饼，凡吃它的必被玷污。}\BibleRef{欧 9:4} 圣神也藉依撒依亚呼喊说：\Quote{祸哉，你们这些叛逆的儿女！上主说：你们图谋，却不出于我；你们缔结盟约，却不出于我的神，以致罪上加罪。}\BibleRef{依 30:1}\par

6. 但是——我因受激而说话；我因悲痛而说话；我因被迫而说话——当一位主教被派往已故者的位置，当他在和平时期由全体人民的投票选出，当他在迫害中得到天主助佑的保护，与所有同僚忠实地联结，在其主教任期内如今已历时四年，被其人民认可；在和平时期遵守纪律；在动乱时期，带着已赐予并附着于他的主教名号而遭放逐；在竞技场中屡次被索要\Quote{喂狮子}；在圆剧场中受到天主屈尊就卑的见证而受尊荣；甚至就在我写这封信给你的这些天，由于通过公告所命令的人民应当举行的祭献，民众在竞技场中再次喊叫着要把他\Quote{喂狮子}——当这样一个人，最亲爱的兄弟，被一些绝望和鲁莽的人以及那些在教会之外的人所攻击时，是谁在攻击他，就显然了：肯定不是基督，祂或是任命或是保护祂的司祭；而是那作为基督的仇敌和祂教会的对头，正是为此目的，以其恶意迫害教会之首，以便当舵手被移除时，他能够更凶暴、更猛烈地肆虐，以期分散教会。\par

7. 我极亲爱的兄弟，任何忠信并铭记福音、遵守那预先警告我们的宗徒命令的人，都不应该为此感到困扰：如果在末世，某些骄傲、悖逆、敌视天主司祭的人，或是离开教会，或是攻击教会；因为主和祂的宗徒们早已预告过会有这样的人。也不要让任何人感到惊奇，那被置于他们之上的仆人竟会被一些人离弃，因为主自己行过如此伟大奇妙的作为，并以祂的作为见证了天主父的属性，却连祂自己的门徒也都离弃了祂。然而，当他们离去时，祂并没有责备他们，甚至没有严厉地威胁他们；反而转身对祂的宗徒们说：\Quote{你们也要走吗？}\BibleRef{若 6:67} 明显遵守了这条法则：人既被置于自由之中，并树立于自己的选择之内，便为自己欲求死亡或救恩。然而，伯多禄，就是同一位主在其上建立了教会的那一位，他代表众人发言，并以教会的声音回答说：\Quote{主！惟祢有永生的话，我们去投奔谁呢？我们相信，而且知道祢是天主的圣者。}\BibleRef{玛 15:13} 无疑地指明并显示出：那些离开基督的人是因自己的过错而灭亡，然而那相信基督、并持守那一次所学之道的教会，绝不离开祂；而那些留在天主的家里的人才是教会；反之，那些不是由天主父所栽种的人，我们看见他们并非稳固如麦子，而是像糠秕一样被仇敌吹散；若望在他的书信中也论到这些人说：\Quote{他们从我们中间出去了，却不是属于我们的；因为如果他们是属于我们的，他们必会留在我们中间。}\BibleRef{若一 2:19} 保禄也警告我们：当恶人从教会中灭亡时，不要感到困扰，也不要因无信者的离去而削弱我们的信德。\Quote{即使有些人不忠信，他们的不忠信能使天主的忠信失效吗？断乎不能！因为天主是真实的，而人都是虚谎的。}\BibleRef{罗 3:3}\par

8. 就我们自己而言，最亲爱的弟兄，我们的良心理应努力，不让人因我们的过失而\Emph{离开}教会丧亡；但如果有人自愿并因自己的罪恶而丧亡，且不愿悔改并回到教会，那么，我们这些关切他们得救的人，在审判之日将无可指摘，而只有那些拒绝借我们忠告的健全功效而得到医治的人，将留在惩罚中。丧亡者的辱骂也不应使我们有任何偏离正路和确定规条，因为宗徒也教导我们说：\Quote{如果我取悦于人，我就不是基督的仆人了。}\BibleRef{迦 1:10} 一个人是期望获得人还是天主的赞许，这有很大的区别。如果我们寻求取悦人，主就会不悦。但如果我们努力并劳苦，为能取悦天主，我们就应轻视人的辱骂和责备。\par

9. 但是，最亲爱的弟兄，我没有立即写信给你关于福图纳图斯的事，那个由少数且是顽固异端者所立的假主教，此事并非那种应该立刻且仓促地引起你注意的事，仿佛它是重大或可畏惧的；尤其是因为你对福图纳图斯的名字已经足够熟悉，他是五个长老之一，不久以前离弃了教会，并最近被我们众多且极有权威的同僚主教的判决所绝罚，他们去年曾就此事写信给你。你也会认出费利奇西穆斯，那叛乱的旗手，他本人也包含在同一些主教们很久以前写给你的同一封信中，他不仅在这里被他们绝罚，而且最近还在那里被你逐出教会。由于我确信这些事为你所知，并确信它们常留在你的记忆和纪律中，我认为没有必要立即急迫地将异端者的愚行告诉你。因为关心异端者和分裂者的狂妄在彼此中间所尝试的，实在不应属于天主教会的威严或尊严。因为据说诺瓦提安的党羽现在也立了长老马克西穆斯——他最近作为诺瓦提安的使者被派到我们这里，并拒绝了与我们的共融——在那里作他们的假主教；然而我还没有写信告诉你这事，因为这一切都被我们轻视了；而且我最近已寄给你那里所任命的主教们的名字，他们以健全而正常的纪律治理天主教会中的弟兄。因此，经我们众人商议，决定写信给你，这样便可找到一种摧毁错误、发现真理的简便方法，使你和我们的同僚知道应该写信给谁，以及相应地应当从谁那里收信；但若有人，除了我们在信中所列出的那些人以外，胆敢写信给你，你就知道，他或是被祭献或接受证书所玷污，或是异端者之一，因而乖戾邪恶。然而，藉着一位伟大的朋友和一位文书的机会，我已通过副祭费利奇亚努斯写信给你，他是你与我们的同僚佩尔修斯一同派来的，在其它应从他那一党引起你注意的事项中，也提到那福图纳图斯。但正当我们的弟兄费利奇亚努斯或因风浪受阻在那里，或因收我们其它的信而耽搁时，费利奇西穆斯抢先赶到了你那里。因为邪恶总是如此匆忙，仿佛凭借其速度就能胜过无辜。\par

10. 但我通过费利奇亚努斯向你，我的兄弟，说明，有一个古老的异端者普里瓦图斯来到了迦太基，他曾在兰贝萨殖民地，许多年前因许多严重的罪行被九十位主教的审判所定罪，并被我们的前任法比安和多纳图斯的书信严厉斥责，这对你并非隐晦。他在我们于刚过去的五月中旬举行的会议中，声称要向我们申诉他的案件，但不被允许，便为自己立了那福图纳图斯作假主教，堪配他的同伙。还有一个名叫费利克斯的人与他同来，这人原是他在教会之外、在异端中立的假主教。但约维努斯和马克西穆斯也与这被证实的异端者同来，他们因邪恶的祭献和经我们九位同僚主教的审判所证实的罪行而被定罪，并在去年的会议中被我们许多人再次绝罚。与这四个人并列的还有苏图尔尼卡的雷波斯图斯，他不仅在迫害中自己跌倒，而且以亵渎的劝说使他的大部分人民跌倒。这五个人，加上几个要么曾献祭、要么良心有愧的人，合意为自己希望福图纳图斯作假主教，以便他们的罪行一致，统治者如同被统治者一样。\par

11. 因此，最亲爱的兄弟，你现在可以知道那些绝望而堕落的人在那里散布的其他谎言了：尽管献祭者或异端者中只有不到五位假主教来到\Place{迦太基}，任命\Person{福图纳图斯}为他们疯狂的同伙；然而，他们作为魔鬼的子女，充满谎言，竟敢如你所写，夸口说有二十五位主教在场；他们在这里也曾在我们的弟兄中吹嘘这个谎言，说将有二十五位主教从\Place{努米底亚}来为他们祝圣一位主教。当他们的这个谎言被揭露和驳斥后（只有五位沉船的人聚集，且已被我们绝罚），他们带着谎言的报酬航行到\Place{罗马}，仿佛真理不能航行到他们后面，并以确凿的证据揭露他们的说谎之舌。而这，我的弟兄，是真正的疯狂：不认为也不知晓谎言不能长久欺骗，黑夜只持续到白天放明；但当白天清明、太阳升起时，黑暗和阴霾就让位于光明，夜间进行的抢劫也就停止了。总之，如果你向他们索取名字，他们甚至一个也提不出来。因为在他们中间，连恶人都是如此匮乏，以致无论是献祭者还是异端者，都无法为他们凑齐二十五人；然而，为了欺骗单纯的人和不在场者的耳朵，他们用谎言夸大数目，仿佛即使这个数目是真的，教会会被异端击败，或正义会被不义者击败。\par

12. 最亲爱的兄弟，我也不应以同样方式对待他们，并在我的论述中历数他们已经犯下和仍在犯下的事，因为我们必须考虑什么是天主的司铎所宜于说出和写下的。在我们中间，羞愧而非恼怒应当发言，我不应似乎是被激怒而堆积辱骂而非罪行和罪过。因此，我对教会的欺诈保持沉默。我略过阴谋、奸淫和各种罪行。然而，他们的邪恶中只有这一点，其原因不在于我，也不在于人，而在于天主，我认为不可隐瞒：从迫害的第一天起，当罪人的新近罪行仍然炙热，不仅魔鬼的祭坛，而且失足者的手和口都仍在冒着可憎祭献的烟雾时，他们并未停止与失足者共融，并干涉他们的补赎。天主说：\BibleQuote{凡祭献与众神，惟独主除外，必被消灭。} \BibleRef{出 22:20} 在福音中，主说：\BibleQuote{凡否认我的，我也要否认他。} \BibleRef{玛 10:33} 在另一处，天主的义怒和怒气也未沉默，说：\BibleQuote{你们向他们倒上奠祭，向他们献上素祭。我岂不因这些事发怒？主说。} \BibleRef{依 57:6} 而他们却干预，使天主不被恳求——祂自己声明祂发怒；他们阻拦，使基督不被祈祷和赔补所恳求——祂宣称祂要否认那否认祂的人。\par

13. 在迫害正盛时，我们就此事写了信，却未被理会。召开全体会议后，我们颁令——不仅凭我们的同意，也凭我们的警告——要弟兄们悔改，且不可轻易将平安给与不悔改的人。而那些亵圣者，以不敬的疯狂冲撞天主的司铎，离开教会；他们举起弑亲的手臂攻击教会，为使魔鬼的恶意完成其工作，竭力阻止天主的仁慈在祂的教会内医治受伤的人。他们以谎言的诡诈腐化可怜之人的悔改，使这悔改不能平息被冒犯的天主——使那先前以身为基督徒为耻或畏惧的人，后来不再寻求基督——他的主，也使那离开教会的人不再归回教会。他们竭力阻止罪过以合宜的补赎和哀恸得以赎偿，阻止创伤以眼泪得以洗去。真正的平安被假平安的虚假废除了；慈母的健全胸怀因继母的干预而关闭，使背教者胸中和口内的哭泣与呻吟不得听闻。除此之外，背教者被迫在从前他们犯罪的\Place{卡皮托利丘}上，用口舌辱骂司铎——以谩骂和恶言攻击宣信者、贞女，以及那些因信德赞美而最杰出、在教会中最光荣的义人。这些事所打击的，其实更多不是我们子民的谦逊、谦卑和羞耻，而是他们自己的希望和生命被撕裂。因为并非那听到的人，而是那出口辱骂的人，才是有祸的；也并非那被弟兄击打的人，而是那击打弟兄的人，才是律法下的罪人；当有罪的人对无辜者行不义时，自以为在施加伤害的人，反是承受了伤害。最后，他们的心因此受到打击，精神迟钝，正义感疏离：这是天主的义怒，使他们不察觉自己的罪，以免悔改随之而来——正如经上记载：\BibleQuote{天主给了他们沉睡的精神，}意思是，他们不能回头而被治愈，不能以合宜的祈祷和补赎在犯罪后得到痊愈。宗徒\Person{保禄}在他的书信中规定，说：\BibleQuote{他们没有接受对真理的爱，为能得救。为此，天主将一种强烈的谬误降到他们身上，使他们相信谎言，为使一切不信真理、反而喜爱不义的人都被定罪。}\BibleRef{得后2:10} 最大的福分是不犯罪；其次，是承认我们的罪过。前者中，无罪的纯洁涓涓流淌，纯全无玷以保守我们；后者中，有药物前来医治我们。他们由于冒犯天主，两者都失去了：既失去了圣洗圣化中所领受的恩宠，悔改也不来帮助他们——而悔改本可治愈罪过。兄弟，你以为他们对抗天主的恶行是微末的吗？他们的罪小而且轻吗？——因为借由他们，愤怒天主的威严不被祈求，主的日子里的愤怒和烈火不被惧怕——当假基督临近时，战斗子民的信德因被夺去基督的德能和敬畏而被解除武装？让平信徒自行留意如何改正此事。司铎们有更重的责任去宣告和维护天主的威严，以免我们在这一方面似乎有所忽略——当天主告诫我们，说：\BibleQuote{如今，司祭们，这条诫命是给你们说的。如果你们不听，不把光荣我名这事放在心上——主说——我必使诅咒降到你们身上，必诅咒你们的祝福。}\BibleRef{拉2:1} 那么，当天主的威严和法令被如此蔑视时——当祂宣告对那些献祭的人愤怒和发怒，并威胁永罚和永久的惩罚时——亵圣者却提议并说：\Quote{不要顾虑天主的义怒，不要惧怕主的审判，任何人都不要叩击基督的教会；不但废除悔改，不认罪，藐视并践踏主教，还要由长老们以欺诈的言语宣告平安；免得背教者站起来，或那些在教会外的人归回教会，竟让共融给予那些不在共融中的人？}\par

14. 对这些人而言，他们从福音中退去、剥夺了失足者获得补偿与悔改的希望、剥夺了那些卷入欺诈或被奸淫玷污、或被祭献的致命毒素污染之人向天主恳求或在教会中告解其罪行的机会——这一切还不够。他们还在外面——在教会之外，与教会对立——为自己建立一个废弃派系的聚会处，那里聚集了一群良心邪恶、不愿恳求并补偿天主的人。此外，他们竟还敢——由异端者为他们任命了一个假主教——扬帆起航，将分裂分子与亵渎之人的信件带到\Person{伯多禄}的宝座和那司铎职统一源头的首席教会；他们没有想到，这些罗马人的信德在宗徒的宣讲中受到赞扬，不忠信无法接近他们。但他们前来宣布反对主教们而设立假主教的原因是什么？要么他们对自己所做之事感到满意，坚持其邪恶；要么，若他们不满而退却，他们知道该回到哪里。因为，正如我们所有人都已定断——且同样公平合理——每个人的案件应在罪行发生之处受审；每个牧人分得一群羊的一部分，由他治理和管辖，并向主交代他的行为；那么，受我们管辖的人当然不应四处奔走，也不应以狡猾欺诈的鲁莽破坏主教们的一致和谐，而应在他们能找到对其罪行的控告者和见证人的地方提出申诉——除非在非洲设立的主教们的权威对少数绝望和堕落的人而言似乎不足，这些人已受到他们的审判，并且最近因其判断的严肃而定了他们被许多罪链束缚的良心的罪。他们的案件已被审查，判决也已宣布；司铎的尊严不应因反复无常的轻率而受指责，因为主教导说：\BibleQuote{你们的话应当：是就说是，非就说非。}\BibleRef{玛 5:37}\par

15. 如果将去年审判他们的人数与司铎和执事一起计算，那么出席审判和听审的人比现在似乎与\Person{福图纳特斯}同伙的人还要多。因为，最亲爱的弟兄，你应当知道，自从他被异端者立为假主教后，几乎所有人都立即离弃了他。因为那些在过去被欺哄、被虚假话语蒙骗，以为他们将一起回到教会的人，当他们看到那里立了一个假主教，便知道自己被愚弄和欺骗了，于是天天回来敲教会的门；与此同时，我们必须向主交代，正在焦虑地权衡和仔细检查谁该被接纳进入教会。因为有些人被他们的罪行阻碍到如此程度，或者他们受到弟兄们如此固执和坚定的反对，以致他们只能以冒犯和危及许多人的方式被接纳。因为不可将某些腐败之物收集并聚拢，以致损害健全的部分；那位牧人若不智地混入病羊以致因恶疾感染整个羊群，也毫无用处。（不要在意他们的数目。因为敬畏天主者胜过一千个不虔敬的儿子，正如主通过先知所说：\Quote{孩子，不要因不虔敬的儿子而喜悦，即便他们增多，除非上主的敬畏与他们同在。}）哦，最亲爱的弟兄，若你能与我们一起在这里，当那些邪恶悖逆的人从分裂中返回时，你就会看到我付出何等辛劳劝说我们的弟兄忍耐，使他们平息心中的悲伤，并同意接纳和医治恶人。因为正如他们欢迎和喜乐那些可容忍、罪过较轻的人回来，同样，每当那些不可救药、暴戾、被奸淫或祭献污染，且加上骄傲的人回到教会，以致败坏内部的善德时，他们就抱怨和不满。我几乎无法说服民众；不，我是从他们那里强行要求他们允许接纳这样的人。弟兄们的悲伤更有道理，因为一个又一个被我轻易接纳的人——尽管人民反对和抗议——结果比以前更坏，并且未能持守悔改的信德，因为他们没有带着真正的悔改而来。\par

16. 但我该说什么有关那些现已与\Person{Felicissimus}一起航海到你们那里的人呢？他们犯有各种罪行，是由假主教\Person{Fortunatus}派遣的使者，带给你们像他本人一样虚假的信件，他们的信件正如其良心充满罪恶、生命可憎、可耻一样；所以，即使他们在教会内，这样的人也应被逐出教会。此外，由于他们认识自己的良心，他们不敢到我们这里来，也不敢接近教会的门槛，而是在教会之外四处游荡，穿越省区，为的是欺骗和欺诈弟兄们；如今，他们已为众人所熟知，并因他们的罪行被各处排斥，他们又航行到你们那里去。因为他们无法有脸面向我们靠近，或站在我们面前，因为弟兄们指控他们的罪行极其严重和严重。如果他们愿意接受我们的审判，就让他们来。最后，如果他们能找到任何借口或辩护，让我们看看他们对于赔补有什么想法，他们带来了什么悔改的果实。教会在这里对任何人都不关闭，主教也不对任何人拒绝。我们的耐心、宽容和仁爱随时预备迎接那些来的人。我恳求所有人回到教会。我恳求我们所有的战友被收纳在基督的营垒内，并成为天主圣父的居所。我宽恕一切。我闭目不看许多事，存着聚集弟兄们的愿望和期盼。即使是那些冒犯天主的事，我也不以完全的宗教判断来追究。在赦免超过我本分的罪过时，我几乎自己犯罪。我以迅速而圆满的爱拥抱那些带着悔改回来、以谦卑和真诚的赔补承认自己罪过的人。\par

17. 但如果有些人认为他们可以不是以祈祷而是以威胁回到教会，或者认为他们能以恐怖而非哀痛和赔补为自己开路，那么让他们确信，主的教会对这些人是关闭的；基督的营垒，因主的保护而不可征服、坚定，不会向威胁屈服。天主的司铎，坚守福音，遵守基督的诫命，可以被杀，但不能被征服。天主的司铎\Person{匝加利亚}给我们提供了勇气和信德的榜样，他无法被威胁和石头砸死吓倒，在天主的圣殿中被杀害，同时呼喊并说着——我们同样也向异端者呼喊并说——\Quote{上主这样说：你们离弃了上主的道路，上主也要离弃你们。}\BibleRef{编下 24:20} 因为少数鲁莽和邪恶的人离弃了天上和健康的道路，不做圣洁的事而被圣神离弃，我们也不应因此忘记神圣的传统，以致认为疯狂之人的罪行大于司铎的判断；或以为人的努力能更有力地攻击，而天主的保护却无助于捍卫。\par

18. 最亲爱的弟兄，难道天主教的尊严要被抛弃？那忠诚无瑕的子民威严、司铎的权威与权柄，全都要被抛弃，使那些在教会之外的人可以说他们想审判教会内的主教？异端者审判基督徒？受伤的人审判完整的人？残缺的人审判健全的人？跌倒的人审判站定的人？有罪的人审判他们的审判者？亵圣者审判司铎？还剩下什么，除非教会向\Place{卡比托利欧}屈服，并且司铎们离开并移除主的祭台，而雕像和偶像连同它们的祭台进入我们神圣可敬的圣职人员的集会？并且给\Person{诺瓦提安}提供更大更充分的材料来控告和辱骂我们，如果那些祭过偶像并在公开场合否认基督的人，不仅开始被恳求和接纳而不做补赎，而且反过来以他们的恐怖权势来支配？\par

19. 如果他们想要平安，就让他们放下武器。如果他们做赔补，为何威胁？或者如果他们威胁，就让他们知道，天主的司铎并不惧怕他们。因为即使假基督开始来临，也不会因他的威胁而进入教会；我们也不会向他的武力和暴力屈服，因为他宣称如果我们抵抗就要毁灭我们。异端者武装我们，当他们以为我们被他们的威胁吓倒时；他们并没有使我们瞩目结舌，反而提升我们、激发我们，当他们使平安本身比迫害对弟兄们更可怕时。我们确实希望，他们不会以罪行来完成他们疯狂中所说的话，使那些以背信和残忍言语犯罪的人也不在行为上犯罪。我们祈祷并恳求天主——他们不断激怒的天主——愿他软化他们的心，使他们放下疯狂，恢复理智；愿他们被罪恶的黑暗所覆盖的胸膛，能认识悔改的光明，并宁愿寻求司铎的祈祷和恳求为他们倾流，而不是他们自己倾流司铎的鲜血。但如果他们继续疯狂，残忍地坚持他们这些弑亲的欺骗和威胁，那么没有一个天主的司铎如此软弱、如此沮丧、如此卑微，或如此因人性软弱的无力而无效，以至于不能因来自高天的力量而被兴起对抗天主的仇敌和攻击者；也不能找到他的谦卑和软弱被保护他的主的活力与力量所激发。我们被谁或何时杀死无关紧要，因为我们将从主那里接受我们死亡和鲜血的报酬。他们的\Emph{割裂}应被哀悼和悲叹，魔鬼使他们如此盲目，以致不顾革赫纳的永恒惩罚，企图模仿那正在临近的假基督的来临。\par

20. 最亲爱的弟兄，虽然我知道，由于我们彼此之间所欠并表现出的互爱，你常常把我的书信读给与你一同管理那里的极有声望的圣职人员，以及你那至圣且人数众多的会众；然而如今我既劝勉又请求你，按照我的要求去做你平时出于自愿和礼貌所做的事：这样，通过宣读我这封信，若任何毒害性的言语和瘟疫般的传播的感染已经潜入那里，它可以完全从弟兄们的耳中和心中被清除，而善者那健全而真挚的情感得以从异端诽谤的一切污秽中重新洁净。\par

21. 至于其余的事，让我们最亲爱的弟兄们坚决拒绝并躲避那些人的言语和交谈，因为他们的言语如同毒疮般蔓延；正如宗徒所说：\Quote{邪恶的交往败坏良好的品行。}\BibleRef{格前 15:33} 又说：\Quote{异端者，在第一次警告后，就应拒绝；因为你知道，这样的人已经走偏，而且犯罪，自己定了自己的罪。}\BibleRef{铎 3:10} 圣神也通过撒罗满说话：\Quote{乖僻的人口中携带败亡；在他的嘴唇里隐藏着火。}\BibleRef{箴 16:27} 此外，他又警告我们说：\Quote{用荆棘围住你的耳朵，不要听邪恶的舌头。} 又说：\Quote{作恶的人听从不义者的舌头；但义人不听虚谎的嘴唇。}\BibleRef{箴 17:4} 虽然我知道那里的弟兄们，必定因你的远见而得到坚固，而且因他们自己的警觉而足够谨慎，不会被异端的毒药所擒获或欺骗，并且天主的教导和诫命在他们心中之所以有力量，只是因为天主的敬畏在他们里面；然而，即使可能不必要，我的挂虑或爱心也促使我写这些给你，要我们不可与这样的人交往；不可与恶人共餐或交谈；而应与他们分离，如同他们离弃教会一样；因为经上记载：\Quote{若他拒绝听从教会，就把他当作外邦人和税吏一样看待。}\BibleRef{玛 18:17} 并且蒙福的宗徒不仅警告，而且命令我们远离这样的人。他说：\Quote{我们奉我们主耶稣基督的名命令你们，要远离每一个行为不端正、不按从我们所领受的传授而行的弟兄。} 信德与不信之间不能有相通。谁不与基督在一起，谁是基督的敌人，谁敌视基督的合一与平安，就不可与我们交往。如果他们带着祈祷和补赎而来，就应听他们；如果他们堆砌咒骂和威胁，就应拒绝他们。最亲爱的弟兄，祝你永远安康。\par

\SourceRecord{Translated by Robert Ernest Wallis.；Ante-Nicene Fathers；Vol. 5.；Edited by Alexander Roberts, James Donaldson, and A. Cleveland Coxe.；Buffalo, NY: Christian Literature Publishing Co.,；1886.；<http://www.newadvent.org/fathers/050654.htm>.}

% structure: p0001 source=h1 target=section reason=page-title

\section{\WorkTitle{第五十五封书信}}

致\Place{提巴里斯}人民：劝勉殉道。\par

\Emph{西彼廉首先向\Place{提巴里斯}人解释他未能访问他们的原因，并警告即将来临的迫害，然后提供欣然接受殉道的动机。}\par

1. \Person{西彼廉}致寄居在\Place{提巴里斯}的人民，问候。亲爱的弟兄们，我确实曾经想过并热切祈求——如果事态与时机允许，正如你们时常所愿——亲自到你们那里去，并亲临你们当中，以我所拥有的那点微弱的劝勉之力坚固弟兄团体。然而，我被如此紧迫的事务所羁绊，无法远行此地，也无法长久离开那蒙天主仁慈托付给我的人民，因此我先写了这封信，代我来到你们当中。因为，蒙主屈尊教导我，我屡受激励与警戒，我理应将我的警戒之忧也带给你们的良心。因为你们应当知道、相信并确知：苦难之日已开始悬于我们头顶，世界末日与\OriginalTerm{假基督}{Antichrist}的日子临近，因此我们都必须准备好作战；除了永生之荣耀与承认主的冠冕，不要思虑任何事物；也不要看那将要来临的如同已逝的那样。一场更严厉、更激烈的战斗正在威胁，基督的士兵应当以纯全的信德与刚毅的勇气预备自己，因为他们每日畅饮基督的血杯，为的是他们自己也能为基督流血。因为这就是愿意与基督同在，就是效法基督所教导和所行的，正如宗徒若望所说：\BibleQuote{那说自己住在祂内的，就应当照祂所行的去行。}此外，蒙福的宗徒保禄劝勉并教导说：\BibleQuote{我们既是子女，便是承继者，是天主的承继者，是基督的同承继者；只要我们与祂一同受苦，也必要与祂一同受光荣。}\par

2. 这一切现在都必须为我们所深思，好叫没有人贪求那正在死去之世界的任何事物，而是要跟随那永远活着并使祂的仆人复活的基督，这些仆人因信祂的名而得以坚固。因为时候到了，亲爱的弟兄们，我们的主早已预告并教导我们临近了，祂说：\BibleQuote{时候要到，凡杀害你们的，还以为是尽恭敬天主的义务。他们这样做，是因为没有认识父，也没有认识我。我告诉你们这一切，是为叫你们到时候想起我早就告诉了你们。}任何人都不必惊讶我们常受迫害，并不断被加重的磨难所考验，因为主早已预言这些事将在末世发生，并借祂话语的教导与劝勉装备我们作战。祂的宗徒伯多禄也教导说，迫害的发生是为了考验我们，我们也应藉着以往义人的榜样，借着死亡与苦难与天主的爱结合。因为他在书信中写道：\BibleQuote{亲爱的，你们不要因为在你们中间的火炼而惊异，好像遭遇了什么新奇的事；反而要喜欢，因为你们分受了基督的苦难，这样，在祂光荣显现时，你们也能欢欣踊跃。如果你们为了基督的名字而受辱骂，便是有福的，因为光荣的神，即天主的神，就安息在你们身上。}宗徒们教导我们的，正是他们自己从主的训示和天主的命令中所学到的，主亲自这样坚固我们，说：\BibleQuote{没有一个人为了天主的国舍弃了房屋、田地、父母、兄弟、姊妹、妻子或儿女，不在今世得到百倍，并在来世获得永生的。}祂又说：\BibleQuote{几时，为了人子的缘故，人恼恨你们，并弃绝你们，并且以你们的名字为恶的而加以辱骂，你们才是有福的。在那一天，你们要欢喜跳跃，因为你们的赏报在天上是丰厚的。}1\par

3. 主愿意我们在迫害中欢喜跳跃，因为当迫害发生时，信德的冠冕被赐予，天主的士兵被考验，天堂为殉道者敞开。因为我们参与这战争并非只应想到和平而退缩拒绝战争，因为在这战争本身中，主首先行在前——祂是谦逊、坚忍与受苦的师傅——为使我们所行与所劝勉所受的，祂先自己为我们忍受了。亲爱的弟兄们，你们眼前要看见：那唯独从父接受审判、并将要来审判的那一位，已经宣告了祂审判与将来承认的谕令，预告并见证：在父前承认祂的，祂也要承认；否认祂的，祂也要否认。如果我们能逃脱死亡，我们就有理由惧怕死亡。但既然反之，凡有死的人必然要死，我们就当拥抱那借着天主应许与屈尊而来的时机，以不朽的酬报完成死亡所定的结局；也不要惧怕被杀，因为我们确知被杀时必得冠冕。\par

4. 任何一个人，亲爱的弟兄们，当他看到我们的人民因迫害的恐惧而被驱散流亡时，不要因看见弟兄们聚集、主教们讲论而惊慌。\newline{}并非所有人都能同时在场，他们不会杀害人，却必须被杀。无论何时，在那几天，因时势所迫，每位弟兄的身体暂时与羊群分离，但心灵不离；不要为那逃亡的恐惧所动摇；即使他退避隐藏，也不要因旷野的孤寂而惊慌。他并不孤单，因为\Person{基督}与他一同逃亡；他并不孤单，因为他无论在哪里都保存着天主的殿宇，并不缺少天主。若在孤寂中或山野里，一个逃犯遇上了强盗；若野兽袭击你；若饥渴寒冷困扰你，或风暴暴雨在海上急速航行中倾覆你，\Person{基督}处处注视着祂的士兵作战；为了迫害的缘故，为了祂圣名的荣耀，当士兵死去时祂必奖赏他，正如祂在复活中所应许的。\OriginalTerm{殉道}{martyrdom}的荣光并不因他不是当众、在众人面前死去而减少，因为死的原因是为基督而死。那位\OriginalTerm{见证人}{Witness}证明并加冕\OriginalTerm{殉道者}{martyr}，祂对殉道者所作的见证便已足够。\par

5. 亲爱的弟兄们，让我们效法义人\Person{亚伯尔}，他开创了\OriginalTerm{殉道}{martyrdom}，首先为了正义而被杀。让我们效法\Person{亚巴辣罕}，天主的朋友，他毫不犹豫地亲手献上自己的儿子为祭，以虔敬的\OriginalTerm{信德}{faith}服从了天主。让我们效法三个青年\Person{阿纳尼雅}、\Person{阿匝黎雅}和\Person{米沙耳}，他们既不为年少的年龄所吓倒，也不为被掳所压垮——\Place{犹大}被征服、\Place{耶路撒冷}被攻取——却在自己的国土上以信德之力胜过了君王；当奉命朝拜\Person{拿步高}王所立的金像时，他们挺身而出，比君王的威胁和火焰更坚强，以这些话语呼喊并证明了他们的信德：\BibleQuote{拿步高王，我们无须为此事回答你。因为我们所事奉的天主，能将我们从烈火窑中救出；祂也必从你手中救出我们，大王。即或不然，你当知道，我们决不事奉你的神，也不朝拜你所立的金像。}他们相信凭着信德可以逃脱，却又说：\BibleQuote{即或不然}，好让王知道他们也能为所事奉的天主而死。这便是勇气和信德的力量：相信并知道天主能从现世的死亡中解救，却不惧怕死亡，也不退缩，好使信德更受考验。\OriginalTerm{圣神}{Holy Spirit}那无玷、不可战胜的力量从他们的口中迸发出来，以致主在福音中所说的话被证明是真实的：\BibleQuote{当他们捉拿你们时，不要思虑说什么；因为在那时刻，必赐给你们当说的话。因为说话的不是你们，而是你们父的圣神在你们内说话。}祂说，我们所能说和回答的，在那时刻是从天上赐给我们的，并供应我们；那时不是我们在说话，而是我们父天主的圣神——祂从不离开也不分离于那些承认祂的人——祂亲自在我们内说话并受加冕。同样，\Person{达内尔}奉命朝拜当时人民和君王所敬拜的\Person{贝尔}偶像时，他彰显了天主的荣耀，以圆满的信德和自由呼喊说：\BibleQuote{我只朝拜我的天主上主，祂创造了天与地。}\par

6. 关于\BibleBook{玛加伯}{书}中可敬的\OriginalTerm{殉道者}{martyr}所受的残酷折磨，以及七位兄弟的多种苦难，还有那位母亲在痛苦中安慰她的儿子们，并最终与儿子们一同死去，我们还有什么可说的呢？它们岂非见证了伟大的勇气和\OriginalTerm{信德}{faith}，并以他们的苦难激励我们走向殉道的光荣？至于\OriginalTerm{先知}{prophet}们，圣神激发他们预知未来之事，又当如何？至于\OriginalTerm{宗徒}{Apostle}们，主所拣选的，又当如何？这些义人为了正义而被杀，岂不也教导我们赴死？\Person{基督}的诞生同时见证了婴孩们的殉道，凡两岁及以下的婴儿都因祂的圣名而被杀。一个尚不适合作战的年龄，却已堪当冠冕。为显明因基督而被杀者是无辜的，无辜的婴孩为祂的圣名而被处死。这显明无人能免于迫害的危险，因为连这些婴孩也完成了殉道。然而，一个基督徒若作为仆人，竟不愿受苦，而他的主人首先受苦，这是何等严重的事！我们不愿为自己的罪受苦，而祂自己本无罪，却为我们受苦！天主圣子受苦，为使我们成为天主的儿女；而人子若不愿受苦，怎能继续成为天主的儿女？如果我们因世界的憎恨而受苦，基督首先忍受了世界的憎恨。如果我们在这世上忍受辱骂、流放、酷刑，世界的创造者和主宰经历了比这些更艰苦的事，祂也警戒我们说：\BibleQuote{世界若恨你们，你们要知道，它恨你们以前，已经恨了我。你们若是属于世界，世界必爱属于自己的人；但因为你们不属于世界，而是我从世界中拣选了你们，因此世界才恨你们。你们要记得我对你们说过的话：仆人不能大于主人。如果他们迫害了我，也要迫害你们。}我们的主、天主所教导的一切，祂也实行了，好使\OriginalTerm{门徒}{disciple}若只学习而不实行，便无从推辞。\par

7. 你们中任何人，亲爱的弟兄，都不要因对未来迫害的恐惧，或即将来临的威胁性假基督，而如此惊慌，以致没有因福音的劝勉和规诫，以及天上的警告而为一切事武装起来。假基督即将来临，但基督在他之上也来临。1敌人四处游荡并咆哮，但主立刻跟随而来，为我们的痛苦和创伤复仇。对手愤怒并威胁，但有一位能从他的手中解救我们。那位无人能逃避其愤怒的，应当被敬畏，正如祂亲自预先警告并说：\BibleQuote{不要害怕那些杀害身体而不能杀害灵魂的；但更要害怕那能把灵魂和身体都投入地狱里的。}\BibleRef{玛 10:28}1 祂又说：\BibleQuote{爱惜自己生命的，必要丧失生命；在现世憎恨自己生命的，必要保存生命到永生。}\BibleRef{若 12:25}1 祂在\WorkTitle{默示录}中教导并预先警告说：\BibleQuote{若有人朝拜那兽和它的像，并在额上或手上接受了它的印记，这人必要喝天主愤怒的酒，即倒入祂忿怒之杯中的纯酒；他必在圣天使面前和羔羊面前，受火与硫磺的折磨；他们折磨的烟上升，直到永远；朝拜那兽和它的像的，昼夜不得安宁。}\BibleRef{默 14:9-11}2\par

8. 世俗的竞赛，人训练并准备，若能在人民眼前、在皇帝面前受戴冠，便视为荣耀中的极大荣耀。看啊，一个崇高而伟大的竞赛，也因着天上冠冕的赏报而荣耀，因为当我们奋斗时，天主注视着我们，将祂的视线扩展到那些祂屈尊使之成为祂儿子的人，祂享受我们竞赛的景象。天主在战事中注视我们，在信仰的冲突中战斗；祂的天使注视我们，基督也注视我们。这尊严何其大，这荣耀的幸福何其大，在天主面前参加竞赛，并戴冠，以基督为审判者！亲爱的弟兄，让我们以全部力量武装起来，以不败坏的心思、健全的信德、忠诚的勇气，为斗争预备自己。让天主的营队出发到为我们指定的战场。让健全的人武装起来，免得健全的人失去刚站稳的优势；让跌倒的人也武装起来，使跌倒的人可以重获他所失去的：让荣誉激励全体；让悲伤激励跌倒的人投入战斗。宗徒保禄教导我们武装并准备，说：\BibleQuote{我们不是对抗血肉，而是对抗权势，对抗这世界和这黑暗的首领，对抗天界中邪恶的鬼神。所以要穿上全副武装，使你们能在邪恶的日子抵挡，使你们做完一切后仍能站立；用真理束上你们的腰，穿上正义的胸甲；以和平福音的预备穿上你们的脚；拿起信德的盾牌，借此你们能熄灭那恶者的一切火箭；戴上救恩的头盔，和圣神的剑，就是天主的话。}\BibleRef{弗 6:12-17}2\par

9. 让我们拿起这些武器，让我们用这些属神和天上的保障加强自己，使我们在邪恶的日子能够抵挡并抵抗魔鬼的威胁：让我们穿上正义的胸甲，使我们的胸膛得以加强并安全抵挡敌人的箭矢：让我们的脚穿上福音教导并武装，使得当那蛇开始被我们践踏和粉碎时，他不能咬我们并绊倒我们：让我们勇敢地拿起信德的盾牌，借着它的保护，敌人射向我们的任何东西都能被熄灭：让我们也戴上救恩的头盔保护我们的头，使我们的耳朵得以保护，免听致命的法令；使我们的眼睛得以加强，免得看见可憎的偶像；使我们的额头得以加强，以保全天主的标记；2使我们的口得以加强，使得胜的舌头能宣认基督为其主：让我们也用圣神的剑武装右手，使它勇敢地拒绝致命的祭献；使那记住圣体、曾接受主身体的右手，2能拥抱主自己，以后从主那里接受天上冠冕的赏报。\par

10. 哦，那日子来临之际，将是何等何等地大啊，亲爱的弟兄，当主开始清点祂的子民，并借着祂神圣知识的察看，认出每个人的功过，将罪人遣送到地狱，并用惩罚之火的永久燃烧点燃我们的迫害者，却将我们的信德和虔敬的赏报付给我们！那将是何等的荣耀，何等的喜乐，获准看见天主，受光荣与基督、你的主天主一起接受永生的救恩和光明的喜乐——迎接\Person{亚巴郎}、\Person{依撒格}、\Person{雅各伯}，以及所有圣祖、先知、宗徒和殉道者——与义人和天主的朋友们在~天国~中欢欣，借着赐予我们的不死之乐——在那里接受眼所未见、耳所未闻、人心未曾想到的事物！因为宗徒宣告，我们将接受比我们在此所做或所受的一切更大的事物，说：\BibleQuote{现今时代的苦难，不配与将来要在我们身上显明的荣耀相比。}\BibleRef{罗 8:18}2 当那启示来到时，当天主的荣耀照耀我们时，我们将如此幸福和喜乐，因天主的屈尊而受光荣，正如那些作魔鬼之意的，无论是离弃天主还是反抗祂的人，仍将有罪和悲惨，以致他们必须与魔鬼一同在不灭的火中受折磨。\par

11. 亲爱的弟兄们，让这些事占据我们的心；让这成为我们武装的准备，成为我们日日夜夜的默想：将恶人的惩罚与义人的赏报和功德常放在眼前，并不断在心绪和情感中反复思量：主以惩罚威胁那些否认祂的人；另一方面，祂以光荣应许那些承认祂的人。如果我们正思想并默想这些事，而迫害之日来到，那么基督的士兵，因受祂的诫命和警告的教导，就不会害怕战斗，而是为冠冕做好准备。我向你们，最亲爱的弟兄们，诚心告辞，愿你们永远安好。\par

\SourceRecord{Translated by Robert Ernest Wallis.；Ante-Nicene Fathers；Vol. 5.；Edited by Alexander Roberts, James Donaldson, and A. Cleveland Coxe.；Buffalo, NY: Christian Literature Publishing Co.,；1886.；<http://www.newadvent.org/fathers/050655.htm>.}

% structure: p0001 source=h1 target=section reason=page-title

\section{书信 56}

致流放中的\Person{科尔乃略}，论及他的宣认。\par

\Person{西彼廉}称赞\Person{科尔乃略}及其子民因基督之名甚至被流放而作的宣认；并劝勉他们持守恒心，彼此代祷，关乎此生即将来临的争战之日，以及死后亦然。\par

西彼廉致其兄弟科尔乃略，祝安。最亲爱的兄弟，我们已得知你们信德与勇气的光荣见证，并满怀喜悦地领受了你们宣认的尊荣，以致我们自认也是你们功绩与赞美的分享者和同伴。因为我们有同一教会，同一心灵，同一不分裂的和谐，哪一位司铎不因同僚司铎的赞美而如同自己的赞美般自喜？哪一位兄弟团体不因其兄弟的喜乐而欢欣？当我们听闻你们的成功与英勇，得知你们在那里的弟兄中成为宣认的领袖，且领袖的宣认因弟兄们的同意而增强，以致你们先行迈向荣耀，使许多人成为你们荣耀的同伴，并因首先预备为众人宣认而劝服百姓成为宣信者——这一切，我们难以充分表达这里的喜乐和欢欣是何等巨大！以至我们不知首先应称赞你们哪一点：是你们那迅捷而坚定的信德，还是弟兄们那不可分离的爱。在你们中间，主教领先的勇气已获公开证明，跟随的弟兄们的团结也已彰显。如同你们拥有同一心灵、同一声音，整个罗马教会已经宣认。\par

最亲爱的兄弟，蒙福的宗徒曾在你们中间称赞的信德如今大放光芒。他那时已在精神上预见了这勇气的赞美和力量的坚定；并通过称赞你们未来的作为来见证你们的功绩，在赞扬父母时激励子女。你们既如此同心，又如此勇敢，已给其余弟兄树立了同心与勇敢的伟大榜样。你们教导他们要深切敬畏天主，坚定依附基督；使百姓与司铎在危难中联合；使弟兄在迫害中不与弟兄分离；使已建立的和谐绝不能被打败；使众人同时祈求的一切，和平的天主必赐予和平之人。仇敌曾猛然跃起，欲以暴力惊吓基督的营盘；但他以同样猛烈的冲势被打退并征服；他带来了多少恐惧和惊慌，也找到了同样多的勇气和力量。他原以为能再次推翻天主的仆人，并依其惯常方式搅动他们，仿佛他们是新手、无经验、预备不足、不够谨慎。他先攻击一人，如同狼试图将羊从羊群中分离，如同鹰试图将鸽子从飞行队伍中分离；因为那对全体没有足够力量的人，试图从个体的孤立中获取优势。但当被联合军队的信德和活力击退后，他察觉到基督的士兵如今正在警醒，清醒地武装着准备战斗；他们不能被征服，但可以死；而正因他们不惧怕死亡，他们才不可战胜；他们不反过来攻击攻击者，因为连无辜者也不可杀害有罪之人；但他们甘愿献出生命和鲜血；既然如此恶毒和残酷在世上肆虐，他们可以更快地脱离邪恶和残酷。在天主眼前，那是何等光荣的景象！在基督面前，祂的教会何等喜乐——不是单个士兵，而是整个营盘，立即奔赴仇敌企图开始的战斗！因为显然，若他们能听到，所有人都会前来，因为凡听到的人都急忙赶来。有多少背教者借由光荣的宣认得以恢复！他们勇敢地站立，并借着悔改本身所受的苦痛，在战斗中变得更勇敢，这表明他们先前是被突然袭击，因对新颖不惯之事恐惧而战栗，但后来回心转意；真正的信德和从敬畏天主中积聚的力量，已恒常坚定地坚固他们承受一切；如今他们不是为罪过求宽恕，而是为苦难求冠冕。\par

针对这些事，最亲爱的兄弟，诺瓦替安会说什么？他难道还不放下他的错误吗？或者，正如愚人的惯常做法，他因我们的利益和兴旺而变得更加狂怒；并且随着爱与信德的荣耀日益增长，纷争与嫉妒的疯狂在那里重新爆发？那可怜的人非但不医治自己的伤口，反而更严重地伤害自己和朋友，用舌头喧嚣以毁灭弟兄，投掷有毒的雄辩之矛，其凶恶根据世俗哲学的邪恶性比照主智慧的温良更为严厉——他是教会的逃兵，慈悲的仇敌，悔改的毁灭者，傲慢的教师，真理的败坏者，爱德的凶手？他现在承认谁是主的天主司铎，哪一个是教会和基督的房屋，谁是天主的仆人（魔鬼折磨他们），谁是基督徒（假基督攻击他们）？因为那已征服的人，他不去寻找；那已归为己有的人，他不费心去推翻。教会的仇敌和敌人蔑视并绕过那些被他从教会疏远、作为俘虏和征服者带出去的人；他继续骚扰那些他看见基督居住在其中的人。\par

4. 即使这样的人中有谁被抓住了，他也没有理由沾沾自喜，仿佛这是宣认了主名；因为显而易见，如果这类人在教会之外被处死，那不是信德的冠冕，而是背信的惩罚。那些因不协的疯狂而从和平神圣的家中退出的人，我们看见他们已离开了天主殿宇中同心合意的人群。\par

5. 最亲爱的弟兄，因着我们彼此联结的互爱，我们尽量恳切劝勉：既然主的眷顾教导我们，天主的慈悲也以有益的劝告提醒我们，我们争战和奋斗的日子已经临近，我们就当不断以斋戒、守夜、祈祷勉励全体民众。让我们以不断的呻吟和频密的祈祷迫切祈求。因为我们这是天上的武器，使我们坚定站立，勇敢持守。这些是保护我们的属灵防御和神圣武器。让我们在和谐与同心合意中彼此记念。让我们双方时常为彼此祈祷。让我们以互爱减轻负担和痛苦，以便我们中若有人因天主的迅速垂顾而先离世，我们的爱能在主面前持续，我们对弟兄姐妹的祈祷在天父的慈悲前不止息。至亲爱的弟兄，祝你永远衷心平安。\par

\SourceRecord{Translated by Robert Ernest Wallis.；Ante-Nicene Fathers；Vol. 5.；Edited by Alexander Roberts, James Donaldson, and A. Cleveland Coxe.；Buffalo, NY: Christian Literature Publishing Co.,；1886.；<http://www.newadvent.org/fathers/050656.htm>.}

% structure: p0001 source=h1 target=section reason=page-title

\section{书信 57}

致罗马主教路爵\newline{}从放逐归来\par

西彼廉与其同僚恭贺路爵从放逐中归来，提醒他延迟的殉道并不会减损光荣。接着，指出科尔乃略的殉道和路爵的放逐是出于天主的安排，为使诺瓦提安派蒙羞；他预言路爵即将殉道，是天主如此安排，使殉道不在异地完成，而在本族人民中完成。\par

1. 西彼廉与其同僚问候他的弟兄路爵安好。最近，亲爱的弟兄，我们确实也曾恭贺你，当天主的上智以双重的荣誉，在天主的教会的管理中任命你同时为宣信者与司铎时。但现在我们同样恭贺你、你的同伴以及全体弟兄，因为主的仁慈而自由的护佑已带着同样的光荣和赞美将你带回祂自己的羊群中；如此，牧者得以归回喂养他的羊群，舵手得以管理船只，统治者得以治理人民；并且显明你的放逐是如此神圣地安排，并非为了使被放逐和驱逐的主教缺席教会，而是使他比离开时更伟大地回到教会。\par

2. 因为三位青年的殉道荣耀并未因他们的死亡被挫败而从火窑中安然出来而减少；达尼尔也并非因被投给狮子为食却蒙主保护而活着获得荣耀，而在应得的赞美上有所欠缺。在基督的宣信者中，延迟的殉道并不减损宣信的功绩，反而显明天主护佑的伟大。我们在你身上看到了那些勇敢而杰出的青年在国王面前所宣告的：他们确实准备被火焚烧，以免事奉他的神或敬拜他所造的像；但他们所敬拜的天主，也是我们所敬拜的，能够将他们从火窑中拯救出来，从国王的手中以及迫近的苦难中解救出来。我们现在发现，这在你宣信的信德和主的护佑中实现了：当你准备并愿意承受一切惩罚时，主却使你免于惩罚，并为教会保存了你。在你归来时，宣信的光荣并未在主教的身上减损，司铎的权威反而增长；如此，一位司铎在天主的祭台前服务，他鼓励人民拿起宣信的武器，接受殉道，因为敌基督近了，他预备战士上阵，不仅凭他言论和话语的迫切，更凭他信德和勇气的榜样。\par

3. 亲爱的弟兄，我们以心灵的全部光明理解并领悟了神圣尊威那救恩且神圣的计划，为何最近在那里突然爆发了迫害——为何世俗权力突然对基督的教会和真福殉道者科尔乃略主教以及你们所有人发作；如此，为使异端者蒙羞并被击倒，主显明了哪一个是教会——哪一个是祂所选定的唯一主教——哪些是与主教在司祭尊荣中联合的司铎——谁是基督联合且真正的子民，在主羊群的爱中相连——谁是敌人要骚扰的人，而魔鬼却饶恕他如同自己的。因为基督的敌人只迫害和攻击基督的营地和士兵；异端者，一旦被征服并成为他的，他就藐视并放过。他试图打倒那些他看见站立的人。\par

4. 亲爱的弟兄，我愿如今能赐予我们力量，在你归来时与你同在，好使我们这些彼此相爱的人，能亲身与其余人一起，也领受你归来那极为喜乐的果实。所有弟兄们在那里如何欢腾！每个人到来时如何奔跑和拥抱！你几乎无法满足于那些紧抱你的人的亲吻；人们的面容和双眼几乎无法因看见你而满足。因你归来的喜乐，那里的弟兄们已开始认识到，当基督来临时，将有怎样伟大的喜乐随之而来。因为祂的来临即将临近，一种预演如今已在你身上先行了；正如若翰，祂的前驱和预备道路者，来宣讲基督已经来临，同样，如今一位主教作为主的宣信者和司铎归来，似乎主也正在归来。但我、我的同工以及全体弟兄，以此信代替我们寄给你，亲爱的弟兄；并在信中向你表达我们的喜乐，同时在此也在我们的祭献和祈祷中表明我们爱的忠诚心意，不住地感谢天主圣父及祂的圣子、我们的主基督；并祈祷恳求祂，那完美且使人完美的，在你身上保持并完善你宣信的光荣冠冕，祂或许正是为此召你归来，以免你的荣耀在异地完成殉道而被隐藏。因为那为弟兄们提供勇气与信德榜样的牺牲，应当在弟兄们临在时奉献。亲爱的弟兄，我们衷心祝愿你常安好。\par

\SourceRecord{Translated by Robert Ernest Wallis.；Ante-Nicene Fathers；Vol. 5.；Edited by Alexander Roberts, James Donaldson, and A. Cleveland Coxe.；Buffalo, NY: Christian Literature Publishing Co.,；1886.；<http://www.newadvent.org/fathers/050657.htm>.}

% structure: p0001 source=h1 target=section reason=page-title

\section{第五十八封书信}

\Emph{致\Person{菲杜斯}：论婴儿圣洗}\par

\Emph{论据——在此书信中，\Person{西彼廉}并非制定任何新法令；而是坚定持守教会的信德，为纠正那些认为婴儿在出生后第八天之前不得受洗的人，他与几位同僚主教一同裁定：婴儿一出生便可合法地受洗。然而，他借此机会拒绝撤销已赐予一位名叫\Person{维克多}之人的平安，尽管该平安是违反有关背教者的会议法令而赐予的；但他禁止\Person{德拉庇乌斯}主教在其他情况下这样做。}\par

1. \Person{西彼廉}，以及当时出席会议的其他同僚，共六十六位主教，致我们的兄弟\Person{菲杜斯}，问安。最亲爱的兄弟，我们已收到你的书信，你在信中提及关于原先的司铎\Person{维克多}一事：我们的同僚\Person{德拉庇乌斯}鲁莽地过早、过于急切地赐予他平安，而此人尚未完全悔改，也未向他所犯罪过的上主天主做出补偿；此事颇令我们不安，因为这是背离我们法令的权威，在合法的和完全的补赎时期之前，且未经民众请求和知情——并无疾病催促，亦无必要迫使——就赐予他平安。但经过长时间权衡，我们认为足以谴责我们的同僚\Person{德拉庇乌斯}行事鲁莽，并指示他不得再对任何其他人做同样的事。然而，我们认为不应撤销一位天主的司铎一旦赐予的平安，因此允许\Person{维克多}享有已赐予他的共融。\par

2. 然而，关于婴儿一事，你说他们不应在出生后第二或第三天受洗，并且应遵守古代割损礼的法律，以致你认为刚出生的人不应在第八天之前受洗和圣化，我们的议会在这一点上意见截然不同。因为在你所认为应当采取的做法上，无人赞同；我们全体反倒认为，天主的仁慈和恩宠不应拒绝给任何出生于人的人。因为正如主在福音中所说：\BibleQuote{人子来，不是要毁灭人的性命，而是要拯救他们。}\BibleRef{路 4:56}，我们应尽可能努力，务使任何灵魂都不至丧亡，若可能的话。因为那在母腹中一经天主亲手形成的人，还缺少什么呢？对我们而言，且在我们眼中，按照世界日子的进程，出生的人似乎得到增长。但凡是天主所造的一切，都因天主造物主的威严和工程而得以完成。\par

3. 此外，对圣经的信仰向我们宣告，在所有的人中，无论是婴儿还是较年长者，天主的恩赐都有同等的平等。\Person{厄里叟}恳求天主，这样俯伏在那寡妇的儿子身上，那孩子已死躺着：他的头对着孩子的头，他的脸对着孩子的脸，\Person{厄里叟}的四肢伸展并贴合在孩子每一个肢体上，他的脚对着孩子的脚。若从我们出生和身体的不平等来考虑这件事，一个婴儿不能与一个长大成熟的人相等，他细小的肢体也无法适应并等于较大的肢体。但这里所表达的是神圣的、属灵的平等，即所有的人都相似且平等，因为他们都是一次由天主所造；我们的年龄在身体增长上可能有差异，但这是按世界，而非按天主；除非那赐予受洗者的恩宠本身也按照接受者的年龄而或多或少，然而圣神不是按分量赐予的，而是因圣父的爱和仁慈同等地赐予所有人。因为天主，正如不偏袒人，也不偏袒年龄；祂以均衡的平等向所有人显示自己为父，以获得天上的恩宠。\par

4. 至于你所说，婴儿在出生初期的面容不洁净，以致我们中任何人仍会因亲吻它而畏惧，我们认为这不应成为获得天上恩宠的障碍。因为经上记载：\BibleQuote{为洁净者，一切都是洁净的。}\BibleRef{铎 1:15} 我们中谁也不应畏惧天主所屈尊创造的事物。因为虽然婴儿刚出生，但并非如此不堪，以致有人会在赐予恩宠和缔造和平时因亲吻它而畏惧；因为在亲吻婴儿时，我们每个人都应为着宗教本身的缘故，思量那仍崭新的天主亲手，我们在某种意义上正在亲吻这些手，当我们拥抱那由天主所造、最近形成并新出生的人时。关于犹太人的肉身割损礼中第八天的遵守，在影子和用法中预先赐下了一个圣事；但当基督来到时，这圣事在真理中得以实现。因为第八天，即安息日之后的第一天，正是主要复活、要使我们活过来、并赐予我们心神割损的那一天，所以第八天，即安息日之后的第一天，亦即主的日子，在预象中先行出现；当真理随后到来，并赐予我们属灵的割损时，这预象便止息了。\par

5. 因此我们认为，不应因那已制定的法律而阻碍任何人获得恩宠，肉身的割损也不应阻碍心灵的割损，而是绝对应当让每一个人都领受基督的恩宠，因为伯多禄在\WorkTitle{宗徒大事录}中也说：\Quote{主对我说：你们不要把任何人称为凡俗或不洁的。} \BibleRef{宗 10:28} 但若有什么能阻碍人们获得恩宠，那么更严重的罪过反而更会阻碍那些成熟、长大和年长的人。然而，如果连最大的罪人，以及那些曾严重冒犯天主的人，在他们后来相信时，也获得了罪的赦免——没有人被阻止领受圣洗和恩宠——那么，我们更应避免阻止一个婴儿，他刚出生，尚未犯罪，只是因按照亚当血肉而生，在最初出生时就感染了古时死亡的毒害；正因如此，他更容易获得罪的赦免——因为赦免给他的不是他自己的罪，而是别人的罪。\par

6. 因此，最亲爱的弟兄，这正是我们在公议中的意见：我们不应阻碍任何人领受圣洗和天主的恩宠，祂对众人是仁慈、良善和慈爱的。既然这应在所有人身上遵守和维护，我们认为更应在婴儿和新出生的人身上遵守，正因如此，他们更应得到我们的帮助和天主的怜悯，因为他们一出生就哭泣哀号，除此之外别无他求。最亲爱的弟兄，祝你永远安康。\par

\SourceRecord{Translated by Robert Ernest Wallis.；Ante-Nicene Fathers；Vol. 5.；Edited by Alexander Roberts, James Donaldson, and A. Cleveland Coxe.；Buffalo, NY: Christian Literature Publishing Co.,；1886.；<http://www.newadvent.org/fathers/050658.htm>.}

% structure: p0001 source=h1 target=section reason=page-title

\section{第五十九封书信}

\Emph{致努米底亚主教们，论从蛮族手中赎回被掳的弟兄。}\par

\Emph{前言——\Person{西彼廉}首先哀悼他听努米底亚主教们说起的弟兄被掳之事，并说他将寄送由弟兄、姐妹和同工们捐献的十万塞斯特斯。}\par

\Person{西彼廉}致\Person{雅努阿里乌}、\Person{马克西穆斯}、\Person{普罗库鲁}、\Person{维克托}、\Person{莫迪安努}、\Person{内梅西安努}、\Person{南普卢斯}和\Person{霍诺拉图斯}诸位弟兄，安。最亲爱的弟兄们，我怀着极度悲痛的心情，并非没有眼泪，读了你因爱心关怀而写给我关于我们弟兄姊妹被掳的信。因为谁不为这类不幸悲伤呢？谁不视弟兄的悲伤为自己的悲伤呢？正如宗徒保禄所说，\BibleQuote{若一个肢体受苦，所有的肢体都一同受苦；若一个肢体得荣耀，所有的肢体都一同欢乐；}\BibleRef{格前 12:26}又在另一处说，\BibleQuote{谁软弱，我不软弱呢？}\BibleRef{格后 11:29}因此，现在我们也必须把弟兄的被掳看作是被掳，那些身处危险之人的悲伤也应视为我们的悲伤，因为我们的结合确实是一个身体；不仅爱，而且宗教也应激励并坚固我们去赎回弟兄的肢体。\par

2. 因为宗徒保禄又说：\BibleQuote{你们不知道你们是天主的宫殿，天主圣神住在你们内吗？}\BibleRef{格前 3:16}——即使爱较少地催促我们去帮助弟兄，在此我们也必须考虑，被掳的是天主的宫殿，我们不应因长期怠惰和忽视他们的苦难而让天主的宫殿长久被掳，而应尽我们所能努力，并在我们的服从上迅速行动，以赢得基督我们的审判者、主和天主的喜悦。因为正如宗徒保禄所说：\BibleQuote{凡你们因基督受洗的，都穿上了基督}\BibleRef{迦 3:27}，所以基督应在被掳的弟兄身上被默观，祂要从被掳的危险中被赎出，祂曾从死亡的危险中救赎了我们；因此，那把我们从魔鬼的口中取出、现在仍在我们内居住和停留的祂，现在应当用金钱从蛮族手中被解救和赎回——祂曾用十字架和血救赎了我们——祂允许这些事发生，是为了考验我们的信德，看我们每个人是否会为他人做他愿意为自己做的事，如果他自己被蛮族俘虏的话。因为谁若是念及人性，并被提醒彼此相爱，如果他是父亲，难道不立刻认为他的儿子在那里？如果他是丈夫，难道不想到他的妻子被俘虏在那里，带着同等的悲伤和婚姻束缚的羞耻？但我们所有人共同的悲伤和对处于危险中的贞女的忧虑是何等巨大！为她们，我们不只要哀叹自由的丧失，还要哀叹贞操的丧失；我们更要哀叹蛮族的锁链少于勾引者和可憎之处的暴力，以免那些奉献给基督、并因贞洁的美德永远在贞洁的光荣中献身的肢体，受到侮辱者的亵渎和污秽的玷污。\par

3. 我们的兄弟团体根据你的信考虑这一切，并悲伤地审查，都迅速、情愿、慷慨地为弟兄们聚集了金钱供应，他们始终按照信德的力量倾向于天主的工作，但现在更因如此巨大的苦难的考虑而受到激励去从事有益的善工。因为主在福音中说：\BibleQuote{我病了，你们来探望了我}\BibleRef{玛 25:36}，现在祂将说，\Quote{我被掳，你们赎了我！}这工作将带来多大的报酬啊！又因为祂说：\Quote{我在监里，你们来看我}，更何况当祂开始说：\Quote{我被囚禁在俘掳的牢房，躺在蛮族中受捆绑，你们从那为奴的监牢中释放了我}，你们将在审判之日从主那里得到报酬！最后，我们衷心感谢你，因为你愿意让我们分享你的忧虑，并在如此伟大而必要的工作中——你为我们提供了肥沃的田地，我们可以在那里播下希望的种子，期待从这天国和救赎的行动中结出最丰硕的果实。我们于是寄给你十万塞斯特斯，这是在此地由我们凭天主仁慈主持的教会中，通过与我们同在的圣职人员和教友的捐献收集的，你可尽你所能谨慎地那里分发。\par

4. 我们真的希望此类事不再发生，我们的弟兄因主的威严保护而安全免于这类危险。然而，如果为了搜求我们心内的爱，并为考验我们心内的信德而发生类似的事，请不要延迟在你们的来信中告诉我们，要确信我们的教会和这里的全体兄弟团体恳求祈祷这些事不再发生；但如果发生，他们将情愿并慷慨地提供帮助。为了使你们在祈祷中纪念那些为这必要工作如此迅速慷慨地劳动的弟兄姊妹，使他们常劳动；并为了回报他们的善工，你们在他们的祭献和祈祷中纪念他们，我将每个人的名字附在下面；此外，我还加上了我的同工和同事司铎们的名字，他们自己当时也按其能力以他们自己和其教友的名义捐献了一些。除了我们自己的金额，我也通知并寄送了他们的少量捐款，所有这些人，你们理应根据信德和爱德的要求在他们的恳求和祈祷中纪念他们。最亲爱的弟兄们，祝你们永远衷心安康，并纪念我们。\par

\SourceRecord{Translated by Robert Ernest Wallis.；Ante-Nicene Fathers；Vol. 5.；Edited by Alexander Roberts, James Donaldson, and A. Cleveland Coxe.；Buffalo, NY: Christian Literature Publishing Co.,；1886.；<http://www.newadvent.org/fathers/050659.htm>.}

% structure: p0001 source=h1 target=section reason=page-title

\section{第六十封书信}

\Emph{论点——他禁止一名演员，若他继续从事其可耻的职业，在教会内共融。他也不允许以下情况成为他的借口：他本人并未实践戏剧艺术，只要他将其教授给他人；他也不因缺乏生计而原谅他，因为必需品可从教会的资源中供应给他；因此，若当地教会的资源不足，他建议他来\Place{迦太基}。}\par

1. \Person{西彼廉}致其弟兄\Person{欧克拉提乌斯}，问候。因我们彼此的爱德和你对我的敬意，你以为我应当被咨询，最亲爱的弟兄，关于一名演员的意见。他住在你们中间，却仍坚持他那同样可耻的技艺；并且作为师傅和教师，不是为教导，而是为败坏少年们，他将自己不幸学到的也传授给别人：你问这样的人是否应当与我们共融。我认为，这既不适合天主的威严，也不适合福音的纪律，即教会的端庄与信誉被如此可耻和臭名昭著的污染所玷污。因为，在法律中，男人被禁止穿女人的衣服，那些如此犯罪的人被视为可诅咒的；那么，不仅穿女人衣服，而且通过不道德的技艺的教导，表现出下流、柔弱和淫荡的姿态，罪过何等更大。\par

2. 也不要让任何人以他已放弃剧院为借口，而他仍在将这门技艺教给他人。因为一个人若用别人代替自己，并且不单单自己一人，而是提供多人作为替代，他不能显得已经放弃了；他违背天主的安排，教导和传授如何使一个男人被瓦解成女人，如何用技艺改变性别，以及那亵渎天主肖像的魔鬼如何因败坏和软弱的身体的罪恶而欢喜。但如果这样的人以贫穷和微薄生计的必要性为理由，他的需要也可以得到教会的供养，与那些靠教会支持维生的人一样；不过他必须满足于非常节俭但无罪的食物。他不要以为通过一笔津贴就能赎回自己，使他停止犯罪，因为这益处不是为我们，而是为他本人。他若想要更多，就必须从那使人离开\Person{亚巴郎}、\Person{依撒格}和\Person{雅各伯}的宴席，并可悲而有害地在此世养肥，最终引向永远饥渴的折磨的收益中寻求；因此，尽你所能，将他从这种堕落和耻辱中召回，走向无罪的道路，并走向永生的希望，使他满足于教会的供给，虽然简朴，却有益健康。但如果你们那里的教会不足以供养需要帮助的人，他可以转移到我们这里，在此获得他所需要的食物和衣服，而不是在教会之外向别人教授致命的东西，而是自己在教会内学习有益健康的东西。我祝你，最亲爱的弟兄，永远衷心地平安。\par

\SourceRecord{Translated by Robert Ernest Wallis.；Ante-Nicene Fathers；Vol. 5.；Edited by Alexander Roberts, James Donaldson, and A. Cleveland Coxe.；Buffalo, NY: Christian Literature Publishing Co.,；1886.；<http://www.newadvent.org/fathers/050660.htm>.}

% structure: p0001 source=h1 target=section reason=page-title

\section{第六十一封书信}

\Emph{致庞波尼乌斯，论某些童贞女。}\par

\Emph{论证。——西彼廉与其数位同僚答复其同僚庞波尼乌斯：那些曾决心以节德和坚毅保守其身份，但后来被发现与男人同床的童贞女，若仍被证实为童贞，则应接纳其共融并准其进入教会。否则，由于她们对基督行了奸淫，应强迫她们做完整的补赎；那些顽固坚持者，则应被逐出教会。}\par

1. 西彼廉、凯基利乌斯、维克多、塞达图斯、特尔图鲁斯，以及与他们同在的司铎们，致其弟兄庞波尼乌斯，祝安。最亲爱的弟兄，我们已阅读了你托我们弟兄帕科尼乌斯送来的信函，你请求并希望我们回复你，告知我们对那些曾决定保持其身份并坚定持守贞洁，但后来被发现与男人同床共枕的童贞女有何看法；你提到其中一人是执事；然而这些与男人同寝过的童贞女却声称自己是贞洁的。关于这些事，既然你渴望我们的意见，要知道我们并未偏离福音与宗徒的传统，而是坚定恒久地为我们的弟兄姊妹筹划，并以一切有益与安全的方式维护教会的纪律，因为主说：\BibleQuote{我要赐给你们合我心的牧者，他们要以纪律牧养你们。}\BibleRef{耶 3:15}经上又记载：\BibleQuote{轻视纪律的人是有祸的。}\BibleRef{智 3:11}在圣咏中，圣神也劝诫并教训我们说：\Quote{守住纪律，免得上主发怒，你们从正路上灭亡，当祂的怒气迅速向你们燃烧。}\par

2. 因此，最亲爱的弟兄，首先，无论是监督还是民众，最应渴求的莫过于我们敬畏天主的人以对纪律的全然遵守来持守天主的诫命，不容许我们的弟兄迷失，凭自己的幻想和私欲生活；而应忠信地为每个人的生命着想，不让童贞女与男人同住——我且不说同寝，就是同住也不可——因为她们软弱的性别和仍然危险的年龄，应当由我们在一切事上加以约束和管束，免得给那诱骗我们、企图象猛兽般侵害我们的魔鬼留下机会，使她们受到伤害，因为宗徒也说：\BibleQuote{不要给魔鬼留余地。}\BibleRef{弗 4:27}船必须警觉地脱离险境，以免在岩石和峭壁间撞碎；行囊必须迅速从火中取出，免得被蔓延的火焰烧毁。凡接近危险的人不会长久安全，天主的仆人也无法逃脱魔鬼，若他已陷入魔鬼的网罗。我们必须立即干预这样的人，使他们尚能在清白中分离；因为过不了多久，当他们因极重的罪疚结合后，便无法再藉我们的干预而分离。此外，我们从这里看到了多少严重的祸害，目睹了多少童贞女因这种非法而危险的结合而堕落，令我们心中极为悲痛！若她们已忠信地将自己奉献给基督，就让她们在谦逊和贞洁中持守，不招任何恶名，以勇气和坚定等待贞洁的赏报。若她们不愿或不能坚持，那么结婚比因罪行而陷入火中更好。当然，她们不可使弟兄姊妹跌倒，因为经上记载：\BibleQuote{如果食物使我的弟兄跌倒，我就永远不吃肉，免得使我的弟兄跌倒。}\BibleRef{格前 8:13}\par

3. 任何人不可认为能用这样的借口为自己辩护：她可以被检查并证实是否为童贞；因为产婆的手和眼也常常被欺骗；即使她在女人可为童贞的那方面被证明是童贞，她仍可能在身体的其他部分犯了罪，那部分可能被玷污而无法被检查。的确，仅仅是同寝、仅仅是拥抱、仅仅是交谈、亲吻，以及两人同卧的羞耻与污秽的睡眠，这承认了多少不光彩与罪！若一个丈夫撞见他的妻子与别的男人同寝，他岂不怒不可遏，在愤怒的冲动下取剑在手吗？那么，我们的主基督和审判者，当祂看见那奉献给祂、归于祂圣洁的童贞女与别人同寝时，祂会怎么想？祂是何等愤怒与生气，对这样不洁的结合威胁何等惩罚！为了使每一位弟兄都能逃脱祂的属灵之剑和将来的审判之日，我们应当以一切筹划去预备和努力。既然所有人都应以各种方式持守纪律，那么监督和执事更应当为此谨慎，好能在他们的言行和品格上为他人提供榜样和教导。因为如果他们自身的腐化和罪的教训开始出现，他们怎能引导他人的正直和节制呢？\par

4. 因此，最亲爱的弟兄，你们明智而果断地行事，将那常与一位贞女同住的执事，以及那些惯与贞女同寝的其他人，予以绝罚。然而，若他们已悔改自己这种非法的同寝，并彼此分离，则暂时让贞女们由产婆仔细检查；若她们被发现仍是贞女，便应让他们重领共融，并被接纳进入教会；但须加此警告：若他们日后重归于同一男子，或与同一男子同住一屋或同一屋顶下，便当以更严厉的处分予以驱逐，且此后不易再被接纳进入教会。若其中一人被发现已受玷污，则让她彻底悔改，因为犯此罪者，不只是对丈夫，而是对基督行了奸淫；因此，指定一段适当的时间后，在她作了告解后，再回归教会。但若他们执迷不悟，不肯彼此分离，就让他们知道，以这种无耻的固执，他们将永不能被我们接纳进入教会，以免他们开始为他人树立榜样，因自己的罪行走向毁灭。也不可让他们以为，若拒绝服从主教和司铎，生命或救恩之路仍为他们敞开，因为申命纪中上主天主说：\BibleQuote{若有人擅敢行事，不听从那日在那里的司祭或判官，那人就必治死；众百姓都要听见害怕，不再擅敢行事。}\BibleRef{申 17:12} 天主曾命令那不听从祂司祭的人，以及那不听从当时所设立之判官的人，当被处死。那时，肉身的割礼尚在施行，他们确实被刀剑所杀；但如今，在天主忠信的仆人当中，割礼已始于心灵，骄傲和悖逆之人便被圣神的剑所杀，即被逐出教会。因为他们不能活在教会之外，因为天主的家只有一个，任何人除非在教会内，否则不能得救。然而，圣经证实，不受管教的人必灭亡，因为他们不听从不服从有益的训诫；因为经上说：\Quote{不受管教的人不爱那责备他的人；但憎恶责斥的人必被羞辱所消灭。}\par

5. 因此，最亲爱的弟兄，要努力不让那些不受管教的人被消灭而丧亡，要尽你所能，以你有益的劝导治理弟兄团体，并为着每个人的得救而为他们出主意。进入生命的道路是狭窄而艰难的，但当我们进入光荣时，赏报却是卓越而伟大的。那些曾为天国自阉的人，当在一切事上悦乐天主，并不可因他们邪恶的丑闻而得罪天主的司祭或主的教会。若眼下我们的一些弟兄似乎因我们而悲伤，我们仍当坚守我们有益的信服，知道宗徒也曾说过：\BibleQuote{那么，我因对你们说实话，就成了你们的仇敌吗？}\BibleRef{迦 4:16} 但若他们听从我们，我们便赢得了我们的弟兄，并以我们的劝导使他们既得救恩，又获尊严。但若有些悖逆的人拒绝听从，就让我们追随同一宗徒，他说：\BibleQuote{若我还求人的欢心，我就不是基督的仆役。}\BibleRef{迦 1:10} 若我们不能使一些人欢心，好让他们悦乐基督，那么就让我们确实尽我们所能，遵守祂的诫命，以悦乐我们的主天主基督。我祝你，我所爱的、极其渴望的弟兄，在主内全心平安。\par

\SourceRecord{Translated by Robert Ernest Wallis.；Ante-Nicene Fathers；Vol. 5.；Edited by Alexander Roberts, James Donaldson, and A. Cleveland Coxe.；Buffalo, NY: Christian Literature Publishing Co.,；1886.；<http://www.newadvent.org/fathers/050661.htm>.}

% structure: p0001 source=h1 target=section reason=page-title

\section{\WorkTitle{第六十二封书信}}

\Emph{\Person{凯齐略}，论主之杯的圣事。}\par

\Emph{辩论。——\Person{西彼廉}教导，反对那些在主的晚餐中使用水的人，指出不可只献水，而应献搀水之酒；在圣经中，水象征圣洗，而非圣体。借旧约预表，阐明在主的身体圣事中用酒；并宣告水之象征代表基督徒集会。}\par

1. \Person{西彼廉}致其兄弟\Person{凯齐略}，安好。虽知最亲爱的弟兄，许多由天主俯就派定管理普世各处主之教会的主教，皆持守福音真理及主之传统，不因人而新设的规矩偏离我主基督所命所行；但亦有人或因无知、或因简朴，在祝圣主之杯并分施于民时，不按耶稣基督，我们的主和天主——此祭献之创立者与导师——所行所教而行，故我认为，写此信于你，既是虔诚亦属必要，好使若有人仍陷此错，得见真理之光，归回主之传统的根源与起源。最亲爱的弟兄，你切勿以为我在写自己的或人的思想，或凭己意妄自担当，因我常以谦逊温和持守我的平庸。但当有任何事由天主的默感和命令所规定时，忠信的仆人须服从主，如此便可免于妄自尊大之责，因他若不行所命，则畏惧得罪主。\par

那么要知道，我已受提醒，在献杯时，必须遵守主的传统，且我们不可行任何主未先为我们所行之事，即那为纪念祂而献的杯，应搀酒而献。因为当基督说：\BibleQuote{我即是真葡萄树}\BibleRef{若 15:1}，基督的血确实不是水，而是酒；我们藉之得救赎与复活的祂的血，不可能显在杯中，若杯中无酒显明基督的血，而此乃全部圣经的圣事与见证所宣告的。\par

因为在创世记中，我们也发现，就\Person{诺厄}的圣事而言，同样的事为他们预兆并预表了主的苦难；他喝酒；他醉了；他在家中赤身；他躺卧，大腿赤裸暴露；他的次子见了父亲的裸体，并宣扬出去，但长子和幼子将其遮盖；以及其他无须详述之事，因我们只须把握这一点就够了：\Person{诺厄}展露未来真理的预表时，喝的并非水，而是酒，如此便表达了主受难的预像。\par

此外，在司祭\Person{默基瑟德}身上，我们看到主祭献的圣事被预表，神圣圣经如此作证说：\BibleQuote{默基瑟德，撒冷王，带了饼和酒来。}\BibleRef{创 14:18}他原是至高天主的司祭，祝福了\Person{亚巴郎}。\Person{默基瑟德}预表基督，圣神在圣咏中从圣父对圣子说：\BibleQuote{我在晨星之前生了祢；祢按照默基瑟德的品位，永为司祭。}这品位确实来自那祭献而从此延续：\Person{默基瑟德}是至高天主的司祭；他献了酒和饼；他祝福了\Person{亚巴郎}。因为谁比我们的主耶稣基督更是至高天主的司祭呢？祂向天主圣父献了祭，并献上同样之物，即\Person{默基瑟德}所献的——饼和酒，亦即祂的体和血。至于\Person{亚巴郎}，那先前的祝福属于我们的人民。因为若\Person{亚巴郎}信了天主，且这算为他的义，那么凡信天主并生活在信德中的人，便被算为义，且已在有信德的\Person{亚巴郎}内蒙受祝福，并被显示为成义；正如蒙福的宗徒\Person{保禄}所证明的，他说：\BibleQuote{亚巴郎信了天主，这算为他的义。所以你们要知道，那以信德为本的人，就是亚巴郎的子孙。圣经预见天主将由信德使外邦人成义，便预先向亚巴郎传报：万民都要因他获得祝福。所以那以信德为本的人，与有信德的亚巴郎同受祝福。}\BibleRef{迦 3:6}因此在福音中我们读到：\BibleQuote{从石头中给亚巴郎兴起子孙，}\BibleRef{玛 3:9}即从外邦人中聚集。当主称赞\Person{匝凯}时，回答说：\BibleQuote{今天救恩临到了这一家，因为他也是亚巴郎之子。}\BibleRef{路 19:9}因此，在创世记中，为使司祭\Person{默基瑟德}对\Person{亚巴郎}的祝福得以适当举行，基督祭献的预像先行，即预定于饼和酒中；主完成并满全此事，献了饼及搀酒的杯，如此，祂既是真理的满全，便满全了预先描绘之像的真理。\par

5. 此外，圣神藉撒罗满预先展示了主祭献的预像，提到了被宰杀的牺牲、饼和酒，还有祭台和宗徒，并说：\BibleQuote{智慧建造了自己的房屋，立了七根柱子；她宰杀了牲畜，将酒调和在爵中；又摆设了筵席；并派遣了她的仆人，用高声宣告召聚人到她的杯前来，说：\InnerQuote{谁是愚蒙的，让他转向我；对缺乏明智的人，她说：来，吃我的饼，喝我为你们调和了的酒。}}祂宣示了调和了的酒，也就是说，用先知的声音预言了主用酒和水调和的杯，这要表明在我们主的受难中所成就的，正是先前所预言的。\par

6. 在犹大的祝福中，也预示了同一件事，那里也表达出基督的预像：祂应得自祂兄弟的赞美和朝拜；祂要用那背负了十字架、战胜了死亡的双手，压服那屈服而逃跑的仇敌的背脊；并且祂自己是犹大支派的狮子，应在受难中躺卧安眠，然后复活，并亲自成为外邦人的希望。与此相合，圣经又加上说：\BibleQuote{祂将在酒中洗净祂的衣服，在葡萄的血中洗净祂的袍衣。}\BibleRef{创 49:11} 但当提到葡萄的血时，除了主的血之杯中的酒之外，还预示了什么呢？\par

7. 依撒意亚也为同一件事，论到主的受难，圣神作证说：\BibleQuote{为什么你的衣服是红的，你的外衣好像从榨酒池踩踏过的一样？}\BibleRef{依 63:2} 水能使衣服变红吗？或者，在酒榨中被脚踩踏、被压榨出来的是水吗？所以，确实，这里提到酒，是为了让主的血被理解，并且后来在主的杯中显明的事，由宣告它的先知们预言了。并且反复强调酒榨的踩踏和压榨；因为正如要喝酒，就必须先踩踏和压榨葡萄串，同样，我们也不能喝基督的血，除非基督先被践踏和压榨，并且祂自己先饮了那杯，就是祂也应给信徒们喝的杯。\par

8. 但是，每当圣经中单独提到水时，指的是圣洗，正如我们在依撒意亚中所见的暗示：\BibleQuote{不要记念，他说，先前的事，也不要思念古时的事。看哪，我要做一件新事，如今就要发生，你们要知道。我必在旷野开道路，在干旱之地开江河，给我选民喝，就是我为自己所赎的子民，好叫他们宣扬我的赞美。}\BibleRef{依 43:18} 在那里，天主藉先知预言，在列邦中，在先前干旱的地方，以后要有丰沛的河流，并为天主的选民——也就是藉圣洗的再生而成为天主子女的人——提供水。另外，又再次预言和预告，犹太人若渴了，寻求基督，就要与我们同饮，也就是说，要获得圣洗的恩宠。他说：\BibleQuote{如果他们口渴，他说，祂必引导他们经过旷野，从磐石中给他们出水；磐石裂开，水就涌出，我的百姓要喝。}\BibleRef{依 48:21} 这在福音中应验了，当基督——祂是磐石——在受难中被长矛刺穿时；祂也提醒先知先前所宣告的，呼喊说：\BibleQuote{若有人渴了，让他到我这里来喝。信我的人，就如经上所说，从他腹中要流出活水的江河。}并且为使更清楚，主在那里说的不是杯，而是圣洗，圣经加上说：\BibleQuote{但祂这话是指着信祂的人将要领受的圣神说的。}因为藉圣洗领受圣神；因此，那些受洗并得到了圣神的人，就得到了主的杯的饮用。并且不要让任何人感到困扰，因为当圣经论到圣洗时，说我们渴了并饮用，因为主在福音中也说：\BibleQuote{饥渴慕义的人是有福的，}\BibleRef{玛 5:6} 因为怀着贪婪和渴望的心所接受的，就会被更充分、更丰富地饮用。同样，在另一处，主对撒玛黎雅妇人说：\BibleQuote{凡喝这水的，还要再渴；但若喝了我所赐的水，就永远不渴。}\BibleRef{若 4:13} 这也预表了拯救之水的圣洗，这圣洗确实只领受一次，不再重复。但主的杯在教会中总是被渴求并饮用。\par

9. 最亲爱的弟兄，无需许多论证即可证明：圣洗总以\Quote{水}之名来表明，我们理当这样理解它。因为主来临之时，藉着命令将那忠信的水、永生之水在圣洗中赐给信徒，并以自己权柄的榜样教导，杯应掺和酒与水而成为合一的，从而彰显了圣洗与杯的真理。因为主在受难前夕拿起杯，祝福了，递给祂的门徒说：\BibleQuote{你们大家都要喝这个，因为这是我的血，新约的血，为多人流出来，以赦免罪过。我告诉你们：从今以后，我不再喝这葡萄树的果实，直到我在我父的国里与你们同喝新酒的那一天。}\BibleRef{玛 26:28}在这段经文中我们发现，主所献的杯是调和的，而祂称为自己血的正是酒。由此可见，若杯中无酒，基督的血便不被献上；除非我们的祭献与牺牲回应祂的苦难，主的祭献便不能以合法的祝圣来庆祝。然而，若我们在天主父及基督的祭献中不献酒，也不按主自己的传统调和主的杯，我们怎能与基督在天父的国里喝葡萄树果实的新酒呢？\par

10. 此外，蒙福的宗徒保禄，蒙主拣选并派遣，被立为福音真理的宣讲者，在他的书信中阐述了同样的事，说：\BibleQuote{主耶稣在被出卖的那一夜，拿起饼，祝谢了，擘开说：\InnerQuote{这是我的身体，为你们而舍的，你们要这样做，以纪念我。}晚餐后，祂也照样拿起杯来说：\InnerQuote{这杯是用我的血所立的新约；你们每次喝它，要这样做，以纪念我。}因为你们每次吃这饼，喝这杯，就是宣告主的死，直到祂来临。}\BibleRef{格前 11:23}但是，如果主既如此命令，祂的宗徒又确认并传授了同样的事——即我们每次喝的时候，要照主所做的去行以纪念主——那么我们就会发现，我们所遵守的并非所命令的，除非我们也做主所做过的；并且我们如此调和主的杯，并未偏离神圣的教导；但我们绝不可偏离福音的规诫，门徒们也应当遵守并践行师傅所教导和所做的一切。蒙福的宗徒在另一处更恳切有力地教导说：\BibleQuote{我很惊奇，你们这么快就离开了那藉恩宠召叫你们的那一位，而去归从别的福音；其实并不是别的福音，只是有些扰乱你们的人，想要改变基督的福音。但是，即使是我们，或是从天上来的天使，若给你们宣讲与我们先前所传给你们的不同福音，当受诅咒。我们先前说过，如今我再重复：若有人给你们宣讲与你们所领受的不同福音，当受诅咒。}\BibleRef{迦 1:6}\par

11. 既然连宗徒自己或从天上来的天使都不能宣讲或教导与基督一次所教导及祂的宗徒所宣布的不同，我非常惊异，这做法从何而起：竟然违背福音与纪律，在某些地方主的杯中只献水；单凭水无法表达基督的血。圣神在\WorkTitle{圣咏}中论及这事的圣事时也并非沉默，祂提及主的杯说：\BibleQuote{令人陶醉的杯，何其美好！}那令人陶醉的杯必定是搀了酒的，因为水不能使人醉。主的杯如此使人沉醉，正如\WorkTitle{创世纪}中诺厄喝酒而醉。但因主的杯与血的沉醉不同于世俗之酒的沉醉，圣神在\WorkTitle{圣咏}中说：\BibleQuote{令人陶醉的杯}，又加上\BibleQuote{何其美好}，无疑因为主的杯如此使饮者沉醉，以致他们清醒；使他们的心灵回复属灵的智慧；使每个人从那世俗的风味转向认识天主；并且，正如那普通的酒使心神涣散、灵魂松弛、一切忧伤被放下，同样，当主的血与救恩之杯被饮下时，旧人的记忆被放下，对先前属世生活的遗忘生起，那被折磨人的罪恶所压迫的忧伤愁苦的胸襟，因神圣慈悲的喜乐而舒缓；因为唯有那在教会中饮者能得喜乐，当它被饮时，它保留着主的真理。\par

12. 然而，何其悖逆与相反：尽管主在婚宴上将水变成了酒，我们却将酒变成水，而连那事的圣事也本应劝诫并教导我们在主的祭献中献上酒。因为犹太人缺乏属灵的恩宠，酒也缺乏。因为万军之主的葡萄园原是以色列家；但基督教导并显明外邦人要接替他们，我们因信德的功劳将随后获得犹太人丧失的位置，祂将水变成了酒；也就是说，祂表明在基督与教会的婚宴上，犹太人失败之后，万民应更涌流聚集：因为神圣圣经在\WorkTitle{默示录}中宣告水象征人民，说：\BibleQuote{你所看见的水，那淫妇坐在上面的，乃是诸民族、群众、外邦人的百姓和异语的人。}\BibleRef{默 17:15}我们显然看到这一点也包含在杯的圣事中。\par

13. 因为基督背负了我们众人，也背负了我们的罪，我们便看到，在水中预表的是人民，而在酒中显示的是基督的血。但当水在杯里与酒相混合时，人民便与基督成为一体，信众的集会便与它所信奉的那位联合并连结在一起；水与酒的这种联合与连结在主杯里如此融合，以致这混合物再也无法分离。因此，再没有什么能使教会——即建立在教会中、忠信且坚定地持守他们所相信的人民——与基督分离，以致阻止他们那不可分割的爱常存并紧密相系。因此，在祝圣主的杯时，不能只献水，也不能只献酒。因为若有人只献酒，基督的血就与我们分离；但若只有水，人民就与基督分离；但当两者混合，并藉紧密的结合彼此连结，便完成了属灵属天的圣事。因此，主的杯确实不是水单独，也不是酒单独，除非两者互相混合；同样地，主的身体也不能只是面粉或只是水，除非两者结合、连结并凝聚成一个饼团；在这同一圣事中，显示我们的人民成为一体，正如许多谷粒被收集、磨碎并混合在一起成为一个团，做成一饼；同样在基督——这天上的食粮——内，我们可以知道有一个身体，我们的数目与它相连并联合。\par

14. 因此，最亲爱的弟兄，没有理由认为应当追随某些人的习惯，他们曾以为在过去只应在主的杯里献水。因为我们必须查问他们自己追随了谁。因为在基督所献的祭中，除了基督之外无人可追随，我们理当服从并实行基督所行和所命令的，因为祂自己在福音中说：\Quote{你们若遵行我所吩咐你们的，从此我不再称你们为仆人，而是朋友。}\BibleRef{若 15:14} 并且只有基督应当被听从，圣父也从天上作证，说：\Quote{这是我的爱子，我所喜悦的；你们要听从他！}\BibleRef{玛 17:5} 因此，如果只有基督必须被听从，我们就不应留意我们之前的别人可能认为应做的事，而应留意那先于万有的基督首先所行的事。效法人的做法是不相宜的，而应效法天主的真理；因为天主藉依撒依亚先知说话，说：\Quote{他们拜我是徒然的，所教导的是人的诫命和道理。}\BibleRef{依 29:13} 主在福音中又重申了同样的话，说：\Quote{你们废弃天主的诫命，为要保守你们的传统。}\BibleRef{谷 7:13} 此外，在另一处祂肯定了这一点，说：\Quote{凡废掉这些诫命中最小的一条，并这样教导人，他在天国里要称为最小的。}\BibleRef{玛 5:19} 但是，如果我们连主最小的诫命也不可废，那么，犯下这些如此重要、如此伟大、如此关乎我们主的苦难和我们救恩这圣事本身的诫命，或者藉人的传统将其改变为任何不同于神圣所命定的事物，那就更是被禁止了！因为如果耶稣基督、我们的主和天主，祂自己是天主圣父的大司祭，并且首先将自己作为祭献献给圣父，又命令我们如此行以纪念祂，那么，那位效法基督所行的司祭，就是真正履行基督的职务；当他按照他所看见的基督自己所献的而行时，他就在教会中向天主圣父献上真实且圆满的祭献。\par

15. 但是，若不忠信地持守属灵所规定的，一切宗教和真理的纪律就被颠覆了；除非的确有人在早晨的祭献中害怕，以免因尝到酒而发出基督的血的香气。因此，弟兄团体竟开始因在献祭中学习对祂的血和祂的流血感到不安，而在迫害中甚至被阻止参与基督的苦难。然而，主在福音中说：\Quote{凡以我为耻的，人子也必以他为耻。} 宗徒也说话，说：\Quote{若我讨人的喜欢，我就不是基督的仆人了。}\BibleRef{迦 1:10} 但我们这些以喝基督的血为羞耻的人，怎能为我们流血呢？\par

16. 或许有人以这种想法来自我安慰：虽然在早晨只献水，但当我们来到晚餐时，我们献上混合的杯？但是，当我们晚餐时，我们无法召集人民到我们的筵席，以便在全体弟兄面前庆祝圣事的真理。然而，主不是在早晨，而是在晚餐后献上混合的杯。那么，我们是否应当在晚餐后庆祝主的杯，以便通过不断重复主的晚餐来献上混合的杯呢？基督应当在傍晚时分献祭，以便这祭献的时刻显示世界的日落和黄昏；正如在出谷纪中记载：\Quote{以色列全体会众应在黄昏时分宰杀它。}\BibleRef{出 12:6} 又在圣咏中：\Quote{愿我双手的举献成为晚上的祭献。} 但我们是在早晨庆祝主的复活。\par

17. 因为我们一切祭献中都提及祂的受难（因为主的受难就是我们献上的祭献），我们理当只做祂所做的事。因为圣经说：\BibleQuote{你们每次吃这饼，喝这杯，就是宣告主的死，直到祂再来。}\BibleRef{格前 11:26}所以我们每次纪念主和祂的受难而献杯时，要行主所行之事。最亲爱的弟兄，要达成此结论：如果我们先辈中有人因无知或单纯未遵守主以榜样和教导训诲我们的事，他可能因主的仁慈，因其单纯而得宽恕。但我们这些如今受主告诫和训诲，要按照主所献的方式，调酒献上主的杯，并就此给我们的同僚写信，使福音的律法和主的传统各处保持，不偏离基督所教导和所行的事，我们就不能得宽恕。\par

18. 若再忽略这些事，并坚持以前的错误，这岂非落在主的谴责之下？祂在第五十篇圣咏中谴责说：\BibleQuote{你怎么胆敢传述我的诫命，你的口怎敢提及我的盟约？因为你憎恨管教，把我的话抛在脑后。你见了盗贼，就与他同伙，又与奸夫有分。}因为传述主的正义和盟约，却不做主所做的事，岂非抛弃祂的话，藐视主的训诲，犯下不是属世，而是属灵的盗窃和奸淫吗？凡是从福音真理中窃取我们主的话和行为，就是在败坏和亵渎神圣的诫命，正如耶肋米亚所记载的。祂说：\BibleQuote{糠秕与麦子有何相干？因此，上主说：看，我要攻击那些窃取我的言语、各向邻舍偷窃的假先知，他们以谎言和轻浮使我的百姓走迷。}在同一先知书的另一处，祂说：\BibleQuote{她与木头和石头行了奸淫，但尽管如此，她仍不归向我。}为使这种盗窃和奸淫也不临到我们身上，我们应焦虑谨慎，战战兢兢地以虔敬之心警醒。因为如果我们是天主和基督的司祭，我不知道除了天主和基督，我们还应跟随谁，因为祂亲自在福音中强调说：\BibleQuote{我是世界的光；跟随我的，决不在黑暗中行走，必有生命的光。}\BibleRef{若 8:12}因此，为使我们不走在黑暗里，我们应跟随基督，遵守祂的诫命，因为祂在另一处派遣宗徒时亲自告诉他们：\BibleQuote{天上地下的一切权柄都给了我。所以你们要去使万民成为门徒，因父、及子、及圣神之名给他们授洗，教训他们遵守我所吩咐你们的一切。}\BibleRef{玛 28:18}因此，如果我们愿意行走在基督的光中，就不要偏离祂的诫命和训诫，感谢祂，一方面教导我们未来应行之事，另一方面宽恕我们过往因单纯所犯的错误。并且，因为祂的第二次来临已临近我们，祂的仁慈和宽厚屈尊更以真理之光光照我们的心灵。\par

19. 因此，最亲爱的弟兄，我们的宗教信仰、敬畏、这地方本身以及我们司祭的职务，都要求我们在调和并献上主的杯时，遵循主传统的真理，并因主的警告，纠正某些人看来错误之处；这样，当祂开始以光明和天上的威能来临时，祂会看到我们遵守了祂所告诫我们的，履行了祂所教导的，实行了祂所做的事。最亲爱的弟兄，祝你永远由衷安康。\par

\SourceRecord{Translated by Robert Ernest Wallis.；Ante-Nicene Fathers；Vol. 5.；Edited by Alexander Roberts, James Donaldson, and A. Cleveland Coxe.；Buffalo, NY: Christian Literature Publishing Co.,；1886.；<http://www.newadvent.org/fathers/050662.htm>.}

% structure: p0001 source=h1 target=section reason=page-title

\section{第六十三封书信}

\Emph{致\Person{埃皮克泰图斯}与\Place{阿苏拉}教会，关于\Person{福图纳提安努斯}，他们前任主教。}\par

\Emph{论题：他警告\Person{埃皮克泰图斯}与\Place{阿苏拉}教会，不得允许\Person{福图纳提安努斯}——一个失足者，但曾是他们的前任主教——重归其主教职位，原因有多方面，其中之一是已规定失足主教不得恢复其原有品秩。}\par

1. 西彼廉致其兄弟埃皮克泰图斯，及驻阿苏拉的子民，问安。我最亲爱的弟兄们，我闻悉福图纳提安努斯，你们以前的主教，在可悲的失足跌落后，如今竟想装作健康，并开始为自己要求主教职位，这令我深感忧伤和痛苦。这件事使我难过：首先是为了他本人，这个可怜的人，要么完全被魔鬼的黑暗所蒙蔽，要么被某些人的亵渎劝说所欺骗；当他本应做补赎，日夜以眼泪、哀求、祈祷恳求主时，竟仍敢为自己要求被他背叛的司铎职，仿佛从魔鬼的祭台接近天主的祭台是正当的。或者仿佛他不会在审判之日为自己招致主更大的愤怒和义怒；他不能在信德和德行上作弟兄们的向导，反倒成为背信、狂妄、鲁莽的教师；他没有教导弟兄们在战斗中勇敢站立，反而教导那些被击败、匍匐在地的人连\Emph{求宽恕}也不；尽管主说：\BibleQuote{你们向他们奠了酒，向他们献了素祭。我岂不为这些事发怒？主说。}\BibleRef{依 57:6}又在另一处说：\BibleQuote{向任何神献祭，惟独不向主献祭的，应被消灭。}\BibleRef{出 22:20}再者，主又发言说：\BibleQuote{他们敬拜了自己手指所造的：卑微人屈膝，尊贵人自谦，我必不饶恕他们。}\BibleRef{依 2:8}在默示录中，我们也读到主的愤怒的威胁，说：\BibleQuote{若有人朝拜那兽和它的像，并在额上或手上接受了它的印记，他必要喝天主愤怒的酒，即斟在他义怒杯中的酒；他要在圣天使面前和羔羊面前受火与硫磺的折磨；他们受折磨的烟上升，直到永远；那些朝拜兽和兽像的，昼夜不得安息。}\BibleRef{默 14:9}\par

2. 因此，既然主在审判之日以这些折磨、这些惩罚威胁那些服从魔鬼并向偶像献祭的人，那么一个曾经服从并服事魔鬼的司铎的人，怎能认为他可以充当天主的司铎？或者，他那已被亵渎和罪恶所掳的手，怎能用来奉献天主的祭献和主所悦纳的祈祷？因为在圣经中，天主禁止司铎即使犯了较轻的过失也近前献祭，并在肋未记中说：\BibleQuote{凡有任何残疾或污秽的人，不得近前向天主奉献礼品。}\BibleRef{肋 21:17}又在出谷纪中说：\BibleQuote{凡近前到上主天主面前的司铎，应圣化自己，免得上主舍弃他们。}\BibleRef{出 19:22}又说：\BibleQuote{当他们近前服事圣者的祭台时，不可带罪，免得死亡。}\BibleRef{出 28:43}因此，那些将重罪加于自身的人，即那些通过向偶像献祭而做了亵渎祭献的人，不能为自己要求天主的司铎职，也不能在天主面前为弟兄们作任何祈祷；因为福音中记载：\BibleQuote{天主不听罪人；但若有人是敬拜天主的，并遵行他的旨意，天主就听他。}\BibleRef{若 9:31}然而，堕落黑暗的深重阴霾已蒙蔽了一些人的心，以至于他们从有益的诫命中得不到任何光明，而是背离了真道的直路，被自己罪恶的黑暗和错误所裹挟，猝然坠入深渊。\par

3. 如果那些否认主的人现在拒绝我们的劝告或主的诫命，这并不奇怪。他们贪求礼物、奉献和收益，从前他们曾为此贪得无厌地守望。他们还渴望宴席和筵席，他们最近消化不良地呕吐出这些宴席的狂欢，极其明显地证明了他们以前并非服务于宗教，而是服务于自己的肚腹和收益，怀着亵渎的贪欲。由此我们也看出并相信，这斥责来自天主的探查，好使他们不再继续站立在祭台前；且更进一步，作为不洁之人，不得与贞洁沾边；作为背信之人，不得与信德相干；作为亵渎之人，不得与宗教有关；作为属世之人，不得接触神圣之事；作为渎圣之人，不得接触圣物。为阻止这样的人再次亵渎祭台、玷污弟兄们，我们必须用尽全力，竭尽所能，在我们能力范围内阻止他们这种邪恶的妄为，使他们不再试图以司铎的身份行事；这些人已被投入最低的死亡深渊，其毁灭之重远远超过平信徒的失足。\par

4. 但是，在这些疯狂者中间，如果他们不可救药的疯狂持续下去，并且随着圣神的撤回，那已开始的盲目仍留在其深沉的黑暗中，我们的忠告将是：将个别弟兄从他们的欺骗中分离出来，免得有人陷入他们错误的罗网，要把他们从他们的传染中隔离开来。因为圣神不在的地方，祭献不能祝圣；并且，主也不能藉着那自己曾侮辱主的人的祈祷和恳求，帮助任何人。但如果\Person{福尔图尼亚努斯}，或因魔鬼导致的盲目而忘记了自己的罪行，或成为魔鬼欺骗弟兄的仆人和奴仆，仍坚持他的疯狂，那么，你们要尽你们所能，在这魔鬼狂怒的黑暗中，努力将弟兄们的心灵从错误中召回，使他们不要轻易同意另一个人的疯狂，不要使自己成为堕落之人罪恶的同伴。相反，让他们保持健全，持守他们救恩的恒常轨迹，以及他们所保持并守护的完整。\par

5. 然而，那些承认自己罪过重大的失足者，不要离开向主的恳求，也不要抛弃主所指定独一的天主教会；反之，要持续做补赎并恳求主的仁慈，让他们敲教会的\Emph{门}，为能重新进入他们曾经所在的地方，并回到他们曾离开的基督那里，不要听信那些以欺骗和致命的诱惑误导他们的人；因为经上记载：\BibleQuote{不要让人用虚浮的话欺骗你们，因为为了这些事，天主的忿怒必临到悖逆之子身上；所以你们不要与他们同伙。}\BibleRef{弗 5:6} 因此，不要让任何人自己与那些倨傲不恭、不敬畏天主、并完全离开教会的人交往。但如果有人不愿恳求那被冒犯的主，也不肯听从我们，反而追随那些绝望和堕落的人，那么当审判之日来临时，他必须自己承担责任。因为，一个既在之前否认过基督、如今又否认基督的教会、并且不服从健全、有益并活着的诸位主教、而让自己成为将死之人的同伙和共犯的人，在那一天怎能恳求主呢？我祝愿你们，最亲爱的、长久渴望的弟兄们，常常在衷心告别中安好。\par

\SourceRecord{Translated by Robert Ernest Wallis.；Ante-Nicene Fathers；Vol. 5.；Edited by Alexander Roberts, James Donaldson, and A. Cleveland Coxe.；Buffalo, NY: Christian Literature Publishing Co.,；1886.；<http://www.newadvent.org/fathers/050663.htm>.}

% structure: p0001 source=h1 target=section reason=page-title

\section{第六十四封书信}

\Emph{致\Person{罗加提亚努斯}，关于与主教争辩的执事。}\par

\Emph{概述——西彼廉告诫罗加提亚努斯主教要抑制那位以侮辱激怒他的执事的骄傲，并迫使他为自己的鲁莽忏悔；并借此机会重申前信中关于司铎权或主教权的一切论述。}\par

1. 西彼廉致其弟兄罗加提亚努斯，安好。我和与我同在的同道们在阅读你的信时深感悲痛和忧虑，最亲爱的弟兄，你在信中抱怨你的执事，忘记了你司铎的身份，也不顾及他自己的职务和职分，以侮辱和伤害激怒了你。而你确实行事相称，以你一贯对我们的谦逊，宁愿向我们抱怨他，尽管你有权依据主教的威能和你的教座权威，立即惩治他，确信我们所有同道都会满意，无论你以司铎权能对傲慢的执事做什么，因为你对这类人有天主的命令。上主天主在\BibleBook{}{申命纪}中说：\Quote{若有人擅敢行事，不听从那司铎或审判官，那人应处死；众百姓听见，必害怕，再不敢行恶。} \BibleRef{申 17:12} 为使我们知道这一天主的话语为尊崇和报复他的司铎是出自他的真实至极威严，当三位仆人——科辣黑、达堂和阿彼辣——胆敢傲慢行事，抬头对抗亚郎司铎，并与被立在他们之上的司铎平等时，他们被大地裂开吞没，立即承受了亵渎的胆大妄为的处罚。不仅他们，还有与他们同伙的二百五十人，也被上主降下的火吞噬，以证明天主的司铎是由那位设立司铎的报仇。在\BibleBook{}{列王纪}中，当撒慕尔司铎被犹太人民因年老而轻视，正如你现在一样，上主发怒喊叫说：\Quote{他们不是弃绝了你，而是弃绝了我。} 为报复此事，他立撒乌耳为他们的王，使他以严重伤害折磨他们，践踏百姓，以各种侮辱和惩罚压服他们的骄傲，使被轻视的司铎通过神圣报复在傲慢的人民身上得报仇。\par

2. 此外，在圣神中建立的撒罗满也见证并教导司铎的权威和权能，说道：\Quote{你当全心敬畏上主，并尊敬他的司铎；} \BibleRef{德 7:29} 又说：\Quote{当全心尊敬天主，并尊敬他的司铎。} \BibleRef{德 7:31} 铭记这些诫命，蒙福的宗徒保禄，按照我们在\BibleBook{}{宗徒大事录}中读到，当有人对他说：\Quote{你竟辱骂天主的大司祭？} 他回答说：\Quote{弟兄们，我不知道他是大司祭；因为经上记着：不可毁谤你百姓的官长。} \BibleRef{宗 23:4} 此外，我们的主耶稣基督他自己，我们的君王、审判者和天主，甚至在受难之日也尊重司铎和大司祭的荣誉，尽管他们既不敬畏天主也不承认基督。因为他治好痲疯病人时，对他说：\Quote{去，给司祭察看，并献上礼物。} 他以那教导我们也要谦卑的谦逊，仍然称他为司祭，尽管他知道他是亵渎的；甚至在受难的痛苦中，当他挨了一巴掌，有人对他说：\Quote{你竟这样回答大司祭吗？} 他并没有对大司祭本人说什么辱骂的话，而是维护自己的清白，说：\Quote{我若说得不对，你指证不对之处；若说得对，你为什么打我？} \BibleRef{若 18:23} 所有这一切都是他谦卑忍耐地做的，为使我们有谦卑和忍耐的榜样；因为他教导真正的司祭应该合法并完全受尊敬，通过表明他自己在假司祭面前的行为。\par

3. 但执事们应当记住：主选择了宗徒，即主教和监督者；而宗徒们在主升天之后为自己设立了执事，作为他们主教职和教会的事奉者。但我们若敢对那设立主教的天主有所放肆，执事们也就敢对那设立他们的人放肆；因此你所写的那个执事应当为他的鲁莽忏悔，承认司铎的尊荣，并以完全的谦卑满足那位被立在他之上的主教。因为这些都是异端者的开端，以及存心不良的分裂者的根源和意图——自以为是，以傲慢的骄矜轻视那在他们之上的人。如此他们离开教会——如此在外面设立亵渎的祭台——如此反叛基督的和平，以及天主的安排与合一。但如果他进一步以侮辱来骚扰激怒你，你必须运用你尊严的权能对付他，或是罢免他，或是开除他。因为如果宗徒保禄写信给弟茂德说：\Quote{不可叫人小看你年轻，} \BibleRef{弟前 4:12} 那么你的同僚们更该如何对你说：不可叫人小看你年迈？既然你写信说，有一个人已与你这执事联合，并分享他的骄傲和鲁莽，你也可以制止或开除他，以及其他任何似乎有类似倾向并反对天主司铎的人。除非，正如我们劝告和建议的，他们最好认识到自己犯了罪而做出补赎，并让我们保持我们自己的意向；因为我们宁可借着仁慈和忍耐来克服个别人的斥责和伤害，而不是用我们的司铎权能来惩罚他们。最亲爱的弟兄，愿你始终在主内衷心安康。\par

\SourceRecord{Translated by Robert Ernest Wallis.；Ante-Nicene Fathers；Vol. 5.；Edited by Alexander Roberts, James Donaldson, and A. Cleveland Coxe.；Buffalo, NY: Christian Literature Publishing Co.,；1886.；<http://www.newadvent.org/fathers/050664.htm>.}

% structure: p0001 source=h1 target=section reason=page-title

\section{第六十五封书信}

\Emph{致住在\Place{弗尔尼}的圣职人员和子民，关于\Person{维克托}，他曾指定司铎\Person{福斯蒂努斯}为监护人。}\par

\Emph{论题。——由于\Person{杰米尼乌斯·维克托}违反主教会议的决定，在其遗嘱中指定司铎\Person{杰米尼乌斯·福斯蒂努斯}为其监护人或代管人，因此作者禁止为他献祭，也不应为他的安息举行祭献，并顺便从肋未支派的例子推论出圣职人员不应卷入世俗事务。}\par

1. \Person{西彼廉}向住在\Place{弗尔尼}的司铎、执事和子民问安。我们最亲爱的弟兄们，我与当时在场的同僚们，以及与我们同坐的诸位司铎同僚，得知我们的弟兄\Person{杰米尼乌斯·维克托}在离世时指定司铎\Person{杰米尼乌斯·福斯蒂努斯}为其遗嘱的执行人，深感不安。因为早在主教会议中已有规定，任何人不得指定任何圣职人员或天主的仆人为其遗嘱的执行人或监护人，因为凡蒙受神圣司祭职尊荣、在圣职服务中被祝圣的人，只应服务于祭台和祭献，并有余暇从事祈祷和恳求。因为经上记载：\Quote{没有人能服兵役为天主而战，而将自己纠缠在世俗的事务中，以求得那招募他的人的欢心。}既然这话是对众人说的，那么那些忙于神圣和灵性事务、不能脱离教会、无暇顾及世俗事务的人，岂不更不应被世俗的焦虑和牵挂所束缚！关于这种祝圣和约束的形式，肋未人从前在律法之下也遵守过，以致当十一个支派分地分产业时，肋未支派——他们被留作自由以服务圣殿、祭台和神圣职务——并未从那份产业中分得任何东西；但其他人耕种土地时，他们只耕获天主的恩宠，从十一个支派那里收取什一奉献，以作食物和给养，取自生长的果实。这一切都是出于天主的权威和安排，以使那些从事神圣服务的人绝不会被呼召离开，也不被迫考虑或处理世俗事务。这一计划和规则如今对圣职界仍保持，使得那些在主教会中经由圣职祝圣而被提升的人，绝不被从神圣管理职务中召离，也不被世俗的焦虑和事务缠累；而是因着捐助的兄弟们的尊敬，收取如同果实的十分之一，他们不应离开祭台和祭献，却要日夜服务天上和灵性的事物。\par

2. 我们的前任主教们虔敬地考虑此事并有益地予以规定，决定任何离世的弟兄不得指定圣职人员为执行人或监护人；若有人如此行，则不得为他献祭，也不得为他的安息举行祭献。因为那曾企图将司铎和职员从祭台召离的人，不配在天主的祭台上被司铎的祈祷所纪念。因此，既然\Person{维克托}违反司铎们在最近会议中所立的规定，胆敢指定司铎\Person{杰米尼乌斯·福斯蒂努斯}为其执行人，你们不得为他安息而献祭，也不得在教会中以他的名字祈祷，这样我们才能遵守司铎们虔诚而必要制定的法令；同时，给其余弟兄们一个榜样，叫任何人不得将那些忙于天主的祭台和教会服务的司铎及职员呼召到世俗忧虑中。因为如果现在所行之事受到惩罚，将来或许会注意不再对圣职人员发生此类情形。我最亲爱的弟兄们，愿你们永远在主内衷心平安。\par

\SourceRecord{Translated by Robert Ernest Wallis.；Ante-Nicene Fathers；Vol. 5.；Edited by Alexander Roberts, James Donaldson, and A. Cleveland Coxe.；Buffalo, NY: Christian Literature Publishing Co.,；1886.；<http://www.newadvent.org/fathers/050665.htm>.}

% structure: p0001 source=h1 target=section reason=page-title

\section{第六十六书函}

\Emph{致主教\Person{斯德望}，论及\Place{阿尔勒}的\Person{玛尔西安}，他加入了\Person{诺瓦提安}派。}\par

\Emph{论旨——由于\Place{阿尔勒}主教\Person{玛尔西安}跟随\Person{诺瓦提安}派，已迷惑多人，并因他的分裂而与其余主教的共融分离，\Person{西彼廉}警告\Person{斯德望}，他应宣布开除该犯错者，由\Place{罗马}和\Place{迦太基}共同执行，使\Place{阿尔勒}教会能另选他人；如此，对跌倒者及被其迷惑者，在他们悔改后，可赐予和平，并准其回归教会。}\par

1. \Person{西彼廉}问候他的弟兄\Person{斯德望}。我们的同僚\Person{福斯提诺}，驻留在\Place{里昂}，曾多次写信给我，最亲爱的弟兄，告知我那些事，我也确知他及其他同省的主教同仁已告知你：即住在\Place{阿尔勒}的\Person{玛尔西安}已与\Person{诺瓦提安}联合，脱离了天主教会的合一，以及我们团体和司祭职的共识，持守着那种异端狂妄中最极端的堕落，即拒绝将天主之爱和父性温柔的安慰与帮助给予那些悔改、哀恸、以泪、呻吟和忧伤叩击教会之门的仆人，受伤者不得进入以安抚其伤口，反被抛弃，无望获得和平与共融，成为豺狼的猎物和魔鬼的战利品。最亲爱的弟兄，这事正是我们应提供建议和帮助的，因为我们顾念天主的仁慈，并在治理教会时持守平衡，对罪人展现严厉的斥责，同时也不拒绝以天主的良善和慈悲之药来扶起跌倒者、医治受伤者。\par

2. 因此，你应写信给我们派驻\Place{高卢}的主教同仁，一封非常详细的信函，绝不再容忍\Person{玛尔西安}——他倔强傲慢，敌视天主的仁慈和弟兄们的得救——侮辱我们的集会，因为他似乎尚未被我们开除教籍；而他早已长期夸耀并宣称，他依附\Person{诺瓦提安}并跟随其倔强，已脱离我们的共融；尽管他所跟随的\Person{诺瓦提安}本人早先已被开除教籍，并被判定为教会的敌人；当他派遣使者到\Place{非洲}我们这里，请求被接纳进入我们的共融时，从当时在场多位司铎的会议中得到的回话是：他本人已将自己排除，我们任何人都不能接纳他进入共融，因为他曾试图设立一个亵渎的祭台，建立一个非法的宝座，献上违背真司铎的渎神祭品；而\Person{科尔内略}主教是因天主的判断，并藉圣职人员和民众的投票，在天主教会中被祝圣的。因此，如果他愿意回到正念，醒悟过来，就应忏悔，以恳求者的身份回归教会。最亲爱的弟兄，当\Person{诺瓦提安}最近已被天主的司铎们全世界地击退、拒绝并开除教籍，我们居然还容忍他的谄媚者现在嘲笑我们，并判断教会的威严和尊荣，这是多么徒劳！\par

3. 应由你致信给该省及住在\Place{阿尔勒}的民众，开除\Person{玛尔西安}教籍后，另选他人接替他的职位，使基督的羊群——至今仍被他轻视、驱散、伤害——得以聚集。近年来我们许多弟兄在那些地区未得和平而去世，这已经够了；务必要帮助其余仍活着的人，他们昼夜呻吟，恳求天父的仁慈，并求我们援助的安慰。因为，最亲爱的弟兄，司祭团体如此广阔，彼此以共融的纽带和合一的联系相结合，正是为此：若我们同僚中有人妄图开创异端，撕裂和摧毁基督的羊群，其他人可以帮忙，如同仁慈有用的牧人，将主的羊群聚集到羊栈中。试想，若海中的某个港口因堤防破裂而开始对船只造成危害和危险，水手们难道不会将船只转向其他相邻的港口，那里有安全可行的入口和稳固的停泊处？或者，若路上的某家旅店开始被强盗包围占据，以致进去的人都会被埋伏者的攻击所捕捉，旅人们一旦发现这一情况，难道不会寻找路上更安全的其他客栈，那里住宿可靠，旅店对旅行者安全？现在我们也应当如此，最亲爱的弟兄，我们应以乐意和慈爱的人文精神接纳我们的弟兄，他们被\Person{玛尔西安}的礁石颠簸，正在寻求教会安全的港口；我们应为旅客提供如福音中所说的那种客栈，在那里被强盗击伤致残的人可以被接纳、照顾，并受主人的保护。\par

4. 因为还有什么比监督者更伟大或更尊贵的关怀，莫过于以殷勤的担忧和有益的良药来滋养和保护羊群呢？主曾说：\Quote{病弱的，你们没有扶助；患病的，你们没有医治；受伤的，你们没有包扎；被赶散的，你们没有领回；迷失的，你们没有寻找。我的羊因没有牧人而四散，成了田野一切野兽的食物，并且无人寻找，无人追寻。因此，主说：看，我要攻击那些牧人，从他们手中追讨我的羊，使他们不再牧放羊群；他们也不再牧放，因为我要从他们口中救出我的羊，并要以正义牧养他们。} 既然主这样以如此严厉的话威胁那些忽略并使主的羊丧亡的牧人，最亲爱的弟兄，我们除了以全部的热忱来聚集和挽回基督的羊，并以为父的慈爱之良药来医治失足者的创伤，还当做什么呢？因为主在福音中也警告说：\Quote{健康的人不需要医生，而是有病的人需要。} \BibleRef{玛 9:12} 虽然我们是许多牧人，但我们牧养同一群羊，应当聚集并珍视基督以自己的血和苦难所寻找的所有羊；我们不应容许我们那些哀求与哀伤的弟兄被某些人傲慢的自负所残忍地蔑视和践踏，因为经上记载：\Quote{骄傲自大的人，必一无所成，他使自己的灵魂宽广如地狱。} \BibleRef{希 2:5} 主在福音中也责备并谴责这类人，说：\Quote{你们在人前自充义人，但天主知道你们的心；因为人在众人中尊崇的，在天主面前却是可憎的。} \BibleRef{路 16:15} 祂说那些自满自足、膨胀自负、傲慢地自以为是的人是可憎可厌的。既然\Person{玛尔喜安}已经成了这样的人，并与\Person{诺瓦提安}结盟，成为仁慈与相爱之敌，就让他不要宣判，而要接受宣判；让他不要以为自己是司铎团体的审判者，因为他自己正被所有司铎所审判。\par

5. 必须维护我们先辈——真福殉道者\Person{科尔乃略}和\Person{路爵}——的光荣尊荣。我们尊敬他们的纪念，而最亲爱的弟兄，你既然成了他们的代理人和继任者，就更应凭着你的权威和份量尊崇和珍视他们。因为他们充满圣神，并在光荣的殉道中立定，断定应赐予失足者和平，并认为在完成补赎之后，不应拒绝和平与共融的赏报；他们以书信证明了这一点，而我们在所有地方、全体一致地作出了同样的判断。因为在我们中间，既然只有一个圣神，就不可能有不同的意见；因此，显然地，凡我们所见其思想不同的人，并不与众人一同持守圣神的真理。请明白告知我们，在\Place{阿尔勒}谁接替了\Person{玛尔喜安}，以便我们知道该派我们的弟兄去谁那里，又该给谁写信。最亲爱的弟兄，祝您永远安康。\par

\SourceRecord{Translated by Robert Ernest Wallis.；Ante-Nicene Fathers；Vol. 5.；Edited by Alexander Roberts, James Donaldson, and A. Cleveland Coxe.；Buffalo, NY: Christian Literature Publishing Co.,；1886.；<http://www.newadvent.org/fathers/050666.htm>.}

% structure: p0001 source=h1 target=section reason=page-title

\section{书函 67}

\Emph{致西班牙的圣职人员及信众——关于\Person{巴西里德}与\Person{玛尔提亚}}\par

\Emph{论旨。——\Person{巴西里德}与\Person{玛尔提亚}两位主教曾背教并受偶像崇拜证书的玷污，\Person{西彼廉}与其同僚主教赞扬西班牙的圣职人员及信众以合法选举推选\Person{撒比诺}与\Person{斐理斯}取代之；尤其按\Person{高乃略}及其同僚的决定，背教的主教固然可接受忏悔，但被禁止享有司铎尊荣。此外，他顺便提及关于主教选举古礼的若干事项。上下文显示此信写于\Person{斯德望}主教任内。}\par

1. \Person{西彼廉}、\Person{凯基略}、\Person{普利莫}、\Person{波利卡普}、\Person{尼各默德}、\Person{路西利安}、\Person{苏克塞苏}、\Person{塞达图}、\Person{福图纳图}、\Person{雅努雅略}、\Person{塞昆丁}、\Person{蓬波纽}、\Person{奥诺拉图}、\Person{维克多}、\Person{奥勒留}、\Person{萨提乌}、\Person{伯多禄}、\Person{另一位雅努雅略}、\Person{撒图宁}、\Person{另一位奥勒留}、\Person{文南修}、\Person{基耶图}、\Person{罗加提安}、\Person{特纳克斯}、\Person{斐理斯}、\Person{福斯提诺}、\Person{奎因图}、\Person{另一位撒图宁}、\Person{路齐奥}、\Person{文森修}、\Person{利博苏}、\Person{格弥纽}、\Person{玛尔策禄}、\Person{杨布}、\Person{阿德尔斐}、\Person{维克多里克}与\Person{保禄}，致在\Place{莱吉奥}与\Place{阿斯图里卡}居住的司铎\Person{斐理斯}及众信友，以及在\Place{埃梅里塔}居住的执事\Person{莱里奥}及众信友，在主内诸位弟兄，安好。亲爱的弟兄们，当我们聚集时，我们读了你们的信函，你们本着信德的纯全与对天主的敬畏，经由我们同为主教的\Person{斐理斯}与\Person{撒比诺}寄给我们，告知\Person{巴西里德}与\Person{玛尔提亚}因沾染偶像崇拜证书并被邪恶罪行的意识所束缚，不应持有主教职并管理天主的司铎职；你们希望就这些事再得到回信，并以我们裁决的安慰或帮助来解除你们既有理又必要的忧心。然而，对你们的这一愿望，回应的并非我们的意见，而是神圣的诫命，其中早已由天主的法律藉上天之声所命令和规定：谁及何种人应服务祭台并举行神圣祭献。因为在\BibleBook{出谷}{纪}中，天主对\Person{梅瑟}说话并警告他说：\BibleQuote{到主天主面前来的司铎应祝圣自己，免得主抛弃他们。} \BibleRef{出 19:22} 又说：\BibleQuote{当他们走近圣者的祭台服务时，不得带罪，免得死亡。} \BibleRef{出 28:43} 也在\BibleBook{肋未}{纪}中，主命令并说：\BibleQuote{凡身上有斑点或瑕疵的，不可上前向天主奉献礼物。} \BibleRef{肋 21:17}\par

2. 既然这些事已向我们宣布并阐明，我们的服从必须遵从神圣的诫命；在此类事情上，人的纵容不能看人情面，也不能向任何人让步，因为神圣的规定介入并制定了法律。因为我们不应忘记主藉先知\Person{依撒意亚}对犹太人所说的话，责备并愤怒他们轻视神圣的诫命而追随人的道理：\BibleQuote{这民族，}祂说，\BibleQuote{用嘴唇尊敬我，心却远离我；但他们的敬拜是徒然的，因为他们教导人的道理和命令。} \BibleRef{依 29:13} 这也在福音中由主重复并说：\BibleQuote{你们废弃了天主的命令，为要坚守你们的传统。} \BibleRef{谷 7:13} 我们以此为鉴，谨慎而虔诚地予以考虑，在任命的司铎时，应只选择无玷污、正直的侍奉者，他们圣洁而相称地向天主献祭，他们为保全主的子民所作的祈祷才能蒙垂听，因为经上记载：\BibleQuote{天主不听罪人，但若有人敬拜天主并遵行祂的旨意，祂便垂听他。} \BibleRef{若 9:31} 因此，应当以充分的勤勉和真诚的查考，选择那些显然天主会垂听的人来担任天主的司铎职。\par

3. 也不要让百姓自欺，以为与有罪的司铎共融、赞同监督者不义不合法的监督权，就能免于罪的传染，因为天主藉\Person{欧瑟亚}先知的责罚警告并说：\BibleQuote{他们的祭品如同哀悼之饼，凡吃它的都要受玷污；} \BibleRef{欧 9:4} 明确教导并显示凡被亵渎不义司铎的祭品所玷污的，都绝对被罪所束缚。此外，我们也发现这在\BibleBook{户籍}{纪}中得到表明，当时\Person{科辣黑}、\Person{达堂}与\Person{阿彼兰}反对司铎\Person{亚郎}，为自己要求献祭的权力。那里主也藉\Person{梅瑟}命令百姓与他们分开，免得与恶人结合，自己也同样被紧紧束缚于邪恶之中。\BibleQuote{你们要从这些邪恶顽固之人的帐棚中分离出来，}祂说，\BibleQuote{不要接触属于他们的东西，免得你们在他们的罪中一同灭亡。} \BibleRef{户 16:26} 因此，一个顺从主的诫命、敬畏天主的百姓，应当与有罪的监督分离，不与亵渎的司铎的祭品结合，尤其因为他们自己有权选择相称的司铎或拒绝不相称的。\par

4. 我们也观察到这件事出自天主的权威，就是司铎应在人民面前、在众人眼下被拣选，并应通过公众的判断和见证被认可为相称和合适；正如在\BibleBook{}{户籍纪}中，主命令\Person{梅瑟}说：\BibleQuote{你要带着你的兄弟\Person{亚郎}和他的儿子\Person{厄肋阿匝尔}，把他们带到山上，在全会众面前，脱去亚郎的衣冠，给厄肋阿匝尔穿上；亚郎就死在那里，归到他本族的人那里。}\BibleRef{户 20:25} 天主命令司铎应当在全会众面前被任命；也就是说，祂教导并表明，司铎的晋秩不应在附近的人民不知情的情况下举行，而应在人民面前，要么揭露恶人的罪行，要么宣告善人的功绩，并且经过所有人的投票和审查的晋秩才是正义和合法的。这后来在\BibleBook{宗徒}{大事录}中也按照天主的指示得到遵守，当\Person{伯多禄}向人民讲话，要任命一位宗徒代替\Person{犹达斯}时。经上说：\BibleQuote{伯多禄在门徒中间站起来，那时众人都在一处。}\BibleRef{宗 4:2} 我们也观察到，宗徒们不仅在主教和司铎的晋秩中遵守了这一点，而且在执事的晋秩中也如此，关于这件事本身，在\BibleBook{宗徒}{大事录}中也记载着：\BibleQuote{他们就召集了门徒全体会众，对他们说：}\BibleRef{宗 4:2} 这件事做得如此勤劳和谨慎，召集了全体人民，毫无疑问是为了这个原因：使任何不相称的人都不能潜入祭台的服务或司铎的职务。因为不相称的人有时被晋秩，不是按照天主的旨意，而是按照人的僭越，并且那些不是出自合法和正义的晋秩的事令天主不悦，天主亲自通过\Person{欧瑟亚}先知表明，说：\BibleQuote{他们自立了君王，却没有我。}\BibleRef{欧 8:4}\par

5. 因此，你们必须勤勉地遵守并持守那源自神圣传统和宗徒惯例的做法，这做法也在我们中间和几乎所有省份中得以保持：即为了正确举行晋秩，同一省份的所有邻近主教应与即将被祝圣主教的教区的人民一同聚集。并且主教应在人民面前被拣选，这些人最充分地了解每个人的生活，并考察了每个人的行为与其惯常的品行。我们也看到，你们在祝圣我们的同僚\Person{撒比诺}时也是这样做的；因此，通过全体弟兄的投票和聚集在场的众主教的判决（他们曾就他写信给你们），主教职被授予了他，并且在他身上覆手，以代替\Person{巴西利德}。\Person{巴西利德}在罪行被揭露、良心甚至通过他自己的告解而暴露之后，前往罗马，欺骗了我们的同僚\Person{斯德望}——他远在天边，对已发生之事和真相一无所知——以便进行游说，使他能不正当地恢复他被正义地罢免的主教职。这结果并非抹除了巴西利德的罪，反而增加了，因为在他先前的罪上又加上了欺骗和蒙蔽的罪行。因为因疏忽而被诡计欺骗的人，其罪责不如那以诡计欺骗他的人那样可憎。但是，如果巴西利德能欺骗人，却不能欺骗天主，因为经上记载：\BibleQuote{天主是轻慢不得的。}但欺骗也不能使\Person{马尔提亚利斯}得利，使他虽然也深陷重罪，却仍能保持他的主教区，因为宗徒也警告说：\BibleQuote{主教必须是无可指摘的，如同天主的管家。}\BibleRef{铎 1:7}\par

6. 因此，亲爱的弟兄们，正如你们所写的，以及我们的同僚\Person{费利克斯}和\Person{撒比诺}所证实的，还有\Place{凯撒奥古斯塔}的另一位\Person{费利克斯}——一位信德的维护者和真理的捍卫者——在他的信中所指明的，\Person{巴西利德}和\Person{马尔提亚利斯}已被可憎的偶像崇拜证书所玷污；而且巴西利德，除了证书的污点之外，当他卧病在床时，还亵渎了天主，并承认他亵渎了；因自己良心的创伤，他自愿放弃主教职，转向悔改，恳求天主，并认为自己如果甚至能以平信徒的身份领受共融，已是足够幸福的了：此外，马尔提亚利斯，除了长期参与外邦人可耻和污秽的宴席（在他们的社团中），并将其子置于同一社团中，按外邦人的方式，置身于亵渎的坟墓中，并与陌生人一同埋葬他们，他已通过公开在一位\OriginalTerm{杜塞纳里乌斯级检察官}{ducenarian procurator}面前所做的行为，确认自己已屈服于偶像崇拜，并否认了基督；既然还有许多其他严重的罪行，巴西利德和马尔提亚利斯被认为与之有牵连；这样的人企图为自己主张主教职是徒劳的；因为显而易见，这类人既不能统治基督的教会，也不应向天主献祭，更何况我们的同僚\Person{科尔内略}，一位和平而正义的司铎，且更因主的屈尊就卑以殉道为荣耀，早已与我们以及全世界所有被任命的主教一同规定，这种人确实可以被接纳悔改，但被禁止受圣职的晋秩和司铎的尊荣。\par

7. 最亲爱的弟兄们，若在这些末世里，有些人的信德摇摆不定，对天主的敬畏毫无虔诚地动摇，或和平的和谐无法继续，你们不必为此困扰。这些事在世界的终结时已被预言将要发生；主的口舌和宗徒们的见证早已预言：如今世界衰败，假基督临近，一切美善都将衰微，而邪恶与敌对之事却将兴旺。\par

8. 然而，尽管在这末世，天主的教会中福音的严厉并未完全消逝，基督徒的德行或信德的力量也未如此萎靡，以致没有留下一部分司铎，他们在这些事物的废墟和信德的残骸面前决不退缩；反而勇敢坚定，以完全的敬畏之心维护天主威严和司铎尊荣。我们记得并注视着，虽然其他人屈服退让，\Person{玛塔提雅}却勇敢地捍卫了天主的法律；当犹太人退却并背离神圣宗教时，\Person{厄里亚}屹立并高尚地争战；\Person{达尼尔}既不被异国的孤寂所吓倒，也不被不断迫害的骚扰所阻，屡次光荣地忍受殉道；还有那三位青年，既不为年幼所制，也不为威胁所屈，忠信地抵抗了巴比伦的烈火，甚至在他们的俘虏生活中征服了那胜利的君王。让那些背信者或叛徒的数量去操心吧，他们现在开始在教会内兴起反对教会，败坏信德和真理。在很多人中间仍然存在真诚的心志和坚固的宗教，以及只献身于主及其天主的灵。别人的背信弃义并不能将基督徒的信德压入毁灭，反而激发它、提升它至光荣，正如蒙福的宗徒\Person{保禄}所劝勉并说的：\Quote{纵然有些人背弃了信德，难道他们的不信能使天主的信德失效吗？绝对不可！天主是真实的，众人都虚假。}\BibleRef{罗 3:3}但如果人人都虚假，唯独天主真实，那么我们这些天主的仆人，尤其是司铎，除了抛弃人的错误和谎言，持守天主的真理、遵守主的诫命，还能做什么呢？\par

9. 因此，最亲爱的弟兄们，虽然在我们的同僚中发现有些人以为虔诚的纪律可以忽略，并与\Person{巴西利德}和\Person{玛提雅利斯}轻率共融，这样的事不应扰乱我们的信德，因为圣神在\BibleBook{}{圣咏}中如此警告说：\Quote{你恼恨管教，把我的话抛在脑后：你见了盗贼，就与他同谋，又与犯奸者同流合污。}这表明，与罪犯交往的人成了他人罪过的同谋和共犯。而且，宗徒\Person{保禄}也写道，并说了同样的话：\Quote{爱说闲话的，背后诋毁的，憎恨天主的，侮辱人的，傲慢的，自夸的，捏造恶事的；他们虽然明知天主的判决，却不明白行这样事的人是当死的，不仅行这样事的人，连赞同他们的人也是如此。}\BibleRef{罗 1:30}既然祂说，行这样事的人当受死亡，祂就表明并证明：不仅作恶的人当受死亡、进入惩罚，连赞同他们的人也当如此——他们因与邪恶和罪人及不悔改者不法共融，被罪人的接触玷污，且共担罪过，因此在惩罚中也不得分离。为此，最亲爱的弟兄们，我们不仅赞同，而且赞赏你们纯洁和信德的虔诚关注，并尽我们所能用信件劝勉你们：不要与亵渎和受污染的司铎进行亵圣的共融，而要以虔诚的敬畏保持你们信德的健全和坚定。我祝愿最亲爱的弟兄们始终在主内安康。\par

\SourceRecord{Translated by Robert Ernest Wallis.；Ante-Nicene Fathers；Vol. 5.；Edited by Alexander Roberts, James Donaldson, and A. Cleveland Coxe.；Buffalo, NY: Christian Literature Publishing Co.,；1886.；<http://www.newadvent.org/fathers/050667.htm>.}

% structure: p0001 source=h1 target=section reason=page-title

\section{第六十八封书信}

致弗洛伦提乌斯·普皮亚努斯，论诽谤者。\par

《论题》——西彼廉向弗洛伦提乌斯·普皮亚努斯澄清自己被其指控的各项罪行，并指出他轻信诽谤者之浮躁心态。\par

1. 西彼廉，亦名塔西乌斯，致弗洛伦提乌斯，亦名普皮亚努斯，其弟兄，问候。弟兄，我原以为你如今终于悔改，为从前曾如此轻率地听闻或相信了关于我的那些在外邦人中也算邪恶、可耻、可憎之事。然而，即便在你的来信中，我仍看出你与从前无异——你对我仍持相同的看法，且坚持你所信之事；惟恐你那显赫的地位和殉道因与我共融而受玷污，你便仔细查究我的品性；并且，在天主这位任命司铎的审判者之后，你竟要审判——我不说是审判我（我算什么？）——而是审判天主和基督的审判。这并非信从天主，而是起来反抗基督及其福音；因为祂说：\BibleQuote{两只麻雀不是卖一个铜钱吗？没有你们天父的旨意，它们中一只也不会掉在地上。}\BibleRef{玛 10:29} 祂的尊威和真理证明，即使微小之事，若无天主的认知和许可，也不会发生；而你竟认为，天主的司铎是在教会中未经祂的知晓而被任命的。因为相信被任命的人是不配和邪淫的，这难道不是相信天主的司铎不是由天主、也非藉着天主在教会中任命的吗？\par

2. 你以为我自己的见证比天主的更好吗？主亲自教导并说，若有人为自己作证，他的见证便不真实，因为每人必然偏袒自己。没有人会提出对自己有害和不利之事；然而，若关于我们的事由他人宣告和见证，则可能有纯朴的真实。祂说：\BibleQuote{我若为我作证，我的见证不真实；但有一位为我作证。}\BibleRef{若 5:31} 但若主自己，祂即将审判一切，都不愿凭自己的见证被相信，而宁愿藉着天主父的审判和见证被认可，那么祂的仆人们岂不更应遵守此事？他们不仅藉着天主的审判和见证被认可，而且以此夸耀。然而，在你那里，敌对和恶人的捏造竟胜过了天主的谕令和基于信仰之坚毅的我们的良心，仿佛在那些身处教会之外、堕落和亵渎的人中（圣神已从他们胸中离去），还能有什么，除了败坏的心、诡诈的舌、毒害的仇恨和亵渎的谎言？谁相信这些，审判之日来临时，他必与他们一同被找到。\par

3. 至于你所说，司铎应当谦卑，因为主和祂的宗徒都是谦卑的；所有弟兄和外邦人都深知并喜爱我的谦卑；你从前仍在教会中、与我共融时，也曾知道并喜爱它。但我们之中谁远离谦卑呢？是我，每日服事弟兄，善意喜乐地接待每一个来到教会的人；还是你，自封为主教的主教、审判由天主按时赐下的审判者？虽然主天主在申命纪中说：\BibleQuote{若有人擅敢行事，不听从当日的司铎或审判官，那人必死；众百姓听见了便害怕，不再擅敢行事。}\BibleRef{申 17:12} 祂又对撒慕尔说：\BibleQuote{他们不是厌弃你，而是厌弃了我。}\BibleRef{撒上 8:7} 此外，主在福音中，当人对祂说：\Quote{你竟这样答复大司祭吗？}时，祂维护司铎的尊威，教导它应当被持守，对大司祭什么也不说，只澄清自己的无辜，回答说：\BibleQuote{我若说得不好，你指证那不好；若说得好，你为什么打我呢？}\BibleRef{若 18:23} 蒙福的宗徒，当人对他说：\Quote{你辱骂天主的司祭吗？}时，他没有对司祭说任何辱骂的话，尽管他本可大胆地起来反对那些钉死主、已牺牲天主和基督、圣殿和司祭职的人；但即使在虚假和堕落的司铎身上，他仍顾念司铎之名的空影，说：\BibleQuote{弟兄们，我不知道他是大司祭；因为经上记着说：不可毁谤你百姓的首领。}\BibleRef{宗 23:4}\par

4. 除非，也许在迫害之前，当你与我共融时，我是你的司铎；而迫害之后，我就不再是司铎了！因为迫害到来时，把你提升到殉道的最高荣耀；但它却以流放的负担压垮了我，因为公开宣告：\Quote{凡持有或拥有基督徒的主教凯齐留斯·西彼廉的任何财产者……}，以致那些不相信天主任命主教的人，仍能相信魔鬼流放了主教。我并非以此自夸，而是忧伤地提出这些事，因为你自封为天主和基督的审判者，而祂对宗徒们，并藉此对一切以继承方式继任宗徒的首领们说：\BibleQuote{听从你们的，就是听从我；听从我的，就是听从那派遣我来的；轻视你们的，就是轻视我，以及那派遣我来的。}\BibleRef{路 10:16}\par

5. 因为由此就产生了，并且仍在产生分裂和异端，就在于那位唯一的、统治教会的主教，被某些人傲慢的狂妄所藐视；那位因天主的屈尊就卑而受尊敬的人，竟被人判定为不配。这是何等骄傲的膨胀，何等灵魂的傲慢，何等心智的夸大，竟要召见主教和司铎来受自己认可，除非我在你眼中被宣告为清白、被你判决宣布无罪，看啊，如今六年来，弟兄团体既没有主教，百姓没有领袖，羊群没有牧人，教会没有管理者，基督没有代表，天主没有司铎！\Person{普皮安}必须来救援，做出判决，宣告天主和基督的决定已被接纳，免得这么多在我治理下被召走的信友，似乎没有救恩和平安的希望就离世了；免得新的一批信徒被认为因我的职务而未能获得圣洗圣事和圣神的任何恩宠；免得我赐予这么多堕落和悔罪者的平安，以及经由我审查而赐予的共融，被你的判决权威所废除。请你屈尊一次，俯允对我们做出宣判，并以你认可的权威确立我的主教职，好使天主和祂的基督感谢你，因为藉着你，一位代表和治理者得以重新被赋予他们的祭台和他们的百姓。\par

6. 蜜蜂有王，牲畜有首领，而它们对首领保持忠诚。强盗以充满谦卑的服从服从他们的头目。粗野的牲畜和哑巴动物，以及强盗，尽管血腥，在刀剑和兵器中狂怒，他们比你更单纯、更优秀！他们中间的头目被承认和惧怕，虽然没有任何神圣的审判任命他，而只是一个被弃的派系和有罪的团伙同意了他。\par

7. 你确实说，你陷入的疑虑应当从你心中除去。你陷入它，但那是由于你不虔敬的轻信。你陷入它，但那是由于你自己亵渎的性情和意志，轻易听信不洁的、不敬的、不可言说的事反对你的兄弟，反对一位司铎，并且乐意相信它们，捍卫他人的谎言，仿佛它们是你自己的和私有的；并且不记得经上记载说：\Quote{你要用荆棘围上你的耳朵，不要听邪恶的舌头；} 又说：\Quote{行恶的人听从不义之人的舌头；但义人不理会说谎的嘴唇。} 为何充满圣神、且已因苦难临近天主和祂基督面前的殉道者们没有陷入这疑虑？他们从地牢里写信给\Person{西彼廉}主教，承认天主的司铎，并为他作证。为何如此多的主教，我的同僚，没有陷入这疑虑？他们或在我们中间离去时被放逐，或被逮捕投入监狱并戴上锁链；或被流放，以光辉的道路走向主；或在某些地方被判死刑，因主的荣耀而接受了天上的荣冠。为何从我们当中与我们同在、且因天主的屈尊就卑委托给我们的百姓中，如此多的宣信者——他们曾受审讯和折磨，因光荣伤疤的记忆而荣耀；如此多纯洁的贞女，如此多可称赞的寡妇；最后，普世所有在团结的纽带中与我们联合的教会——都没有陷入这疑虑？除非所有这些与我共融的人，如你所写，都被我嘴唇的玷污所污染，并因我共融的传染而失去了永生的望德。只有\Person{普皮安}，健全、无玷、圣洁、谦逊，不愿与我们联合，将独自居住在乐园和天国里。\par

8. 你还写道，由于我，教会如今有一部分处于分散状态，尽管全体的教会百姓都聚集、联合并团结在不可分割的和谐中；只有他们留在外面，而即使他们在里面，也本应被赶出去。主，祂百姓的保护者和守护者，不容许麦子从祂的禾场上被夺走；只有糠秕可以从教会中分离出去，因为宗徒也说：\BibleQuote{即便他们中有一些人背离了信德？他们的不信会使天主的信德失效吗？绝对不可！其实天主是真实的，众人都是虚谎的。} \BibleRef{罗 3:3} 主在福音中也说过，当门徒因祂的话离开祂时，祂转向那十二人说：\BibleQuote{你们也要走吗？} 然后\Person{伯多禄}回答祂说：\BibleQuote{主！我们要投奔谁呢？祢有永生的话，我们相信，并且确知祢是天主的圣者。} \BibleRef{若 6:67} 那里\Person{伯多禄}发言，教会将要建立在他之上，他以教会的名义教导并表明，虽然一群悖逆和傲慢、不愿听从和服从的人可能离开，但教会并不离开基督；而教会就是那些与司铎联合的百姓，以及依附其牧人的羊群。因此你应当知道，主教在教会内，教会也在主教内；若有人不与之在一起，他就不在教会内，那些偷偷潜入、不与天主的司铎保持平安、却自以为与某些人秘密共融的人，是徒然自欺；因为教会，她是大公的、唯一的，不会被切割或分裂，而是确实由彼此紧密相连的司铎的粘合剂连接并捆绑在一起。\par

9. 所以，兄弟，若你想到天主委任司铎的尊威，若你愿意有一次尊重基督——祂凭自己的命令和言语，并藉自己的临在，既治理主教们自身，又藉着主教们治理教会；若你关于主教的纯洁，宁愿信赖天主的审判，而非人间的仇恨；若你开始为你的鲁莽、骄傲和傲慢作迟来的悔改；若你极其丰盛地向天主和祂的基督——我所事奉的那一位——赔补，并以纯洁无瑕的嘴唇，不仅在和平时期，也在迫害中，不断向祂献上祭献；那么我们就有某种与你共融的基础，即使在我们中间仍存有对天主审判的敬畏与恐惧；如此，我应先向我的主咨询，祂是否允许赐予你们和平，并允许你们藉祂亲自的显示和劝告，被接纳到祂教会的共融中。\par

10. 因为我记得那已经向我启示的——不，是由我们的主、天主的权威，向一个顺从而敬畏的仆人吩咐的；在其他祂屈尊显示和启示的事中，祂还加上了这句：\Quote{谁若不相信那立定司铎的基督，以后必将开始相信那报复司铎的基督。}虽然我知道，对某些人来说，梦境似乎可笑，异象显得愚蠢；但这确是对那些宁愿反对司铎而不相信司铎的人而言。这并不奇怪，因为\Person{若瑟}的兄弟们曾论及他说：\InnerQuote{看，那做梦的人来了；来，我们杀了他。}后来那做梦的人得到了他所梦见的；杀他、卖他的人都蒙受了羞辱，以致他们起初不信话语，后来却信了事迹。至于你们在迫害中或和平中所行的事，我若假装审判你们，那是愚蠢的，因为你们反倒自命为审判我们的人。这些事，出于我良心的纯洁，以及对我主、我天主的信心，我已详细写下。你收到了我的信，我也收到了你的信。在审判之日，在基督的审判台前，二者都将被宣读。\par

\SourceRecord{Translated by Robert Ernest Wallis.；Ante-Nicene Fathers；Vol. 5.；Edited by Alexander Roberts, James Donaldson, and A. Cleveland Coxe.；Buffalo, NY: Christian Literature Publishing Co.,；1886.；<http://www.newadvent.org/fathers/050668.htm>.}

% structure: p0001 source=h1 target=section reason=page-title

\section{第六十九封书信}

致\Person{雅努阿里乌斯}及其他\Place{努米底亚}主教：论为异端者施洗。\par

论旨——本函及下一函的要旨见于后来致\Person{斯德望}的一封书信；\Quote{异端者所用的并非洗礼；凡反对基督者，无人能从基督的恩宠中获益；此点近日在致我们在\Place{毛里塔尼亚}的同僚\Person{昆图斯}的书信中已有详细阐述；此外，我们的同僚此前致书\Place{努米底亚}的主教们，亦论及此事；兹将这两封书信的副本附后。}\par

1.\Person{西彼廉}、\Person{利贝拉利斯}、\Person{卡尔多尼乌斯}、\Person{尤尼乌斯}、\Person{普里穆斯}、\Person{凯基利乌斯}、\Person{波利卡普}、\Person{尼科美德斯}、\Person{费利克斯}、\Person{马鲁提乌斯}、\Person{苏克塞苏斯}、\Person{卢基安努斯}、\Person{奥诺拉图斯}、\Person{福尔图纳图斯}、\Person{维克托}、\Person{多纳图斯}、\Person{卢基乌斯}、\Person{赫尔库拉努斯}、\Person{蓬波尼乌斯}、\Person{德米特里乌斯}、\Person{昆图斯}、\Person{撒图尔尼努斯}、\Person{雅努阿里乌斯}、\Person{马尔库斯}、另一位\Person{撒图尔尼努斯}、另一位\Person{多纳图斯}、\Person{罗加提安努斯}、\Person{塞达图斯}、\Person{特尔图鲁斯}、\Person{霍尔滕西安努斯}、又一位\Person{撒图尔尼努斯}、\Person{萨提乌斯}，致他们的弟兄\Person{雅努阿里乌斯}、\Person{撒图尔尼努斯}、\Person{马克西穆斯}、\Person{维克托}、另一位\Person{维克托}、\Person{卡西乌斯}、\Person{普罗库鲁斯}、\Person{莫迪安努斯}、\Person{基提努斯}、\Person{加尔吉利乌斯}、\Person{欧提基安努斯}、另一位\Person{加尔吉利乌斯}、另一位\Person{撒图尔尼努斯}、\Person{内梅西安努斯}、\Person{南普卢斯}、\Person{安托尼安努斯}、\Person{罗加提安努斯}、\Person{奥诺拉图斯}，问安。最亲爱的弟兄们，我们一同在会议中读了你们写给我们的信，关于那些似乎由异端者和分裂者施洗的人，问他们来到唯一的天主教会时，是否应当受洗。关于此事，虽然你们自己已持有天主教规则的真理和确定性，但由于你们认为出于我们彼此的爱德，我们应当被咨询，我们提出我们的意见，并非新意见，而是与你们一致赞同，这是我们前辈早已议定并由我们遵守的意见——即判断并确信，无人能在教会之外、在外地受洗，因为在圣教会中只有一种洗礼。主的话如此记载：\Quote{他们离弃了我这活水的泉源，为自己凿出破裂的水池，不能蓄水。}\BibleRef{耶 2:13} 再者，圣经警告说：\Quote{你要提防陌生之水，不要饮用陌生泉源的水。} 因此，必须由司铎先洁净并祝圣那水，好能藉其洗礼洗去受洗者的罪；因为主藉厄则克耳先知说：\Quote{那时，我要在你们身上洒清水，洁净你们脱离你们的一切污秽，并洁净你们脱离你们的一切偶像；我还要给你们一颗新心，在你们内注入一种新的精神。}\BibleRef{则 36:25} 然而，他自己不洁且圣神不在其中的人，怎能洁净并祝圣水呢？因为主在户籍纪中说：\Quote{凡不洁之人所触摸的，必成为不洁。}\BibleRef{户 19:2} 或者，那施洗的人，自己既在教会之外而不能除去自己的罪，又怎能给予他人罪的赦免呢？\par

2. 此外，在洗礼中所提的问题本身便是真理的见证。因为我们说：\Quote{你藉着圣教会相信永生和罪的赦免吗？}意思是，罪的赦免除非在教会内，否则不予赐予；而在异端者中，没有教会，罪不能被除去。因此，主张异端者可以施洗的人，必须要么改变问题，要么维护真理；除非他们也将教会归于那些他们声称拥有洗礼的人。受洗者也必须受傅油；这样，领受了\OriginalTerm{圣化圣油}{chrism}，即傅油，他就被天主所傅油，并在他内拥有基督的恩宠。此外，正是由圣体圣事，受洗者得以用祭台上祝圣过的油傅油。但是，一个既没有祭台也没有教会的人，不能祝圣被造物的油；因此，在异端者中不可能有属神的傅油，因为显然油不能祝圣，圣体圣事也完全不能在他们中间举行。但我们应当知道并记住经上所载：\Quote{不要让罪人的油傅在我的头上}（圣神在圣咏中早已警告），以免有人偏离正路、迷失真理之道，被异端者和基督的仇敌傅油。此外，一个不虔敬且有罪的司铎，能为受洗者献上什么祈祷呢？因为经上记载：\Quote{天主不垂听罪人；但若有人是敬畏天主的，并承行祂的旨意，天主才垂听他。}\BibleRef{若 9:31} 再者，谁能给予他自己所没有的呢？或者，他自己已失掉了圣神，又怎能执行属神的职务呢？因此，那未经训练来到教会的人必须受洗并更新，以便在教会内由圣洁的人圣化自己，因为经上记载：\Quote{你们应当是圣的，因为我是圣的，上主说。}\BibleRef{肋 19:2} 这样，那被错谬引诱、在教会之外受洗的人，甚至应当在真实而教会的洗礼中放下这一点，即：一个来到天主面前的人，在寻求司铎时，却因错误的欺骗落到了一个亵渎者身上。\par

三、但赞同异端者和分裂者的圣洗，就是承认他们确实付了洗。因为在其中，一部分不能无效，而另一部分有效。若一人能付洗，他也能赐予圣神。但他若不能赐予圣神，因为那从外部被任命的人并未被赋予圣神，他就不能给前来的人付洗；因为圣洗是一个，圣神也是一个，并且由主基督建立在\Person{伯多禄}之上的教会，凭着合一的根源和原则，也是一个。由此得出，既然在他们那里一切都是虚妄和虚假的，他们所行的一切都不应被我们认可。因为主称之为祂的敌人和反对者的人所行的事，怎能被天主批准和确立呢？主在祂的福音中教导说：\BibleQuote{不随同我的，就是反对我；不与我收集的，就是分散。}\BibleRef{路 11:23} 有福的宗徒若望也遵守主的诫命和训导，在他的书信中规定说：\BibleQuote{你们听说假基督要来，现在已经有许多假基督，因此我们就知道这是最后的时期。他们从我们中间出去，却不是属于我们的，因为如果他们属于我们，必然留在我们中间。}\BibleRef{若一 2:18} 因此，我们也应当推断和思考，那些主的反对者，被称为假基督的人，能否赐予基督的恩宠。所以，我们这些与主同在、保持主的合一、并按祂的屈尊就卑在教会中施行祂司祭职的人，应当弃绝、拒绝并视祂的反对者和假基督所行的一切为亵渎；对于那从错谬和邪恶中出来、承认唯一教会真实信仰的人，我们应当借着一切神圣恩宠的圣事，赐予他们合一与信仰的真理。最亲爱的弟兄们，我们祝你们永远衷心平安。\par

\SourceRecord{Translated by Robert Ernest Wallis.；Ante-Nicene Fathers；Vol. 5.；Edited by Alexander Roberts, James Donaldson, and A. Cleveland Coxe.；Buffalo, NY: Christian Literature Publishing Co.,；1886.；<http://www.newadvent.org/fathers/050669.htm>.}

% structure: p0001 source=h1 target=section reason=page-title

\section{第七十封书信}

\Emph{致\Person{克文托}，论异端者的圣洗。}\par

\Emph{论点——答复\Place{摩里塔尼亚}主教\Person{克文托}，他曾就异端者的圣洗征询意见。}\par

1. \Person{西彼廉}致其兄弟\Person{克文托}，问安。我们的同事司铎\Person{路先}向我报告，最亲爱的弟兄，说您希望我向您宣示我对那些似乎由异端者和分裂者施洗之人的看法；关于此事，为使您知道我们几位同僚主教与在座的司铎弟兄们最近在会议中决定了什么，我已将同一书信的副本寄给您。因为我不知道，是怎样的狂妄使我们的某些同僚认为，那些由异端者浸洗的人来投奔我们时不应受洗，其理由是：只有一个圣洗，而它之所以是一个，正是因为教会是唯一的，并且在教会之外不可能有任何圣洗。既然不能有两个圣洗，如果异端者真的施洗，他们自己就拥有这个圣洗。而任何凭自己权威将这一益处授予他们的人，便是向异端者让步并同意他们，使基督的仇敌和对手似乎拥有洗濯、洁净和圣化人的能力。但我们的说法是：从那里来的人在我们当中不是重洗，而是受洗。因为他们确实在那里什么也没有领受，因为那里一无所有；但他们来我们这里，为要在这里领受，因为这里有恩宠和全部真理，因为恩宠和真理都是唯一的。然而，我们的某些同僚宁愿尊敬异端者而不愿与我们一致；他们藉着主张只有一个圣洗而不愿给来者施洗，这样他们要么自己制造了两个圣洗，因为他们说异端者之中有圣洗；要么，更严重的是，他们竭力将异端者污秽不洁的洗置于天主教真实、唯一、合法的圣洗之前并优先考虑，没有考虑到经上记载：\BibleQuote{由死人施洗的人，他的洗有什么用处？}\BibleRef{德 34:25}如今显然，那些不在基督教会内的人被列在死人之中；一个人不能由自己本身不活着的人而得以生活，因为只有一个教会，她既获得了永生的恩宠，就永远活着，并使天主子民生活。\par

2. 他们又说，在此事上他们遵循古老传统；但在古人中，异端和分裂尚属初萌，故那些脱离教会的人最初是在教会内受洗的，因此当他们怀着悔改返回教会时，无需再施洗。这一点我们在今日也遵守：对于那些已知在教会内受洗、却从我们这里投奔了异端者的人，若他们后来承认自己的罪，摒弃错误，返回真理和他们的母亲，为他们施行悔改的覆手便已足够；因为那羊曾是羊，牧人可以将迷失流浪的羊领回自己的羊栈。但如果从异端者来的人先前并未在教会内受洗，而是作为一个陌生人且全然亵渎地前来，他就必须受洗，好成为羊，因为在圣教会内有那唯一能使羊成为羊的水。因此，既然虚假与真理、黑暗与光明、死亡与永生、假基督与基督之间毫无共通之处，我们无论如何都应当维护天主教会的合一，决不在任何方面向信仰和真理的敌人让步。\par

3. 我们也不应从传统来规定这事，而应以理性克服\Emph{相反的传统}。因为主首先拣选并在其上建立祂教会的\Person{伯多禄}，当\Person{保禄}后来与他争论割损问题时，并没有傲慢地为自己要求什么，也没有妄自尊大，以自己拥有首席权而声称新来的人和初信者更应服从他。他也没有因为\Person{保禄}先前是教会的迫害者而轻视他，反而接受了真理的建议，轻易地顺从了\Person{保禄}提出的合理理由，从而既为我们提供了和谐与忍耐的榜样，使我们不要固执地爱自己的意见，而应将我们弟兄和同僚在任何时候有益且健康地提出的建议——只要它们是真实合法的——当作自己的意见来采纳。此外，\Person{保禄}预见此事，并忠信地为和谐与平安着想，在他的书信中立下了这个规则：\Quote{至于先知们，可以两个人或三个人说话，其余的人要审断。但若在座的另一位得到了启示，那先说话的，就应该缄默。}在这里，他教导并表明，许多更好的事情会启示给个人，并且每人不应固执地为自己一度接受和坚持的主张而争辩；如果有什么显得更好、更有益，他应当乐意接受。因为当更好的事情呈现给我们时，我们不是被战胜，而是受到教诲，尤其是在那些关乎教会合一以及我们希望和信仰真理的事上；这样，我们作为天主的司铎和祂教会的监督，因祂的俯就，应当知道：罪过的赦免只能在教会内赐予，基督的仇敌不能将属于祂恩宠的任何东西据为己有。\par

4. 事实上，堪受纪念的\Person{阿格里皮诺}，与他当时管理\Place{非洲}和\Place{努米底亚}地区主教会堂的其他同僚主教，也如此议定，并通过共同会议的审慎考量确立了此规。我们亦遵循了他的见解，因其虔诚、合法、有益，且与天主教信仰和教会和谐一致。为使您知道我们就此事所写的信件，为我们的互爱，我已寄送了一份副本，供您本人和当地同僚主教知晓。我祝您，最亲爱的弟兄，永远衷心平安。\par

\SourceRecord{Translated by Robert Ernest Wallis.；Ante-Nicene Fathers；Vol. 5.；Edited by Alexander Roberts, James Donaldson, and A. Cleveland Coxe.；Buffalo, NY: Christian Literature Publishing Co.,；1886.；<http://www.newadvent.org/fathers/050670.htm>.}

% structure: p0001 source=h1 target=section reason=page-title

\section{第七十一封书信}

\Emph{致\Person{斯德望}，论会议。}\par

\Emph{论据——\Person{西彼廉}与其同僚在一次会议上通知罗马主教\Person{斯德望}，他们已经决定：凡从异端回归教会者均应受洗，且来自异端的主教或司铎，除非以平信徒身份共融，不得以其他条件被接纳。公元255年。}\par

1. \Person{西彼廉}与其他同僚，致他们的兄弟\Person{斯德望}，问候。为安排某些事项并进行共同审议，我们以为有必要召集并举行一次会议，最亲爱的兄弟，许多司铎同时聚集；会上提出并处理了许多事情。但尤其需要我们写信给您并与您的庄重智慧商讨的主题，更特别涉及司铎的权威以及天主教会的合一与尊严，这些源于神圣任命的神职。即：那些在教会外受洗、在异端者和分裂者中被亵渎之水玷污的人，当他们来到我们这里，来到唯一的教会时，应当受洗，因为\Quote{为他们覆手使他们领受圣神}是小事，除非他们也领受教会的圣洗。因为只有当他们由两件圣事而生时，才能完全被圣化，成为天主的子女；正如经上记载：\BibleQuote{人除非由水和圣神重生，不能进入天主的国。}因为我们在\WorkTitle{宗徒大事录}中也发现，宗徒们持守并保守救恩信德的真理：当圣神降临在百夫长\Person{科尔乃略}家的外邦人身上时，他们因信德的热忱而炽热，全心信靠主；当他们充满圣神，以各种言语赞美天主时，蒙福的宗徒\Person{伯多禄}仍然记念神圣的诫命和福音，命令那些已经充满圣神的人也受洗，以免在宗徒训导中似乎忽略了在任何事上遵守神圣诫命和福音的法律。至于异端者所用的不是圣洗，以及与基督对抗的人不能因基督的恩宠获益，这一点最近已在写给我们在\Place{毛里塔尼亚}的同僚\Person{昆图斯}的信中详细阐述；同样，我们同僚此前写给在\Place{努米底亚}主持的主教的信中也已说明。这两封信的副本我已附上。\par

2. 然而，最亲爱的兄弟，我们以共同的同意和权威，补充并联系前文所说：如果任何司铎或执事，或先前在天主教会中受祝圣，后来成为教会的叛徒和反叛者；或由假主教和假基督之手在异端中通过亵圣的祝圣而被擢升，违反基督的任命，并试图在唯一神圣的祭台之外，献上虚假亵圣的祭献；那么，当他们回归时，也应被接纳，但条件是：他们以平信徒的身份共融，且认为在成为和平之敌后，能被接纳到平安中已足够；他们回归时不应保留那些曾用以反叛我们的祝圣和荣誉的武器。因为侍奉祭台和祭献的司祭和仆人应当是健全无瑕的；主天主在\WorkTitle{肋未纪}中发言说：\BibleQuote{凡有污点或瑕疵的人，不可上前向主奉献礼物。}\BibleRef{肋 21:21}此外，在\WorkTitle{出谷纪}中，祂也规定同样的事说：\BibleQuote{凡前来接近主天主的司祭，要自洁，免得主放弃他们。}\BibleRef{出 19:22}又说：\BibleQuote{当他们前来在圣所祭台前供职时，不可在身上带有罪愆，免得死亡。}\BibleRef{出 28:43}然而，有什么罪愆更大，或有什么污点更可憎，比得上对抗基督、分散祂用血所赎买和建立的教会、忘记福音的和平与仁爱、以敌对纷争的疯狂攻击一致和睦的天主子民呢？这样的人，即使他们自己回归教会，也不能恢复和召回那些被他们诱惑、在外死亡而无共融和平安的人；他们的灵魂在审判之日，必从那些成为他们毁灭的作者和领袖之人手中被追索。因此，当他们回归时，赦免已足够；因为悖逆确实不该在信德之家得到升迁。如果我们尊重那些离开我们并与教会对抗的人，那么我们为那些良善无辜、不远离教会的人保留什么呢？\par

3. 最亲爱的兄弟，为着我们彼此间的尊敬和真诚的友爱，我们将这些事告知您的尊意；深信根据您宗教与信德的真理，那些既虔敬又真实的事必得您的认可。但我们知道，有些人不会放弃他们一旦吸收的东西，也不轻易改变他们的意图；但他们紧紧持守同僚间和平与和谐的联系，保留某些在他们中间一度采纳的独特做法。在此事上，我们既不强迫任何人，也不向任何人强加法律，因为每位主教在治理教会时，均可自由运用自己的意志，因为他将向主交账。最亲爱的兄弟，愿您时常在主内安康。\par

\SourceRecord{Translated by Robert Ernest Wallis.；Ante-Nicene Fathers；Vol. 5.；Edited by Alexander Roberts, James Donaldson, and A. Cleveland Coxe.；Buffalo, NY: Christian Literature Publishing Co.,；1886.；<http://www.newadvent.org/fathers/050671.htm>.}

% structure: p0001 source=h1 target=section reason=page-title

\section{第七十二封书信}

\Emph{致犹巴伊安：论异端者的圣洗。}\par

\Emph{提要。——西彼廉驳斥犹巴伊安附寄给他的信件，并极其审慎地收集他认为有助于辩护其主张的一切。此外，他将致努米底亚人及昆图斯的信函副本寄给犹巴伊安，或许还包括上次主教会议的决议。}\par

1. 西彼廉祝他的兄弟犹巴伊安安好。你写信给我，最亲爱的兄弟，希望我向你表达我的心意，关于我对异端者圣洗的看法；他们置身教会之外，自诩有权柄和能力行一件既不在他们权利之内也不在他们能力之内的事。我们不能认为这种圣洗有效或合法，因为在异端者中它显然是不合法的；既然我们已经在信中表达了对这件事的看法，我以便捷的方式，将相同的信函副本寄给你，即我们许多人在场时在会议上所决定的，以及后来我回复我们同僚昆图斯关于同一问题的信函。现在，当我们再次聚集时，非洲省和努米底亚省的主教共七十一位，我们再次以我们的判决确立了同一件事，决定在天主教会中只设立一种圣洗；并且凡是从那奸淫和不洁之水前来，要被救恩之水的真理洗净和圣化的人，不是由我们重洗，而是由我们施洗。\par

2. 你所描述的信件并未扰乱我们，最亲爱的兄弟，即诺瓦提安人重洗那些被他们从我们这里引诱走的人，因为教会的敌人做什么与我们毫无关系，只要我们自身持守我们的权柄以及理性与真理的坚定。因为诺瓦提安模仿猿猴——虽然它们不是人，却模仿人的行为——企图为自己夺取天主教会之权威与真理，然而他自己并不在教会内；不仅如此，他至今仍作为叛徒和敌人对抗教会。因为他知道只有一个圣洗，却将这一个圣洗据为己有，以便声称教会与他同在，并使我们成为异端者。但我们这些持守唯一教会之头和根的人知道并确切相信，在教会之外没有任何事是合法的，并且那唯一的圣洗是在我们中间，而他本人从前也曾在那里受洗，当时他还持守神圣合一的智慧和真理。但如果诺瓦提安认为，在教会内受洗的人必须在教会外——即教会之外——重洗，那么他应当从自身开始，先以一种外来的、异端的圣洗重洗自己，因为他认为在教会之后，甚至违背教会，人们应在教会之外受洗。但这算什么呢？因为诺瓦提安敢做这事，我们就必须认为我们不可做！那么呢？因为诺瓦提安也篡夺司铎宝座的尊荣，我们就应当放弃我们的宝座吗？或者因为诺瓦提安不正当地设立祭台并奉献祭献，我们就应当停止我们的祭台和祭献，以免显得与他举行相同或类似的事吗？极其虚妄和愚蠢的是，因为诺瓦提安在教会之外自称拥有真理的影像，我们就应当放弃教会的真理。\par

3. 但在我们中间，判断那些从异端者来到教会的人应当受洗，这并不是新的或突然的事情，因为许多年前，在值得纪念的阿格里皮努斯领导下，许多主教聚集在一起已经决定了此事；从那时直到今日，我们各省中成千上万的异端者已归向教会，他们既不轻视也不延迟，反而合理而喜乐地拥抱获得生命之洗礼和救恩圣洗之恩宠的机会。因为对于一个教师来说，将真实和合法的事情灌输给一个人的思想并不困难，这人既已谴责异端的邪恶，又发现了教会的真理，他为此而来，目的就是学习，并且他学习是为了生活。我们不应因我们的赞同而助长异端者的愚顽，当他们甘愿且顺从地服从真理的时候。\par

4. 诚然，既然我在你寄给我的那封信的副本中看到，它写着：\Quote{不应问是谁施洗，因为受洗者可能按他所信的获得罪的赦免}，我认为这个话题不可忽略，尤其是因为我在同一封信中还提到玛尔基翁，说：\Quote{甚至那些从他那里来的人也不需要受洗，因为他们似乎已在耶稣基督之名内受洗了}。因此，我们应当考虑那些在教会之外相信的人的信德，他们是否因着同一信德能获得任何恩宠。因为如果我们和异端者有同一信德，我们也可能有同一恩宠。如果父受苦派、人形派、华伦提努斯派、阿佩莱斯派、蛇派、玛尔基翁派以及其他为颠覆真理而存在的祸害、刀剑和毒药——这些异端——与我们承认同一圣父、同一圣子、同一圣神、同一教会，那么他们若也有同一信德，便也可能有同一圣洗。\par

5. 为了避免逐一审视所有异端、枚举它们各自的愚行或狂妄之劳（因为谈论那使人畏惧或羞于知道的事并非愉快之事），让我们暂且只审查关于\Person{马尔基昂}的情形——你们转交给我们的信函中已提及此人——看他圣洗的依据是否成立。因为主在复活后派遣门徒时，教导并指示他们应以何种方式施洗，说：\Quote{\BibleQuote{天上地下的一切权柄都交给了我。所以你们要去，使万民成为门徒，因圣父、圣子、圣神之名给他们授洗。}}祂提示了\OriginalTerm{天主圣三}{Trinity}，万民将在\OriginalTerm{圣事}{sacrament}中藉此而受洗。那么，\Person{马尔基昂}是否持守天主圣三？他是否像我们一样宣称同一位圣父、造物主？他是否认识同一位圣子、由\Person{童贞玛利亚}所生的\OriginalTerm{基督}{Christ}，祂即\OriginalTerm{圣言}{Word}成了血肉，背负了我们的\OriginalTerm{罪}{sin}，藉死亡战胜了死亡，自己首先开启了肉身的\OriginalTerm{复活}{resurrection}，并向门徒显示祂已在这同一肉身中复活？\Person{马尔基昂}的\OriginalTerm{信德}{faith}相去甚远，而且其他\OriginalTerm{异端者}{heretic}也是如此——不，他们所有的只是背信、亵渎和争辩，这与圣洁和真理为敌。那么，一个在他们中间受洗的人，怎能因他的信德而看似获得了罪的赦免和神圣怜悯的\OriginalTerm{恩宠}{grace}呢？因为他根本没有真实的信德本身。因为，如果像某些人所猜测的，人能在\OriginalTerm{教会}{Church}之外按自己的信德领受某种东西，那么他确实领受了所信的；但若他信的是虚假，就不能领受真实的；反而他所领受的是淫乱和亵渎之物，正如他所信的。\par

6. 关于这亵渎淫乱的圣洗，\Person{耶肋米亚}先知明确斥责说：\Quote{\BibleQuote{为什么迫害我的人得胜？我的创伤严重，我怎能痊愈？它实在成了我欺骗的水，没有忠诚。}} \OriginalTerm{圣神}{Holy Spirit}藉先知提到欺骗的水，没有忠诚。这欺骗而无信的水是什么？当然是那虚假地冒充圣洗外表、以虚假的幻觉挫败信德恩宠的水。但若按照歪曲的信德，人能在外面受洗并获得罪的赦免，那么按照同样的信德，他也能获得圣神；这样当他前来时，无需覆手便可获得圣神并被印封。要么他能凭信德在外面获得两者，要么在外面的人两者都未获得。\par

7. 然而，显然在哪里、由谁可以给予罪的赦免，即在圣洗中所给予的。因为首先，主将这权柄赐给了\Person{伯多禄}——祂在伯多禄身上建立了\OriginalTerm{教会}{Church}，并由此指定并显示了合一之源——即凡他在地上释放的，在天上也要被释放。复活后，祂也对宗徒们说：\Quote{\BibleQuote{就如父派遣了我，我也同样派遣你们。说了这话，就向他们嘘了一口气，并说：你们领受\OriginalTerm{圣神}{Holy Ghost}！你们赦免谁的罪，就给谁赦免；你们保留谁的罪，就给谁保留。}}因此我们看出，只有那些被安置在教会之上、立于\OriginalTerm{福音}{Gospel}律法和主的诫命之上的人，才获准施洗并给予罪的赦免；而在外面，什么也不能被束缚或释放，因为那里没有人能束缚或释放任何事物。\par

8. 最亲爱的弟兄，我们提出这一点并非没有神圣\OriginalTerm{圣经}{Scripture}的权威；我们说，万事万物都是按一定的规律和特定的诫命由神圣旨意所安排，没有人能擅自篡夺不属于自己权利和权柄的、违背主教和司铎的任何事物。因为\Person{科辣黑}、\Person{达堂}和\Person{阿彼兰}曾企图违背\Person{梅瑟}和司铎\Person{亚郎}擅自举行\OriginalTerm{祭献}{sacrifice}，他们非法胆敢做的事，没有不受惩罚的。\Person{亚郎}的儿子们也将凡火放在\OriginalTerm{祭台}{altar}上，立即在发怒的\OriginalTerm{主}{Lord}眼前被焚灭；这惩罚也临到那些藉虚假\OriginalTerm{圣洗}{baptism}引入凡水的人身上，好让神圣的报复惩罚\OriginalTerm{异端者}{heretic}，当他们违背\OriginalTerm{教会}{Church}做那唯独教会才被允许做的事时。\par

9. 至于有些人关于在\Place{撒玛黎雅}受洗之人的主张，即当\Person{伯多禄}和\Person{若望}两位宗徒来到时，只给他们覆手以使他们领受\OriginalTerm{圣神}{Holy Ghost}，却没有重新为他们施洗；最亲爱的弟兄，我们看出那处经文并不适用于现在的情形。因为在\Place{撒玛黎雅}信了的人是以真\OriginalTerm{信德}{faith}相信的，并且是在内里——在唯一的\OriginalTerm{教会}{Church}中，唯独她蒙赐予圣洗的\OriginalTerm{恩宠}{grace}并赦免\OriginalTerm{罪}{sin}——由\Person{斐理伯}执事为他们施洗，而这位执事正是同一些宗徒所派遣的。因此，因为他们已获得合法且合乎教会的圣洗，无需再受洗；只需由\Person{伯多禄}和\Person{若望}完成必要的事，即为他们祈祷、覆手，呼求\OriginalTerm{圣神}{Holy Spirit}并倾注于他们身上。这如今在我们中间也照样做：凡在教会中受洗的人，被带到教会的首领面前，藉我们的祈祷和覆手获得圣神，并以\OriginalTerm{主}{Lord}的印玺得以成全。\par

10. 最亲爱的弟兄，因此没有任何理由认为我们应该向异端者让步，以至于考虑将只授予唯一教会的那种圣洗也交给他们。好士兵的职责是保卫将军的营寨抵御叛徒和敌人。杰出领袖的职责是守护托付给他的旗帜。经上记载：\Quote{上主你的天主是忌邪的天主。}\BibleRef{申 4:24} 我们这些领受了圣神的人，理应为神圣的信德怀有忌邪之心，就如丕乃哈斯忌邪那样，既悦乐了天主，又在祂发怒、百姓丧亡时，合理地平息了祂的义怒。我们为何接纳一个不贞洁的、外邦的、与神圣合一为敌的教会为合法，而我们知道只有一位基督和祂的唯一教会呢？教会描绘出乐园的样式，在她的围墙内包括结果实的树，其中不结好果子的被砍下丢进火里。她以四条河流浇灌这些树，即四部福音，借此藉天降的洪水，分赐拯救的圣洗恩宠。谁若不在教会内，岂能从教会的泉源中取水呢？谁若自己败坏、已被定罪、被逐出乐园的泉源之外、且因永渴之干涸而枯萎失败，又怎能将乐园中那些有益而拯救的饮料赐给别人呢？\par

11. 主大声呼喊说：\Quote{凡口渴的，都应来喝从祂怀中涌出的活水江河。}口渴的人该到哪里去呢？到异端者那里去吗？那里根本没有活水的泉源和江河；还是到那唯一的、建立在由主的声音领受钥匙的那一位身上的教会？唯有她持有并拥有她新郎及主的一切权能。我们在她内主持；我们为她尊荣和合一而战斗；我们以忠诚的热忱捍卫她的恩宠，也捍卫她的光荣。我们得天主许可，灌溉口渴的天主子民；我们守护活水的泉源边界。因此，如果我们持有我们拥有的权利，如果我们承认合一的圣事，为何我们被视为真理的犯法者？为何我们被判断为合一的背叛者？教会那忠实、拯救而圣洁的水不能受败坏和掺杂，正如教会本身也是无玷、贞洁而端庄的。如果异端者属于教会，在教会内，他们就可以使用她的圣洗和其他拯救的恩惠。但如果他们不在教会内，甚至反对教会，他们怎能用教会的圣洗施洗呢？\par

12. 因为这不是一件小事或无足轻重的事，当我们承认异端者的圣洗时，就等于将重要的东西让给了他们；因为由此产生了信仰的整个根源和通往永生望德的拯救之门，以及天主为洁净和使祂的仆人复活的屈尊就卑。因为如果有人在异端者中能受圣洗，他当然也能获得罪的赦免。如果他获得了罪的赦免，他就被圣化了。如果他被圣化了，他也成了天主的圣殿。我问：是哪个天主的圣殿？如果是创造者的，他不能，因为他没有相信祂。如果是基督的，他不能成为祂的圣殿，因为他否认基督是天主。如果是圣神的，因为三者是一体，那与圣子或圣父为敌的人，圣神怎能与他平安相处呢？\par

13. 因此，有些被理由挫败的人徒然以习俗反对我们，好像习俗大于真理；或者在属灵的事上，圣神所启示的更好的东西，反而不去追求。因为由于单纯而犯错的人可以蒙赦免，正如蒙福的宗徒保禄论自己所说：\Quote{我先前曾是亵渎者、迫害者和施暴者，但我蒙了怜悯，因为我是在无知中做的。}\BibleRef{弟前 1:13} 但在领受了灵感和启示之后，明明知道却仍故意坚持他犯错时的作法，就因无知而不可赦免地犯罪。因为他被理由战胜时，还带着某种傲慢和固执来抵抗。也不要有人说：\Quote{我们遵循我们从宗徒领受的。}但宗徒只传授了唯一的教会和唯一的圣洗，而它只有在同一教会内才被设立。我们找不到任何人，当他由异端者付洗后，宗徒在同一个圣洗中接纳了他，并以这样的方式与他共融，以致宗徒似乎认可了异端者的圣洗。\par

14. 至于有些人所说的话，似乎有利于异端者，就是宗徒保禄宣称：\Quote{无论如何，无论是假意还是真诚，只要基督被传扬就好了。}\BibleRef{斐 1:18} 我们发现，这对那些支持和赞扬异端者的人也毫无益处。因为保禄在书信中并非谈论异端者，也不是他们的圣洗，以致可以证明有任何与此相关的事被提及。他是在谈论弟兄们，无论他们是行为不检、违反教会纪律，还是以敬畏天主之心保持福音真理。他说他们中有些人以坚定和勇敢宣讲天主的话，有些人则出于嫉妒和纷争而行；有些人对他保持慈爱的善意，有些人则怀有恶意的纷争之灵；但他仍耐心忍受这一切，只要，无论是出于真理还是假意，保禄所宣讲的基督之名能传到许多人那里；并且话语的播种，那时还是崭新而不规则的，能通过说话者的宣讲而增长。此外，在教会内谈论基督之名是一回事，在教会外、反对教会的人以基督之名付洗又是另一回事。因此，那些支持异端者的人，不要提出保禄关于弟兄们所说的话，而应证明他认为有什么可以迁就异端者，或者他认可他们的信仰或圣洗，或者他命定那些背信和亵渎的人能在教会外获得罪的赦免。\par

15. 然而，如果我们思考宗徒们对异端者的看法，就会发现他们在所有的书信中都咒骂并憎恶异端者的亵渎邪恶。因为当他们说\BibleQuote{他们的言语如同毒疮蔓延}\BibleRef{弟后 2:17}，这样的言语如何能给予罪过的赦免，它如同毒疮蔓延到听者的耳中？当他们说\BibleQuote{义和不义没有相通，光明和黑暗没有共融}\BibleRef{格后 6:14}，黑暗如何能照亮，不义如何能成义？当他们说\BibleQuote{他们不是出于天主，而是出于假基督的灵}\BibleRef{若一 4:3}，这些天主的敌人、内心被假基督的灵占据的人，如何能施行属灵和神圣的事？因此，如果我们放下人性争论的错误，以真诚和宗教般的信德回归福音的权威和传统，我们就会明白，那些分散并攻击基督的教会、被基督亲自称为仇敌、被祂的宗徒称为假基督的人，在授予教会和救恩的恩宠方面无能为力。\par

16. 再者，任何人为了欺骗基督真理，都不应以基督之名反对我们，说：\Quote{凡在任何地方、以任何方式因耶稣基督之名受洗的，都获得了圣洗的恩宠}——而基督亲自说：\BibleQuote{不是凡向我说\InnerQuote{主啊，主啊}的人，就能进入天国}\BibleRef{玛 7:21}。祂又预先警告并教导说，不要轻易被假先知和假基督以祂的名迷惑。祂说：\BibleQuote{将有许多人假冒我的名来，说：我是基督，并要迷惑许多人。}随后祂又补充说：\BibleQuote{你们要当心；看，我已经预先告诉了你们一切}\BibleRef{玛 24:5}。由此可见，并非所有以基督之名夸耀的事都当立即接受和采纳，只有那些在基督的真理中所行的事才当接受。\par

17. 因为在福音和宗徒书信中，基督之名被引用于罪过的赦免，但这并非意味着圣子可以离弃圣父或反对圣父而对任何人有益；而是为了向那些夸耀自己拥有圣父的犹太人表明，除非他们相信祂所派遣的圣子，否则圣父对他们毫无益处。因为认识天主圣父造物主的人，也应当认识基督圣子，免得他们因只承认圣父而不承认祂的圣子便自夸自耀；圣子也曾说：\BibleQuote{除非经过我，谁也不能到父那里去}\BibleRef{若 14:6}。然而，同一圣子指出，认识两者才能得救，祂说：\BibleQuote{永生就是：认识你，唯一的真天主，和你所派遣的耶稣基督}\BibleRef{若 17:3}。因此，既然根据基督自己的宣讲和见证，必须先认识派遣的父，然后认识被派遣的基督，并且除非同时认识两者，否则没有得救的希望；那么，当天主圣父不被认识，甚至还被亵渎时，那些在异端者中所谓因基督之名受洗的人，怎么能被视为获得了罪过的赦免呢？因为宗徒时代犹太人的情况是一回事，而外邦人的处境是另一回事。前者既然已经领受了法律和梅瑟的最古老的圣洗，就还应当因耶稣基督之名受洗，正如伯多禄在\BibleBook{宗徒}{大事录}中对他们所说的：\BibleQuote{你们悔改，并因主耶稣基督之名受洗，好使你们的罪得赦免，并领受圣神的恩惠。因为这恩许是为你们和你们的子女，以及所有远方的人，凡是上主我们的天主所召叫的}\BibleRef{宗 2:38}。伯多禄提及耶稣基督，并非要省略圣父，而是使圣子也与圣父相连。\par

18. 最后，当宗徒们复活后由主派遣到外邦人那里时，他们奉命\BibleQuote{因父、及子、及圣神之名}为外邦人施洗。那么，有些人怎能说，一个在教会之外、甚至反对教会、只凭耶稣基督之名、在任何地方、以任何方式受洗的外邦人，可以获得罪过的赦免？而基督亲自命令外邦人要在完整而合一的天主圣三内受洗。除非否定基督的人被基督否定，但否定基督所承认的圣父的人却不被否定；而亵渎基督称为\Quote{主}和\Quote{天主}的那一位的人，反倒被基督赏赐并获得罪过的赦免和圣洗的圣化！然而，那否定天主造物主、即基督之父的人，如何在洗礼中获得罪过的赦免？因为基督正是从同一圣父那里领受了那使我们受洗和圣化的权能；圣父是基督称为比自己\Quote{更大}的那一位，是基督愿意受其光荣的那一位，是基督顺从至死、饮杯服从其旨意的那一位。因此，想要坚持并断言，那亵渎并严重得罪基督之父、主和天主的人，能因基督之名获得罪过的赦免，这岂不是与亵渎的异端者同流合污吗？此外，那又是什么，何等情形：否认天主圣子的人就没有圣父，而否认圣父的人却被认为拥有圣子？尽管圣子亲自作证说：\BibleQuote{除非蒙我父恩赐，谁也不能到我这里来}\BibleRef{若 6:65}。由此可见，在洗礼中不能从圣子那里获得任何罪过的赦免，除非显然圣父已经赐予。尤其因为祂又重复说：\BibleQuote{凡不是我天父所栽种的植物，必要连根拔除}\BibleRef{玛 15:13}。\par

19. 但如果基督的门徒不愿从基督学习对圣父的名号应有的敬礼和尊荣，至少让他们从世俗和尘世的例子学习，并知道基督曾带着最严厉的斥责声明：\Quote{今世之子应付自己的世代，比光明之子更为精明。}\BibleRef{路 16:8} 在我们现世中，若有人冒犯了某人的父亲，若因恶意和悖逆，以恶毒的舌头伤害了他的名誉和尊荣，儿子便会愤怒、恼怒，并尽己所能为父亲所受的伤害报仇。你以为基督会容忍不敬和亵渎祂圣父的人，并在圣洗中赦免他们的罪，而这些人显然在受洗时仍对圣父口出恶言，以亵渎之舌不断作恶吗？一个基督徒，一个天主的仆人，能这样想、这样信，或这样宣讲吗？神圣法律的诫命说：\Quote{应孝敬你的父亲和你的母亲}\BibleRef{出 20:12}，又会变成什么？如果那命令人在人中应受尊敬的\Quote{父亲}之名，在天主身上被违反而不受惩罚，那么基督自己在福音中规定的\Quote{咒骂父亲或母亲的，应处死刑}\BibleRef{玛 15:4}，又该如何？祂命令那咒骂肉身父母的人应受惩罚和处死，却赐予生命给那些辱骂天上和属灵的圣父、敌视教会——他们的母亲——的人？有人竟主张一件可憎可恶的事：那威胁亵渎圣神者将永远有罪的主，竟屈尊以救恩的圣洗圣化那些亵渎天主圣父的人。现在，那些认为必须与未受洗而来到教会的人共融的人，没有考虑到，他们接纳未受洗者时，便成了他人——甚至是永恒之罪——的共犯，因为那些人除非藉圣洗，否则无法除去他们亵渎的罪。\par

20. 此外，多么虚妄和乖谬的事：当异端者自己已经弃绝和抛弃了他们先前所处的错误或邪恶，承认教会的真理时，我们却要阉割这同一真理的权利和圣事，对那些来到我们这里并悔改的人说，他们已获得罪过的赦免，因为他们承认自己犯了罪，并为此来寻求教会的宽恕！所以，最亲爱的弟兄，我们应当坚定持守天主教会的信德和真理，教导并藉一切福音的训诫阐述天主上智安排和合一的计划。\par

21. 当一个人在众人前宣认基督，并用自己的血受洗时，圣洗的能力难道能比宣认和苦难更大或更有效吗？然而，即使是这样的圣洗，对异端者也无益处，虽然他宣认了基督，并在教会外被杀，除非异端者的庇护者和辩护者宣称，那些在虚假的基督宣认中被杀的异端者是殉道者，并将殉道的荣耀和冠冕归给他们，这与宗徒的见证相反，宗徒说：即使他们被焚被杀，也毫无益处。\BibleRef{格前 13:3} 但如果连公开宣认和流血的圣洗都不能使异端者得救恩——因为在教会之外没有救恩——那么，若他在一个隐藏处和贼窝里，被奸淫之水的传染所玷污，不仅没有除去旧罪，反而堆积了更新更大的罪，这圣洗对他又有何益呢？因此，圣洗不能为我们和异端者所共有，因为与他们既没有共同的天主圣父，也没有共同的基督圣子，没有共同的圣神，没有共同的信德，甚至没有共同的教会本身。所以，那些从异端来到教会的人理应受洗，这样，他们藉神圣教会的合法、真实、唯一的圣洗，通过神圣的再生，为天主的国而预备，就能由两件圣事而生，因为经上记载：\Quote{除非由水和圣神而生，不能进天主的国。}\par

22. 关于这一点，有些人仿佛能凭人的推理使福音宣告的真理失效，便以慕道者的情况来反对我们：倘若他们中任何一人在教会中受洗之前，因宣认主名而被捕并被杀，他是否会失去救恩的希望和宣认的赏报，因为他先前没有由水重生？让这些异端者的帮助者和支持者知道：首先，这些慕道者持守着教会纯正的信德和真理，从神圣的营垒出发，以对圣父、基督和圣神圆满而真诚的承认，与魔鬼作战；其次，那些用最光荣、最伟大的血洗而受洗的人，确实没有被剥夺圣洗的圣事，主也曾论及这血洗，说祂有\Quote{另一种当受的圣洗}。但同一主在福音中宣告，那些用自己的血受洗、并通过苦难而成圣的人，得以成全并获得神圣应许的恩宠，正如祂在受难时对相信并宣认的强盗所说，应许他应与自己同在乐园里。因此，我们这些被立於信德和真理之上的人，不应欺骗和误导那些来到信德和真理、悔改并请求赦罪的人，而应教导他们，由我们矫正，藉天上的纪律为天国而改造。\par

23. 但有人说：\Quote{那么，那些从前从异端来到教会、未受圣洗就被接纳的人，将会怎样呢？}主能以其仁慈给予宽赦，并不断开那些因纯朴而被接纳进入教会、并在教会中安眠的人脱离祂教会的恩赐。然而，不能因为曾经有过错误，就必须永远有错误；因为明智而敬畏天主的人，在真理被揭示并认识时，欣然且毫不迟延地服从真理，要比顽固而执拗地为异端者同弟兄和同僚司铎争辩更为适宜。\par

24. 也不要有人认为，因为向他们提出圣洗，异端者就会因第二次洗礼之名而感到冒犯，因而被阻止前来教会；不，恰恰相反，正是由于这个原因，他们更因向他们展示并证明的真理的见证，而被驱使到不得不前来的地步。因为如果他们看到根据我们的判断和决定，他们在那里所受的洗礼被视为公正而合法的，他们就会认为自己也在公正而合法地拥有教会和教会的其他恩赐；那么，既然他们有了圣洗，似乎也就有了其余的一切，就没有理由到我们这里来了。但是，更进一步，当他们知道外面没有圣洗，教会之外不能给予罪的赦免时，他们就会更加迫切而欣然地投奔我们，恳求教会我们母亲的恩赐和恩惠，确信他们除非先来到教会的真理中，否则绝不能获得天主恩宠的真正应许。当异端者从我们这里得知，连保禄也是在已领受了若翰的洗礼之后才受洗的（就如我们在\WorkTitle{宗徒大事录}中所读到的），他们也不会拒绝在我们中间领受教会合法而真正的圣洗。\par

25. 如今，我们当中有些人断言异端者的洗礼占据同等的地位，并且仿佛出于对重洗的某种厌恶，认为在天主的仇敌之后施洗是不合法的。\Emph{而这事}，尽管我们发现那些曾被若翰施洗的人又受了洗：若翰，被尊为先知中最大的；若翰，甚至在母胎中就充满了天主的恩宠；他以厄里亚的精神和能力被扶持；他不是主的对头，而是祂的前驱和宣告者；他不仅用言语预言了我们的主，甚至以双眼指给祂看；他为自己受洗的基督施洗。但如果因此一个异端者能获得洗礼的权利，因为他先行施洗，那么洗礼就不属于拥有它的人，而属于夺取它的人。既然洗礼与教会绝不可彼此分离、分开，那么谁能先夺得了洗礼，也就同样夺得了教会；并且你在他眼中就变成了异端者，因为你被抢先而成了后来的，并且由于让步和退让，放弃了你已领受的权利。然而，在神圣事务中，一个人放弃自己的权利和能力是多么危险，圣经在\WorkTitle{创世纪}中宣告：厄撒乌由此失去了长子的名分，并且之后无法重新获得他一度放弃的。\par

26. 最亲爱的弟兄，我照我的能力，将这些事简略地写信给你，不加规范于任何人，也不预断任何人，以免阻碍任何一位主教做他认为适当的事，并自由运用他的判断。我们尽其所能，不为了异端者而与我们保持神圣和谐与主平安的同僚及主教弟兄们争辩；尤其宗徒说：\Quote{但如果有人想强辩，我们没有这样的习惯，天主的教会也没有。}我们以忍耐和温良持守精神的仁爱、团体的荣誉、信德的联系和司铎间的和谐。为此，我们还尽我们微薄的能力，蒙主的允准和启迪，撰写了一篇关于\WorkTitle{忍耐的益处}的论文，为了我们彼此的爱而传送给你。最亲爱的弟兄，祝你永远衷心安康。\par

\SourceRecord{Translated by Robert Ernest Wallis.；Ante-Nicene Fathers；Vol. 5.；Edited by Alexander Roberts, James Donaldson, and A. Cleveland Coxe.；Buffalo, NY: Christian Literature Publishing Co.,；1886.；<http://www.newadvent.org/fathers/050672.htm>.}

% structure: p0001 source=h1 target=section reason=page-title

\section{第七十三封书信}

\Emph{致庞培，驳斥斯蒂芬关于异端者圣洗的书信。}\par

\Emph{提要：——本信之要旨见于圣奥斯定的\WorkTitle{驳多纳徒派}卷五第二十三章。他说：\Quote{西彼廉又就同一事写信给庞培，明示斯蒂芬（如我们所知，他当时是罗马教会的主教）不仅在这些点上与他意见不合，甚至写信并责令反对他。}}\par

1. 西彼廉致其弟兄庞培，问安。最亲爱的弟兄，虽然我在寄给你的信件中已充分涵盖了关于异端者圣洗所应说的一切，但因你希望斯蒂芬弟兄对我信件的回复也为你知道，我便将他的回信副本寄给你；阅读之后，你将越发看出他的错误：他试图维护异端者反对基督徒和天主的教会的立场。因为在他傲慢地假定、或与事情无关、或自相矛盾（他笨拙而无预见地写下的）的其他事中，他又加上了这句话：\Quote{因此，若有人从任何异端来到你们这里，不要创新（或做什么）未曾传下来的事，即：为他悔改而覆手；因为异端者本身，按其本有特性，并不为彼此前来的人施圣洗，只接纳他们进入共融。}\par

2. 他禁止从任何异端来的人在教会中受圣洗；也就是说，他判定所有异端者的圣洗都是正当合法的。而尽管各特殊异端有其特殊圣洗和不同罪恶，他却与所有异端者的圣洗共融，将所有人的罪恶收集起来，堆积在自己的怀中。他责令不要创新，只应持守传下来的；好像那保持合一、为唯一教会主张唯一圣洗的是创新者，而不是那忘记合一、采用亵渎洗涤的谎言和传染的。他说：\Quote{不要创新，\Emph{不要持守}，只应持守传下来的。} 这传统从何而来？是从主和福音的权威而来，还是从宗徒的命令和书信而来？因为那些写下来的事必须遵行，天主以向农的儿子若苏厄所说的话作证并劝诫说：\BibleQuote{这法律书不可离开你的口，但你要日夜默思，好能谨慎遵行其中所写的一切话。}\BibleRef{苏 1:8} 主派遣宗徒时，也命令万民应受圣洗，并遵守祂所吩咐的一切。因此，若福音中规定，或宗徒书信或\BibleBook{}{宗徒大事录}中载明，从任何异端来的人不应受圣洗，只应覆手使其悔改，那么就当遵守这神圣而传统的规条。但若异端者到处只被称为仇敌和假基督，被视为应躲避、自趋败坏、自定己罪的人，那么我们为何不该认为他们当受我们谴责？因为宗徒的见证清楚显示他们自定己罪。所以无人应当诽谤宗徒，仿佛他们曾认同异端者的圣洗，或未以教会的圣洗就与他们共融，因为宗徒对异端者写了这样的话。况且，那时异端更可怕的祸患尚未爆发；本都的马尔西翁尚未从本都出现，其师傅切尔东来到罗马——当时希吉诺仍是主教，他是该城第九位主教——马尔西翁随从他，并以更大的无耻为他的罪恶增添其他加剧，更大胆地开始亵渎天主圣父、造物主，并以亵渎的武器武装那以更大恶毒和决心反叛教会的异端疯狂。\par

3. 但若显然此后异端变得更多更坏；若过去从未规定或写道，异端者只应覆手悔改，然后与之共融；若只有一种圣洗，是在我们中、在内里，并唯由天主的恩赐授予教会，那么有何顽固或傲慢，竟将人的传统置于天主的诫命之上，而不理会当天主藉依撒意亚先知呼喊说：\BibleQuote{这民族用嘴唇尊敬我，他们的心却远离我。但他们恭敬我也是徒然，因为他们教导人的教训和命令。}\BibleRef{依 29:13} 时，祂是如何愤慨和发怒？主在福音中也同样斥责并谴责说：\BibleQuote{你们废弃天主的诫命，为遵守你们的传统。}\BibleRef{谷 7:13} 有福的宗徒保禄自己也记念这诫命而警告教导说：\BibleQuote{若有人讲异端道理，不赞同我们主耶稣基督健全的话和祂的道理，他是骄傲无知，你们要避开这样的人。}\BibleRef{弟前 6:3}\par

4. 诚然，我们的兄弟\Person{斯德望}的教导为我们提出了一条卓越而合法的传统，它可为我们提供合适的权威！因为在他书信的同一处，他补充并继续说道：\Quote{那些特别属于异端的人，并不为彼此之间来投奔者施洗，而仅接纳他们共融。} 天主教会、基督的净配竟已发展到如此邪恶的地步，以至于效法异端者的榜样；为了庆祝天上的圣事，光明竟向黑暗借用其纪律，基督徒竟做反基督者所做的事。然而，那是怎样的灵魂盲目，怎样的信德堕落——竟拒绝承认来自天主圣父、以及来自主\Person{耶稣基督}、我们的天主的传统所赋予的合一！因为如果教会与异端者无份，因为教会是唯一的，不能分裂；如果圣神也不在那里，因为祂是唯一的，不能存在于亵圣者和外人在内之中；那么，同样，存在于同一合一中的圣洗，也不能存在于异端者中，因为它既不能与教会分离，也不能与圣神分离。\par

5. 或者，若他们将圣洗的效果归于圣名的尊威，以致凡在任何地方、以任何方式奉\Person{耶稣基督}之名受洗者，都被视为得以更新与圣化；那么，为何不在他们中间为受洗者行覆手礼以领受圣神呢？为何同一圣名的同一尊威在覆手礼上无效，而他们却主张它在圣洗的圣化中有效？因为如果任何生于教会之外的人能成为天主的殿，为何圣神不能倾注在这殿上呢？因为那已在圣洗中除去罪恶、被圣化、并属灵地更新为新造的人，已适合领受圣神；正如宗徒所说：\Quote{凡你们因基督受洗的，都穿上了基督。}\BibleRef{迦 3:27} 凡能在异端者中受洗而穿上基督的，更能领受基督所派遣的圣神。否则，被派遣者将大于派遣者；如此，在外受洗者或能开始穿上基督，却不能领受圣神，仿佛基督可以没有圣神而被穿上，或圣神可与基督分离。再者，说第二重生虽属灵——我们藉重生之洗在基督内重生——却能在异端者中属灵地重生（他们声称那里没有圣神），这是愚蠢的。因为单凭水不能洗净罪恶、圣化人，除非他也有圣神。因此，他们必须承认：凡他们主张有圣洗之处，圣神也在那里；否则，没有圣神之处就没有圣洗，因为没有圣神便不能有圣洗。\par

6. 然而，断言并争辩说未在教会内出生的人能成为天主的子女，这是何等的事！因为蒙福的宗徒阐明并证明：圣洗是旧人死去、新人生成之处，他说：\Quote{祂藉着重生的洗，拯救了我们。}\BibleRef{铎 3:5} 但若重生在于洗，即圣洗，那么异端（不是基督的净配）怎能藉基督为天主生子呢？因为只有教会与基督结合、联合，才属灵地生子；同一宗徒又说：\Quote{基督爱了教会，并为她舍命，为要圣化她，用水洗洁净她。}\BibleRef{弗 5:25} 那么，如果她是蒙爱的和配偶，唯独她由基督圣化、唯独她由祂的洗洁净，那么显然，异端——既非基督的配偶，也不能被祂的洗所洁净或圣化——不能为天主生子。\par

7. 此外，人不是因领受圣神时的覆手礼而生，而是在圣洗中出生，这样，既然已经出生，他就可以领受圣神，正如在最初的人亚当身上所发生的。因为天主先塑造了他，然后向他鼻孔中吹入生命的气息。因为圣神不能被领受，除非领受者先存在。但既然基督徒的出生是在圣洗中，而圣洗的重生与圣化唯独属于基督的配偶（她才能属灵地怀孕并为天主生子），那么，那不是教会之子的人，在何处、由谁、向谁出生，才能以天主为父，在此之前他已有教会为母？然而，既然任何异端，同样任何分裂，由于在外，都不能拥有拯救性圣洗的圣化，为何我们的兄弟\Person{斯德望}的苦毒固执竟发展到如此地步，以致主张从\Person{马尔西翁}、甚至\Person{瓦伦廷}、\Person{阿培肋斯}以及其他亵渎天主圣父者的圣洗中，能生出天主的子女；并说在亵渎圣父及主\Person{基督}、天主之名之处，奉\Person{耶稣基督}之名可赦免罪过？\par

8. 在此处，最亲爱的弟兄，为了我们所履行的司祭职务的信德与宗教，我们必须考虑：在审判之日，一位天主的司铎，若他维护、赞同并默许亵渎者的圣洗，岂能给出令人满意的解释？主曾警告说：\Quote{现在你们司祭，这条诫命是为你们的：如果你们不听从，如果不放在心上，将光荣归于我的名——万军的上主说——我必使诅咒临于你们，且诅咒你们的祝福。}与马尔西翁的圣洗相通的人，是光荣天主吗？判断那些亵渎天主者之间能得赦罪的人，是光荣天主吗？声称从奸夫淫妇而生、在教会之外的人能成为天主的子女，是光荣天主吗？不持守源于神圣法律的合一与真理，却维护反对教会的异端，是光荣天主吗？作为异端者的朋友、基督徒的敌人，认为支持基督真理与教会合一的天主司铎应被逐出教会，是光荣天主吗？若这样就是光荣天主，若天主的敬畏与纪律就这样被祂的敬拜者与司铎所保存，那我们就放下武器吧！我们自甘被掳，将福音的任命、基督的委托、天主的威严交给魔鬼吧！让神圣争战的圣事松懈，让天军营垒的旗帜被出卖，让教会屈服于异端，光明向黑暗，信德向背信，望德向绝望，理智向谬误，不朽向死亡，爱德向仇恨，真理向虚假，基督向敌基督投降吧！这正该如此，异端与分裂日益频繁地出现，更丰饶地滋长，以蛇发之势向天主的教会喷吐毒液的毒素；而某些人的辩护，既给予它们权威又提供支持；他们的圣洗被辩护，信德、真理被出卖；在教会之外反对教会的事，竟在教会内部被辩护。\par

9. 但是，至爱的弟兄，若我们心中有对天主的敬畏，若信德的持守得胜，若我们遵守基督的诫命，若我们捍卫祂净配的无玷无损的圣洁，若主的话存留于我们的思想与心中，当祂说：\Quote{人子来临时，你认为祂会在世上找到信德吗？}\BibleRef{路 18:8}那么，既然我们是天主忠信的士兵，为天主的信德与纯正宗教而战，就让我们以忠勇之姿守护天主托付给我们的营垒。也不应让在某些人中习以为常的惯例阻碍真理得胜并征服；因为无真理的惯例乃是错误的古老。因此，让我们抛弃错误，追随真理，知道在厄斯德拉书中也是真理得胜，如经上所载：\Quote{真理永存并日益强大，直到永远，且永远活着并得胜。在她面前没有偏袒或区别；她只做正义之事；在她的审判中没有不义，却有万世的权能、国度、威严与能力。愿上主真理的天主受赞美！}这真理基督在祂的福音中启示给我们，并说：\Quote{我是真理。}\BibleRef{若 14:6}所以，如果我们是在基督内，基督也在我们内，如果我们居于真理中，真理也居于我们内，就让我们坚守那些真实的事。\par

10. 但由于对傲慢与执拗的喜爱，人宁愿坚持自己邪恶虚假的立场，也不赞同属于他人的正确真实的事。有鉴于此，蒙福的宗徒保禄写信给弟茂德，警告他主教不可\Quote{好争辩，也不可好胜，而应温良、善于教导。}善于教导的人，是那温良而耐心学习的人。因为主教不仅应教导，也应学习；因为每日通过学习而长进的人，也能更好地教导。同样的事，宗徒保禄也教导说：\Quote{如果一位旁边坐着的得到了什么更好的启示，那先发言的就应静默。}\BibleRef{格前 14:30}但有捷径使虔敬而纯朴的心灵既可弃绝错误，又可发现并引出真理。因为如果我们回到神圣传统的源头与泉源，人的错误就会止息；看清了天上圣事的理由，凡在黑暗的阴云与阴影中隐藏的，就光照在真理之中。如果一条供水渠道，从前丰沛自由地流水，突然中断，我们岂不是要回到泉源，查明缺水的缘故：是由于源头的水源干涸，还是虽然从那里自由涌出充沛，却在中间途中断绝？如此，若由于渠道中断或渗漏的毛病，致使常流之水不能持续不断地涌流，那么渠道被修补加固后，积蓄的水就能以同样丰裕和充足供应城市的用度与饮用，正如它从泉源流出一样？现在天主的司铎也应这样做，如果他们愿遵守神圣的诫命：若在任何方面真理有所动摇或踌躇，我们就应回到我们的本源与主那里，回到福音与传统；那作为我们行动根据的，正是我们的秩序与起源所从出的根源。\par

11. 因为曾传授给我们：只有一个天主，一个基督，一个望德，一个信德，一个教会，一个圣洗，这圣洗惟独在独一的教会中订立；任何人如偏离这合一，就必定与异端者同流；当他支持异端者反对教会时，就是抨击神圣传统的圣事。这合一的圣事我们在《雅歌》中也看到，由基督所说：\Quote{我的妹妹，我的新娘，是关闭的花园，是封锁的泉源，是活水的井，是结有苹果的果园。} \BibleRef{歌 4:12} 但如果祂的教会是关闭的花园、封锁的泉源，那不在教会里的人怎能进入那花园，或饮其泉源呢？此外，\Person{伯多禄}本人也在指明并维护这合一时，命令并警告我们：除非藉着独一教会的唯一圣洗，我们不能得救。他说：\Quote{在\Person{诺厄}的方舟里，只有八个人藉着水得了救，同样，圣洗也要这样拯救你们。} \BibleRef{伯前 3:20} 他在何等简短而属神的概要中陈述了合一的圣事！因为，正如在那以水洗除世界古老罪污的洗礼中，凡不在\Person{诺厄}方舟内的人不能藉水得救；同样，凡未在按照独一方舟的圣事而建于主的合一中的教会内受洗的人，也不能被认为藉圣洗得救。\par

12. 因此，最亲爱的兄弟，在探索并看到了真理之后，我们遵守并持守：凡从任何异端归向教会的人，都必须由教会唯一合法的圣洗受洗，但那些先前已经在教会内受洗、而后转到异端者那里的人除外。因为这些人悔改返回时，只须藉覆手礼接纳，并由牧人领回他们迷失的羊栈。最亲爱的兄弟，祝你永远衷心安康。\par

\SourceRecord{Translated by Robert Ernest Wallis.；Ante-Nicene Fathers；Vol. 5.；Edited by Alexander Roberts, James Donaldson, and A. Cleveland Coxe.；Buffalo, NY: Christian Literature Publishing Co.,；1886.；<http://www.newadvent.org/fathers/050673.htm>.}

% structure: p0001 source=h1 target=section reason=page-title

\section{第七十四封书信}

\Emph{\Person{斐米利安}，\Place{加帕多家的凯撒勒雅}主教，致\Person{西彼廉}，驳斥\Person{斯德望}的书信。公元256年。}\par

\Emph{论点——本信论点与前信完全相同，但写得更多几分激昂与尖刻，超出主教应有；主要理由，正如可以猜测的，是因为\Person{斯德望}也曾写了一封信给\Person{斐米利安}、\Person{赫勒努斯}及其他该地区的某些主教。}\par

1. 致主内的弟兄\Person{西彼廉}，问候。我们藉着我们可爱的执事\Person{若加提安}，收到了你寄给我们、写给我们的信，至爱的弟兄；我们向主献上最大的感谢，因为那在身体上彼此分离的我们，竟如此在精神上合一，仿佛我们不仅同住一地，更是共居一室。这也正是我们适宜说的，因为天主的属灵房屋确实只有一个。先知说：\BibleQuote{到末日，上主的山必要矗立，天主的殿宇必要高耸于群山之上。}\BibleRef{依 2:2}来到这房屋的人，都喜乐地合一，正如圣咏中向主所求的：一生居住在上主的殿宇中。因此，在另一处也显明，圣者们对聚会有极大而渴望的爱情。他说：\BibleQuote{看，兄弟们同居共处，多么美好，多么快乐！}\par

2. 因为合一、和平与和谐不仅给相信并认识真理的人带来极大的喜悦，也给天上的天使们自己带来喜悦；天主圣言说，当一个罪人悔改并回到合一的联系时，天使们就欢喜。但若天使们自己是在天上生活，这话肯定不会对他们说，除非他们自己与我们合一，因我们的合一而喜乐；正如另一方面，当他们看见某些人分歧的心意和分裂的意志时，他们肯定悲伤，仿佛这些人不仅不是一起呼求同一个天主，而且彼此分离、分裂，不能有共同的交往或谈论。只是我们在此要感谢\Person{斯德望}，因为现在藉着他的不仁慈，我们得以领受你信德与智慧的证明。但虽然我们因\Person{斯德望}而领受了这恩惠的好处，\Person{斯德望}却一定不配得到仁慈与感谢。因为\Person{犹达斯}也不能因他的背信弃义与诡诈，就是他以邪恶对待救主的方式，被认为配得如此大的好处，好像藉着他，世界及外邦人民藉着主的受难而得救。\par

3. 但让\Person{斯德望}所做的这些事暂时搁过一边，免得我们记住他的胆大妄为与骄傲时，因他所行的恶事而给自己带来更持久的悲伤。我们知道，关于现今讨论的这个问题，你是按照真理的准则和基督的智慧所定的，我们就大大欢喜，感谢天主，因为在相隔如此遥远的弟兄们身上，我们找到了与己如此一致的信仰与真理。因为天主的恩宠是大能的，能将那些看来被相当大的地域间隔所分隔的事物，也在爱德与合一的联系中结合起来，正如古时天主的德能也将\Person{厄则克耳}和\Person{达尼尔}——他们虽是较晚的时代，且与最早的\Person{约伯}和\Person{诺厄}相隔很长的时间——结合在同心合意的联系中；这样，虽然他们被漫长的时期隔开，却因天主启示而感到同样的真理。现在我们在你身上也看到这一点：你虽与我们相隔最辽阔的地区，却证明自己与我们心相印、灵相通。这一切都源自天主的合一。因为居住在咱们内的主是同一的，祂处处在合一的联系中联合并连接祂的子民，因此他们的声音传遍普世——他们被主差遣，在合一的精神中迅速奔跑；而另一方面，有些人虽然身体上很近且联结在一起，若在精神与心意上相异，则毫无益处，因为灵魂若使自己与天主的合一分离，便根本不能合一。经上说：\BibleQuote{因为，请看，远离你的人必致灭亡。}但这样的人将按他们的所作所为受天主的审判，因为他们离开了那位为合一向父祈祷的主的话，祂说：\BibleQuote{父啊，愿他们在我们内合而为一，如同你在我内，我在你内一样。}\BibleRef{若 17:21}\par

4. 我们接受你所写的，如同我们自己的；我们并非草草阅读，而是再三诵读，记在心里。这并不妨碍救恩的益处，无论是为确认真理而重复同样的话，还是为累积证据而添加一些内容。但若我们添加了什么，并非因为你的话有所不足；而是因为天主的话语超越人性，灵魂不能领会或把握那完整而圆满的话语，所以先知的数量如此众多，为使多样性的天主智慧通过许多人分施。因此，在先知讲预言时，若启示给予第二人，那先发言的应静默。为此，我们中间每年都有必要聚会，长老与主教们共同安排委托给我们照管的事务；若有更重大的事，可由共同商议来指导。此外，\Emph{我们这样做}，是为跌倒的弟兄以及那些在得救的洗礼后受到魔鬼伤害的人，通过悔改寻求某种补救：并非他们从我们这里获得罪的赦免，而是藉着我们，他们得以转化，认识自己的罪，并被促使向主做更圆满的补赎。\par

5. 但因你派来的使者急于返回你处，且冬季紧迫，我们只能就你的来信尽我们所及予以答复。的确，关于\Person{斯德望}所说的，仿佛宗徒们禁止为从异端而来的人施洗，并将此作为传统交付给他们的后继者，你已极为充分地答复说：没有人会愚蠢到相信宗徒们交付了此事，因为众所周知，这些可憎可恶的异端本身是后来才出现的；甚至\Person{马西昂}——\Person{塞尔多}的门徒——是在宗徒之后，且相隔漫长岁月之后，才引入了他那亵渎天主的传统。\Person{阿佩莱斯}，也赞同他的亵渎，添加了许多其他新的、更重要的、敌视信仰与真理的东西。此外，\Person{瓦伦廷}和\Person{巴西里德}的时代也很明显，他们也是在宗徒之后，过了很长时期，以他们的邪恶谎言反叛了天主的教会。显然，其他异端后来也引入了他们邪恶的派别和乖谬的发明，各人被错误引领；所有这些人都显然自定其罪，在审判之日以前就对自己宣判了不可避免的刑罚。而谁若确认这些人的洗礼，他除了是与他们一同受审、定自己的罪、使自己成为这样的同伙之外，还能做什么呢？\par

6. 然而，在\Place{罗马}的人并非在所有事情上都遵守从起初传下来的那些东西，却徒然妄称宗徒的权威；任何人也可以从这一点知道：关于庆祝逾越节，以及关于神圣事物的许多其他圣事，他可以看到他们中有一些差异，并非所有事情都像在\Place{耶路撒冷}所遵守的那样被他们同样遵守，正如在许多其他省份中，许多事情也因地名不同而有所变化。然而，为此并没有丝毫偏离天主教会和平与合一，\Person{斯德望}现在却胆敢这样做；他破坏了针对你们的和平，而他的前任们一直以互爱互敬保持着与你们的和平，甚至在此事上诽谤蒙福的宗徒\Person{伯多禄}和\Person{保禄}，仿佛正是这些人交付了那藉其书信咒骂异端、并警告我们要远离他们的人所交付的东西。由此可见，那维护异端、并断言他们有洗礼——这洗礼唯独属于教会——的传统，是人的传统。\par

7. 此外，你已很好地答复了\Person{斯德望}在他的信中所说的部分：即异端者自己在洗礼上也意见一致，他们并不给彼此来投靠的人施洗，而只是与他们共融，仿佛我们也应当这样做。在此处，虽然你已经证明，跟随那些在错误中的人是足够荒谬的，但我们还要补充这一点：异端者如此行并不奇怪，他们虽然在某些次要事情上有分歧，但在那最重要的事上却持有同样的共识，即亵渎造物主，为自己臆想出一个未知天主的某些梦幻和幻象。确实，这些人自然会在拥有一种不真实的洗礼上达成一致，正如他们在否认天主性的真理上一致那样。答复他们的种种恶言或愚昧陈述是冗长的，所以仅需简要概括：他们若不能拥有真主——圣父，就不能拥有圣子或圣神的真理；据此，那些被称为弗吕家人、并试图声称自己拥有新预言的人，也不能拥有圣父、圣子、\Emph{也不能有圣神}。若问他们宣讲什么基督，他们会回答说，他们宣讲那位派遣那藉着\Person{孟他努}和\Person{普黎斯加}发言的圣神的那一位。我们看到，在他们身上的不是真理之神，而是错误之神，因此我们知道，那些坚持他们反对基督信仰的虚假预言的人，不能拥有基督。此外，所有其他异端者，如果他们已经与天主的教会分离，就不能拥有任何能力或恩宠，因为一切能力和恩宠都建立在教会内，其中长老们主持，他们拥有施洗、按手和祝圣的权力。正如一个异端者不可合法地祝圣或按手，同样他也不能施洗，或做任何神圣属灵的事，因为他是异于属灵和使人成圣的圣洁者。这一切是我们先前在\Place{依科尼雍}（\Place{夫黎基雅}的一个地方）所证实的，当时我们与来自\Place{迦拉达}、\Place{基里基雅}及其他邻近地区的人聚集在一起，确定并坚决维护反对异端者的立场，当时有些人心对此事存有疑虑。\par

8. 既然\Person{斯德望}及其赞同者争辩说，罪过的赦免和重生可能来自异端者的洗礼，而他们自己又承认在异端者中没有圣神；那么让他们思考并明白：属灵的出生不能没有圣神。与此一致，蒙福的宗徒\Person{保禄}也曾以属灵的洗礼重新为那些在圣神由主派遣之前已受\Person{若翰}洗礼的人施洗，然后给他们按手，使他们领受圣神。这是什么样的事呢？当我们看到\Person{保禄}在\Person{若翰}的洗礼之后再次为他的门徒施洗，我们却迟疑着不为那些从异端来到教会的人施洗，尽管他们的浸洗是不洁和亵渎的。难道\Person{保禄}还不如当世的主教们，以至于这些主教仅凭按手就能给来投靠的异端者圣神，而\Person{保禄}却不能用按手给那些受\Person{若翰}洗礼的人圣神，除非他先用教会的洗礼给他们施洗吗？\par

9. 此外，荒谬的是，他们认为无需查问是谁施洗，因为受洗者借着呼求天主圣三，即呼求圣父、圣子、圣神之名，已获得了恩宠。那么，这将是\Person{保禄}所说的在成全者中拥有的智慧。但教会中谁是成全而智慧的人，能够辩护或相信：仅仅呼求名字便足以获得罪赦和圣洗圣事的圣化？因为只有当施洗者拥有圣神，并且圣洗本身也不是没有圣神而施行时，这些才有效。但他们说，凡是在外面以任何方式受洗的人，都可能凭自己的意愿和信德获得圣洗的恩宠，这本身无疑可笑，好像邪恶的意愿能从天吸引正义的圣化，或虚假的信德能吸引信徒的真理。但并非所有呼求基督之名的人都蒙垂听，他们的呼求也不能获得任何恩宠，主亲自揭示了这一点，说：\Quote{将有许多人冒我的名来，说：我是基督，并且要迷惑许多人。}\BibleRef{谷 13:6}因为假先知与异端者之间并无区别。前者以天主或基督的名欺骗，后者以圣洗圣事欺骗。两者都以谎言企图迷惑人的意愿。\par

10. 但我愿向你们讲述一些发生在我们中间、与此事有关的事实。大约二十二年前，在\Person{亚历山大皇帝}之后的什一税期间，这些地区发生了许多纷争和困难，无论是针对所有人，还是私下针对基督徒。此外，还有许多频繁的地震，以至于\Place{卡帕多西亚}和\Place{本都}许多地方被颠覆；甚至有些城市被拖入深渊，被裂开的大地吞噬。因此，一场针对我们基督徒的严厉迫害也随之而来；这场迫害在之前长期的和平之后突然爆发，并以不寻常的灾祸变得更加可怕，扰乱了我们的人民。那时，\Person{塞雷尼亚努斯}是我们省的总督，一个残酷可怕的迫害者。但信徒们处于这种混乱状态，因害怕迫害而四处逃散，离开自己的乡国，前往其他地区——因为当时有迁徙的机会，因为那迫害并非遍及全世界，而是局部性的——在我们中间突然出现了一个女人，她在出神状态中宣称自己是女先知，并装作充满圣神。她被首要的恶魔的冲动所驱动，长期使弟兄们忧虑和受骗，完成了一些奇事和异兆，并承诺要使大地震动。并非恶魔的力量大到足以震动大地或扰乱元素；而是有时邪恶的精灵预知并察觉将有地震，就假装要做出它预见到要发生的事。通过这些谎言和夸耀，它如此征服了个人的心，以致他们服从它，无论它命令和引导到哪里都跟随。它还会使那女人在严冬中赤脚走在冰冻的雪地上，而行走丝毫不受困扰或伤害。此外，她还会说自己正赶往\Place{犹太地}和\Place{耶路撒冷}，假装仿佛从那里来。在这里，她也欺骗了一位乡村长老和另一位执事，以致他们与那女人发生关系，这件事不久后被发觉。因为突然有一位驱魔者出现在她面前，他是位被认可的人，在宗教纪律方面始终品行良好；他受到许多在信德上坚强且值得称赞的弟兄们的激励，起来对抗那邪恶的精灵以战胜它；而那个精灵此前以它的狡猾阴谋曾预言过，说某个敌对和不信的试探者将要到来。然而，那位驱魔者受天主恩宠激励，勇敢地抵抗，并表明那先前被认为神圣的其实是一个极邪恶的精灵。那女人先前凭借恶魔的诡计和欺骗试图做许多事来迷惑信徒，在她已迷惑许多人的其他事情中，也经常胆敢这样做：以不可轻蔑的呼求祝圣饼并举行圣体圣事，向主献上祭献，且不缺少惯常言辞的圣事；她也用惯常且合法的询问辞为许多人施洗，以便在一切事上似乎与教会规矩无异。\par

11. 那么，关于这女人的圣洗，即一个极邪恶的精灵藉一个女人所施行的圣洗，我们该说什么呢？\Person{斯德望}和与他意见一致的人也赞同这个吗？尤其当她既不缺少天主圣三的象征，也不缺少合法且教会的询问辞时？能否相信罪赦已被赐予，或拯救之层的重生已适当完成，尽管一切虽按真理的形象，却是由一个恶魔所做的？除非，也许，那些为异端者的圣洗辩护的人争辩说，恶魔也以圣父、圣子、圣神之名赋予圣洗的恩宠。在他们中间，无疑存在着同样的错误——这正是魔鬼的欺骗，因为他们中间根本没有圣神。\par

12. 再者，斯德望所坚持的，即基督的亲临和圣洁与那些在异端者中受洗者同在，这有何意义呢？因为如果宗徒所言非虚，\BibleQuote{你们凡受洗归于基督的，就是穿上了基督} \BibleRef{迦 3:27}，那么，在异端者中受洗归于基督的人，确实穿上了基督。但如果他穿上了基督，他也可以领受由基督所派遣的圣神，那么，当他来到我们这里领受圣神时，为他覆手就是徒劳的；除非，也许他并没有从基督那里穿上圣神，以致基督果然与异端者同在，但圣神却不同在他们身上。\par

13. 但让我们也简要地过一遍你所详尽阐述的其他事项，尤其因为我们敬爱的执事洛加提亚努斯正赶往你处。因为接着我们需要质问那些为异端者辩护的人：他们的圣洗是属血气的还是属神的？如果是属血气的，那么它与犹太人的圣洗毫无区别，后者被使用的方式如同普通而世俗的洗濯，只洗去外在的污秽。但如果是属神的，那么在那些没有圣神的人中，圣洗怎能是属神的呢？因此，他们用以洗濯的水对他们来说只是属血气的洗濯，而不是圣洗圣事。\par

14. 但是，如果异端者的圣洗能带来第二次出生的重生，那么在他们当中受洗的人就不应被视为异端者，而是天主的子女。因为在圣洗中发生的第二次出生，产生了天主的子女。然而，如果基督的净配只有一位，那就是至公教会，唯独她生出天主的子女。因为基督的净配并非许多，正如宗徒所说：\BibleQuote{我已把你们许配给一个丈夫，如同把贞洁的童女献给了基督} \BibleRef{格后 11:2} 以及 \BibleQuote{女子，请听，请看，请侧耳细听：忘记你的民族和你的父家，因为君王恋慕你的美貌} \BibleRef{咏 45:11-12} 和 \BibleQuote{我的新娘，请从黎巴嫩与我同去，从黎巴嫩，从阿玛纳山巅，从色尼尔和赫尔孟山巅，从狮子的巢穴，从豹子的山岭下来} \BibleRef{歌 4:8} 以及 \BibleQuote{我进入了我的花园，我的妹妹，我的新娘} \BibleRef{歌 5:1}。我们看到，处处只提到一位，因为净配只有一位。但异端者的会堂与我们并不合一，因为净配不是淫妇和娼妓。因此，她也不能生出天主的子女；除非，如斯德望所认为的，异端确实生出他们并抛弃他们，而教会则收养被抛弃的，并养育那些并非自己所生的子女，但教会不能成为陌生子女的母亲。因此，我们的主基督表明祂的净配只有一位，并宣示祂合一的奥迹，说：\BibleQuote{不与我相合的，就是反对我；不同我收集的，就是分散} \BibleRef{路 11:23}。因为如果基督与我们同在，而异端者不与我们在同一侧，那么异端者确实是反对基督的；如果我们与基督一同收集，而异端者不与我们一同收集，那么他们无疑是分散的。\par

15. 但我们也不应忽略你所必然提到的：教会根据 \BibleBook{雅}{歌} 是一座封闭的花园，一个封锁的泉源，一个结有苹果的乐园。\BibleRef{歌 4:12} 正如宗徒伯多禄所奠定的，说：\BibleQuote{这水所预表的圣洗，现在也拯救你们} \BibleRef{伯前 3:21}，表明正如那些不在诺厄方舟内的人，不仅没有被水净化而得救，反而在洪水中立即灭亡；同样，如今凡不在教会内与基督同在的人，也将在外灭亡，除非他们通过补赎皈依教会唯一而救恩的圣洗。\par

16. 但是，他的错误何其大，他的盲目何其深！他竟说罪过的赦免可以在异端者的会堂中赐予，却不住留在那唯一教会的根基上——这教会曾由基督建立在磐石上。这一点可以从以下事实看出：基督单独对伯多禄说：\BibleQuote{凡你在地上所束缚的，在天上也要被束缚；凡你在地上所释放的，在天上也要被释放} \BibleRef{玛 16:19}。此外，在福音中，当基督只向宗徒们嘘气时说：\BibleQuote{你们领受圣神吧！你们赦免谁的罪，谁的罪就得赦免；你们保留谁的罪，谁的罪就被保留} \BibleRef{若 20:22-23}。因此，赦罪的权柄赐给了宗徒们，以及由基督派遣他们所建立的教会，并赐给了通过替代性祝圣继承他们的主教们。但那些反对我们所在的唯一至公教会的敌人，以及反对我们这些宗徒继承者的敌人，他们为自己篡夺不合法的司祭职，设立亵渎的祭台，他们不是别的，正是科辣黑、达堂和阿彼兰，怀着同样的邪恶亵渎，并要遭受与他们同样的惩罚；而那些赞同他们的人，也正如他们的同伙和帮凶一样，将以同样的死亡而灭亡。\par

17. 就此而言，我对于斯德望如此公开而明显的愚昧感到义愤填膺：他如此夸耀自己主教座位的地位，并坚持说自己继承自伯多禄——教会的基础就安放在他之上——却引入许多别的磐石，建立许多新教会的基础；声称凭他的权柄，在他们中有圣洗。因为受洗的人无疑充实了教会的数目。但认可他们圣洗的人，也主张教会与他们同在。他不明白，当他如此背叛并抛弃合一之时，基督磐石的真理就被他遮蔽，并在某种程度上被废除了。宗徒承认，犹太人虽因无知而盲目，且为最严重的邪恶所束缚，却仍有对天主的热情。斯德望宣布自己继承伯多禄的宝座，却不对异端者激起任何热情，竟将恩宠的权能——不是中等的，而是最大的——让给他们：甚至说并断言，借着圣洗圣事，旧人的污秽被他们洗去，他们赦免了先前的大罪，他们借着天上的重生使人成为天主的子女，并借着神圣水泉的圣化更新人获得永生。凡如此将教会伟大而属天的恩赐让给异端者的人，除了与他所维护并为之要求如此多恩宠的人共融之外，还能做什么呢？如今他徒然犹豫不决，不肯同意他们，不肯在其他事上也与他们有份，与他们聚会，平等地与他们一起祈祷，并设立共同的祭台和祭献。\par

18. 但是，他说，基督的名字对于信德和圣洗的圣化大有裨益；因此，无论何人在任何地方因基督的名受洗，就立刻获得基督的恩宠：尽管这个主张可以简单地回应和驳斥：如果因基督的名而举行的洗礼（但不在教会内）对人有益，那么，同样因基督的名而行的按手礼也能使人领受圣神；并且异端者中所行的其他事物，若是因基督的名而行，也将开始显得正当合法；正如你在信中坚持的：基督的名只有在教会内才有效用，因为基督唯独将天上恩宠的权能交给了教会。\par

19. 至于他们对习俗的辩驳——他们似乎以习俗反对真理——有谁如此愚蠢，竟偏爱习俗胜过真理，或看见光明却不弃绝黑暗呢？——除非最古老的习俗在基督（即真理）来临后，对犹太人仍有所裨益，使他们固守旧习，而弃绝真理的新路。你们非洲人确实可以这样对斯德望说：当你们认识真理时，就弃绝了习俗的错误。但我们把习俗与真理结合起来，并以习俗反对罗马人的习俗，但这是真理的习俗；我们从起初就持守基督和宗徒所传授的。我们也不记得这在我们中间有任何开始，因为这里一直遵守着：我们不认识天主的教会以外的任何一个教会，并且只认可神圣教会的圣洗为圣洁。诚然，有些人对于接受新先知、却似乎与我们承认同一位圣父和圣子的人的圣洗有所怀疑；我们许多人在依科尼雍集会，非常仔细地审查了此事，并断定：凡在教会之外安排的洗礼，均应完全弃绝。\par

20. 至于他们为异端者所引用的主张说，宗徒曾言：\Quote{或是假意，或是真心，基督总被传扬} \BibleRef{斐 1:18}，我们无需作答；因为显然，宗徒在提及此事的那封书信中，既未提到异端者，也未提到异端者的洗礼，而只是谈论弟兄们——无论是背信地随声附和，还是持守真诚的信德；此事无需长篇大论，只需阅读书信本身，并从宗徒自己那里领会宗徒所说的话。\par

21. 那么，他们说，那些从异端者而来、未经教会洗礼就被接纳的人，将如何呢？如果他们已辞世，他们就被算在我们当中望教者的数目里，但未受洗而死——他们因弃绝错误而获得的真理和信德并非微不足道，尽管因死亡受阻未能获得恩宠的圆满。但仍在世的人，应当接受教会的洗礼，以获罪赦，免得因他人的妄断而仍陷于旧错，未得恩宠圆满而死。然而，接纳者或被接纳者的罪是何等重大啊！他们未受教会水泉的洗涤，罪过未得赦免，便草率领受共融，触摸主的身体和血，尽管经上记载：\Quote{凡不相称地吃这饼或喝这主杯的，就是干犯主的身体和血！} \BibleRef{格前 11:27}\par

22. 我们认为，那些曾经是天主教会的主教、后来擅自掌握圣职授职权的人所施行的洗礼，也应视为无效。我们中间遵守此规：凡由他们施洗而来到我们这里的人，在我们中间作为陌生人、并未获得任何东西，而要接受天主教会唯一而真实的洗礼，获得生命水泉的重生。然而，因迫害的必要而勉强屈服的人，与怀着亵渎的意志大胆反抗教会、或以不敬的口舌亵渎基督的父和天主、全世界的创造者的人之间，有很大区别。斯德望却不以为耻地断言，罪过的赦免可以由那些自身陷于各种罪中的人赐予，仿佛在死亡之屋里能有救恩的水泉一样。\par

23. 那么，经上所言何解？\Quote{戒饮异水，勿汲异泉}，若你离开教会的封泉，自己取用异水，并以不圣洁的泉源玷污教会？当你与异端者的圣洗相通时，你所做的岂非是饮自他们的污泥和泥浆？当你自己藉教会的圣化被洁净时，却因接触他人的污秽而被玷污？你为异端者作证而反对教会，岂不惧怕天主的审判？尽管经上写道：\Quote{作假见证的必不受惩罚吗？} 但你确实比所有异端者更坏。因为当许多人一旦得知自己的错误，便从他们那里来到你处，为要接受教会真光时，你却援助来者的错误，遮蔽教会真理之光，堆积异端黑夜的黑暗；并且当他们承认自己处于罪中，没有恩宠，因此来到教会时，你却夺走他们在圣洗中所领受的罪赦，说他们已受圣洗，并在教会之外获得了教会的恩宠，你没有觉察到当审判之日来临时，他们的灵魂将从你手中被追讨，因为你拒绝给口渴者教会的饮品，并成为那些渴望生活者致死的原因。经过这一切，你竟愤慨！\par

24. 请想一想，你是多么缺乏判断，竟敢责备那些为真理对抗虚妄的人。因为谁更应正当地对另一方发怒？——是那支持天主仇敌的人，还是那反对支持天主仇敌者、与\Emph{我们}联合为教会真理辩护的人？——只不过显而易见，无知者也会激动发怒，因为他们缺乏谋略和谈论，容易转为怒气；因此，神圣圣经论到谁都不如论到你那样说：\BibleQuote{易怒的人挑起争端，狂暴的人堆积罪恶。}\BibleRef{箴 29:22} 因为你在全世界的教会中挑起了多少争端和分裂！此外，当你将自己从如此众多的羊群中割断时，你为自己堆积了多大的罪！因为是你割断了自己。不要自欺，那真正分裂的人，是使自己成为教会合一共融的背弃者。因为当你以为所有人都能被你绝罚时，你只将你自己从所有人中绝罚了；甚至连宗徒的训诫也未能将你塑造成真理与和平的规范，尽管他警告并说：\BibleQuote{所以我这为主被囚的，劝你们：既然蒙召，行事为人就当与蒙召的恩相称，凡事谦虚、温柔、忍耐，用爱心互相宽容，用和平彼此联络，竭力保守圣神的合一。身体只有一个，圣神只有一个，正如你们蒙召，同有一个指望；一主，一信，一圣洗；一天主，就是众人的父，超乎众人之上，贯乎众人之中，也住在我们众人之内。}\par

25. \Person{Stephen}是何等仔细地遵守了宗徒这些有益的诫命和警告，首先保持了谦虚的心和温柔！因为有什么比与全世界的这么多主教意见不合、以各种纷争与每一位决裂和平更谦虚、更温柔的呢？——有时与东方教会，我们确信你们知道；有时又与你们在南方的人，他从你们那里接待主教作为使者，却如此耐心和温柔，甚至不愿与他们进行普通会谈；不仅如此，他还念念不忘爱德，命令全体弟兄不可接待他们进屋，以致他们来时不仅和平与共融被拒绝，连住宿和款待也被拒之门外！这就是保持圣神在和平中的合一——将自己从爱的合一中割断，在各方面使自己与弟兄们隔绝，并以倔强纷争的疯狂反叛圣事和信德！与这样的人，能有一个圣神和一个身体吗？他里面甚至没有统一的思想，是如此滑溜、多变、不定！\par

26. 但就他而言，我们暂且不提；我们不如来处理那最大的争议。那些主张在异端者中受圣洗的人应被接纳，如同他们已获得合法圣洗的恩宠的人，声称圣洗对他们和我们是一样的，毫无差别。但宗徒保禄怎么说？\BibleQuote{一主，一信，一圣洗，一天主。}\BibleRef{弗 4:5} 如果异端者的圣洗和我们的圣洗是同一而相同的，那么无疑他们的信德也是同一的；但如果我们的信德是同一的，那么我们也确实有同一的主；如果有同一的主，那么我们必须说他是一位。但如果这种不可分离、不可分割的合一本身也在异端者中，那我们为何还要争论？为何称他们为异端者而不称基督徒？此外，既然我们和异端者没有同一的天主，没有同一的主，没有同一的教会，没有同一的信德，甚至没有同一的圣神，没有同一的身体，那么显然，圣洗也不能在我们与异端者之间共有，因为我们之间毫无共同之处。然而，\Person{Stephen}并不羞于对这样的人给予庇护以反对教会，并且为了维护异端者而分裂弟兄团体，此外，他还称呼\Person{Cyprian}为\Quote{假基督、假宗徒、欺诈的工人}。而他明知所有这些品性都在他自己身上，却抢先于你们，错误地将他自己本应聆听的话加于他人。我们全体，以我们所有人的名义，与所有在非洲的主教、全体圣职人员和全体弟兄团体，向你告别；愿你常怀同一心意，思想同一事物，即使相隔遥远，也使我们发现你与我们结合。\par

\SourceRecord{Translated by Robert Ernest Wallis.；Ante-Nicene Fathers；Vol. 5.；Edited by Alexander Roberts, James Donaldson, and A. Cleveland Coxe.；Buffalo, NY: Christian Literature Publishing Co.,；1886.；<http://www.newadvent.org/fathers/050674.htm>.}

% structure: p0001 source=h1 target=section reason=page-title

\section{第七十五封书信}

\Emph{致马格努斯，论为诺瓦提安派施洗及病榻上获得恩宠者。}\par

\Emph{摘要——本函前半部分与前几封信内容相同，唯关于诺瓦提安派，重申了对所有异端者所说过的要点；并顺便暗示，科尔乃略在罗马的合法继承权是人所共知的，正如教会为人所知一样。第二部分（迄今，如标题足够表明的那样，一直被误作为单独信函发表）教导说，在病榻上以洒水礼所领受的圣洗，与在教会中受浸礼一样，都是真实的圣洗。}\par

西彼廉致其子马格努斯，问安。最亲爱的儿子，你以惯常的宗教热诚，就关于其他异端者中，那些来自诺瓦提安派的人，是否在承受了他亵渎的圣洗之后，应在天主教会内以教会合法、真实、唯一的圣洗受洗并得以圣化一事，咨询了我贫乏的才智。关于此事，按我的信德能力以及神圣圣经的圣洁与真理所提示的，我回答说：任何异端者和分裂者都完全没有权能与权利。因此，诺瓦提安既不应、也不能期望被算作对手和假基督，因为他同样是在教会之外，并且与基督的平安与爱相敌对。因为我们的主耶稣基督在祂的福音中作证，那些不与祂同在的人就是祂的敌人时，并未指出任何异端类别，而是表明所有不与祂同在、且不与祂一同收聚而分散祂羊群的人，都是祂的敌人，祂说：\Quote{不与我相合的，就是敌我的；不同我收聚的，就是分散的。} 此外，蒙福的宗徒若望本人也没有区分任何异端或分裂，也没有特别指定任何人；但他称所有从教会出去、与教会敌对的人为假基督，说：\BibleQuote{你们听见说过，假基督要来，现在已经有许多假基督出来了；由此我们知道这是最后的时刻。他们从我们中间出去，却并不是属于我们的；因为如果是属于我们的，他们必会与我们同在。} \BibleRef{若一 2:18} 由此可见，所有被称为离开了爱德和天主教会一体性的人，都是主的敌人和假基督。此外，主在祂的福音中确立并说：\BibleQuote{若他连教会也不听从，你就把他看作外邦人和税吏。} \BibleRef{玛 18:17} 如今，如果那些藐视教会的人被算作外邦人和税吏，那么那些背叛与仇敌——他们伪造祭台、设立不法的司祭职、举行亵渎的祭献、败坏名字——就更必然要被算作外邦人和税吏了；因为那些罪过较轻、仅是藐视教会的人，已被主的判决断定为外邦人和税吏。\par

2. 然而，圣神在 \BibleBook{}{雅歌} 中，以基督的口吻宣称教会是唯一的，说：\BibleQuote{我的鸽子，我的完全人，是唯一的；她是她母亲的独一者，是生养她者的蒙拣选者。} \BibleRef{歌 6:9} 关于此，祂又说：\BibleQuote{我的妹妹，我的新娘，是关闭的花园，是封闭的泉源，是活水的井。} \BibleRef{歌 4:12} 但如果基督的新娘——教会——是关闭的花园，那么关闭的东西就不能向陌生人和亵渎者敞开。如果它是封闭的泉源，那置身于外、无法触及泉水的人，既不能从那里饮水，也不能被封印。同样，活水的井，如果它是在内的同一口井，置身于外的人就不能从那水中获得生命和圣化，因为那水只允许里面的人使用或饮用。伯多禄也显示这一点，他指出教会是唯一的，只有那些在教会里的人才能受洗；他说：\BibleQuote{在诺厄方舟里，只有少数人，即八个灵魂，借着水得了救；这水所预表的圣洗，如今也拯救你们。} \BibleRef{伯前 3:20} 以此证明和证实，诺厄的独一方舟是唯一教会的预像。那么，如果在那次世界被罪罚和净化的洗礼中，不在诺厄方舟里的人能借着水得救，那么现在，不在那唯一被授予圣洗的教会里的人，也能借着圣洗获得生命吗？此外，宗徒保禄更公开、更清楚地阐明同一件事，写信给厄弗所人说：\Quote{基督爱了教会，并为教会舍了自己，为能以水的洗涤圣化并洁净教会。} 但如果教会是基督所爱的那一位，并且唯独由祂的洗涤所洁净，那么一个不在教会内的人怎能被基督所爱，或被祂的洗涤所洗净和洁净呢？\par

3. 因此，既然只有教会拥有活水以及施洗并洁净人的权能，那么任何声称有人能被诺瓦提安施洗并圣化的人，必须首先证明并教导：诺瓦提安是在教会内，或是管理着教会。因为教会是唯一的，且既然她是唯一的，就不能同时在内又在外。因为如果教会与诺瓦提安同在，她就未与科尔乃略同在。但如果教会与科尔乃略同在——他借着合法祝圣继承了法比盎主教之位，并且主除了司祭的尊荣外，还以殉道荣耀了他——那么诺瓦提安就不在教会内；他也不能被算作主教，因为他没有继承任何人，反倒藐视福音与宗徒传统，自己从自身兴起。因为凡未在教会内受祝圣的人，绝不能拥有或持守教会于任何方式。\par

4. 因为圣经的信德昭示，教会不是在外，也不能分裂或自相分割，而是保持一个不可分离且不可分割之家的合一；因为经上论及逾越节的圣事和羔羊（这羔羊预指基督）记载说：\BibleQuote{在一座房屋内应吃；不可将肉带到屋外。}\BibleRef{出 12:46} 这也见于论及\Person{辣哈布}的话，她本人也是教会的预象。她曾接到命令说：\BibleQuote{你应把你的父亲、母亲、兄弟和你父亲的全家，都聚集到你屋内；谁若走出你的屋门到街上去，他的血债应归在他自己头上。}\BibleRef{苏 2:18} 在这奥秘中宣告：凡愿存活并逃脱世界毁灭的，必须被聚集到唯一的房屋内，即教会内；但那些被聚集的人中，谁若走到外面，即如果有人虽然在教会中获得了恩宠，却离开并走出教会，他的血债就归在他自己头上；也就是说，他必须自己承担自己的丧亡。宗徒\Person{保禄}解释这事，教训并命令应躲避异端者，因为他是乖僻的、有罪的、且是自定罪的。因为那未被主教逐出，却自愿离开教会，并因异端者的狂妄而自定罪的人，将对自己的灭亡有罪。\par

5. 因此，主向我们提示一个源自天主权威的合一，祂规定说：\BibleQuote{我与父原是一体。}\BibleRef{若 10:30} 为使祂的教会归于此合一，祂又说：\BibleQuote{将只有一个\Emph{羊群}，一个牧人。}\BibleRef{若 10:16} 但如果羊群只有一个，那不在羊群数目中的，怎能算在羊群里呢？或者，当真正的牧人仍在，并借着连续的祝圣管理天主的教会时，那继承无人、自成一派，成为陌生和亵渎的人、主和平及天主合一的仇敌，不住在天主的家里，即天主的教会中的人，怎能被尊为牧人呢？在教会中，除非是同心合意的人，无人能居住，因为圣神在圣咏中发言说：\Quote{是天主使人同心合意地住在家里。}\par

6. 此外，连主的祭献本身也宣告，基督徒的同心合意是靠坚固且不可分离的爱德彼此相连的。因为主称那由许多谷粒结合而成的饼为祂的身体，祂就指明祂所背负的我们这民族是合一的；祂称那由许多葡萄和果串压榨并收集而成的酒为祂的血，祂也表示我们这羊群是由合一群体的交融而相连的。如果\Person{诺瓦替安}与此主的饼合一，如果他也与此基督的杯相融，那么他似乎也能获得教会唯一圣洗的恩宠，只要他显然持守着教会的合一。最后，合一的圣事多么不可分割，那些制造分裂、离开自己的主教、在外部另立一个假主教的人，为自己从天主的义怒中赚得了何等绝望和过度的毁灭——这在列王纪中宣告了：当十个支派从犹大和本雅明支派分裂，离开他们的王，在外部另立一个王时，经上说：\BibleQuote{上主对以色列所有后裔大发忿怒，将他们赶走，交在抢夺者手中，直到将他们从自己面前抛去；因为以色列从达味家分裂，他们为自己立了王，乃巴特的儿子雅洛贝罕。}\BibleRef{列下 17:20} 经上说上主大发忿怒，将他们交于丧亡，因为他们从合一中分裂，为自己另立了王。上主对那些制造分裂者的义怒如此巨大，以至于当天主的人被派到\Person{雅洛贝罕}那里，指责他的罪，并预言将来的惩罚时，他被禁止与他们一同吃饼或喝水。当他没有遵守这一命令，违反天主的命令而吃了肉时，他立刻被天主审判的威严击打，以至于从那里返回时，被一头袭击他的狮子在路上咬死。谁敢说，与那些连地上的食物世上的饮料都不应共享的分裂者，能共享圣洗的救赎之水与天上的恩宠？此外，主在福音中满足我们，并显示更光明的领悟力：那些从前从犹大和本雅明支派分裂，离开\Place{耶路撒冷}退到\Place{撒玛黎雅}的人，应被视为亵渎者和外邦人。因为当祂初次派遣门徒从事救恩的职务时，祂吩咐他们说：\Quote{不要走外邦人的路，不要进入撒玛黎雅人的任何城市。} 祂先派遣到犹太人那里，命令暂时越过外邦人；但加上甚至要略过\Place{撒玛黎雅}人的城市（那里有分裂者），就显示分裂者应被置于与外邦人同等的地位。\par

7. 但若有人提出反对，说\Person{诺瓦提安}持有与天主教会相同的法规，以我们所用同一标记施洗，认识同一天主圣父、同一基督圣子、同一圣神，并因此他可主张施洗之权，即他在圣洗询问上似乎与我们无异；任何认为可如此反对的人，首先要知道，信经的法规并非同一，我们与分裂者也没有相同的询问。因为他们说：\Quote{你相信藉着圣教会得赦罪和永生吗？}他们在询问中撒谎，因为他们没有教会。此外，他们自己亲口承认，赦罪只能藉圣教会赐予；既然他们没有教会，便表明在他们中间不能赦罪。\par

8. 但说他们与我们有同一天主圣父、认识同一基督圣子、同一圣神，这对这样的人毫无益处。因为\Person{科辣黑}、\Person{达堂}和\Person{阿彼兰}也认识同一天主，如同司祭\Person{亚郎}和\Person{梅瑟}一样。他们生活在同一法律和宗教之下，呼求那应当被呼求和崇拜的独一真天主；然而，因为他们违反自己职位的职务，反对司祭亚郎——他因天主俯允和主的祝圣而接受了合法司祭职——并擅自夺取祭献之权，遭受神罚，立即因他们非法行径而受罚；而且那些不虔诚、不合法、违反神圣任命权利的祭献，不能蒙悦纳，也不能使他们获益。甚至那些非法献香用的香炉，为使司祭不再使用，反要作为神罚和义怒的纪念，以教训后代，奉命被火熔化炼净，打成薄片，钉在祭台上，正如圣经所说：\BibleQuote{作为以色列子民的纪念，使任何非亚郎子孙的外人不得前来在主前献香，免得他像科辣黑一样。}然而那些人并未制造分裂，也未外出，而是敌对天主的司祭无耻地反叛；但这些人现今却分裂教会，作为和平与基督合一的叛逆，设法为自己设立宝座，夺取首位，并主张施洗和献祭之权。他们怎能完成所行之事，或从天主那里因非法企图获得什么？因为他们所行的，是反对天主、是他们不合法的事。因此，那些支持\Person{诺瓦提安}或此类分裂者的人，徒然主张在他们中间能有人被施以救恩之洗而圣化，而施洗者显然没有施洗之权。\par

9. 此外，为更好理解天主对类似大胆行为的审判，我们发现，在如此邪恶中，不仅首领和发起者，连参与者也将受罚，除非他们与恶人的共融分离；如主藉梅瑟命令说：\BibleQuote{你们应从这些最顽固之人的帐棚中退出，不要碰他们的任何东西，免得你们在他们的罪中被消灭。}\BibleRef{户 16:26}而主藉梅瑟所警告的，祂也实现了，即凡未与\Person{科辣黑}、\Person{达堂}和\Person{阿彼兰}分离的，立即因其不敬的共融而受罚。此例表明并证明，凡以不敬的胆量，在反对上司和司祭中与分裂者同流合污的，都将负罪受罚；正如圣神藉\Person{欧瑟亚}先知见证说：\BibleQuote{他们的祭献必像丧饼，凡吃它的都必污秽。}\BibleRef{欧 9:4}无疑地教导并表明，凡被他们的罪行玷污的，全与首领一同受罚。\par

10. 那么，在天主眼中，那些被天主宣布惩罚的人，能有什么功德？或者说，这样的人——他们身为司祭的仇敌，妄图篡夺违法且不属于自己的事，且无任何权利给他们——怎能成义并圣化受洗者？然而，我们并不奇怪，他们根据其邪恶而为此争辩。因为他们每人必须维护自己所行之事；即使失败，也不轻易屈服，尽管他们知道所行不合法。那令人惊奇的事，更确切地说，令人愤慨和痛心的是：基督徒竟支持假基督；信仰的背叛者和教会的出卖者竟留在教会内。这些人，尽管在其他方面顽固不化、不受教导，但仍然至少承认这一点：所有异端者或分裂者都没有圣神，因此他们确实能施洗，但不能赋予圣神；正是在这一点上，我们抓住了他们，因为我们表明，没有圣神的人根本不能施洗。\par

11. 因为在圣洗中每个人都获得了自己罪的赦免，主在祂的福音中证明并宣告，只有那些拥有圣神的人才能赦罪。因为祂复活后，派遣祂的门徒，对他们说：\Quote{就如父派遣了我，我也同样派遣你们。}说了这话，就向他们嘘了一口气，说：\Quote{你们领受圣神吧！你们赦免谁的罪，谁的罪就得赦免；你们保留谁的罪，谁的罪就被保留。}\BibleRef{若 20:21}在此处祂表明，只有拥有圣神的人才能施洗并赐予罪的赦免。再者，连要为我们主基督本人施洗的若翰，也在母胎中就预先领受了圣神，为要确凿而明显地证明，除了拥有圣神的人，无人能施洗。因此，那些袒护异端或分裂者的人必须回答我们：他们是否拥有圣神。如果他们拥有，那么为什么当他们中间受洗的人来到我们这里时，要为他们覆手，使他们领受圣神？因为圣神必定已在彼处被领受，若在那里已被领受，便能赐予。但如果异端者和分裂者施洗时没有圣神，因此在我们这里为他们覆手，为使\InnerQuote{在此处}能领受\InnerQuote{彼处}既没有也不能赐予的圣神；那么显然，那些肯定没有圣神的人也不能赐予罪的赦免。因此，为了使按照天主的安排和福音的真理，那些从仇敌和假基督那里来到基督教会的人能够获得罪的赦免、得以圣化、成为天主的宫殿，他们全都必须绝对用教会的洗礼受洗。\par

12. 最亲爱的儿子，你也问了我对在疾病和虚弱中获得天主恩宠之人的看法：他们是否应被视为合法的基督徒，因为他们不是用救恩之水洗涤，而是被洒水。在这一点上，我的谦逊和谨慎并不预断任何人，以免妨碍任何人按自己的判断感受，并按自己的感受行事。就我浅薄的理解而言，我认为天主的恩惠绝不会被削弱或减损；在那情况下，当施予者和领受者都以完整而全备的信德领受来自天主恩赐的事物时，不会发生任何减少。因为在救恩的圣事中，罪的污染被洗去，并不像身体和皮肤的污垢在肉体和寻常的洗涤中被洗去那样，需要硝石和其他用品，以及浴盆和盆子来清洗和净化这个卑贱的身体。信徒的胸怀乃以另一种方式被洗净；人的心智乃借着信德的功绩被净化。在救恩的圣事中，当需要迫不得已，而天主施予仁慈时，天主的途径将完整的恩惠赐予信徒；也不应困扰任何人，以为病人在获得主的恩宠时似乎只是被洒水或浇水，因为圣经借厄则克耳先知的口说：\Quote{那时我洒清水在你们身上，你们就洁净了；我要洁净你们脱离各种不洁和各种偶像。我还要赐给你们一颗新心，在你们五内注入一种新的精神。}\BibleRef{则 36:25} 此外在户籍纪中：\Quote{那人不洁净直到晚上，第三天和第七天应洁净自己；若他第三天不洁净自己，第七天就不洁净。那人应从以色列中被铲除，因为洒水没有洒在他身上。}又说：\Quote{上主对梅瑟说：你从以色列子民中选出肋未人，洁净他们。你要这样对他们做，以洁净他们：用除罪水洒在他们身上。}\BibleRef{户 8:5} 又说：\Quote{洒水是除罪水。}\BibleRef{户 19:9} 由此看来，洒水同样具有救恩洗涤的功效；当这在教会中施行，且施予者和领受者的信德都纯正时，一切事物都借主的威严和信德的真理得以成全和圆满。\par

13. 此外，关于有些人称那些借着救恩之水和合法信德获得基督平安的人为\Quote{临床者}而非\Quote{基督徒}，我不知道他们从何处取用这名称，除非他们读过更隐晦的典籍，并从希波克拉底或索拉努斯那里借用了这些\Quote{临床者}。因为我在福音中知道一个临床者，就是那个瘫痪而软弱的人，他在自己漫长的生命中卧床不起，但他的软弱并未妨碍他完全获得天上的力量。他不仅因天主的恩宠从床上被扶起，而且以恢复和增强的力量拿起了自己的床。因此，就信德允许我领会和思考而言，我的看法是：任何人只要借信德的法律和权利在教会中获得了天主的恩宠，就应被视为合法的基督徒。如果有人以为那些仅被洒过救恩之水的人一无所获，仍是空虚和无效的，那么请不要受骗，以致如果他们摆脱了疾病的恶果而痊愈，他们\Emph{应当寻求}再受洗。但如果那些已经受过教会洗礼而圣化的人不能再受洗，那么为什么他们要在自己的信德和主的仁慈上受困扰呢？或者他们确实获得了天主的眷顾，但在天主恩赐和圣神的度量上较小、较有限，以致他们确实被视为基督徒，但却不能与其他人同等看待？\par

14. 不，确实，圣神不是按份量赐予的，而是整个倾注在相信者身上。因为如果白昼同样照耀众人，如果太阳以同样均等的光辉普照万物，那么基督——祂是真正的太阳、真正的白昼——在祂的教会中以同样的均等赐予永生之光，岂不更是如此！我们在\WorkTitle{出谷纪}中看到这均等的圣事被庆祝，那时玛纳从天降下，预表将来之事，显示了天上之粮和将要来临的基督的滋养。因为在那里，无分性别或年龄，每人收集的敖默尔都是相等的。由此表明，基督的仁慈和随后将要降临的天上恩宠是平等地分给众人的；无性别之差，无年龄之别，不偏待人，属神的恩宠之赐倾注于天主的全体子民。确实，同样的属神恩宠在圣洗中平等地被相信者领受，随后在我们的言行中或增加或减少；正如在福音中，主的种子被均匀地撒下，但因土壤的不同，有的浪费了，有的则增长为丰盛的收获，结出三十倍、六十倍或一百倍的丰硕果实。但再想一想，当每人被召来领取一个银钱时，为何那由天主平等分施的，却因人的解释而减少了呢？\par

15. 但如果有人因某些患病受洗者仍受不洁之灵的诱惑而感到不安，就让他知道：魔鬼顽固的邪恶甚至能延伸到拯救之水前，但在圣洗中，它便失去了一切恶毒。我们可在法老王身上看到这样的例子，他长期挣扎，在背信中拖延，能够抵抗并坚持直到水边；但当他到了那里，既被征服，又被消灭。而那海是圣洗的预像，蒙福的宗徒保禄宣告说：\BibleQuote{弟兄们，我不愿意你们不知道，我们的祖先都曾在云彩下，都从海中经过，都曾在云中和海中受洗归于梅瑟。}他又加上说：\BibleQuote{这些事都是我们的鉴戒。}同样的事也在今天发生：魔鬼被驱魔者以人的声音和天主的大能鞭打、焚烧、折磨；虽然它常常说自己要出去，要离开天主的子民，但它在所说的话中欺骗，并实行了法老先前同样顽固和欺诈的诡计。然而，当他们来到拯救之水以及圣洗的圣化时，我们理应知道并相信：在那里，魔鬼被击倒，献身于天主的人因天主的仁慈而获得自由。因为如同蝎子和蛇在干地上能逞凶，一旦被投入水中，便不能逞威，也不能保持毒液；同样，恶灵——它们被称为蝎子和蛇，而且我们因主赐予的能力践踏它们——不能再留在一个人身上，当他受洗并被圣化后，圣神开始住在其中。\par

16. 最后，我们在实际中也确实经历到：那些因急病而受洗、并获得了恩宠的人，摆脱了先前搅扰他们的不洁之灵，在教会中过着赞美和尊荣的生活，并且日复一日地因信德的增长而在天上恩宠的增加中不断进步。另一方面，有些健康时受洗的人，如果后来开始犯罪，就因不洁之灵的回归而被摇动，以致于显明：魔鬼在圣洗中是因相信者的信德而被驱逐，如果后来信德消失，它就会回来。除非在某些人看来，那些在教会之外、在敌对者和假基督者中，以亵渎之水受污的人被认为受了洗；而在教会中受洗的人反而被认为获得了更少的天主仁慈和恩宠；并且对异端者如此看重，以致于那些从异端来的人不被问及他们是否被洗或洒水，是\OriginalTerm{病床受洗者}{clinics}还是\OriginalTerm{可走动受洗者}{peripatetics}；但在我们中间，信德的纯正真理却被轻视，而教会圣洗的尊威和圣洁受到贬损。\par

17. 最亲爱的儿子，我已尽我微薄之能答复了你的来信；我尽可能地表明了我的看法；我不向任何人规定，以免阻止任何主教按他认为正确的决定，因为他将为自己的行为向主交账，正如蒙福的宗徒保禄在致罗马人书中写道：\BibleQuote{我们每人都要为自己的事向天主交账：所以我们不可再彼此判断。}\BibleRef{罗 14:12}我愿你，最亲爱的儿子，永远衷心安康。\par

\SourceRecord{Translated by Robert Ernest Wallis.；Ante-Nicene Fathers；Vol. 5.；Edited by Alexander Roberts, James Donaldson, and A. Cleveland Coxe.；Buffalo, NY: Christian Literature Publishing Co.,；1886.；<http://www.newadvent.org/fathers/050675.htm>.}

% structure: p0001 source=h1 target=section reason=page-title

\section{\WorkTitle{第七十六封书信}}

\Emph{\Person{西彼廉}致\Person{乃米西安}及在矿场中的其他殉道者。}\par

\Emph{内容简介——他用美妙的对照，以各种折磨的对应安慰来极尽赞美矿场中的殉道者。}\par

\Person{西彼廉}致\Person{乃米西安}、\Person{费利克斯}、\Person{路济乌斯}、另一位\Person{费利克斯}、\Person{利特乌斯}、\Person{波利安}、\Person{维克多}、\Person{雅德尔}、\Person{达提乌斯}，诸位同为主教者，亦致诸位同为长老与执事者，以及在矿场中的其余弟兄，全能的圣父天主、我们的主耶稣基督、以及保存我们的天主的殉道者，永远的平安。最蒙福、最亲爱的弟兄们，你们的荣耀确实要求我亲自前来探望并拥抱你们，但若我被放逐（因我为了\OriginalTerm{圣名}{the Name}的宣认而被流放）的指定地点的界限没有阻止我。然而，凡我所能的，我使自己临于你们面前；即使不容我以身体和行动来到你们那里，我却以爱情和心神前来，在信中表达我的思念，在这思念中，我因你们的这些德行与赞美而欢欣雀跃，自视为你们的分享者，虽非在身体的受苦中，却在爱德的共融中。当我得知关于我最亲爱的朋友们如此众多而荣耀的事迹——这些事迹是天主纡尊降贵所赐予你们的荣耀，以致你们中一部分人已借着殉道的圆满先走一步，从主那里领受他们应得的荣冠——我岂能缄默无声、抑制我的言语呢？另一部分人仍留在监牢的地牢中，或在矿场与锁链中，借着他们刑罚的延迟，展现出更伟大的榜样，以坚固和装备弟兄们，借着折磨的漫长进至更丰盛的功勋，好按他们如今在刑罚中所计数的日子，领受同样多的天国赏报。最勇敢、最蒙福的弟兄们，我毫不惊讶这些事临到你们，鉴于你们宗教与信德的应得奖赏；主竟如此以祂荣耀的尊威将你们提升至荣耀的高峰，因为你们在祂的教会中始终兴盛，持守信德的原貌，坚定遵守主的诫命；在淳朴中持守纯洁，在爱德中持守和睦，在谦卑中持守端庄，在管理上持守勤勉，在帮助受苦者上持守警醒，在照顾穷人上持守仁慈，在维护真理上持守恒毅，在纪律的严厉上持守判断。为使你们善行的榜样一无所缺，甚至如今，在你们言语的宣认和身体的受苦中，你们激发弟兄们的心思趋向神圣的殉道，借着展示你们作为德行的领袖，好使羊群跟随牧者，效法所见上司所行的事，从而由主获得同样服从的功勋荣冠。\par

2. 然而，你们先被棍杖狠狠击打、受虐待，以这类苦难开始了你们光荣的首批宣认，这并非我们该诅咒的事。因为基督徒的身体并不十分畏惧棍杖，因为它的全部希望在于\OriginalTerm{木十字架}{the Wood}。基督的仆人承认他救恩的圣事：被木头赎回到永生，借木头被提升至荣冠。但有何奇怪，如果你们如同金银器皿被交付给那金银之家的矿场，只是如今矿场的性质改变了，那些先前惯常产出金银的地方开始接收它们？此外，他们把足械放在你们脚上，捆绑你们蒙福的肢体、天主的宫殿以可耻的锁链，好像连心神也能与身体一同被捆绑，或你们的金子能被铁的接触玷污。对于献身于天主、以宗教的勇武见证信德的人，这类事物是装饰，而非锁链；它们不是捆绑基督徒的脚以招耻辱，而是荣耀他们以得荣冠。哦，蒙福被捆绑的脚，它们被释放，不是由铁匠，而是由主！哦，蒙福被捆绑的脚，它们被引导进入乐园，走在救恩的路上！哦，在现世暂时被捆绑的脚，好使它们永远与主同在自由中！哦，脚，在锁链和铁栏间滞留片刻，却要在光荣的道路上迅速奔向基督！任凭嫉妒或恶毒的残酷在这里用它的束缚和锁链扣留你们多久，你们必要迅速从这大地、这些苦难中来到天国。身体在矿场中不以卧榻和靠枕受照料，却以基督的安慰和慰藉受照料。因劳苦而疲惫的躯体卧倒在地，但与基督同卧并非惩罚。你们未沐浴的肢体，因污秽和泥土而肮脏变形；但内心却受着属灵的洁净，尽管外在肉体被玷污。那里食物稀少；但人活着不只靠饼，而是靠天主口中的话。因寒冷而缺少衣服；但穿上基督的人既充分着装又装饰华美。你们半剃的头发似乎令人厌恶；但既然基督是男人的头，则因基督之名而荣耀的头，无论怎样都必须合适。那对外邦人而言可憎可恶的丑陋，将有何等光采的赏报！这暂时的短暂苦难，将如何换来光明永恒荣耀的赏报，正如蒙福宗徒的话：\BibleQuote{主将改变我们卑微的身体，使之相似祂荣耀的身体！}\BibleRef{斐 3:21}\par

3. 但是，最亲爱的弟兄们，关于宗教或信德，如今在那里天主的司铎没有机会奉献并举行神圣祭献，这并不会感到任何损失；是的，你们向天主奉献并举行同样宝贵和光荣的祭献，这将大大有助于你们获得天上赏报的回报，因为圣经说话说：\BibleQuote{天主所要的祭献就是痛悔的精神；天主不轻视痛悔和谦卑的心。} 你们向天主奉献这祭献；你们昼夜不停地举行这祭献，成为天主的牺牲，把自己展示为圣洁无瑕的祭品，正如宗徒劝勉并说：\BibleQuote{所以弟兄们，我以天主的仁慈请求你们，献上你们的身体当作生活、圣洁、悦乐天主的祭品。你们不可与此世同化，反而要以更新的心思变化自己，使你们能辨别什么是天主的旨意，什么是善事、悦乐和成全的事。}\BibleRef{罗 12:1}\par

4. 这正是特别悦乐天主的；正是藉此，我们的工作因更大的功德而成功地赢得天主的善意；正是这，我们信德和敬虔的服从所能单单献给主以回报祂伟大和救恩的恩惠，正如圣神在圣咏中宣告并见证说：\BibleQuote{我该拿什么报谢，}祂说，\BibleQuote{给上主祂赐给我的一切恩惠？我必举起救恩的杯，呼唤上主的名。上主眼中看重的，是祂圣人的死亡。} 谁不乐意而欣然地接受救恩的杯呢？谁不喜乐而愉快地渴望那使他自己也能向他的主有所回报的事呢？谁不勇敢而坚定地接受在主眼中宝贵的死亡，以悦乐祂的眼目呢？祂从高处俯视我们这些为祂的名而处于战斗中的人，赞赏那自愿的，援助那挣扎的，以忍耐、良善和慈爱的赏报加冕那得胜的，在我们身上酬报祂自己所赐予的一切，并尊荣祂所成就的一切。\par

5. 因为我们的得胜，以及我们藉着制服仇敌而获得最艰巨竞赛的棕榈枝，乃是祂的作为，主在祂的福音中宣告并教导说：\BibleQuote{但当人把你们交出时，不要思虑怎样或说什么；因为在那时刻，自会赐给你们该说什么。因为不是你们说话，而是你们父的圣神在你们内说话。} 又说：\BibleQuote{因此，你们心中要定志，不要预先思索怎样回答；因为我要赐给你们口才和智慧，你们的仇敌不能抵抗。}\BibleRef{路 21:14} 在其中，确乎既有信徒的巨大信心，也有不信者的极重过失，因为他们不信任那应许要帮助那些宣认祂的人，反之也不惧怕那以永罚威胁那些否认祂的人。\par

6. 所有这些事，基督最勇敢和忠信的士兵们，你们已经向你们弟兄们提议了，在行动上满全了你们以前在言语上所教导的，日后在天国里将成为最大的，正如主所应许并说：\BibleQuote{谁若这样实行并教导，将在天国里称为最大的。}\BibleRef{玛 5:19} 此外，许多民众效法你们的榜样，与你们一同宣认，一同被加冕，藉着最强爱德的联系与你们结合，既没有被监狱也没有被矿坑与他们的主教们分离；在他们中间，不乏贞女，她们在六十倍果实上加上百倍，双重的光荣将她们推进到天上的冠冕。连男孩们也有超越年龄的勇气，在宣认的赞美中超越他们的岁数，以致每种性别和每种年龄都装饰了你们殉道者的蒙福羊群。\par

7. 最亲爱的弟兄们，你们那得胜的觉悟该有何等的力量，你们的心灵该有何等的崇高，你们的感受该有何等的喜乐，你们的心中该有何等的凯旋，因你们每人已接近天主所应许的赏报，已稳妥于天主的审判，在矿坑中行走，身体虽被囚禁，心却为王，你们知道基督与你们同在，为祂仆人们的坚忍而喜乐，他们正跟随祂的足迹、沿着祂的道路上升至永恒的国度！你们每日喜乐地期盼离世得救的日子；你们即将离开此世，正奔向殉道的赏报，奔向天主的家园，为在此世的黑暗之后看见最纯净的光，并接受超越一切痛苦和战斗的光荣，正如宗徒见证并说：\BibleQuote{现今的苦楚，与将来在我们身上要显示的光荣，是不能比较的。}\BibleRef{罗 8:18} 因为如今你们的祈祷话语更有效力，在患难中的恳求更容易获得所求，所以你们要更热切地寻求并祈求，愿天主的俯就成全我们所有人的宣认；愿天主在此世的黑暗和罗网中，也与你们一同释放我们，使我们健康而光荣；使我们这些在这里藉着爱德与和平之链联合，并一同抵抗异端者的过犯和外邦人的压迫的人，能在天国一同欢欣。最蒙福和最亲爱的弟兄们，我祝愿你们永远在主内平安，并时常、处处记念我。\par

\SourceRecord{Translated by Robert Ernest Wallis.；Ante-Nicene Fathers；Vol. 5.；Edited by Alexander Roberts, James Donaldson, and A. Cleveland Coxe.；Buffalo, NY: Christian Literature Publishing Co.,；1886.；<http://www.newadvent.org/fathers/050676.htm>.}

% structure: p0001 source=h1 target=section reason=page-title

\section{第七十七封书信}

\Emph{内梅西亚努斯、达提乌斯、费利克斯、维克多答复西彼廉。}\par

\Emph{内容提要。—— 本封书信以及接下来两封，除了对上述书信的答复外，别无他物，因为它们既包含对来信所传达的安慰的感谢，也包含对所送援助的感谢。但从三封不同的书信回复西彼廉单一的一封这一事实，我们应推断，写信的主教们被安置在矿场的不同部门。}\par

1. 内梅西亚努斯、达提乌斯、费利克斯、维克多向他们的兄弟西彼廉，在主内致以永远的救恩。亲爱的西彼廉，你总是在你的信中，以适合时代状况的深意说话，借着对这些信的勤勉阅读，恶人得以纠正，有信德的人得以坚定。因为当你不停地在你的著作中揭露隐藏的奥秘时，你就这样使我们在信德中成长，并使世上的人们接近信仰。因为无论你通过你的许多书籍引入了什么美好的事物，你都不自觉地向我们描绘了你自己。因为你在言谈中比所有人都伟大，在演讲中更雄辩，在智慧上更明智，在忍耐中更淳朴，在工作中更丰盛，在节制中更神圣，在服从上更谦卑，在善行中更纯洁。而且你自己知道，亲爱的，我们热切的愿望是，我们能看到你，我们的老师兼挚爱者，达到伟大宣认的花冠。\par

2. 因为，在总督面前的诉讼中，作为一位良好而真实的老师，你首先宣告了那些我们这些跟随你的门徒在长官面前应当说的话。并且，像一面响亮的号角，你已经唤醒了天主装备了天上武器的士兵，投入近距离搏斗；你在前线战斗，用属灵的剑杀死了魔鬼；你还用你的话语从左右两侧命令了弟兄们的队伍，以致对敌人布下了四面埋伏，并且公众敌人的尸体被切断的筋骨被踩在脚下。请相信我们，最亲爱的，你纯洁的灵魂离百倍的赏报不远了，因为它既不惧怕世界的第一次攻击，也不因流放而退缩，也不犹豫离开城市，也不害怕居住在旷野之地；既然它为许多人提供了宣认的榜样，它自己首先说出了殉道的见证。因为它以自身的榜样激发了他人走向殉道；不仅开始成为已经离开世界的殉道者的同伴，而且还与那些将要成为殉道者的人建立了天上的友谊。\par

3. 因此，那些与我们一同被定罪的人，亲爱的西彼廉，在天主面前向你致以最大的感谢，因为你用你的信慰藉了他们痛苦的胸膛；治愈了他们被棍棒击伤的肢体；松开了他们被脚镣捆绑的脚；抚平了他们半剃的头发的发梢；照亮了地牢的黑暗；将矿山的山脉夷为平地；甚至在他们鼻孔前放置了芬芳的花朵，并驱散了烟雾的恶臭。此外，你持续的馈赠，以及我们亲爱的奎里努斯的馈赠，你通过副执事赫伦尼亚努斯、辅祭卢西安、马克西穆斯和阿曼提乌斯送来分发，为他们身体所需的一切提供了供应。那么，让我们在祈祷中互相帮助；并且让我们像你所吩咐的那样祈求，使我们在一切行动中都有天主、基督和天使作为支持者。我们衷心祝愿你，主和兄弟，永远安好，并记住我们。请问候所有与你在一起的人。我们这里所有的人都爱你、问候你，并渴望见到你。\par

\SourceRecord{Translated by Robert Ernest Wallis.；Ante-Nicene Fathers；Vol. 5.；Edited by Alexander Roberts, James Donaldson, and A. Cleveland Coxe.；Buffalo, NY: Christian Literature Publishing Co.,；1886.；<http://www.newadvent.org/fathers/050677.htm>.}

% structure: p0001 source=h1 target=section reason=page-title

\section{第七十八封书信}

\Emph{致路爵与其余殉道者同一问题的答复。}\par

\Emph{论证——本信的论证实质上与前信相同；因此它不是罗马主教路爵的信，而是非洲主教兼殉道者路爵的信。}\par

\Person{路爵}及所有与我在主内同在的弟兄，问候我们的弟兄兼同僚\Person{西彼廉}。最亲爱的弟兄，当我们因天主而欢欣踊跃时，祂武装我们投入战斗，并屈尊使我们成为战场的胜利者，你的信函到达了我们手中；这封信是由副执事\Person{厄勒尼雅诺}、\Person{路齐亚诺}、\Person{玛克西莫}及辅祭者\Person{阿曼提奥}带给我们的。我们读信时，在束缚中得了缓解，苦难中得了安慰，困境中得了支持；我们被激起，更坚强地忍受可能等待我们的任何更多惩罚。因为在我们的苦难之前，你已召唤我们走向光荣，你首先为我们提供了导向宣认基督之名的引导。我们确实追随你们宣认的足迹，希望与你们同享相等的恩宠。因为赛跑中领先者也是奖赏的领先者；你们先行占据了跑道，便从你们的起点将之传递给我们，无疑显明了你们始终爱我们的不分之爱，使我们在和平的束缚中同有一个圣神，能得着你们祈祷的恩宠，以及宣认的唯一冠冕。\par

然而，最亲爱的弟兄，在你们身上，宣认的冠冕之外还加上了你们劳苦的赏报——一份丰盛的度量，你们将在报应的日子从主那里领受，因为你们藉着书信将自己呈现给我们，向我们显示了那纯洁而蒙福的胸怀，我们素来所认识的，并且照其宽广，与我们一同向天主发出赞美，并非我们配得聆听，而是你们所能表达的。因为你们以言语装饰了我们中那些尚未完全受教之处，并坚固我们承受所忍受的苦难，既因你们确信天上的赏报、殉道的冠冕及天主的国，从你们充满圣神的预言中，在书信中向我们许诺了这一切。这一切将会实现，亲爱的，如果你们在祈祷中记念我们，我相信你们如此行，正如我们确实也如此行一般。\par

因此，极渴望的弟兄，我们收到了你们从\Person{基里诺}和你们自己那里带给我们的礼物——从一切洁净之物所献的祭献。正如\Person{诺厄}向天主奉献，天主喜爱馨香，垂顾他的供物；同样，愿祂也垂顾你们的供物，并乐意将这善工的赏报归给你们。但我请求你们命令将我们写给\Person{基里诺}的信转交出去。最亲爱的且热切愿望的弟兄，我祝你们永远衷心安好，并记念我们。问候所有与你们同在的人。安好。\par

\SourceRecord{Translated by Robert Ernest Wallis.；Ante-Nicene Fathers；Vol. 5.；Edited by Alexander Roberts, James Donaldson, and A. Cleveland Coxe.；Buffalo, NY: Christian Literature Publishing Co.,；1886.；<http://www.newadvent.org/fathers/050678.htm>.}

% structure: p0001 source=h1 target=section reason=page-title

\section{第七十九封书信}

\Emph{\Person{斐理斯}、\Person{雅德尔}、\Person{波利亚努斯}及其他殉道者给\Person{西彼廉}的答复。}\par

\Emph{上文所述的殉道者感谢\Person{西彼廉}所派送的援助。}\par

致我们最亲爱、最挚爱的\Person{西彼廉}，\Person{斐理斯}、\Person{雅德尔}、\Person{波利亚努斯}，连同司铎及所有与我们一起住在\Place{西瓜}矿坑的人，在主内永享安康。亲爱的兄弟，我们通过副执事\Person{赫勒尼安}、我们的兄弟\Person{路西安}和\Person{玛克西穆}回复您的问候，因着您祈祷的帮助，他们强壮而安全；我们从他们那里收到了一笔称为献仪的款项，连同您所写的书信，您在信中屈尊用天上的言语安慰我们，如同安慰儿子一般。我们已感谢并继续感谢全能的天主圣父，借着祂的基督，因为我们因您的讲话而得安慰和坚强，我们从您坦诚的心中请求，请您屈尊在您不断的祈祷中记念我们，使主补足您和我们的宣认所缺欠的，祂已屈尊赐予我们这宣认。请问候所有与您同在的人。亲爱的兄弟，我们在主内衷心向您告别。\Signature{我\Person{斐理斯}书写此信；我\Person{雅德尔}签署；我\Person{波利亚努斯}宣读。我问候我的主上\Person{欧提基安}。}\par

\SourceRecord{Translated by Robert Ernest Wallis.；Ante-Nicene Fathers；Vol. 5.；Edited by Alexander Roberts, James Donaldson, and A. Cleveland Coxe.；Buffalo, NY: Christian Literature Publishing Co.,；1886.；<http://www.newadvent.org/fathers/050679.htm>.}

% structure: p0001 source=h1 target=section reason=page-title

\section{第八十封书信}

\Emph{西彼廉致塞尔吉乌斯、洛加提亚努斯以及狱中其他宣信者。}\par

\Emph{内容提要。——他安慰洛加提亚努斯及其同工，即狱中的宣信者，并藉年长的洛加提亚努斯和斐利奇西穆斯的殉道榜样给予他们勇气。这封信本身表明它是在流亡中写成的。}\par

1. 西彼廉致塞尔吉乌斯、洛加提亚努斯以及其余在主内的宣信者，永享康泰。我问候你们，最亲爱的、最蒙福的弟兄们，我自己也渴望能见到你们，若我目前的处境允许我来到你们中间。因为还有什么比现在靠近你们，让你们用那些纯洁无罪、持守主之信德、拒绝了世俗顺服的双手来拥抱我，更令我渴望和喜乐呢？还有什么比现在亲吻你们的嘴唇（它们曾以荣耀的声音宣认了主），并被你们的眼睛（它们藐视世界，已配得瞻仰天主）所注视，更令人愉快和崇高呢？但是由于我没有机会分享这份喜乐，我寄上此信代替我来到你们的耳中和眼中，借此我祝贺并劝勉你们，要坚定稳固地持守对这天上荣耀的宣认；并且既已踏上主降尊就卑的道路，就应在圣神的德能中前行，以领取冠冕，以主为你们的保护者和向导，祂曾说：\BibleQuote{看，我同你们天天在一起，直到今世的终结。}\BibleRef{玛 28:20} 哦，蒙福的监狱，因你们的临在而变得光明！哦，蒙福的监狱，将天主的人送往天堂！哦，比太阳本身更明亮、比这世界的光更清澈的黑暗，那里现在安置了天主的殿宇，你们的肢体将因神圣的宣认而成圣。\par

2. 现在，除了圣神曾激发你们忍受痛苦的那些神圣规诫和天上的命令之外，不要让任何事物在你们的心中和思想里萦绕。不要有人想到死亡，而要想到不死；不要想到暂时的惩罚，而要想到永恒的荣耀；因为经上记载：\BibleQuote{在主的眼中，圣者的死是珍贵的；}又说：\BibleQuote{心灵破碎是献给天主的祭献：天主不轻视痛悔和谦卑的心。}还有，圣经谈到那些使天主的殉道者成圣、并在痛苦考验中圣化他们的酷刑时，说：\BibleQuote{他们在人眼中虽受了苦，他们的希望却充满不死；他们受了少许惩罚，却要获得丰厚赏报；因为天主考验了他们，发现他们配得上祂。祂试验了他们，如同在炉中炼金，然后接纳了他们犹如全燔祭的牺牲；在适当的时候，他们将受到眷顾。义人将发光，像在禾秸中飞溅的火星。他们将审判万民，统治百姓；他们的主要永远为王。}\BibleRef{智 3:4} 因此，当你们想到你们将同基督审判并作王时，你们必然要欢欣，并在未来喜乐的期盼中践踏当前的痛苦；知道从创世之初就已如此预定：义必要在世界的冲突中受苦，因为一开始，甚至最初，义人亚伯尔就被杀害，此后所有被派遣的义人、先知和宗徒也是如此。对于这一切，主也在自身内立了榜样，教导说，除非那些跟随祂走祂自己的路的人，没有人能达到祂的王国，祂说：\BibleQuote{那爱惜自己性命的，必要丧失性命；那恨恶自己在此世性命的，必要保存性命到永生。}\BibleRef{若 12:28}又说：\BibleQuote{你们不要害怕那杀害肉身，而不能杀害灵魂的；但更要害怕那能使灵魂和肉身都丧亡在地狱里的。}\BibleRef{玛 10:28} 保禄也劝勉我们这些渴望获得主之应许的人，应当在一切事上效法主。他说：\BibleQuote{我们是天主的子女；如果子女，便是承继者，是天主的承继者，与基督同为承继者；只要我们与祂一同受苦，也必与祂一同受光荣。}\BibleRef{罗 8:16} 此外，他又将现今与未来的荣耀相比较，说：\BibleQuote{现今的苦楚，与将来在我们身上要显示的光荣，是不能比拟的。}\BibleRef{罗 8:18} 当我们思量这光耀的荣耀时，我们理当承受一切困苦和迫害；因为，虽然义人的困苦众多，但那信赖天主的人必从这一切中获得解救。\par

3. 有福的妇女们，她们与你们一同分享同样的宣认之荣耀，持守主的信德，比她们的性别更为勇敢，不仅自身临近荣耀的冠冕，还以她们的坚忍为其他妇女提供了榜样！且为使你们的数目毫不欠缺荣耀，为使每个性别和每个年龄都能与你们一同获得尊荣，天主的屈尊就卑也使男孩们参与了荣耀的宣认；为我们展现了某种与昔日阿纳尼雅、阿匝黎雅和米沙耳相似的情形，那杰出的青年们，当他们被投入火窑时，火焰为之退避，烈焰反成凉风，因为主与他们同在，证明了对祂的宣认者和殉道者而言，地狱的炎热毫无能力，而凡信赖天主的人，在一切危险中必始终安然无恙。我恳求你们按照你们的宗教信仰，更细致地思量：这些青年身上究竟有何等信德，竟能蒙受主如此充分的垂爱？因为，我们所有人都当准备面对任何命运，他们对国王说：\Quote{拿步高王，我们不必回答你这件事；因为我们所事奉的天主能救我们脱离烈火巨窑；祂必救我们脱离你手，大王！然而若不，大王，请你知悉，我们决不事奉你的神，也不朝拜你所立的金像。}\BibleRef{达 3:16} 尽管他们相信，并按他们的信德知道自己甚至可能从当前的刑罚中被解救出来，他们仍未以此自夸，也不为此妄求，却说：\Quote{然而若不。} 以免他们的宣认之德因缺乏受苦的见证而减弱，他们补充说天主能做一切；但他们并不信赖于此，以致期望在那一刻被解救；他们所思念的，是那永恒的自由与安全之荣耀。\par

4. 你们也持守这信德，日夜默想，全心预备于天主，只思念将来，轻看现在，好能到达永恒国度的果实，并得到主的拥抱、亲吻和瞻仰，使你们在一切事上追随长老罗加弟亚诺，那荣耀的老人，他以虔诚的勇气和天主的屈尊就卑，为我们的时代增添了荣耀，为你们开辟了道路；他与我们的弟兄斐里奇西摩，一贯温良节制，承受了猛悍人群的攻击，首先在监狱中为你们预备了居所，并在某种程度上为你们标画了道路，如今他也在你们前面前行。为使这一切在你们身上得以圆满，我们以不断的祈祷恳求主，愿祂使那些祂已使之宣认的人，从开端进至最高成就，并终获冠冕。我最亲爱的弟兄们，衷心在主内告别；愿你们获得天上荣耀的冠冕。执事维克多以及同我一起的人问候你们。\par

\SourceRecord{Translated by Robert Ernest Wallis.；Ante-Nicene Fathers；Vol. 5.；Edited by Alexander Roberts, James Donaldson, and A. Cleveland Coxe.；Buffalo, NY: Christian Literature Publishing Co.,；1886.；<http://www.newadvent.org/fathers/050680.htm>.}

% structure: p0001 source=h1 target=section reason=page-title

\section{第八十一封书信}

\Emph{致苏克塞苏斯：关于从罗马传来的迫害消息。}\par

\Emph{论题——西彼廉告知苏克塞苏斯主教，在瓦勒良皇帝下令的残酷迫害中，八月八日罗马的西斯笃主教已殉道；他请求苏克塞苏斯将此事告知其他同僚，以便每位主教鼓励自己的羊群勇敢殉道。}\par

1. 西彼廉问候他的兄弟苏克塞苏斯。最亲爱的兄弟，我未能立即写信给你的原因是：所有神职人员正处于争战最激烈之时，根本无法离开此地；他们每个人都依照心灵的虔诚，准备迎接神圣与天上的光荣。但你要知道：我派往罗马城去探明真相的人已经回来了，他们带回了关于我们命运的任何决定的实情。因为人们众说纷纭，有许多不确定的流言。实际情况如下：瓦勒良向元老院发出谕旨，命令立即处罚主教、司铎和执事；而元老、显贵及罗马骑士应失去他们的尊位，并被剥夺财产；如果他们在财产被夺后仍坚持作基督徒，则还要被斩首；至于贵妇，则应被剥夺财产，流放他乡。此外，皇室仆从中，凡先前或现今承认信仰者，其财产将被没收，并被戴上镣铐送往皇帝庄园。瓦勒良皇帝还在这道诏令中附上了一封他写给各省总督关于我们的信函；我们天天盼望着这些信函的到来，按照我们信德的力量等待着受苦的坚忍，并期望因主的助佑和仁慈而获得永生的冠冕。但你要知道：西斯笃于八月八日在墓窟中殉道，与他一同殉道的还有四位执事。此外，罗马城的行政长官每天都催促着这场迫害；以致如果有人被带到他们面前，他们就被处死，财产被充公。\par

2. 我恳求你将此事转告我们其余的同僚，以便各地通过他们的劝勉，使弟兄们得到加强，为属灵的争战做好准备；好让我们每个人都更轻视死亡，更重视不朽；并怀着完满的信德和全部勇气，将自己奉献于主，在宣认信仰时更喜乐而非恐惧；他们知道，天主的和基督的士兵在此宣认中不是被杀害，而是被加冕。最亲爱的兄弟，愿你常在主里衷心安康。\par

\SourceRecord{Translated by Robert Ernest Wallis.；Ante-Nicene Fathers；Vol. 5.；Edited by Alexander Roberts, James Donaldson, and A. Cleveland Coxe.；Buffalo, NY: Christian Literature Publishing Co.,；1886.；<http://www.newadvent.org/fathers/050681.htm>.}

% structure: p0001 source=h1 target=section reason=page-title

\section{第八十二封书信}

\Emph{致圣职人员和民众：关于他殉道前不久的退隐。}\par

\Emph{论点——当\Person{西彼廉}接近生命终点，回到他的花园时，被告知已派人要带他到\Place{乌提卡}受刑，他便退避了。为免有人认为他这样做是出于怕死，他在这封信中说明了原因，即：他愿在\Place{迦太基}、在他自己的子民眼前殉道，而不在其他地方。公元258年。}\par

\Person{西彼廉}致司铎、执事和全体民众，问候。最亲爱的弟兄们，当我们被告知狱卒已被派来带我到\Place{乌提卡}，而我被那些最亲爱之人的劝告说服，暂时从我的花园退避，因有正当理由，我便同意了。因为主教的适宜在于，在他治理主之教会的城市里，在那里宣认主，并且全体子民因他们主教的宣认在他们面前而得光荣。因为在宣认的那一刻，\OriginalTerm{宣认主教}{confessor-bishop}所说的一切，都是藉天主的默感，在众人口中说的。但我们教会的光荣尊荣，若我——一个被置于另一教会的主教——在\Place{乌提卡}接受判决或宣认，便从那里作为殉道者去到主前，这尊荣就会受损。事实上，为了我自身和你们的缘故，我以不断的恳求祈祷，以我全部的心愿祈求，我能在你们中间宣认，在那里受苦，并如我应当的那样从那里去到主前。因此，我在此一隐藏的退修之所，等候返回\Place{迦太基}的总督到来，好从他那里听到皇帝们关于基督徒平信徒和主教的命令，并说出主在那个时刻愿意说的话。\par

但你们，最亲爱的弟兄们，按照你们从主的诫命从我领受的纪律，以及你们从我讲论中非常多次学到的，要保持平安与恬静；也不要任何人为弟兄们激起骚乱，或自愿把自己交给外邦人。因为当被逮捕并交出时，他应当说话，因为主居住在我们内，在那时刻说话，祂愿意我们宣认而非声称。至于其余的事，在总督因宣认天主之名对我宣判之前，我们应遵守什么，我们将在主的训导下共同安排。愿我们的主，最亲爱的弟兄们，使你们在祂的教会中保持安全，并屈尊护佑你们。愿藉祂的慈悯如此成就。\par

\SourceRecord{Translated by Robert Ernest Wallis.；Ante-Nicene Fathers；Vol. 5.；Edited by Alexander Roberts, James Donaldson, and A. Cleveland Coxe.；Buffalo, NY: Christian Literature Publishing Co.,；1886.；<http://www.newadvent.org/fathers/050682.htm>.}
